% !TEX program = xelatex
% !BIB program = biber
%
% 从裸机到操作系统：基于 xv6 的简单实现
% 模板：ElegantBook（https://github.com/ElegantLaTeX/ElegantBook）
% 编译：XeLaTeX -> Biber -> XeLaTeX -> XeLaTeX（latexmk 会自动完成）
%
% 本书必须用 XeLaTeX（或 LuaLaTeX）编译：ctex 的 Fandol 中文字体在 pdfLaTeX 下不可用。
% 在线平台请在“项目设置/编译器”中选择 XeLaTeX。
\RequirePackage{iftex}
\ifPDFTeX
  \errmessage{This book must be compiled with XeLaTeX (or LuaLaTeX), not pdfLaTeX.
    Please change the compiler to XeLaTeX in the project settings.
    请把编译器改为 XeLaTeX}
  \expandafter\stop
\fi
\documentclass[lang=cn,newtx,10pt,scheme=chinese,cyan,bibend=bibtex]{elegantbook}

\title{从裸机到操作系统}
\subtitle{基于 xv6 的简单实现 \ ——\ Raspberry Pi 3 / AArch64}

\author{Jiming Kang}
\institute{xv6-aarch64 项目}
\date{2026 年 10 月}
\version{0.1}
\bioinfo{源码}{xv6-aarch64（分支 rpi3-vfs）}

\extrainfo{从上电后的第一条指令，到能远程 SSH 登录、连 Wi-Fi、拍照的多核内核。}

\setcounter{tocdepth}{2}
\setcounter{secnumdepth}{2}

\cover{cover.jpg}
\logo{logo.png}

\definecolor{customcolor}{RGB}{23,96,120}
\colorlet{coverlinecolor}{customcolor}

% ---------------------------------------------------------------- 宏包
\usepackage{array,tabularx,longtable}
\usepackage{cprotect}
\usetikzlibrary{arrows.meta,positioning,shapes.geometric,fit,calc,backgrounds,matrix,decorations.pathreplacing}

\addbibresource[location=local]{reference.bib}

% ---------------------------------------------------------------- 代码样式
% 模板默认把 listings 设成 LaTeX 语言并高亮一批 LaTeX 词，这里为 C、汇编、
% 终端和纯文本示意图分别定义样式，并清空模板的 emph 列表。
\definecolor{codekw}{RGB}{160,30,60}
\definecolor{codecm}{RGB}{110,120,130}
\definecolor{codestr}{RGB}{20,110,90}
\definecolor{codebg}{RGB}{248,249,250}

\lstdefinelanguage{armasm}{
  sensitive=false,
  morekeywords={mov,movk,movz,ldr,ldrb,str,strb,ldp,stp,adr,adrp,add,sub,subs,
    and,orr,bic,lsr,lsl,mul,cmp,cbz,cbnz,tbz,tbnz,b,bl,br,blr,ret,eret,msr,mrs,
    isb,dsb,dmb,tlbi,ic,dc,sev,wfe,wfi,svc,hvc,smc,ldar,stlr,ldaxr,stxr,nop},
  morekeywords=[2]{.global,.section,.balign,.align,.org,.word,.quad,.ascii,
    .macro,.endm,.ltorg,.p2align,.text,.data},
  morecomment=[l]{//},
  morecomment=[s]{/*}{*/},
  morestring=[b]",
}

\lstdefinestyle{xbase}{
  emph={}, emphstyle={},
  basicstyle=\ttfamily\small,
  columns=fullflexible,
  keepspaces=true,
  breaklines=true,
  frame=single,
  framerule=0.3pt,
  rulecolor=\color{structurecolor!60},
  backgroundcolor=\color{codebg},
  tabsize=4,
  showstringspaces=false,
  aboveskip=6pt, belowskip=6pt,
  xleftmargin=2pt, xrightmargin=2pt,
}
\lstdefinestyle{xc}{style=xbase, language=C,
  keywordstyle=\color{codekw}\bfseries,
  commentstyle=\color{codecm},
  stringstyle=\color{codestr},
  morekeywords={uint,uint8,uint16,uint32,uint64,uchar,pte_t,pagetable_t,
    asm,volatile,static,inline,__attribute__}}
\lstdefinestyle{xasm}{style=xbase, language=armasm,
  keywordstyle=\color{codekw}\bfseries,
  keywordstyle=[2]\color{structurecolor},
  commentstyle=\color{codecm},
  stringstyle=\color{codestr}}
\lstdefinestyle{xsh}{style=xbase, language=bash,
  keywordstyle=\color{codekw},
  commentstyle=\color{codecm},
  stringstyle=\color{codestr}}
\lstdefinestyle{xtxt}{style=xbase, language={},
  basicstyle=\ttfamily\small, columns=fixed, basewidth=0.5em}
\lstdefinestyle{xmake}{style=xbase, language=make,
  keywordstyle=\color{codekw},
  commentstyle=\color{codecm}}

\lstnewenvironment{ccode}[1][]{\lstset{style=xc,#1}}{}
\lstnewenvironment{asmcode}[1][]{\lstset{style=xasm,#1}}{}
\lstnewenvironment{shcode}[1][]{\lstset{style=xsh,#1}}{}
\lstnewenvironment{txtcode}[1][]{\lstset{style=xtxt,#1}}{}
\lstnewenvironment{makecode}[1][]{\lstset{style=xmake,#1}}{}

% 行内代码：\code{...}，下划线等特殊字符请自行转义
\newcommand{\code}[1]{\texttt{\hyphenchar\font=-1 #1}}
\setlength{\emergencystretch}{3em}
\XeTeXlinebreaklocale "zh"
\XeTeXlinebreakskip = 0pt plus 1pt
\newcommand{\file}[1]{\texttt{\color{structurecolor!80!black}#1}}
\newcommand{\reg}[1]{\texttt{\textsc{#1}}}

% TikZ 公共样式
\tikzset{
  box/.style={draw=structurecolor!70!black, rounded corners=2pt, fill=structurecolor!8,
    align=center, minimum height=7mm, inner sep=3pt, font=\small},
  hbox/.style={box, fill=structurecolor!22},
  gbox/.style={box, draw=gray!70, fill=gray!8},
  arr/.style={-{Stealth[length=2mm]}, thick, draw=structurecolor!70!black},
  darr/.style={arr, dashed},
  lab/.style={font=\footnotesize, text=gray!70!black},
}

\begin{document}

\maketitle
\frontmatter
\tableofcontents

\mainmatter

\chapter{序言}

\begin{introduction}[本章内容]
  \item 为什么要从裸机写起
  \item 为什么选 xv6、AArch64 和 Raspberry Pi 3
  \item 这个项目走过的路
  \item 本书的结构与阅读方法
  \item 开发环境与约定
\end{introduction}

\section{从一条指令开始}

一台树莓派插上电源后，GPU 固件把一段二进制文件拷到物理地址 \code{0x80000}，然后让 CPU0 从那里开始取指。此时没有栈，没有堆，没有 \code{printf}，没有页表，甚至不知道自己处于哪个异常级别。本书要回答的问题是：从这条指令出发，要依次搭起哪些东西，才能得到一个可以多核调度、读写 SD 卡、通过 Wi-Fi 被 SSH 远程登录、还能驱动摄像头的操作系统？

教科书通常按“进程、内存、文件、设备”分章讲操作系统原理，每一章都假设别的部分已经存在。真正动手时顺序恰好相反：要先让串口吐出一个字符，才能知道代码跑到了哪里；要先建好页表，内存分配器才有地方放空闲链表；要先有块设备驱动，文件系统才能读到 superblock。本书沿着这条“依赖链”组织内容，每一章都只依赖前面已经建好的东西。

\section{为什么是 xv6、AArch64 和树莓派}

\paragraph{xv6.} xv6 是 MIT 为操作系统课程 6.S081/6.1810 编写的教学内核\cite{xv6book}，用 C 重写了第六版 Unix 的核心思想。它只有一万行左右，却完整包含了进程、虚拟内存、系统调用、日志文件系统、管道和 shell。代码短，意味着每一个机制都能读到底；结构清楚，意味着加新功能时知道该放在哪一层。

\paragraph{AArch64.} 原版 xv6 运行在 RISC-V 上。本项目的起点是一个面向 QEMU \code{virt} 机器的 AArch64 移植。ARMv8-A 的异常级别（EL0～EL3）、两套页表基址寄存器（\reg{TTBR0\_EL1}/\reg{TTBR1\_EL1}）、内存属性间接表（\reg{MAIR\_EL1}）以及显式的缓存维护指令，都比 RISC-V 更接近手机和服务器上真实运行的系统，也更容易踩到“QEMU 正常、真机失败”的坑——而这些坑恰恰最有教学价值。

\paragraph{Raspberry Pi 3.} 树莓派 3B/3B+ 使用 BCM2837 SoC（四核 Cortex-A53），价格低、资料多，Linux 驱动全部开源，可以对照阅读。它的外设足够丰富：Mini UART、SD 卡控制器、DWC2 USB 主机、板载 SDIO Wi-Fi、I2C、CSI-2 摄像头接口……几乎每一类典型驱动都能在这块板子上找到一个真实的例子。同时 QEMU 提供了 \code{raspi3b} 机器模型，可以先在模拟器里调通流程，再上真机验证。

\section{这个项目走过的路}

项目的提交记录大致反映了“从裸机到操作系统”的顺序（表~\ref{tab:history}）。最早的一个月只做了一件事：把 QEMU \code{virt} 上的 xv6 搬到 \code{raspi3b}；此后两周里，几乎每天都在真机上解决一个新问题。

\begin{table}[htbp]
  \centering
  \caption{主要里程碑（取自 git 提交记录）}
  \label{tab:history}
  \small
  \begin{tabularx}{\linewidth}{@{}lX@{}}
    \toprule
    日期 & 内容 \\
    \midrule
    05-27 & 初始提交：QEMU \code{virt} 上的 xv6-aarch64 \\
    09-23 & 从 \code{virt} 转到 Raspberry Pi 3；SD 卡与 FAT32 分区 \\
    09-24 & 真机串口与内存属性问题；D/I-cache 一致性；SD 命令超时重试 \\
    09-27 & 进程间同步原语（mutex、semaphore、event）；网络协议栈与 \code{ping} \\
    09-28 & TTY 中间层；Linux 风格的总线/设备/驱动模型；SDIO Wi-Fi \\
    09-30 & \code{sshd} 远程登录；SDIO 中断；workqueue 下半部 \\
    10-01 & 每设备私有 work、delayed work、NAPI；FAT32 启动分区挂载到 \code{/boot}；epoll；TFTP；信号；USB 键盘 \\
    10-02 & USB URB 接口；input 子系统与 evdev \\
    10-03 & 多核软实时调度；用户线程 \\
    10-04 & POSIX 源码兼容层 \\
    10-05 & CSI-2 摄像头（OV5647 + Unicam） \\
    \bottomrule
  \end{tabularx}
\end{table}

这些功能都附有一篇专题文档（\file{readme/} 目录），记录了实现思路和调试过程。本书就是把这些文档按依赖顺序重新组织、补全前因后果之后得到的。

\section{本书的结构}

\begin{itemize}
  \item \textbf{第 2 章 项目组织}：目录结构、构建流程、镜像如何生成，以及内核初始化顺序。这一章是全书的地图。
  \item \textbf{第 3 章 基于汇编的启动代码}：固件把控制权交给内核之后，\file{entry.S} 如何统一异常级别、点亮早期串口、建立临时页表、打开 MMU，并唤醒其余三个核；异常向量表和中断控制器也在这一章。
  \item \textbf{第 4 章 MMU}：39 位虚拟地址、三级页表、内存属性；内核页表与用户页表；MMIO 与 DMA 地址；以及三个与缓存和 TLB 有关的真机故障。
  \item \textbf{第 5 章 进程创建}：\code{struct proc}、第一个用户进程、\code{fork()} 的每一步、上下文切换和调度器，以及在此基础上加入的同步原语和软实时调度。
  \item \textbf{第 6 章 文件系统}：xv6 日志文件系统的分层，SD 卡驱动，根文件系统的三种存放方式，VFS 中间层（procfs、ext2、FAT32），以及如何把系统烧录到 SD 卡。
  \item \textbf{第 7 章 驱动架构}：设备模型与上下半部；控制台、UART 与 TTY；USB 与输入子系统；I2C；SDIO 与 USB 两种无线网卡；CSI-2 摄像头。
  \item \textbf{第 8 章 用户进程、线程与 libc}：系统调用桩、\code{init} 与 shell、用户线程、POSIX 兼容层和一个能在 xv6 里运行的小型 C 编译器。
\end{itemize}

\section{阅读方法与约定}

本书以源码为准。正文中的文件路径都相对于仓库根目录，例如 \file{kernel/entry.S}。代码片段尽量保持原样，只删去与讨论无关的行。描述某个函数时，一般按“它解决什么问题 → 调用路径 → 关键代码 → 容易出错的地方”的顺序展开。

\begin{note}
  QEMU 的 \code{raspi3b} 模型与真机有不少差别：QEMU 对内存属性和缓存一致性非常宽容，不模拟 SDIO Wi-Fi、摄像头和板载 USB 网卡，固件 mailbox 也只实现了一部分。本书凡是只在 QEMU 上验证过的结论都会注明。一条经验贯穿全书：\textbf{QEMU 能跑通不代表真机正确}，原子操作、缓存维护和设备时序必须在 Cortex-A53 上验证。
\end{note}

开发环境如下：

\begin{itemize}
  \item 交叉工具链：\code{aarch64-linux-gnu-}、\code{aarch64-elf-} 或 \code{aarch64-unknown-elf-} 均可，Makefile 会自动探测；
  \item 模拟器：\code{qemu-system-aarch64 -machine raspi3b}；
  \item 真机：Raspberry Pi 3B 或 3B+，一张 microSD 卡，一根 3.3\,V 的 USB-TTL 串口线（GPIO14/15，115200 8N1）；
  \item 宿主机：macOS（文中的 \code{diskutil}、\code{/dev/rdiskN} 等命令以 macOS 为例，Linux 下换成 \code{fdisk}/\code{dd} 的等价写法即可）。
\end{itemize}

\section{致谢}

本项目从以下工作中获益良多：MIT 的 xv6\cite{xv6book}；Sergey Matyukevich 的 \emph{raspberry-pi-os}\cite{rpios}，它对树莓派裸机异常与中断的讲解非常清楚；rockytriton 的 LLD 裸机教程\cite{lld}，其中的 EMMC 驱动是本项目 SD 驱动的参考；Circle 裸机框架\cite{circle}；以及 Raspberry Pi Linux 内核中的 \code{brcmfmac}、\code{mt7601u}、\code{bcm2835-unicam}、\code{ov5647} 等驱动\cite{rpilinux}。

\chapter{项目组织}
\label{chap:organization}

\begin{introduction}[本章内容]
  \item 仓库目录结构
  \item 内核子系统与源文件
  \item 从源码到镜像的构建流程
  \item 安装、运行与调试目标
  \item 内核初始化顺序：全书的地图
\end{introduction}

\section{目录结构}

仓库的顶层目录如表~\ref{tab:dirs} 所示。它保留了 xv6 的基本布局：\file{kernel/} 是内核，\file{user/} 是用户程序和用户库，\file{mkfs/} 是在宿主机上运行的文件系统镜像生成工具。在此之上，为了支撑 POSIX 源码兼容、真机部署和调试，又增加了 \file{include/}、\file{firmware/}、\file{tools/} 等目录。

\begin{table}[htbp]
  \centering
  \caption{仓库顶层目录}
  \label{tab:dirs}
  \small
  \begin{tabularx}{\linewidth}{@{}lX@{}}
    \toprule
    路径 & 内容 \\
    \midrule
    \file{kernel/} & 内核源码：启动汇编、异常、内存、进程、文件系统、VFS、网络栈、全部驱动 \\
    \file{user/} & 用户程序（\code{sh}、\code{ls}、\code{sshd}、\code{wifi}、\code{camshot} 等）和用户库（\code{ulib}、\code{printf}、\code{umalloc}、线程、POSIX 包装） \\
    \file{include/} & 面向 POSIX 源码的标准头文件：\code{stdio.h}、\code{unistd.h}、\code{pthread.h}、\code{sys/socket.h} 等 \\
    \file{mkfs/} & 宿主机工具 \code{mkfs}，把用户程序打包成 xv6 原生文件系统镜像 \code{fs.img} \\
    \file{readme/} & 30 余篇专题文档，记录每个功能的设计与调试过程 \\
    \file{firmware/} & 无线网卡固件：BCM43430/43455 的 \code{.bin/.txt/.clm\_blob}，MT7601U 的 \code{mt7601u.bin} \\
    \file{tools/} & 调试脚本：串口跟踪的采集、回放和网页导出（\code{trace\_xv6.py} 等） \\
    \file{ports/} & 向其他开发板移植的实验（如全志 T113） \\
    \file{config.txt} & 树莓派固件配置文件，随内核一起拷到启动分区 \\
    \file{Makefile} & 构建、打包、安装、QEMU 运行的全部规则 \\
    \bottomrule
  \end{tabularx}
\end{table}

\section{内核子系统}

内核约 2.8 万行（含头文件），用户空间约 1.1 万行。按功能可以把内核源文件分成几组（表~\ref{tab:kernel-files}）。阅读时建议先看“启动与异常”和“内存与进程”两组，它们是其余一切的基础。

\begin{table}[htbp]
  \centering
  \caption{内核源文件分组}
  \label{tab:kernel-files}
  \small
  \begin{tabularx}{\linewidth}{@{}p{2.2cm}X@{}}
    \toprule
    子系统 & 主要文件 \\
    \midrule
    启动与异常 & \file{entry.S}、\file{start.c}、\file{main.c}、\file{kernel.ld}、\file{trapasm.S}、\file{trap.c}、\file{timer.c}、\file{bcm2837.c} \\
    内存 & \file{kalloc.c}（物理页分配）、\file{vm.c}（页表）、\file{memlayout.h}、\file{aarch64.h}（系统寄存器与页表位） \\
    进程与同步 & \file{proc.c}、\file{proc.h}、\file{sched.h}、\file{swtch.S}、\file{spinlock.c}、\file{sleeplock.c}、\file{sync.c}、\file{workqueue.c}、\file{exec.c}、\file{syscall.c}、\file{sysproc.c} \\
    文件系统 & \file{bio.c}、\file{log.c}、\file{fs.c}、\file{file.c}、\file{sysfile.c}、\file{pipe.c}、\file{vfs.c}、\file{rootdev.c}、\file{fat32.c}、\file{ext2.c}、\file{xjournal.c} \\
    设备核心 & \file{device.c}、\file{device.h}（总线、设备、驱动、字符设备） \\
    字符设备 & \file{uart.c}、\file{console.c}、\file{tty.c}、\file{pty.c}、\file{input.c} \\
    存储 & \file{sdhost.c}（外置 SD 卡）、\file{sd.c}（早期 Arasan SD 驱动） \\
    USB & \file{dwc2.c}（主机控制器）、\file{usb.c}（USB core 与 URB）、\file{usbkbd.c}、\file{usbnet.c} \\
    无线网络 & \file{arasan\_sdio.c}、\file{sdio.c}、\file{brcmfmac.c}、\file{wpa\_crypto.c}、\file{mt7601u.c} \\
    网络协议 & \file{net.c}、\file{net.h}、\file{sysnet.c}、\file{epoll.c} \\
    摄像头 & \file{mbox.c}（固件 mailbox）、\file{i2c.c}、\file{camsensor.h}、\file{ov5647.c}、\file{unicam.c} \\
    其他 & \file{rng.c}（硬件随机数）、\file{printf.c}、\file{string.c} \\
    \bottomrule
  \end{tabularx}
\end{table}

\section{构建流程}

\subsection{编译选项}

内核是一个独立运行的程序：没有 C 运行库，没有动态链接器，也不能假设 CPU 已经打开浮点单元。关键的编译选项如下：

\begin{makecode}
CFLAGS = -Wall -Werror -Os -g -fno-omit-frame-pointer -mcpu=cortex-a53+nofp
CFLAGS += -ffreestanding -fno-common -nostdlib
CFLAGS += -I. -Iinclude -fno-stack-protector -fno-pie -no-pie
LDFLAGS = -z max-page-size=4096
\end{makecode}

\begin{itemize}
  \item \code{-mcpu=cortex-a53+nofp}：生成 A53 指令，但禁止编译器使用浮点/SIMD 寄存器。内核上下文切换只保存通用寄存器，若编译器偷偷用 \code{q0} 做 \code{memcpy}，用户态的浮点状态就会被破坏。
  \item \code{-ffreestanding -nostdlib}：不假设标准库存在，\code{memset}、\code{memmove} 等由 \file{kernel/string.c} 自己实现。
  \item \code{-fno-pie}：内核链接在固定的高虚拟地址，不需要位置无关代码。
  \item \code{-Werror}：所有警告视为错误，真机调试的成本很高，宁可在编译期多拦住一些问题。
\end{itemize}

\subsection{从目标文件到 \code{kernel8.img}}

图~\ref{fig:build} 是完整的构建流水线。

\begin{figure}[htbp]
  \centering
  \begin{tikzpicture}[node distance=6mm and 9mm]
    \node[box] (src) {\file{kernel/*.c}\\\file{kernel/*.S}};
    \node[box, right=of src] (obj) {\code{*.o}};
    \node[box, below=of obj] (bi) {\code{buildinfo.c}\\编译时间/分支};
    \node[hbox, right=of obj] (elf) {\code{kernel/kernel}\\ELF，VMA 高地址};
    \node[hbox, right=of elf] (img) {\code{kernel8-xv6\_wifi.img}\\裸二进制};
    \node[box, above=of elf] (ld) {\file{kernel.ld}};
    \node[box, below=of elf, yshift=-6mm] (ic) {\code{user/initcode}\\（嵌入内核）};
    \draw[arr] (src) -- node[lab, above]{gcc} (obj);
    \draw[arr] (obj) -- node[lab, above]{ld} (elf);
    \draw[arr] (bi) -- (elf);
    \draw[arr] (ld) -- (elf);
    \draw[darr] (ic) -- (elf);
    \draw[arr] (elf) -- node[lab, above]{objcopy} (img);

    \node[box, below=24mm of src] (usrc) {\file{user/*.c}};
    \node[box, right=of usrc] (libc) {\code{libc.a}\\ulib/usys/printf\\umalloc/thread/posix};
    \node[box, right=of libc] (uprog) {\code{\_sh}、\code{\_ls}…\\-Ttext 0};
    \node[hbox, right=of uprog] (fs) {\code{fs.img}\\32 MiB};
    \draw[arr] (usrc) -- (libc);
    \draw[arr] (libc) -- node[lab, above]{ld} (uprog);
    \draw[arr] (uprog) -- node[lab, above]{mkfs} (fs);
  \end{tikzpicture}
  \caption{构建流水线：上方生成内核镜像，下方生成根文件系统镜像}
  \label{fig:build}
\end{figure}

链接脚本 \file{kernel/kernel.ld} 同时给出两个地址：

\begin{txtcode}
. = 0xffffff8000080000;          /* VMA：运行时的虚拟地址 */
.text : AT(0x80000) { ... }      /* LMA：镜像装入的物理地址 */
\end{txtcode}

ELF 中的所有符号都按高虚拟地址解析，但固件并不解析 ELF，它只是把文件内容原样拷到物理地址 \code{0x80000}。因此需要 \code{objcopy -O binary} 去掉 ELF 头，得到“从第一个字节开始就是第一条指令”的裸镜像。在 MMU 打开之前，代码实际运行在物理地址上，这段代码必须只使用 PC 相对寻址（\code{adr}/\code{adrp}），这是第~\ref{chap:boot} 章要处理的问题。

每次链接都会重新生成 \file{kernel/buildinfo.c}，写入当时的日期和 git 分支、提交号，启动时打印出来：

\begin{txtcode}
kernel image built: 2026-10-05 14:02:11 CDT (rpi3-vfs ebefc6c+dirty)
\end{txtcode}

真机调试时经常要确认“SD 卡上的内核是不是刚编译的那个”，这一行能省去很多困惑。

\subsection{用户程序与文件系统镜像}

用户程序先编译成 \code{.o}，再与 \code{libc.a} 静态链接。每个程序都链接在虚拟地址 0、入口为 \code{main}：

\begin{makecode}
LIBC_OBJS = $U/ulib.o $U/usys.o $U/printf.o $U/umalloc.o $U/thread.o \
            $U/posix.o $U/tweetnacl.o $U/ssh_crypto.o
_%: %.o $(ULIB)
	$(LD) $(LDFLAGS) -N -e main -Ttext 0 -o $@ $^
\end{makecode}

\code{usys.S} 由 Perl 脚本 \file{user/usys.pl} 生成，每个系统调用对应三条指令（第~\ref{chap:user} 章）。最后 \code{mkfs} 把所有程序写进一个 32\,MiB 的 xv6 原生文件系统镜像：

\begin{makecode}
fs.img: mkfs/mkfs $(UPROGS)
	mkfs/mkfs fs.img $(UPROGS)
	truncate -s 32M fs.img
\end{makecode}

\begin{note}
  \code{mkfs} 是宿主机程序。如果仓库在 macOS 上编译过，\file{mkfs/mkfs} 就是一个 Mach-O 可执行文件，换到 Linux 构建时必须先删掉它重新编译，否则 \code{make fs.img} 会报 “Syntax error”。
\end{note}

\section{安装、运行与调试}

Makefile 提供了几组目标，覆盖从模拟器到真机的完整工作流（表~\ref{tab:targets}）。

\begin{table}[htbp]
  \centering
  \caption{常用 make 目标}
  \label{tab:targets}
  \small
  \begin{tabularx}{\linewidth}{@{}lX@{}}
    \toprule
    目标 & 作用 \\
    \midrule
    \code{make} & 构建内核与 \code{fs.img}，并把内核同步到本机 TFTP 目录 \\
    \code{make qemu} & 用 \code{raspi3b} 机器启动，\code{fs.img} 作为虚拟 SD 卡 \\
    \code{make qemu-net} & 额外接入 QEMU 的 USB 网卡（CDC-ECM）和用户态 NAT \\
    \code{make qemu-gdb} & 启动后暂停，等待 GDB 连接 \\
    \code{make install-rpi3} & 把内核、\code{config.txt}、Wi-Fi 固件拷到已挂载的启动分区；若给出 \code{RPI3\_XV6\_DEV} 则同时烧写根文件系统 \\
    \code{make install-rpi3-rawfs} & 只把 \code{fs.img} 写入 SD 卡上的 xv6 原始分区 \\
    \code{make install-tftp} & 把内核拷到宿主机 TFTP 目录，真机上用 \code{tftp} 命令下载到 \code{/boot} \\
    \bottomrule
  \end{tabularx}
\end{table}

QEMU 的 \code{raspi3b} 有两个串口：第一个是 PL011，第二个是 AUX Mini UART。本内核使用 Mini UART，因此 QEMU 参数要把第一个串口丢弃：

\begin{shcode}
qemu-system-aarch64 -machine raspi3b -kernel kernel/kernel8-xv6_wifi.img \
    -display none -serial null -serial mon:stdio \
    -drive file=fs.img,if=sd,format=raw
\end{shcode}

真机通过启动分区里的 \file{config.txt} 指定内核文件名和几个关键参数：

\begin{txtcode}
arm_64bit=1                      # 以 AArch64 状态启动 CPU
enable_uart=1                    # 打开 GPIO14/15 上的串口
uart_2ndstage=1                  # 固件自身的启动日志也输出到串口
core_freq=250                    # 固定 VPU 时钟：Mini UART 波特率依赖它
core_freq_min=250
kernel=kernel8-xv6_wifi.img
camera_auto_detect=0             # 不让固件接管摄像头
\end{txtcode}

其中 \code{core\_freq=250} 很重要：Mini UART 的波特率由 VPU 核心时钟分频得到，如果固件动态调频，串口就会乱码（第~\ref{sec:early-uart} 节）。

\section{内核初始化顺序}
\label{sec:init-order}

\file{kernel/main.c} 是整个内核的入口函数，它的初始化顺序就是全书的目录。下面的代码去掉了注释中的细节，保留了调用顺序：

\begin{ccode}
void main() {
  if(cpuid() == 0){
    kinit1(end, P2V(EARLYTOP));  // 只用启动页表覆盖的前 2 MiB    -> 第 4 章
    kvminit();  kvminithart();   // 建立并切换到最终内核页表      -> 第 4 章
    kinit2(P2V(EARLYTOP), P2V(PHYSTOP));
    device_init(); sdio_bus_init(); usb_bus_init();             // -> 第 7 章
    consoleinit(); ttyinit(); printfinit();                     // 从这里起 printf 可用
    procinit(); workqueue_init();                               // -> 第 5 章
    rnginit(); syncinit(); netinit();
    input_init(); usbkbd_driver_init(); dwc2_driver_init();     // USB 与键盘
    trapinit(); trapinithart(); gicv3init(); gicv3inithart();   // -> 第 3 章
    timerinit();
    binit(); iinit(); fileinit();                               // -> 第 6 章
    sd_driver_init(); arasan_sdio_driver_init();
    fat32init(); rootdev_init();                                // 找到根文件系统，挂起 bootfs
    brcmfmac_driver_init(); mt7601u_driver_init();              // 需要读 bootfs 里的固件
    ov5647_init(); camera_init();                               // 摄像头
    ext2init(); vfsinit();
    userinit();                                                 // 第一个用户进程
    start_secondary_cpus();                                     // 唤醒 CPU1~3
    started = 1;
  } else {
    while(started == 0) ;
    kvminithart(); trapinithart(); gicv3inithart(); timerinit();
  }
  scheduler();
}
\end{ccode}

几条顺序约束值得先记住，后面各章会解释原因：

\begin{itemize}
  \item 页表必须在内存分配器之前就绪，但建页表本身又要分配页——所以分配器分两次初始化（\code{kinit1}/\code{kinit2}）。
  \item 设备核心（\code{device\_init}）要在任何驱动注册之前初始化；控制台要尽早，因为之后的错误都要靠 \code{printf} 报告。
  \item 异常向量、中断控制器、定时器依次初始化，但真正打开 CPU 中断要等到 \code{scheduler()}。
  \item Wi-Fi 驱动需要从 FAT32 启动分区读取固件，所以必须排在 \code{rootdev\_init()} 之后——它扫描 MBR 时顺带挂起 FAT32 启动分区。
  \item 其余三个 CPU 最后才被唤醒，它们只做每核相关的初始化，然后直接进入调度器。
\end{itemize}

\chapter{基于汇编的裸机代码}
\label{chap:boot}
\setcounter{secnumdepth}{3}

\begin{introduction}[本章内容]
  \item 上电、链接地址与装载地址
  \item 异常级别：从 ARM9 的七种模式到 AArch64 的 EL0～EL3
  \item 不依赖任何环境的早期串口
  \item 异常和中断：向量表、现场保存、中断控制器与定时器
\end{introduction}

\ref{sec:xv6-startup} 节讲本项目的启动代码：从固件交出控制权、统一异常级别，到点亮早期串口。\ref{sec:irq} 节讲异常和中断：向量表、现场保存、两级中断控制器和通用定时器。清 BSS、建启动页表、打开 MMU 和多核启动与页表关系密切，放在第~\ref{chap:mmu} 章的 \ref{sec:bootpgt}、\ref{sec:smp-boot} 节。这两节的代码全部来自 xv6-aarch64；它要兼容真机固件、QEMU 和多核，比教学内核复杂，所以每一步都先讲清楚背后的原理，这些原理主要参考了 Sergey Matyukevich 的 \emph{raspberry-pi-os}\cite{rpios}（下文简称 RPi OS）第 1～3 课。RPi OS 后三课的原理也按同样的方式融入了后面各章：第 4 课（进程调度器）在第~\ref{chap:process} 章，第 5 课（用户进程与系统调用）在第~\ref{chap:user} 章，第 6 课（虚拟内存管理）在第~\ref{chap:mmu} 章，那些地方同样只用本项目的代码来讲。

\section{xv6-aarch64 的启动代码}
\label{sec:xv6-startup}
% 3.1 xv6-aarch64 的启动代码（融入 raspberry-pi-os 第 1、2 课的原理）
\subsection{上电之后发生了什么}

树莓派的启动与 PC 不同：最先运行的不是 ARM 核，而是 VideoCore GPU。

\begin{enumerate}
  \item GPU 内部 ROM 从 SD 卡的 FAT32 分区读取 \code{bootcode.bin}，再加载 \code{start.elf}（GPU 固件）；
  \item \code{start.elf} 读取 \code{config.txt}，初始化 SDRAM 和时钟，把内核镜像 \code{kernel8-xv6\_wifi.img} 拷到物理地址 \code{0x80000}；
  \item 固件还在物理地址 0 处放一段很小的 \emph{armstub}。它让四个 Cortex-A53 核一起开始执行：CPU0 跳到 \code{0x80000}，CPU1～3 在 \emph{spin-table} 上循环等待——各自读取 \code{0xe0}、\code{0xe8}、\code{0xf0} 处的 64 位地址，非零就跳过去；
  \item 固件自带的 64 位 armstub 会先把核从 EL3 降到 EL2，因此内核通常在 EL2 拿到控制权；QEMU 也是如此。
\end{enumerate}

项目早期曾使用自制的极简 armstub（\file{kernel/armstub.S}），它只做一次跳转、保留在 EL3。后来发现这样做会让固件的 property mailbox 在真机上完全没有响应——而 SD 时钟查询、USB 上电、摄像头供电都依赖 mailbox。所以现在直接使用固件自带的 armstub，内核启动代码则要能同时处理 EL3、EL2 两种入口。

\paragraph{固件、\code{config.txt} 与镜像文件名.} 树莓派的固件是闭源的，ARM 一侧能控制它的只有启动分区里的 \code{config.txt}。从一张空白 SD 卡启动一个裸机内核，至少需要四个文件：GPU 引导程序 \code{bootcode.bin}、GPU 固件 \code{start.elf}、\code{config.txt} 和内核镜像。固件根据镜像文件名决定 CPU 的执行状态：以 \code{8} 结尾的 \code{kernel8.img} 表示 ARMv8，固件会让四个核以 AArch64 状态启动；本项目用 \code{kernel=} 指定了自己的文件名，同时写了 \code{arm\_64bit=1}。

\paragraph{另一种启动方式：\code{kernel\_old}.} 最简单的裸机教程（如 raspberry-pi-os\cite{rpios}）在 \code{config.txt} 中写 \code{kernel\_old=1}。这时固件把镜像装在物理地址 0，并且不安装 armstub：四个核都从地址 0 开始执行同一段内核代码，而且停留在 EL3。内核必须自己读 \reg{MPIDR\_EL1} 的低 8 位区分核号，把非 0 号核挂进死循环，再自己从 EL3 降到 EL1。这种方式的好处是“什么都由自己掌控”，代价是失去了固件 armstub 提供的 spin-table 多核唤醒协议，也失去了 armstub 对 EL3 安全状态的初始化。本项目选择相反的路线：让固件做它擅长的事，内核只在 \code{\_entry} 中兼容两种入口级别。

\subsection{链接地址与装载地址}

内核在高地址 \code{0xffffff8000080000} 链接，却被装到物理地址 \code{0x80000}。\file{kernel/memlayout.h} 用两个宏描述这种对应关系：

\begin{ccode}
#define KERNBASE  0xffffff8000000000L     // 内核虚拟地址空间起点
#define KERNLINK  (KERNBASE + KERNPA)     // 内核链接地址
#define V2P(a) (((uint64)(a)) - KERNBASE)
#define P2V(a) ((void *)(((char *)(a)) + KERNBASE))
\end{ccode}

\code{KERNBASE} 选在 39 位地址空间的高半部分起点：当 \reg{TCR\_EL1.T1SZ}$=25$ 时，\reg{TTBR1\_EL1} 负责的范围恰好是 \code{0xffffff8000000000}～\code{0xffffffffffffffff}。内核把全部物理内存线性映射到 \code{KERNBASE + PA}，所以物理地址与内核虚拟地址之间只差一个常数。

MMU 打开之前，CPU 用物理地址取指。此时如果执行 \code{ldr x0, =symbol} 得到的是高虚拟地址，访问它会立即出错；只有 PC 相对寻址的 \code{adr}/\code{adrp} 能得到正确的物理地址。这一限制决定了 \file{entry.S} 前半部分的写法。

\paragraph{独立环境下的编译.} 内核不能依赖任何运行时环境，编译时要明确告诉编译器这一点：不链接 C 标准库（库中的大多数函数最终都要调用操作系统，而这里下面已经没有操作系统了）；不使用标准启动文件（设置栈、初始化静态数据、调用 \code{main} 这些事都由启动汇编自己完成）；声明为独立（freestanding）环境，使编译器不对 \code{memcpy} 之类的函数做任何假设；只使用通用寄存器，不让编译器生成 NEON/浮点指令，否则上下文切换就必须额外保存这些寄存器。本项目对应的选项是 \code{-ffreestanding -nostdlib -fno-common} 和 \code{-mcpu=cortex-a53+nofp}（第~\ref{chap:organization} 章）。

\paragraph{为什么要把 ELF 变成裸镜像.} 链接器输出的是 ELF 文件，它需要一个“装载器”去解析段表、把各段放到指定地址。固件不做这件事，它只是把文件按字节原样拷到内存，然后从第一个字节开始执行。所以要用 \code{objcopy -O binary} 抽出所有代码段和数据段，拼成一个裸镜像，并且保证启动代码位于镜像的最开头：常见的做法是把启动代码放进单独的 \code{.text.boot} 段并在链接脚本中排在第一位；本项目则让 \file{entry.o} 成为 \code{OBJS} 的第一个目标文件，\code{\_entry} 自然落在 \code{.text} 的起点。

\subsection{统一异常级别}
\label{sec:xv6-el}

AArch64 有四个异常级别：EL0 运行用户程序，EL1 运行操作系统内核，EL2 运行虚拟机监控器，EL3 运行安全监控固件。xv6 只需要 EL0 和 EL1。

\paragraph{为什么需要异常级别.} 异常级别可以看成处理器的几种执行模式，每种模式只能使用一部分指令和寄存器。操作系统要实现进程隔离，就必须让用户程序运行在特权最低的 EL0：在这一级只能使用通用寄存器、栈指针和普通的访存指令，不能修改页表基址等系统寄存器，因此一个进程既看不到也改不了别的进程的内存。操作系统本身运行在 EL1，可以配置虚拟内存和大多数系统寄存器。EL2 留给虚拟机监控器：宿主系统在 EL2 隔离各个运行在 EL1 的客户系统，就像操作系统隔离用户进程一样。EL3 负责在“安全世界”和“非安全世界”之间切换，两边的隔离由硬件保证，非安全世界的软件无论如何都访问不到安全世界的指令和数据。

\paragraph{ARM9 时代：七种处理器模式.} 四个异常级别是 ARMv8 才有的说法。在 ARM9（ARM920T、ARM926EJ-S 等 ARMv4T/v5 内核，S3C2440、AT91 这类板子上最常见）时代，处理器用的是\emph{处理器模式}（processor mode），当前模式记在 \reg{CPSR} 的低 5 位 M[4:0] 里，一共七种（表~\ref{tab:arm-modes}）：用户模式（User）运行应用程序；系统模式（System）是“有特权的用户模式”；其余五种各对应一类异常——复位和 \code{swi} 指令进入管理模式（Supervisor，SVC），未定义指令进入未定义模式（Undefined），取指或数据访问出错进入中止模式（Abort），普通中断和快速中断分别进入 IRQ 模式和 FIQ 模式。

\paragraph{模式不等于特权级.} 七种模式里只有 User 是非特权的，其余六种的特权完全相同，都能改 \reg{CPSR}、访问 CP15 协处理器（MMU、缓存的控制寄存器都在这里）。所以 ARM9 实际上只有两个特权级，模式之间的区别不在于“能做什么”，而在于\emph{用哪一组寄存器}：每个异常模式有自己私有（banked）的 \code{r13}（SP）、\code{r14}（LR）和 \reg{SPSR}，FIQ 还额外私有 \code{r8}～\code{r12}，这样快速中断处理程序不用保存任何通用寄存器就能开工。异常的类型直接决定进入哪个模式；向量表固定在地址 \code{0x00000000}（或 CP15 选择的高端 \code{0xffff0000}），每类异常占一个 4 字节的槽，只放得下一条跳转指令。

这种设计给操作系统带来不少麻烦。内核必须在启动时依次切进 IRQ、Abort、Undefined 各模式，为每个模式单独设置栈指针；而内核真正想在一个统一的上下文里处理所有陷入，所以 Linux 的 ARM 移植在每个异常入口只往该模式的小栈里存三个字（\code{r0}、\code{lr}、\code{spsr}），随即切换到 SVC 模式，在进程的内核栈上继续处理——七种模式在软件里又被折回成“用户态”和“内核态”两种。另外，特权代码可以直接用 \code{msr cpsr\_c, \#0xd3} 改写 \reg{CPSR} 进入任意模式，模式切换并不一定经过异常。

\paragraph{TrustZone 与虚拟化：再加两种模式.} ARMv6K 和 ARMv7 引入安全扩展（TrustZone），处理器多了一个“安全/非安全”状态位，以及负责在两个世界之间切换的监控模式（Monitor），由 \code{smc} 指令进入。ARMv7 的虚拟化扩展（Cortex-A15 起）又加了一个 Hyp 模式，供虚拟机监控器使用，由 \code{hvc} 进入。为了说清楚这些模式之间的高低，ARMv7 手册开始使用\emph{特权级}（PL0：User；PL1：六种原有特权模式加 Monitor；PL2：Hyp）。到这时已经有九种模式，而“特权高低”和“用哪组寄存器”这两件事仍然搅在一起。

\paragraph{AArch64：把模式折叠成异常级别.} ARMv8 的 AArch64 状态把这两件事拆开了。\emph{异常级别}只表示特权的高低：EL0 对应原来的 User，EL1 一次吸收了 SVC、IRQ、FIQ、Abort、Undefined、System 六种模式，EL2 对应 Hyp，EL3 对应 Monitor。异常的\emph{类型}不再决定模式，而由两样东西表达：

\begin{itemize}
  \item 异常发生时跳到目标级别向量表里的哪一项。\reg{VBAR\_ELx} 指向的表有 16 项，按“来自哪一级、用哪个栈 $\times$ 同步/IRQ/FIQ/SError”排列，每项 128 字节，能直接放下一段保存现场的代码（\ref{sec:irq} 节）；
  \item 同步异常的具体原因写进 \reg{ESR\_ELx} 的 EC 字段：\code{svc} 是 \code{0x15}，来自 EL0 的数据中止是 \code{0x24}，未定义指令是 \code{0x00}……原来要靠“当前处于 Abort 模式还是 Undefined 模式”来区分的事情，现在读一个寄存器就知道。
\end{itemize}

私有寄存器也大大精简：31 个通用寄存器不再分组，每一级只保留自己的栈指针 \reg{SP\_ELx}、返回地址 \reg{ELR\_ELx} 和保存状态 \reg{SPSR\_ELx}，FIQ 的快速寄存器组被取消。原来的 System 模式（特权级、但用 User 的寄存器）也有对应物：EL1t 在 EL1 下借用 \reg{SP\_EL0}，EL1h 使用自己的 \reg{SP\_EL1}（见下文 \code{SPSR} 的模式字段）。最重要的一点是：AArch64 里没有任何指令能直接改写当前级别，升级只能靠异常，降级只能靠 \code{eret}。

\begin{table}[htbp]
  \centering
  \caption{ARM9/ARMv7 处理器模式与 AArch64 异常级别}
  \label{tab:arm-modes}
  \small
  \begin{tabularx}{\linewidth}{@{}llllX@{}}
    \toprule
    模式 & \reg{CPSR}.M & 出现于 & AArch64 & 何时进入（旧架构） \\
    \midrule
    User (usr)       & \code{10000} & ARMv4 & EL0 & 应用程序运行 \\
    FIQ (fiq)        & \code{10001} & ARMv4 & EL1 & 快速中断；私有 \code{r8}～\code{r14} \\
    IRQ (irq)        & \code{10010} & ARMv4 & EL1 & 普通中断 \\
    Supervisor (svc) & \code{10011} & ARMv4 & EL1 & 复位、\code{swi}/\code{svc} 系统调用 \\
    Abort (abt)      & \code{10111} & ARMv4 & EL1 & 预取中止、数据中止 \\
    Undefined (und)  & \code{11011} & ARMv4 & EL1 & 未定义指令 \\
    System (sys)     & \code{11111} & ARMv4 & EL1 & 无对应异常，由特权代码直接切入 \\
    Monitor (mon)    & \code{10110} & ARMv6K/v7 & EL3 & \code{smc}，安全/非安全世界切换 \\
    Hyp (hyp)        & \code{11010} & ARMv7VE & EL2 & \code{hvc}、虚拟化陷入 \\
    \bottomrule
  \end{tabularx}
\end{table}

这张对应关系并不只是类比。ARMv8 处理器仍能在 EL0、EL1 上运行 32 位代码（AArch32 状态），这时旧的模式依旧存在，而且严格按表~\ref{tab:arm-modes} 落在对应的异常级别上；那些私有寄存器也被映射到 64 位通用寄存器上（例如 IRQ 模式的 \code{r13}/\code{r14} 就是 \code{X17}/\code{X16}），所以 64 位内核切换到 32 位进程时不需要另外保存它们。对比早期面向树莓派 1（ARMv6）的 xv6 移植：它要在启动时为 IRQ、Abort、Undefined 各设一个栈，并在每个入口切回 SVC 模式；本项目在 AArch64 上只需要每个核一个 \reg{SP\_EL1}，所有异常都在同一个 EL1 里处理。

\paragraph{异常级别只能靠异常来改变.} 没有更高特权软件的参与，程序无法提升自己的级别——否则隔离就形同虚设。级别只在异常发生和返回时改变。发生异常（非法访存、\code{svc} 指令、硬件中断等）并由 EL$n$ 处理时，硬件依次做三件事：把返回地址存入 \reg{ELR\_ELn}，把当时的处理器状态存入 \reg{SPSR\_ELn}，跳到异常向量；处理程序结束时执行 \code{eret}，硬件再从这两个寄存器恢复状态和地址。关键在于这两个寄存器都是可写的，处理程序不一定要“回到原处”。启动时降级正是利用这一点。

调试启动代码时，一个简单而有用的手段是读出 \reg{CurrentEL}：当前级别保存在它的第 [3:2] 位，右移 2 位即得 0～3。用不同的固件或 \code{config.txt} 启动时，读到的值可能是 3 或 2，下面的代码对两者都要处理：

\file{kernel/entry.S} 的第一件事就是读 \reg{CurrentEL}，不论从 EL3 还是 EL2 进入，都降到 EL1：

\begin{asmcode}
_entry:
        mrs x15, mpidr_el1
        and x15, x15, #0xff          // 记下核号，后面区分主核与副核
        mrs x0, CurrentEL
        cmp x0, #0xc                 // CurrentEL[3:2] == 3 ?
        b.ne 2f
        mov x0, #(1 << 31)           // HCR_EL2.RW：EL1 使用 AArch64
        msr hcr_el2, x0
        msr cntvoff_el2, xzr         // 虚拟计数器偏移清零
        mov x0, #3
        msr cnthctl_el2, x0          // 允许 EL1 访问物理计数器/定时器
        mov x0, #0x431               // SCR_EL3：RES1 | RW | NS
        msr scr_el3, x0
        mov x0, #0x3c5               // 目标 EL1h，屏蔽 D/A/I/F
        msr spsr_el3, x0
        adr x0, 3f
        msr elr_el3, x0
        eret
2:      cmp x0, #0x8                 // EL2 路径与上面相同，只是写 *_el2
        ...
        eret
3:      ldr x0, =((3 << 28) | (3 << 22) | (1 << 20) | (1 << 11))
        msr sctlr_el1, x0            // 只保留 RES1 位：MMU、缓存先全部关闭
        isb
        cbnz x15, secondary_wait     // 副核去等待
        bl early_mini_uart           // 主核点亮串口
        b entry
\end{asmcode}

“降级”的方式很有意思：处理器没有一条“切换到 EL1”的指令，只能伪造一次异常返回。把目标状态写进 \reg{SPSR\_ELx}、把返回地址写进 \reg{ELR\_ELx}，然后执行 \code{eret}，CPU 就会“返回”到 EL1 的标号 \code{3} 处。表~\ref{tab:el-regs} 列出了涉及的寄存器。

\begin{table}[htbp]
  \centering
  \caption{降级到 EL1 时设置的系统寄存器}
  \label{tab:el-regs}
  \small
  \begin{tabularx}{\linewidth}{@{}llX@{}}
    \toprule
    寄存器 & 值 & 含义 \\
    \midrule
    \reg{HCR\_EL2} & bit 31 & RW=1：下一级 EL1 以 AArch64 运行，否则会掉进 AArch32 \\
    \reg{CNTVOFF\_EL2} & 0 & 虚拟计数器 = 物理计数器，内核用的是虚拟定时器 \\
    \reg{CNTHCTL\_EL2} & 3 & EL1/EL0 可以访问物理计数器和物理定时器 \\
    \reg{SCR\_EL3} & \code{0x431} & 非安全状态，下一级为 AArch64 \\
    \reg{SPSR\_ELx} & \code{0x3c5} & M[3:0]=\code{0101} 即 EL1h（使用 \reg{SP\_EL1}）；DAIF 四位全部置 1 \\
    \reg{SCTLR\_EL1} & RES1 位 & 不继承固件留下的 MMU/缓存状态 \\
    \bottomrule
  \end{tabularx}
\end{table}

\code{SPSR} 中的 D、A、I、F 四位分别屏蔽调试异常、SError、IRQ 和 FIQ。启动期间必须全部屏蔽：此时既没有异常向量表，也没有栈，任何一个中断都会让 CPU 跳到不可预知的地方。

\paragraph{\code{SPSR} 里保存的处理器状态.} 写入 \reg{SPSR\_ELx} 的值就是 \code{eret} 之后的 PSTATE，它由三部分组成：条件标志 N、Z、C、V（上一次运算结果是否为负、为零、无符号溢出、有符号溢出，条件分支指令据此跳转）；中断屏蔽位 D、A、I、F；以及模式字段 M[3:0]，它决定返回到哪一级、使用哪个栈指针。EL1 有两种栈模式：EL1h 使用 EL1 自己的 \reg{SP\_EL1}，EL1t 借用 EL0 的 \reg{SP\_EL0}。内核总是使用 EL1h，这样陷入内核时自动换到内核栈，用户栈指针原样留在 \reg{SP\_EL0} 中。有的教学内核只屏蔽 A、I、F 三位（值 \code{7<<6}），本项目连同调试异常 D 一起屏蔽（\code{0xf<<6}）。

\paragraph{\code{SCTLR\_EL1} 的初值.} \reg{SCTLR\_EL1} 控制 EL1 下的 MMU、数据缓存、指令缓存和字节序。手册把其中若干位标为 RES1（保留，必须写 1），即上面那个 \code{(3<<28)|(3<<22)|(1<<20)|(1<<11)}；其余位全写 0，含义是：EL1 和 EL0 的数据访问都是小端，指令缓存、数据缓存和 MMU 全部关闭。不继承固件留下的任何状态，后面的代码才有确定的起点。

\paragraph{\code{HCR\_EL2} 与 \code{SCR\_EL3}.} 即使不实现虚拟机监控器，也必须设置 \reg{HCR\_EL2} 的 RW 位，因为它决定 EL1 运行在 AArch64 还是 AArch32；同理，\reg{SCR\_EL3} 的 RW 位决定 EL2 的执行状态，NS 位让所有更低的级别处于非安全状态。漏掉 RW 位，\code{eret} 之后处理器会以 32 位状态执行 64 位指令，表现为“毫无输出地死掉”。

\subsection{不依赖任何环境的早期串口}
\label{sec:early-uart}

普通的 \code{printf} 依赖一长串设施：有效的栈、清零的 BSS、页表、内存分配器、自旋锁和控制台驱动。如果内核在这些设施就绪之前挂掉，串口上什么也看不到，连“死在哪一步”都无从判断。项目在真机上调通的第一个关键是：用纯汇编写一个不需要栈、不需要内存、不需要 C 运行环境的串口输出。

\paragraph{内存映射的设备寄存器.} BCM2837 上所有外设都通过内存映射寄存器访问：树莓派 3 把 \code{0x3F000000} 以上的一段物理地址留给外设，每个寄存器是 32 位，读写某一位就是配置或启动设备的某项功能。Broadcom 手册中的地址以 \code{0x7E} 开头，那是 VideoCore 总线上的地址，ARM 一侧对应的物理地址要换成 \code{0x3F} 开头（第~\ref{chap:mmu} 章）。

树莓派 3 的 GPIO14/15 可以连到 PL011 UART0 或 AUX Mini UART。本内核使用 Mini UART，即 GPIO14/15 的 ALT5 功能。树莓派有两个 UART：PL011 功能完整，但在树莓派 3 上默认被板载蓝牙占用；Mini UART 结构简单，只需几个寄存器就能工作，代价是它的波特率依赖 VPU 核心时钟。初始化分三步：

\paragraph{第一步：选择 GPIO 复用功能.} 大多数 GPIO 引脚可以连接不同的外设，每个引脚有 0～5 共六种“复用功能”，用 3 位编码，十个引脚共用一个 \code{GPFSELn} 寄存器。GPIO14、15 由 \code{GPFSEL1} 的第 14:12 位和 17:15 位控制；它们的 ALT5 分别是 Mini UART 的 TXD1 和 RXD1，ALT5 的编码是 \code{010}。所以要先把这两个 3 位字段清零，再各写入 2。注意 GPIO 编号与排针针脚号不同：GPIO14 是第 8 脚，GPIO15 是第 10 脚，第 6 脚是地。

\paragraph{第二步：关闭内部上下拉.} 一个悬空的输入引脚读到的值是随机的；上拉或下拉电阻让它在悬空时稳定地读到 1 或 0。串口两根线始终接着对端，不需要上下拉，而且上下拉状态在重启后仍然保持，所以使用前要显式撤销。BCM2837 改变上下拉状态需要一个带时钟的握手过程：先在 \code{GPPUD} 中写入目标状态（0 表示都不要），等待约 150 个周期让控制信号建立；再在 \code{GPPUDCLK0} 中把要修改的引脚对应位写 1，把控制信号“打”进这些引脚，再等 150 个周期保持；最后清除 \code{GPPUD} 和 \code{GPPUDCLK0}。只有收到时钟的引脚会改变，其余引脚保持原状。

\paragraph{第三步：配置 Mini UART.} 顺序很重要：先在 AUX 使能寄存器中打开 Mini UART，因为在此之前它的其他寄存器都不可访问；配置期间先关闭收发器，并永久关闭自动流控（它需要额外的 RTS/CTS 引脚，普通 USB-TTL 线不支持）；关闭收发中断（早期阶段没有中断处理）；把数据位设为 8 位——手册写的是第 0 位，实际硬件要第 0、1 两位都置 1，即写 3；RTS 置为常高；写波特率除数；最后打开收发器。

\code{early\_mini\_uart} 依次完成：

\begin{asmcode}
early_mini_uart:
        mov x13, x30                 // 没有栈，用寄存器保存返回地址
        ldr x9, =0x3f200004          // GPFSEL1：GPIO14/15 = ALT5（功能码 2）
        ldr w10, [x9]
        ldr w11, =0x0003f000
        bic w10, w10, w11
        ldr w11, =0x00012000
        orr w10, w10, w11
        str w10, [x9]
        ...                          // GPPUD/GPPUDCLK0 时序：关闭内部上下拉
        ldr x9, =0x3f215000          // AUX 基址
        ...                          // 使能 AUX、8N1、清 FIFO
        mov w10, #270
        str w10, [x9, #0x68]         // AUX_MU_BAUD：115200 @ 250 MHz
        mov w10, #3
        str w10, [x9, #0x60]         // 打开收发
        mov w12, #'E'
        bl early_mini_uart_putc
        ...
\end{asmcode}

这里用的都是物理地址，因为 MMU 还没打开。Mini UART 的波特率公式是
\[
  \text{baud} = \frac{f_\text{core}}{8\,(\text{divisor}+1)}, \qquad
  \text{divisor} = \frac{250\,000\,000}{8 \times 115\,200} - 1 \approx 270.
\]
它依赖 VPU 核心频率，所以 \file{config.txt} 必须固定 \code{core\_freq=250}。

收发都用轮询：\code{AUX\_MU\_LSR} 的第 5 位为 1 表示发送器空闲，可以向 \code{AUX\_MU\_IO} 写一个字符；第 0 位为 1 表示收到数据，可以从 \code{AUX\_MU\_IO} 读出。早期汇编版只需要发送，所以 \code{early\_mini\_uart\_putc} 只有三条指令：读 LSR、测第 5 位、写 IO。

在没有调试器的裸机上，串口输出几乎是了解程序内部状态的唯一途径。所以一个裸机项目通常要做的第二件事，就是在串口之上移植一个小型 \code{printf}——只需要提供一个发送单个字符的函数。本项目把这件事分成两级：汇编阶段只能输出单个标记字符，C 环境就绪后才由控制台驱动提供完整的 \code{printf}。

此后，启动代码在每个关键节点输出一个字符，形成一串“启动指纹”（表~\ref{tab:boot-marks}）。真机串口上看到的最后一个字符，就指明了内核停在哪一步。

\begin{table}[htbp]
  \centering
  \caption{启动标记字符}
  \label{tab:boot-marks}
  \small
  \begin{tabular}{@{}cl@{\hspace{2em}}cl@{}}
    \toprule
    标记 & 含义 & 标记 & 含义 \\
    \midrule
    \code{E} & 进入 EL1，物理地址串口可用 & \code{1} & \code{kinit1()} 完成 \\
    \code{B} & BSS 清零完成 & \code{2} & 最终内核页表建立 \\
    \code{P} & 启动页表与系统寄存器就绪 & \code{3} & 切换到最终页表 \\
    \code{H} & MMU 已开，跳到高地址并建好栈 & \code{4} & \code{kinit2()} 完成 \\
    \code{0} & 进入 \code{main()} & \code{a}～\code{f} & \code{kalloc.c} 内部细分 \\
    \bottomrule
  \end{tabular}
\end{table}

正常启动时串口先打出 \code{E}，然后是 \code{BPH0abdefc1234}，接着才是 \code{printf} 的第一行 “xv6 kernel is booting”。第~\ref{chap:mmu} 章会讲到一次停在 \code{BPH0abde} 的故障：它说明第一次 \code{kfree()} 已经完成 \code{memset}，却没能拿到 \code{kmem.lock}。

MMU 打开之后，物理地址 \code{0x3f215000} 不再可用。\code{boot\_uart\_mark()} 改用高地址 \code{KERNBASE + 0x3f215000}，所以启动页表里专门为 Mini UART 所在的 2\,MiB 块加了一个设备映射。


% 3.2 异常和中断（全部以 xv6-aarch64 源码讲解）
\section{异常与中断}
\label{sec:irq}

与设备打交道的大多数场合都是异步的：内核向设备发出一条命令，设备并不立即完成，而是在工作结束后“通知”处理器。这种通知打断了正在执行的指令流，迫使处理器转去执行一段处理程序，所以叫“中断”。对操作系统来说，最重要的中断源是定时器：调度器靠周期性的定时器中断度量每个进程运行了多久，并在时间片用完时把 CPU 交给别人。

本节沿着一次中断从发生到返回的路径，依次讲解 xv6-aarch64 的实现：\file{kernel/trapasm.S}（向量表与现场保存）、\file{kernel/trap.c}（分发）、\file{kernel/bcm2837.c}（中断控制器）、\file{kernel/timer.c}（定时器）和 \file{kernel/spinlock.c}（中断屏蔽的嵌套）。

\subsection{异常与中断区别}

在 ARMv8 中，中断只是“异常”的一种。异常分为四类，表~\ref{tab:exc-kinds} 列出了它们在 xv6-aarch64 中的用途。

\begin{table}[htbp]
  \centering
  \caption{四类异常在 xv6-aarch64 中的用途}
  \label{tab:exc-kinds}
  \small
  \begin{tabularx}{\linewidth}{@{}llX@{}}
    \toprule
    类型 & 性质 & 在 xv6-aarch64 中 \\
    \midrule
    同步异常 & 由当前指令引起 & 用户态 \code{svc \#0} 即系统调用；用户态访存错误时打印诊断并杀死进程；内核态同步异常直接 \code{panic} \\
    IRQ & 外部硬件产生，与当前指令无关 & 通用定时器、Mini UART、DWC2 USB、Arasan SDIO、Unicam CSI、核间中断（IPI） \\
    FIQ & 高优先级的“快速中断” & 不使用 \\
    SError & 外部产生的系统错误 & 不处理，进入即停机 \\
    \bottomrule
  \end{tabularx}
\end{table}

同步异常的具体原因记录在 \reg{ESR\_EL1}（异常综合征寄存器）中：高 6 位 [31:26] 是异常类别 EC，低 25 位 ISS 是该类别下的细节；出错的指令地址在 \reg{ELR\_EL1}，出错的数据地址在 \reg{FAR\_EL1}。xv6 只识别一种“正常”的同步异常——EC $=$ \code{0x15}（21），即来自 AArch64 EL0 的 \code{svc}；其余同步异常都按错误处理。内核态的同步异常不尝试恢复：

\begin{ccode}
void kerneltrap()
{
  uint64 esr = r_esr_el1();
  w_esr_el1(0);
  if(intr_get() != 0)
    panic("kerneltrap: interrupts enabled");
  uint64 ec = (esr >> 26) & 0x3f;
  uint64 iss = esr & 0x1ffffff;
  printf("esr %p %p %p\n", esr, ec, iss);
  printf("elr=%p far=%p\n", r_elr_el1(), r_far_el1());
  panic("kerneltrap");
}
\end{ccode}

\subsection{异常向量}

每类异常都要有自己的处理入口；而且同一类异常，从不同的执行状态进入时也要用不同的入口。以运行在 EL1 的内核为视角，进入异常时的执行状态有四种：EL1t（EL1 使用 \reg{SP\_EL0}）、EL1h（EL1 使用自己的 \reg{SP\_EL1}，内核正是这种模式）、来自 64 位 EL0、来自 32 位 EL0。四类异常乘以四种状态，共 16 个入口，它们排成一张\emph{异常向量表}：表必须 2\,KiB 对齐，每个入口占 \code{0x80} 字节（表~\ref{tab:vectors}）。

\begin{table}[htbp]
  \centering
  \caption{\reg{VBAR\_EL1} 向量表布局}
  \label{tab:vectors}
  \small
  \begin{tabular}{@{}lllll@{}}
    \toprule
    偏移 & 来源 & 同步 & IRQ & 本内核 \\
    \midrule
    \code{0x000} & 当前 EL，使用 \reg{SP\_EL0} & \code{+0x000} & \code{+0x080} & 未用（\code{b .}） \\
    \code{0x200} & 当前 EL，使用 \reg{SP\_ELx} & \code{el1sync} & \code{el1irq} & 内核态异常/中断 \\
    \code{0x400} & 低 EL，AArch64 & \code{el0sync} & \code{el0irq} & 系统调用、用户异常/中断 \\
    \code{0x600} & 低 EL，AArch32 & \code{+0x600} & \code{+0x680} & 未用 \\
    \bottomrule
  \end{tabular}
\end{table}

每组的第三、四个入口是 FIQ 和 SError。\file{kernel/trapasm.S} 用 \code{.balign} 保证对齐，每个入口只放一条跳转指令，真正的处理放在表外（\code{0x80} 字节只够放 32 条指令）：

\begin{asmcode}
.global alltraps
.balign 0x800
alltraps:
// Current EL with sp0
  b .
.balign 0x80
  b .
.balign 0x80
  b .
.balign 0x80
  b .

// Current EL with spx
.balign 0x80
  b el1sync
.balign 0x80
  b el1irq
.balign 0x80
  b .
.balign 0x80
  b .

// Lower EL using aarch64
.balign 0x80
  b el0sync
.balign 0x80
  b el0irq
.balign 0x80
  b .
.balign 0x80
  b .

// Lower EL using aarch32
  ...                       // 四个 b .
\end{asmcode}

只有四个入口真正被使用。其余 12 个入口是 \code{b .}，一旦进入就原地停住。这在调试初期足够了，但也有缺点：FIQ 或 SError 发生时，现象只是“毫无输出地死机”。更好的做法是让这些入口也保存现场，然后打印入口编号、\reg{ESR\_EL1} 和 \reg{ELR\_EL1} 再停机。

\subsection{保存寄存器状态}

处理完异常回到原来的代码时，所有通用寄存器必须和异常发生前一模一样，否则一个与当前代码毫不相干的中断就会悄悄改变它的行为。进入异常时，硬件只做三件事：把返回地址存入 \reg{ELR\_EL1}，把原来的 PSTATE 存入 \reg{SPSR\_EL1}，换到 \reg{SP\_EL1}。通用寄存器要由软件自己保存。

\file{trapasm.S} 为两种来源准备了两对宏。来自内核的异常用 \code{savereg}/\code{restorereg}，现场压在当前内核栈上：

\begin{asmcode}
.macro savereg
        sub sp, sp, #272              // 17 组 × 16 字节
        stp x0, x1, [sp, #16 * 0]
        stp x2, x3, [sp, #16 * 1]
        ...
        stp x28, x29, [sp, #16 * 14]
        mrs x9, elr_el1
        mrs x10, spsr_el1
        add x11, sp, #272             // 异常发生前的内核 sp
        stp x30, x9, [sp, #16 * 15]
        stp x10, x11, [sp, #16 * 16]
.endm
\end{asmcode}

来自用户态的异常用 \code{usavereg}/\code{urestorereg}，区别是最后一个字保存的是用户栈指针 \reg{SP\_EL0}，返回时也要写回：

\begin{asmcode}
.macro usavereg
        sub sp, sp, #272
        stp x0, x1, [sp, #16 * 0]
        ...
        stp x28, x29, [sp, #16 * 14]
        mrs x9, elr_el1
        mrs x10, spsr_el1
        mrs x11, sp_el0               // 用户栈指针
        stp x30, x9, [sp, #16 * 15]
        stp x10, x11, [sp, #16 * 16]
.endm

.macro urestorereg
        ldp x30, x9, [sp, #16 * 15]
        ldp x10, x11, [sp, #16 * 16]
        msr elr_el1, x9
        msr spsr_el1, x10
        msr sp_el0, x11
        ldp x0, x1, [sp, #16 * 0]
        ...
        ldp x28, x29, [sp, #16 * 14]
        add sp, sp, #272
.endm
\end{asmcode}

这 272 字节的布局与 \code{struct trapframe}（\code{x0}～\code{x30}、\code{elr}、\code{spsr}、\code{sp}）完全一致。\code{allocproc()} 把每个进程的 trapframe 放在它内核栈的最顶端，而用户态陷入时内核栈恰好是空的，所以 \code{usavereg} 压栈的位置正好就是 \code{p->trapframe}。\file{trap.c} 开头的注释画出了两种现场同时存在时一个内核栈的样子：

\begin{txtcode}
低地址   p->kstack
           |  普通内核栈空间
           |  +---------------------------+  <- 内核态中断后 sp = S - 272
           |  | el1irq savereg 临时现场   |
           |  | x0..x30, elr, spsr, sp    |
           |  +---------------------------+  <- 中断前 sp = S
           |  函数调用帧、局部变量……
           +------------------------------
           |  struct trapframe            |  <- p->trapframe（usavereg 的位置）
           |  x0..x30, elr, spsr, 用户 sp |
高地址   p->kstack + PGSIZE
\end{txtcode}

也就是说：进程在内核里执行系统调用时又被中断，会在同一个内核栈上再压一层 \code{savereg} 现场；如果中断发生在调度器中，用的则是这个 CPU 自己的调度器栈（\code{stack0}）。

有了调度器之后，“以一个进程的身份进入中断、以另一个进程的身份离开”是完全合法的：\code{eret} 依赖的 \code{elr} 和 \code{spsr} 都保存在当前进程自己的内核栈上，所以 \reg{ELR\_EL1}、\reg{SPSR\_EL1} 必须和通用寄存器一起保存和恢复。新进程第一次回到用户态走的是 \code{usertrapret}，它直接把 \code{sp} 指向 trapframe，借用 \code{el0irq} 的后半段返回：

\begin{asmcode}
el0irq:
        usavereg
        mov x0, sp
        bl userirq
trapret:
        urestorereg
        eret

.global usertrapret
usertrapret:                    // void usertrapret(struct trapframe *tf)
        mov sp, x0
        b trapret
\end{asmcode}

\subsection{设置向量表}

向量表准备好了，还要告诉处理器它在哪里：把地址写入 \reg{VBAR\_EL1}（向量基址寄存器）。这是每个核各自的系统寄存器，所以每个 CPU 都要执行一次：

\begin{ccode}
void trapinithart(void)
{
  w_vbar_el1((uint64)alltraps);
}

static inline void w_vbar_el1(uint64 x)     // kernel/aarch64.h
{
  asm volatile("msr vbar_el1, %0" : : "r" (x) );
}
\end{ccode}

\code{alltraps} 是链接时确定的高虚拟地址，所以这一步必须在 MMU 打开、跳到高地址之后进行。主核在 \code{main()} 中依次调用 \code{trapinit()}（只初始化 \code{tickslock}）和 \code{trapinithart()}；副核被唤醒后只调用 \code{trapinithart()}。

\subsection{屏蔽与解除屏蔽中断}

有些代码绝不能被中断打断。例如 \code{usavereg} 执行到一半时若又来一个中断，现场就会被覆盖。所以硬件在进入任何异常时都会自动屏蔽中断；保存好现场之后再解除屏蔽、允许嵌套，也是完全合法的。

PSTATE 中有四个屏蔽位：D（调试异常）、A（SError，也叫异步中止）、I（IRQ）、F（FIQ），用 \code{msr daifset/daifclr, \#imm} 设置或清除，立即数是一个 4 位掩码，从高到低对应 D、A、I、F。xv6 的封装在 \file{kernel/aarch64.h} 中：

\begin{ccode}
static inline void intr_on()  { asm volatile("msr daifclr, #0xf" ::: "memory"); }
static inline void intr_off() { asm volatile("msr daifset, #0xf" ::: "memory"); }

static inline int intr_get()          // IRQ 是否打开？
{
  uint64 x = daif();
  return ((x >> 6) & 0x2) == 0;       // DAIF 寄存器中 I 在 bit 7
}
\end{ccode}

注意这里用的是 \code{\#0xf}，一次改变全部四位，而不只是 I 位（\code{\#2}）。对 xv6 来说效果相同，因为它只用 IRQ；更精确的写法应只操作 I 位，以免意外打开调试异常或 SError。

\subsubsection{自旋锁与中断的嵌套}


如果一个 CPU 持有某把自旋锁时被中断，而中断处理程序又去获取同一把锁，这个 CPU 就会永远自旋下去。所以 xv6 规定：持有任何自旋锁期间，本 CPU 的中断都是关闭的。\code{acquire()} 调用 \code{push\_off()}，\code{release()} 调用 \code{pop\_off()}，二者成对嵌套计数：

\begin{ccode}
void push_off(void)
{
  int old = intr_get();
  intr_off();
  if(mycpu()->noff == 0)
    mycpu()->intena = old;          // 只记录最外层进入前的状态
  mycpu()->noff += 1;
}

void pop_off(void)
{
  struct cpu *c = mycpu();
  if(intr_get())
    panic("pop_off - interruptible");
  if(c->noff < 1)
    panic("pop_off");
  c->noff -= 1;
  if(c->noff == 0 && c->intena)     // 最外层释放、且原来是开着的，才重新打开
    intr_on();
}
\end{ccode}

\subsubsection{中断在哪里被打开}


\begin{itemize}
  \item 启动时，\file{entry.S} 通过 \code{SPSR = 0x3c5} 进入 EL1，四位全部屏蔽；
  \item 每个 CPU 第一次真正接收中断，是在 \code{scheduler()} 循环中调用 \code{intr\_on()}；
  \item 用户态陷入后，\code{usertrap()} 识别出系统调用，在调用 \code{syscall()} 之前打开中断，让长时间的系统调用也能被设备中断和定时器打断；返回用户态前再 \code{intr\_off()}，因为 \code{urestorereg} 期间不能被打断；
  \item \code{kerneltrap()} 进入时检查 \code{intr\_get()}，若中断是开着的就说明某处逻辑出错，直接 \code{panic}。
\end{itemize}

\subsection{配置中断控制器}

设备一般不直接中断 CPU，而是经过中断控制器：控制器负责打开或关闭各个中断源，并告诉 CPU 当前是谁在请求中断。BCM2837 没有 ARM 标准的 GIC，而是两个级联的模块（图~\ref{fig:irq}）：

\begin{itemize}
  \item \textbf{legacy 中断控制器}（\code{0x3f00b000}）管理 SoC 共享外设：Mini UART（IRQ 29）、DWC2 USB（IRQ 9）、Arasan SDIO（IRQ 62）、Unicam CSI1（IRQ 39）等。中断号 0～31 由 \code{ENABLE\_IRQS\_1}/\code{IRQ\_PENDING\_1} 管理，32～63 由 \code{ENABLE\_IRQS\_2}/\code{IRQ\_PENDING\_2} 管理。所有外设中断汇总成一根 “GPU IRQ”；
  \item \textbf{ARM-local 中断路由}（\code{0x40000000}）位于 CPU 簇旁，管理每个核私有的中断源：通用定时器、核间 mailbox、PMU，以及从 legacy 控制器汇总来的 GPU IRQ，最终连到四个核的 IRQ/FIQ 输入。
\end{itemize}

\begin{figure}[htbp]
  \centering
  \begin{tikzpicture}[node distance=5mm and 12mm]
    \node[box, minimum width=28mm] (uart) {Mini UART (29)};
    \node[box, minimum width=28mm, below=2mm of uart] (usb) {DWC2 USB (9)};
    \node[box, minimum width=28mm, below=2mm of usb] (sdio) {Arasan SDIO (62)};
    \node[box, minimum width=28mm, below=2mm of sdio] (csi) {Unicam CSI1 (39)};
    \node[hbox, right=of usb, yshift=-3mm, minimum height=22mm, text width=24mm] (legacy) {legacy\\中断控制器\\\code{0x3f00b000}};
    \foreach \n in {uart,usb,sdio,csi} \draw[arr] (\n.east) -- (legacy.west |- \n.east);
    \node[box, above=12mm of legacy, minimum width=26mm] (timer) {通用定时器 CNTV};
    \node[box, above=2mm of timer, minimum width=26mm] (mbox) {mailbox 0（IPI）};
    \node[hbox, right=of legacy, minimum height=40mm, text width=26mm, yshift=12mm] (local) {ARM-local\\中断路由\\\code{0x40000000}\\[2mm]\footnotesize bit 3：虚拟定时器\\bit 4：mailbox 0\\bit 8：GPU IRQ};
    \draw[arr] (legacy.east) -- node[lab, below]{GPU IRQ} (legacy.east -| local.west);
    \draw[arr] (timer.east) -- (timer.east -| local.west);
    \draw[arr] (mbox.east) -- (mbox.east -| local.west);
    \node[box, right=of local, text width=16mm] (cpus) {CPU0～3\\IRQ 输入};
    \draw[arr] (local) -- (cpus);
  \end{tikzpicture}
  \caption{BCM2837 的两级中断结构}
  \label{fig:irq}
\end{figure}

为了少改 \file{trap.c}，项目保留了原 GIC 版本的函数名 \code{gicv3init()}、\code{gicv3inithart()}、\code{gic\_iar()}、\code{gic\_eoi()}，实现换成了 \file{kernel/bcm2837.c}。全局初始化先关掉 legacy 控制器的全部中断，再只打开已经有驱动的几个：

\begin{ccode}
void gicv3init(void)
{
  irqwrite(DISABLE_IRQS_1, ~0U);
  irqwrite(DISABLE_IRQS_2, ~0U);
  irqwrite(DISABLE_BASIC_IRQS, ~0U);
  irqwrite(ENABLE_IRQS_1, (1U << UART0_IRQ) | (1U << DWC2_IRQ));
  irqwrite(ENABLE_IRQS_2, 1U << (SDIO_IRQ - 32));
}
\end{ccode}

后来加入的驱动（如摄像头）在 probe 成功后调用 \code{bcm2837\_enable\_irq(irq)} 打开对应位，按中断号小于还是大于 32 选择 \code{ENABLE\_IRQS\_1} 或 \code{\_2}。

每个核还要各自配置 ARM-local 路由：把本核的虚拟定时器接到 IRQ，并打开 mailbox 0 作为“重新调度”核间中断：

\begin{ccode}
#define CORE_TIMER_CONTROL(c) (0x40 + 4 * (c))
#define CORE_MBOX_CONTROL(c)  (0x50 + 4 * (c))
#define CORE_IRQ_SOURCE(c)    (0x60 + 4 * (c))
#define CORE_MBOX0_SET(c)     (0x80 + 0x10 * (c))   // 写 1 发出
#define CORE_MBOX0_CLR(c)     (0xc0 + 0x10 * (c))   // 写 1 清除
#define CORE_VIRTUAL_TIMER    (1 << 3)
#define CORE_MBOX0            (1 << 4)

void gicv3inithart(void)
{
  int cpu = cpuid();
  localwrite(CORE_TIMER_CONTROL(cpu), CORE_VIRTUAL_TIMER);
  localwrite(CORE_MBOX0_CLR(cpu), ~0U);     // 先清掉遗留的 mailbox 位
  localwrite(CORE_MBOX_CONTROL(cpu), 1);
}

void send_resched_ipi(int cpu)              // 让 cpu 尽快重新调度
{
  asm volatile("dsb sy" ::: "memory");      // 先让 need_resched 的写入可见
  localwrite(CORE_MBOX0_SET(cpu), 1);
}
\end{ccode}

一个外设中断要真正到达 CPU，三道开关缺一不可：设备自己的中断使能位（例如 Mini UART 的 \code{MU\_IER}）、legacy 控制器的 \code{ENABLE\_IRQS} 位，以及 CPU 的 PSTATE.I $=0$。

\subsection{通用 IRQ 处理程序}

所有 IRQ 都进入同一个入口，按来源分为两条路径：

\begin{asmcode}
el1irq:                    // 内核执行时来的 IRQ
        savereg
        mov x0, sp
        bl kernelirq
        restorereg
        eret

el0irq:                    // 用户执行时来的 IRQ
        usavereg
        mov x0, sp
        bl userirq
trapret:
        urestorereg
        eret
\end{asmcode}

两个 C 处理函数都先调用 \code{devintr()} 找出并服务中断源，再决定是否让出 CPU：

\begin{ccode}
void userirq(void)
{
  struct proc *p = myproc();
  devintr();
  if(resched_pending())      // 时间片用完，或更高优先级进程就绪（定时器或 IPI）
    yield();
  if(p->killed)
    exit(-1);
  intr_off();
}

void kernelirq()
{
  devintr();
  // 内核态也可以被抢占：只有不持任何自旋锁时中断才是开着的
  if(myproc() != 0 && myproc()->state == RUNNING && resched_pending())
    yield();
}
\end{ccode}

\code{devintr()} 先问中断控制器“是谁”，再调用对应设备的处理函数，返回 2 表示定时器、1 表示其他设备、0 表示无法识别：

\begin{ccode}
int devintr()
{
  uint32 iar = gic_iar();
  uint32 irq = gic_iar_irq(iar);
  int dev = 0;

  if(irq == UART0_IRQ){         uartintr();        dev = 1; }
  else if(irq == DWC2_IRQ){     dwc2_irq();        dev = 1; }
  else if(irq == SDIO_IRQ){     arasan_sdio_irq(); dev = 1; }
  else if(irq == CSI1_IRQ){     unicam_irq();      dev = 1; }
  else if(irq == TIMER0_IRQ){
    int logical_tick = timerintr();
    if(cpuid() == 0 && logical_tick)
      clockintr();
    sched_tick();
    dev = 2;
  } else if(irq == IPI_RESCHED_IRQ){
    resched_ipi_ack();          // need_resched 已由发送方设置
    dev = 1;
  } else if(irq == 1023){
    // 没有可处理的中断
  } else if(irq){
    printf("unexpected interrupt irq=%d\n", irq);
  }
  if(dev)
    gic_eoi(iar);
  return dev;
}
\end{ccode}

“是谁”的判断在 \code{gic\_iar()} 中。它先查本核的 ARM-local 源寄存器（定时器和 IPI 是每核私有的），再查 legacy 控制器的挂起寄存器，把两套硬件的状态翻译成一个类似 GIC 的中断号：

\begin{ccode}
uint32 gic_iar(void)
{
  int cpu = cpuid();
  uint32 local = localread(CORE_IRQ_SOURCE(cpu));
  if(local & CORE_VIRTUAL_TIMER)  return TIMER0_IRQ;        // 27
  if(local & CORE_MBOX0)          return IPI_RESCHED_IRQ;   // 伪中断号 1000
  if(irqread(IRQ_PENDING_1) & (1U << UART0_IRQ))       return UART0_IRQ;
  if(irqread(IRQ_PENDING_1) & (1U << DWC2_IRQ))        return DWC2_IRQ;
  if(irqread(IRQ_PENDING_2) & (1U << (SDIO_IRQ - 32))) return SDIO_IRQ;
  if(irqread(IRQ_PENDING_2) & (1U << (CSI1_IRQ - 32))) return CSI1_IRQ;
  return 1023;
}
\end{ccode}

BCM2837 没有 GIC 那样的 EOI 握手，\code{gic\_eoi()} 是空函数：中断要在\emph{源设备}处清除——定时器靠重新装载，UART 靠读空接收 FIFO，USB 靠写 \code{HCINT}，IPI 靠写 mailbox 清除寄存器。设备的挂起位清除之后，legacy 控制器中对应的位才会消失。多个中断可能同时挂起，\code{devintr()} 每次只处理优先级最高的一个；其余的会在 \code{eret} 之后立即再次触发中断，所以不会丢失。

\subsubsection{同步异常：系统调用路径}


同步异常走 \code{el0sync}/\code{el1sync} 入口，保存现场的方式与 IRQ 相同，只是调用的 C 函数不同。用户程序执行 \code{svc \#0} 时，系统调用号放在 \code{x7}，参数在 \code{x0}～\code{x5}：

\begin{txtcode}
user:  mov x7, #SYS_write ; svc #0
  -> VBAR_EL1 + 0x400 -> el0sync -> usavereg（现场 = trapframe）
  -> usertrap(): ESR_EL1.EC == 0x15
       -> intr_on(); syscall(): syscalls[x7]()，返回值写回 trapframe->x0
       -> 若唤醒了更高优先级进程：yield()
  -> intr_off(); urestorereg -> eret -> 回到 svc 的下一条指令
\end{txtcode}

\begin{ccode}
void usertrap(void)
{
  struct proc *p = myproc();
  uint64 esr = r_esr_el1();
  uint64 ec = (esr >> 26) & 0x3f;
  w_esr_el1(0);
  if(ec == 21){                         // SVC from AArch64
    if(p->killed) exit(-1);
    intr_on();
    syscall();
    if(resched_pending()) yield();
  } else {
    printf("usertrap(): unexpected esr=%p ec=%p pid=%d\n", esr, ec, p->pid);
    printf("            elr=%p far=%p\n", r_elr_el1(), r_far_el1());
    ...
    uvmdump(p->vm->pagetable, p->pid, p->name, "usertrap");
    p->killed = 1;
  }
  if(p->killed) exit(-1);
  intr_off();
}
\end{ccode}

硬件保存在 \reg{ELR\_EL1} 中的已经是 \code{svc} 的下一条指令地址，所以不需要像 RISC-V 那样手工给 PC 加 4。

\begin{note}
  早期版本的 \code{usertrap()} 先读出 EC 再清零 \reg{ESR\_EL1}，打印错误信息时又重新读一次寄存器，结果永远显示 \code{ec=0}。现在的代码先把完整的 ESR 存进局部变量 \code{esr}。调试信息本身的正确性，同样需要检查。
\end{note}

\subsection{初始化定时器}

xv6-aarch64 没有使用 BCM 自己的 System Timer，而是使用 ARM 架构定义的通用定时器（Generic Timer）。它的好处是：每个核都有自己的比较器，中断经 ARM-local 的 bit 3 直接送到本核，不经过 legacy 控制器；寄存器是系统寄存器，用 \code{mrs}/\code{msr} 访问，不需要 MMIO 映射。内核使用虚拟定时器那一组：

\begin{itemize}
  \item \reg{CNTFRQ\_EL0}：计数器频率（由固件设置）；
  \item \reg{CNTVCT\_EL0}：当前虚拟计数值，所有核共享，可用来测量时间（\code{clock\_us()}、\code{delay()} 都基于它）；
  \item \reg{CNTV\_TVAL\_EL0}：距离触发还剩多少个计数，写入即开始倒计时；
  \item \reg{CNTV\_CTL\_EL0}：bit 0 ENABLE、bit 1 IMASK、bit 2 ISTATUS。
\end{itemize}

\code{timerinit()} 的顺序是“关闭—装载—打开”，避免改写比较值的过程中产生一个中间中断：

\begin{ccode}
#define CNTV_CTL_ENABLE   (1<<0)
#define CNTV_CTL_IMASK    (1<<1)

void timerinit()
{
  disable_timer();
  reload_timer();
  enable_timer();
}

static void disable_timer()
{
  uint64 c = r_cntv_ctl_el0();
  c &= ~CNTV_CTL_ENABLE;
  c |= CNTV_CTL_IMASK;
  w_cntv_ctl_el0(c);
}

static void reload_timer()
{
  uint64 interval = 10000;                                   // 微秒，即 10 ms
  uint64 interval_clk = interval * (r_cntfrq_el0() / 1000000);
  w_cntv_tval_el0(interval_clk);
}

static void enable_timer()
{
  uint64 c = r_cntv_ctl_el0();
  c |= CNTV_CTL_ENABLE;
  c &= ~CNTV_CTL_IMASK;
  w_cntv_ctl_el0(c);
}
\end{ccode}

\code{timerinit()} 由每个核各自调用。执行完之后定时器已经开始倒计时，但 PSTATE.I 仍然屏蔽着 IRQ，要等进入 \code{scheduler()} 调用 \code{intr\_on()} 时 CPU 才会真正接收它。

\subsection{处理定时器中断}

定时器到期后，中断经 ARM-local 路由到本核，\code{gic\_iar()} 返回 \code{TIMER0\_IRQ}，\code{devintr()} 调用 \code{timerintr()}：

\begin{ccode}
int timerintr()
{
  int cpu = cpuid();
  int logical_tick;

  disable_timer();
  reload_timer();                       // 重新装载即清除了中断条件
  timer_divider[cpu]++;
  logical_tick = timer_divider[cpu] == 10;
  if(logical_tick)
    timer_divider[cpu] = 0;
  if(cpu == 0)
    workqueue_timer_tick();             // 推进 delayed-work 时钟
  enable_timer();
  return logical_tick;                  // 每 10 次（100 ms）返回 1
}
\end{ccode}

一次 10\,ms 的硬件中断在 \code{devintr()} 中引出三件事：

\begin{enumerate}
  \item \textbf{系统时间}：每 10 次中断构成一个 xv6 逻辑 tick。只有 CPU0 调用 \code{clockintr()} 把全局 \code{ticks} 加一并 \code{wakeup(\&ticks)}，避免四个核把时间推进四倍。\code{sleep()}、\code{uptime()} 因此保持 xv6 原来 100\,ms 一个 tick 的语义：
\begin{ccode}
void clockintr()
{
  acquire(&tickslock);
  ticks++;
  wakeup(&ticks);
  release(&tickslock);
  net_udp_timeout_tick();       // 检查带超时的 UDP 接收是否到期
}
\end{ccode}
  \item \textbf{下半部时钟}：CPU0 推进 workqueue 的 delayed-work 时钟，TCP 重传、Wi-Fi 连接超时等到期的工作被放进各自的工作队列，由内核线程执行，而不是在中断里直接运行；
  \item \textbf{调度节拍}：每次中断都调用 \code{sched\_tick()}，扣减当前进程的时间片、累计实时任务的运行时间，必要时设置 \code{need\_resched}（第~\ref{chap:process} 章）。随后 \code{userirq()}/\code{kernelirq()} 检查 \code{resched\_pending()}，决定是否 \code{yield()}。
\end{enumerate}

\code{reload\_timer()} 写的是相对倒计时，所以下一次中断发生在“本次处理程序重新装载的时刻 $+$ 10\,ms”，处理程序的延迟会累积成时钟漂移。若要更稳定的时基，可以改写绝对比较值 \reg{CNTV\_CVAL\_EL0}，每次在上一个截止时间上加固定间隔。

\subsection{小结}

把本节的代码按 \code{main()} 中的调用顺序串起来：

\begin{ccode}
trapinit();        // 初始化 tickslock
trapinithart();    // VBAR_EL1 = alltraps                   （每个核）
gicv3init();       // legacy 控制器：全关，再打开 UART/USB/SDIO
gicv3inithart();   // ARM-local：本核虚拟定时器 + mailbox 0  （每个核）
timerinit();       // 10 ms 虚拟定时器开始倒计时            （每个核）
...
scheduler();       // intr_on()：第一次真正接收中断
\end{ccode}

顺序不能随意交换：先安装向量表，保证中断来了有入口；再配置控制器和路由；再启动定时器；最后才解除 PSTATE 的屏蔽。要收到第一个定时器中断，下面这些条件必须同时满足：

\begin{itemize}
  \item 处于 EL1h，且 EL2 已允许 EL1 访问计数器和定时器（\reg{CNTHCTL\_EL2}）；
  \item \reg{VBAR\_EL1} 指向 2\,KiB 对齐的 \code{alltraps}；
  \item \code{CORE\_TIMER\_CONTROL(cpu)} 打开了虚拟定时器路由；
  \item \reg{CNTV\_TVAL\_EL0} 装入了非零值，\reg{CNTV\_CTL\_EL0} 的 ENABLE $=1$、IMASK $=0$；
  \item PSTATE.I $=0$，且当前内核栈有效，有空间保存 272 字节的现场。
\end{itemize}

调试中断问题时，可以在启动阶段打印 \reg{CurrentEL}、\reg{VBAR\_EL1}、\reg{DAIF}、\reg{CNTFRQ\_EL0} 和 \reg{CNTV\_CTL\_EL0}；对未知中断打印 \code{IRQ\_PENDING\_1/2} 和 \code{CORE\_IRQ\_SOURCE}。但不要在定时器中断里持续 \code{printf}：串口远比中断慢，会改变调度时序，甚至带来新的锁问题。


\setcounter{secnumdepth}{2}

\chapter{MMU}
\label{chap:mmu}

\begin{introduction}[本章内容]
  \item 39 位虚拟地址与三级页表
  \item 硬件如何遍历页表、块映射
  \item 描述符格式与内存属性（MAIR）
  \item 清 BSS、启动页表（恒等映射加高地址映射）、TCR/MAIR/SCTLR、跳到高地址
  \item 多核启动：两道闸门
  \item 两阶段的物理内存初始化
  \item 最终内核页表与用户页表
  \item 第一个进程、\code{fork} 与地址空间切换；TLB；缺页与按需分配
  \item MMIO、总线地址与 DMA
  \item 三个缓存与 TLB 故障
\end{introduction}

\section{AArch64 的地址翻译}

\subsection{地址空间划分}

本内核使用 4\,KiB 页、39 位虚拟地址。\reg{TCR\_EL1} 中 \code{T0SZ}$=$\code{T1SZ}$=25$，即每个地址空间有 $2^{64-25}=2^{39}$ 字节（512\,GiB）。AArch64 用虚拟地址的高位选择页表基址寄存器：

\begin{itemize}
  \item 高位全 0：\code{0x0000000000000000}～\code{0x0000007fffffffff}，由 \reg{TTBR0\_EL1} 翻译，用于\textbf{用户空间}，每个进程一张；
  \item 高位全 1：\code{0xffffff8000000000}～\code{0xffffffffffffffff}，由 \reg{TTBR1\_EL1} 翻译，用于\textbf{内核空间}，所有进程共享一张；
  \item 中间的地址不合法，访问即产生翻译错误。
\end{itemize}

两个基址寄存器分开带来一个很大的好处：进程切换时只需改写 \reg{TTBR0\_EL1}，内核映射保持不变，也不用像 xv6-riscv 那样在每个用户页表里复制一份内核映射和 trampoline 页。

\subsection{三级页表}

39 位地址用 4\,KiB 粒度时，从第 1 级开始遍历，共三级（没有第 0 级）。每级页表占一页，含 512 个 8 字节表项，用虚拟地址中的 9 位索引（图~\ref{fig:va}）：

\begin{figure}[htbp]
  \centering
  \begin{tikzpicture}[x=3.3mm,y=8mm]
    \foreach \x/\w/\t/\c in {0/6/{高位: 全0或全1}/gray!15, 6/9/{L1 索引 VA[38:30]}/structurecolor!15,
        15/9/{L2 索引 VA[29:21]}/structurecolor!25, 24/9/{L3 索引 VA[20:12]}/structurecolor!35,
        33/12/{页内偏移 VA[11:0]}/gray!10}{
      \draw[fill=\c] (\x,0) rectangle ++(\w,1);
      \node[font=\scriptsize, text width=\w*3.1mm, align=center] at (\x+\w/2,0.5) {\t};
    }
    \node[lab, below, text width=140mm, align=center] at (22.5,-0.1) {每级 9 位 = 512 项；L1 一项管 1\,GiB，L2 一项管 2\,MiB，L3 一项管 4\,KiB};
  \end{tikzpicture}
  \caption{39 位虚拟地址的划分}
  \label{fig:va}
\end{figure}

\file{kernel/aarch64.h} 用两个宏取出各级索引：

\begin{ccode}
#define PXMASK          0x1FF                    // 9 位
#define PXSHIFT(level)  (39-(level)*9)           // L1: 30, L2: 21, L3: 12
#define PX(level, va)   ((((uint64)(va)) >> PXSHIFT(level)) & PXMASK)
\end{ccode}

第 1、2 级表项可以是\emph{表描述符}（指向下一级页表）或\emph{块描述符}（直接映射 1\,GiB/2\,MiB 的连续物理内存）；第 3 级表项是\emph{页描述符}。启动页表用 2\,MiB 块，最终页表全部用 4\,KiB 页。

\subsection{硬件如何遍历页表}
\label{sec:hw-walk}

虚拟内存给每个进程一个错觉：整个地址空间都归它一个人用。进程发出的每个地址都是虚拟地址，由 MMU（内存管理单元）在访存的路上翻译成物理地址，进程本身完全察觉不到。这样做同时解决了两个问题：一是\emph{隔离}，进程只能访问自己页表里映射了的内存，碰不到别人的数据，也碰不到内核；二是\emph{分配简单}，每个进程的程序都可以从地址 0 开始链接，不必关心物理内存里哪些地方已经被别人占了。

MMU 翻译一个地址的过程叫\emph{页表遍历}（table walk），完全由硬件完成，软件只负责事先把页表在内存里摆好、再把根页表的物理地址写进 \reg{TTBR0\_EL1} 或 \reg{TTBR1\_EL1}。以用户地址为例：

\begin{enumerate}
  \item 从 \reg{TTBR0\_EL1} 取出 L1 页表的\textbf{物理}地址；
  \item 用 VA[38:30] 作下标读出 L1 表项。若它是表描述符，其中 [47:12] 位就是 L2 页表的物理地址；
  \item 用 VA[29:21] 在 L2 页表中读出表项，同理得到 L3 页表的物理地址；
  \item 用 VA[20:12] 在 L3 页表中读出页描述符，得到物理页框号；
  \item 物理页框号拼上 VA[11:0] 的页内偏移，就是最终的物理地址。
\end{enumerate}

每一步读到的地址都是物理地址——MMU 遍历页表时不可能再去翻译页表自己的地址。这就是为什么 \code{walk()} 往表项里写的是 \code{V2P(pagetable)}，读出来以后又要 \code{P2V()} 才能被内核软件访问（\ref{sec:walk-mappages} 节）。途中任何一级遇到无效表项（bit~0 为 0），遍历就停下来，CPU 产生一个\emph{翻译错误}（translation fault），并把出错的级别记在 \reg{ESR\_EL1} 里、出错的虚拟地址记在 \reg{FAR\_EL1} 里。

\paragraph{块映射.} 有时候要映射一大片连续的物理内存，例如启动时的整个内核镜像。这时可以让 L2 表项直接是一个\emph{块描述符}，指向一个 2\,MiB 对齐的物理块，遍历在第 3 步就结束：省去一级页表，页内偏移变成 VA[20:0] 共 21 位（$2^{21}=2$\,MiB）。同理，L1 块描述符映射 1\,GiB。描述符最低两位区分这几种情况：在 L1/L2 上 \code{01} 是块、\code{11} 是表；在 L3 上只有 \code{11} 合法，表示一个 4\,KiB 页——同样是 \code{11}，含义随所在的级别而变。xv6-aarch64 只在启动页表里用块映射（\ref{sec:bootpgt} 节）。

\subsection{描述符的属性位}

表~\ref{tab:pte-bits} 是项目中用到的属性位。

\begin{table}[htbp]
  \centering
  \caption{页表项属性位（\file{kernel/aarch64.h}）}
  \label{tab:pte-bits}
  \small
  \begin{tabularx}{\linewidth}{@{}llX@{}}
    \toprule
    宏 & 位 & 含义 \\
    \midrule
    \code{PTE\_VALID} / \code{PTE\_TABLE} & [1:0] & L1/L2：\code{01} 块，\code{11} 表；L3：\code{11} 页（\code{PTE\_V}） \\
    \code{PTE\_INDX(i)} & [4:2] & AttrIndx：选择 \reg{MAIR\_EL1} 中的第 $i$ 个内存类型 \\
    \code{PTE\_U}、\code{PTE\_RO} & [7:6] & AP：EL0 是否可访问、是否只读 \\
    \code{PTE\_SH\_INNER} & [9:8] & 共享属性：\code{11} 为 Inner Shareable \\
    \code{PTE\_AF} & [10] & Access Flag：必须置 1，否则首次访问触发 Access Fault \\
    \code{PTE\_PXN}、\code{PTE\_UXN} & 53、54 & 内核态 / 用户态禁止执行 \\
    \bottomrule
  \end{tabularx}
\end{table}

\subsection{内存属性间接表}

页表项里并不直接写“可缓存”“设备内存”之类的属性，而是写一个 3 位的索引，真正的属性放在 \reg{MAIR\_EL1} 的 8 个字节里。项目定义了三种类型：

\begin{ccode}
#define AI_DEVICE_nGnRnE_IDX  0x0
#define AI_NORMAL_NC_IDX      0x1
#define AI_NORMAL_IDX         0x2
#define MT_DEVICE_nGnRnE  0x00    // 设备内存：不聚合、不重排、无提前写确认
#define MT_NORMAL_NC      0x44    // 普通内存，不可缓存
#define MT_NORMAL         0xff    // 普通内存，Inner/Outer Write-Back

#define PTE_DEVICE    PTE_INDX(AI_DEVICE_nGnRnE_IDX)
#define PTE_NORMAL_NC PTE_INDX(AI_NORMAL_NC_IDX)
#define PTE_NORMAL   (PTE_INDX(AI_NORMAL_IDX) | PTE_SH_INNER)
\end{ccode}

用途划分很简单：RAM、内核、用户页面和页表本身用 \code{PTE\_NORMAL}；UART、GPIO、SD 控制器等外设寄存器用 \code{PTE\_DEVICE}；摄像头 DMA 缓冲区用 \code{PTE\_NORMAL\_NC}，这样就不必在每帧前后做缓存维护（第~\ref{chap:drivers} 章）。

\section{启动页表}
\label{sec:bootpgt}

第~\ref{chap:boot} 章的启动代码把处理器降到 EL1、点亮串口之后，内核还运行在物理地址上，MMU 和缓存都是关着的。本节接着讲从这里到跳进高地址的三步：清 BSS、建临时页表、打开 MMU。

\subsection{清 BSS 与早期栈}

C 语言假定未初始化的全局变量为 0，这些变量位于 \code{.bss} 段。镜像文件里并不包含 BSS 的内容，所以要先清零。链接脚本导出了 \code{bss\_start} 和以 8 字节为单位的 \code{bss\_size}：

\begin{asmcode}
entry:
        adrp x1, bss_start          // PC 相对：MMU 未开时得到物理地址
        ldr w2, =bss_size
1:      cbz w2, 2f
        str xzr, [x1], #8
        sub w2, w2, #1
        b 1b
2:      ...                         // 输出 'B'
\end{asmcode}

把需要清零的数据单独放在 \code{.bss} 段，是为了节省镜像空间：ELF 只记录这个段的大小，不存它的内容。代价是装载之后必须由启动代码自己清零，所以链接脚本要导出段的起止位置，并让起点按 8 字节对齐。对齐在这里不只是为了效率：MMU 关闭时所有数据访问都按 Device 内存处理，非对齐访问会直接触发对齐异常，因此一次写 8 字节的 \code{str xzr} 要求地址必须是 8 的倍数。

栈也是同样的道理：C 代码需要一个可用的栈指针。最简单的做法是把 \code{sp} 设到内核镜像上方一个“足够远”的固定地址（例如内核装在 0 时设为 4\,MiB），栈向下增长也碰不到内核。本项目要支持四个核同时运行，所以在 BSS 中定义了 \code{stack0[4096 * NCPU]}，每个核按核号取其中一页（见 \ref{sec:mmu-on} 节末尾的 \code{\_start}）。

\subsection{两棵临时页表}


\file{entry.S} 在 MMU 打开前用汇编构建两棵只有两级的小页表，四张页表页（\code{l1entrypgt}、\code{l2entrypgt}、\code{l1kpgt}、\code{l2kpgt}）静态定义在 \file{start.c} 中、按页对齐。

\begin{figure}[htbp]
  \centering
  \begin{tikzpicture}[node distance=4mm and 10mm]
    \node[box, text width=36mm] (lo) {低地址 VA\\\code{0x0}～\code{end}};
    \node[box, text width=36mm, below=8mm of lo] (hi) {高地址 VA\\\code{KERNBASE}～\code{KERNBASE+end}};
    \node[box, text width=36mm, below=4mm of hi] (hu) {\code{KERNBASE+0x3f200000}};
    \node[hbox, right=22mm of lo, yshift=-7mm, text width=34mm] (pa) {物理内存\\\code{0x0}～\code{end}（2\,MiB 块）};
    \node[gbox, right=22mm of hu, text width=34mm] (mmio) {GPIO/Mini UART\\\code{0x3f200000} 2\,MiB 块};
    \draw[arr] (lo) -- node[lab, above, sloped]{TTBR0 恒等} (pa);
    \draw[arr] (hi) -- node[lab, below, sloped]{TTBR1} (pa);
    \draw[arr] (hu) -- node[lab, above]{Device} (mmio);
  \end{tikzpicture}
  \caption{启动页表：同一段物理内存同时映射到低地址和高地址}
  \label{fig:bootpgt}
\end{figure}

如图~\ref{fig:bootpgt}，恒等映射保证“打开 MMU 的那条指令之后还能取到下一条指令”，高地址映射保证“跳到链接地址之后一切照常”。两者都用 2\,MiB 块描述符，所以每棵树只要一张 L1 和一张 L2：

\begin{asmcode}
        adrp x0, l2entrypgt
        mov  x1, #0                       // 起始 PA
        ldr  x2, =V2P_WO(end)-1           // 结束 PA
        lsr  x3, x1, #PXSHIFT(2)
        and  x3, x3, #PXMASK              // 起始 L2 索引
        lsr  x4, x2, #PXSHIFT(2)
        and  x4, x4, #PXMASK              // 结束 L2 索引
        mov  x5, #(PTE_AF | PTE_NORMAL | PTE_VALID)
        orr  x6, x1, x5                   // 第一个块描述符
l2epgt_loop:
        str  x6, [x0, x3, lsl #3]         // l2entrypgt[idx] = 块描述符
        add  x3, x3, #1
        add  x6, x6, #0x200000            // 下一个 2 MiB
        cmp  x3, x4
        b.ls l2epgt_loop
\end{asmcode}

启动页表只覆盖内核镜像所在的前 2\,MiB（\code{EARLYTOP}$=$\code{0x200000}）。\code{main()} 一开始就检查内核镜像连同 BSS 没有越过这个边界——否则后面的分配器会把内核自己的内存当成空闲页发出去：

\begin{ccode}
if(V2P(end) > EARLYTOP){
  for(char *m = "\nkernel image extends past EARLYTOP\n"; *m; m++)
    boot_uart_mark(*m);
  for(;;) asm volatile("wfe");
}
\end{ccode}

\subsection{打开 MMU：TCR、MAIR 与 SCTLR}
\label{sec:mmu-on}

页表摆好以后，还要告诉 MMU 怎样解释它们。主核在 \code{\_entry}、副核在 \code{entryothers} 里执行的是同一段代码：

\begin{asmcode}
        adrp x0, l1entrypgt
        adrp x1, l1kpgt
        msr  ttbr0_el1, x0        // 低地址：恒等映射
        msr  ttbr1_el1, x1        // 高地址：内核
        ldr  x0, =(TCR_T0SZ(25)|TCR_IRGN0(1)|TCR_ORGN0(1)|TCR_SH0(3)|TCR_TG0(0)| \
                   TCR_T1SZ(25)|TCR_IRGN1(1)|TCR_ORGN1(1)|TCR_SH1(3)|TCR_TG1(2)|TCR_IPS(0))
        msr  tcr_el1, x0
        ldr  x1, =((MT_DEVICE_nGnRnE<<(8*AI_DEVICE_nGnRnE_IDX)) | \
                   (MT_NORMAL_NC<<(8*AI_NORMAL_NC_IDX)) | (MT_NORMAL<<(8*AI_NORMAL_IDX)))
        msr  mair_el1, x1
        isb
        ldr  x1, =_start          // 链接地址（高地址）
        dsb  ish
        tlbi vmalle1              // 丢掉固件可能留下的翻译
        dsb  ish
        isb
        ic   iallu
        dsb  sy
        isb
        mrs  x0, sctlr_el1
        orr  x0, x0, #1           // M：打开 MMU
        orr  x0, x0, #(1 << 2)    // C：数据缓存
        orr  x0, x0, #(1 << 12)   // I：指令缓存
        msr  sctlr_el1, x0
        isb
        br   x1                   // 跳到高地址
\end{asmcode}

\reg{TCR\_EL1}（Translation Control Register）的各个字段决定了前面所有的“规则”：

\begin{itemize}
  \item \code{T0SZ}$=$\code{T1SZ}$=25$：两个地址空间都是 39 位，于是遍历从 L1 开始、一共三级；
  \item \code{TG0}$=0$、\code{TG1}$=2$：两边都是 4\,KiB 粒度。两个字段的编码不一样（TG0 中 \code{00} 是 4\,KiB，TG1 中 \code{10} 才是 4\,KiB），抄错一个数字，高地址的页表就会被按 64\,KiB 或 16\,KiB 粒度解释；
  \item \code{IRGN}/\code{ORGN}$=1$、\code{SH}$=3$：MMU 遍历页表时本身也要读内存，这几位规定这些读取走 Write-Back 缓存、属于 Inner Shareable 域。它们必须和页表页本身的映射属性（\code{PTE\_NORMAL}）一致，否则 CPU 写进缓存的新表项，MMU 可能看不到（第~\ref{sec:fault1} 节的故障与此同源）；
  \item \code{IPS}$=0$：物理地址 32 位（4\,GiB），对 BCM2837 足够。
\end{itemize}

\reg{MAIR\_EL1} 把三种内存类型写进第 0、1、2 字节，页表项的 AttrIndx 就是这里的下标。\reg{SCTLR\_EL1} 是总开关：M 位打开地址翻译，C、I 位打开数据和指令缓存。打开前的 \code{tlbi vmalle1} 和 \code{ic iallu} 清掉固件阶段可能残留的 TLB 和指令缓存；写 \reg{SCTLR\_EL1} 之后的 \code{isb} 保证后续取指都已经按新的翻译进行。此时 PC 仍是低地址，靠 \reg{TTBR0} 的恒等映射继续执行，直到 \code{br x1} 跳进高地址——这就是启动页表要同时建两棵树的原因。

\subsection{跳到高地址}

\code{ldr x1, =\_start} 从文字池中取出链接时确定的高虚拟地址；这里不能写 \code{b \_start}：打开 MMU 并不会改变 PC，而 \code{b} 和 \code{adr} 都是相对当前 PC 计算目标的，会落到“没开 MMU 时 \code{\_start} 所在的低地址”。必须事先用 \code{ldr} 从文字池取出链接器算好的高地址绝对值；这条 \code{ldr} 本身也是 PC 相对的，所以要在打开 MMU 之前执行。\code{br x1} 之后，CPU 就运行在 \reg{TTBR1\_EL1} 映射的高地址上了。\code{\_start} 为每个核在 \code{stack0} 中划出 4\,KiB 栈，然后跳进 C 函数 \code{main()}：

\begin{asmcode}
_start:
        ldr x0, =stack0
        mov x1, #1024*4
        mrs x2, mpidr_el1
        and x2, x2, #0x3
        add x2, x2, #1
        mul x1, x1, x2
        add x0, x0, x1               // sp = stack0 + (cpuid + 1) * 4096
        mov sp, x0
        mov w0, #'H'
        bl boot_uart_mark
        b main
\end{asmcode}

\section{多核启动}
\label{sec:smp-boot}

四个核在上电后都会执行 \code{\_entry}——CPU0 由 armstub 直接跳入，CPU1～3 则要等主核在 spin-table 里写入地址。如果像 \code{kernel\_old=1} 那样没有 armstub，四个核会同时从同一个入口开始执行，内核要立刻用 \reg{MPIDR\_EL1} 的低 8 位区分核号，把除 0 号核以外的所有核挂进死循环——这是单核教学内核的标准做法，但被挂起的核就再也没有机会醒来了。本项目需要四核同时调度，所以副核不能死循环，而要在一个可以被唤醒的地方等待。整个过程有两道“闸门”：

\paragraph{第一道闸门：启动页表。} 副核到达 \code{\_entry} 后，统一好异常级别（\ref{sec:xv6-el} 节）就进入 \code{secondary\_wait}，用 \code{ldar}（带 acquire 语义的加载）读取标志 \code{secondary\_release}，为 0 就 \code{wfe} 休眠。主核建好启动页表后，用 \code{stlr}（带 release 语义的存储）把标志置 1，再执行 \code{sev} 唤醒所有在 \code{wfe} 上休眠的核：

\begin{asmcode}
        adrp x10, secondary_release            // 主核
        add x10, x10, :lo12:secondary_release
        mov w11, #1
        stlr w11, [x10]
        sev

secondary_wait:                                // 副核
        adrp x10, secondary_release
        add x10, x10, :lo12:secondary_release
4:      ldar w11, [x10]
        cbnz w11, entryothers
        wfe
        b 4b
\end{asmcode}

acquire/release 配对保证：副核看到标志为 1 时，主核在此之前写入的页表内容也一定可见。

\paragraph{第二道闸门：内核初始化。} 主核在 \code{main()} 的最后调用 \code{start\_secondary\_cpus()}，把 \code{\_entry} 的物理地址写进三个 spin-table 槽，用 \code{dc cvac} 把这几行缓存写回内存（此时副核的 MMU 和缓存都没打开，只能看到内存里的值），再 \code{sev}：

\begin{ccode}
static void start_secondary_cpus(void) {
  volatile uint64 *spin = (volatile uint64*)P2V(0xe0);
  uint64 entry = V2P((uint64)_entry);
  for(int cpu = 1; cpu < NCPU; cpu++){
    spin[cpu - 1] = entry;
    asm volatile("dc cvac, %0" :: "r"(&spin[cpu - 1]) : "memory");
  }
  asm volatile("dsb sy\n\tsev" ::: "memory");
}
\end{ccode}

副核进入 \code{main()} 后在 \code{while(started == 0)} 处等待主核完成全部初始化，然后只做每核相关的部分：切换到最终页表、安装异常向量、打开本核的定时器中断路由，最后进入调度器。真机串口上会看到：

\begin{txtcode}
PH0hart 3 starting
hart 2 starting
hart 1 starting
\end{txtcode}

三个副核几乎同时出发，谁先抢到 \code{printf} 的锁谁先打印，所以顺序不固定；\code{PH0} 是不经过锁的早期标记，会与其他输出交错。

\section{两阶段的物理内存初始化}

建立最终页表需要分配页表页，分配页表页又需要物理页分配器——这是一个“先有鸡还是先有蛋”的问题。xv6 的解决办法是把分配器分两次喂：

\begin{ccode}
kinit1(end, P2V(EARLYTOP));          // 阶段一：内核末尾到 2 MiB，启动页表已覆盖
kvminit();                           // 用这些页建最终内核页表
kvminithart();                       // 切换 TTBR1，作废 TLB
kinit2(P2V(EARLYTOP), P2V(PHYSTOP)); // 阶段二：2 MiB 到 112 MiB，最终页表已覆盖
\end{ccode}

分配器本身就是 xv6 的空闲链表：每个空闲页的前 8 字节存放下一页的地址，\code{kalloc()} 从表头取，\code{kfree()} 往表头放，全程由 \code{kmem.lock} 保护。为了尽早暴露“使用已释放内存”的错误，\code{kfree()} 把页填成 \code{0x01}，\code{kalloc()} 把页填成 \code{0x05}——第~\ref{chap:user} 章会看到，用户态异常地址 \code{0x0505050505050505} 正是一页刚分配、尚未初始化的内存。

\code{PHYSTOP} 设为 \code{0x07000000}（112\,MiB）。上面保留给 QEMU 时代的 ramdisk（\code{0x07000000}）和摄像头 DMA 区（\code{0x07200000}），分配器永远不会把它们发出去。

\section{最终内核页表}

\subsection{\code{kvmmake()}}

\begin{ccode}
pagetable_t kvmmake(void) {
  pagetable_t kpgtbl = (pagetable_t) kalloc();
  memset(kpgtbl, 0, PGSIZE);
  // BCM2837 外设窗口 16 MiB，以及 ARM-local 外设一页：设备内存，不可执行
  kvmmap(kpgtbl, PERIPHERAL_BASE, PERIPHERAL_BASE_PA, 0x01000000, PTE_DEVICE | PTE_XN);
  kvmmap(kpgtbl, LOCAL_BASE, LOCAL_BASE_PA, PGSIZE, PTE_DEVICE | PTE_XN);
  kvmmap(kpgtbl, RAMDISK, RAMDISK_PA, RAMDISK_SIZE, PTE_NORMAL | PTE_XN);
  // 摄像头 DMA：不可缓存的普通内存
  kvmmap(kpgtbl, CAMDMA, CAMDMA_PA, CAMDMA_SIZE, PTE_NORMAL_NC | PTE_XN);
  // 物理第 0 页：spin-table 槽 0xe0..0xf0，唤醒副核时要写
  kvmmap(kpgtbl, KERNBASE, 0, PGSIZE, PTE_NORMAL | PTE_XN);
  // 内核代码：只读、可执行
  kvmmap(kpgtbl, KERNLINK, V2P(KERNLINK), (uint64)etext-KERNLINK, PTE_NORMAL | PTE_RO);
  // 内核数据和全部可用 RAM：可读写、不可执行
  kvmmap(kpgtbl, (uint64)etext, V2P(etext), (uint64)P2V(PHYSTOP)-(uint64)etext,
         PTE_NORMAL | PTE_XN);
  return kpgtbl;
}
\end{ccode}

代码段只读可执行，数据段可写不可执行（W\^{}X），这样即使内核出现缓冲区溢出，也很难把数据当指令执行。

\subsection{\code{walk()} 与 \code{mappages()}}
\label{sec:walk-mappages}

所有映射最终都落到这两个函数上。\code{walk()} 从 L1 走到 L3，遇到缺失的中间页表就（按需）分配一页：

\begin{ccode}
pte_t *walk(pagetable_t pagetable, uint64 va, int alloc) {
  if(va >= MAXVA) panic("walk");
  for(int level = 1; level < 3; level++) {
    pte_t *pte = &pagetable[PX(level, va)];
    if((*pte & PTE_VALID) && (*pte & PTE_TABLE)) {
      pagetable = (pagetable_t)P2V(PTE2PA(*pte));
    } else {
      if(!alloc || (pagetable = (pde_t*)kalloc()) == 0)
        return 0;
      memset(pagetable, 0, PGSIZE);
      *pte = PA2PTE(V2P(pagetable)) | PTE_TABLE | PTE_VALID;
    }
  }
  return &pagetable[PX(3, va)];
}

int mappages(pagetable_t pagetable, uint64 va, uint64 size, uint64 pa, uint64 perm) {
  uint64 a = PGROUNDDOWN(va), last = PGROUNDDOWN(va + size - 1);
  for(;;){
    pte_t *pte = walk(pagetable, a, 1);
    if(pte == 0) return -1;
    if(*pte & PTE_AF) panic("mappages: remap");
    *pte = PA2PTE(pa) | perm | PTE_AF | PTE_V;
    if(a == last) break;
    a += PGSIZE; pa += PGSIZE;
  }
  return 0;
}
\end{ccode}

注意页表项里存的是\textbf{物理地址}（硬件遍历页表时用物理地址），而内核软件读写页表时要先 \code{P2V()} 转成虚拟地址。页表页本身也位于被直接映射的 RAM 中，所以内核总能通过 \code{KERNBASE + PA} 访问任何一张页表——这不是递归页表技巧，只是直接映射的自然结果。

\subsection{页表的规模}

用户命令 \code{vmmap} 会调用 \code{kvmdump()} 打印内核页表，每个 2\,MiB 窗口汇总一行，再逐个打印各进程的用户页表。表~\ref{tab:kpgt} 是一次快照的统计（当时还保留着 ramdisk 映射）。

\begin{table}[htbp]
  \centering
  \caption{最终内核页表的组成（一次运行时快照）}
  \label{tab:kpgt}
  \small
  \begin{tabularx}{\linewidth}{@{}lllX@{}}
    \toprule
    L1 项 & 虚拟地址 & L2 有效项 & 内容 \\
    \midrule
    \code{L1[0]} & \code{KERNBASE}$+$0～1\,GiB & 65 & 内核和 RAM（56 项）、ramdisk（1 项）、BCM2837 MMIO（8 项） \\
    \code{L1[1]} & \code{KERNBASE}$+$1～2\,GiB & 1 & ARM-local 外设 \code{0x40000000} \\
    \code{L1[255]} & 接近 \code{MAXVA} & 1 & 进程内核栈及其间的保护页 \\
    \midrule
    \multicolumn{4}{@{}l}{合计：L1 一张、L2 三张、L3 67 张，共 71 页 $=$ 284\,KiB；有效 L3 叶子 32\,961 个} \\
    \bottomrule
  \end{tabularx}
\end{table}

如果改用 2\,MiB 块映射 RAM，L3 表可以省掉绝大部分。项目选择全部用 4\,KiB 页，是为了让内核栈的保护页、代码段只读等细粒度属性能用同一套 \code{mappages()} 实现。

\section{用户地址空间}
\label{sec:uvm}

\subsection{内核空间与用户空间}

每个进程有自己的 L1 页表，进程切换时写入 \reg{TTBR0\_EL1}：

\begin{ccode}
void switchuvm(struct proc *p) {
  w_ttbr0_el1(V2P(p->vm->pagetable));
  flush_tlb();
}
\end{ccode}

\subsection{\code{exec()} 之后的布局}

\code{exec()} 建立的用户地址空间从虚拟地址 0 开始（图~\ref{fig:uas}）：先是 ELF 的代码和数据，然后是一页不可访问的保护页和 4 页（16\,KiB）用户栈，栈之上是 \code{sbrk()} 扩展的堆。用户页面的属性是 \code{PTE\_NORMAL | PTE\_U}，代码页可执行。

\begin{figure}[htbp]
  \centering
  \begin{tikzpicture}[x=1cm,y=6mm]
    \foreach \y/\h/\t/\c in {0/1.4/{text / data / bss}/structurecolor!15, 1.4/0.8/{保护页（\code{uvmclear} 清除 U 位）}/gray!25,
        2.2/1.6/{用户栈 16\,KiB（\code{USTACKPAGES}=4）}/structurecolor!25, 3.8/1.2/{堆（\code{sbrk}）}/structurecolor!8}{
      \draw[fill=\c] (0,\y) rectangle (6,\y+\h);
      \node[font=\small] at (3,\y+\h/2) {\t};
    }
    \node[lab, left] at (0,0) {VA 0};
    \node[lab, left] at (0,5) {\code{p->vm->sz}};
    \draw[arr] (6.3,3.8) -- (6.3,2.4) node[midway, right, lab]{栈向下增长};
  \end{tikzpicture}
  \caption{\code{exec()} 之后的用户地址空间}
  \label{fig:uas}
\end{figure}

进程的内核栈不在用户页表里，而在内核页表高端的 \code{KSTACK(slot)} 处，每个栈一页，相邻栈之间隔一页不映射的保护页。内核栈溢出会立刻触发翻译错误，而不是悄悄覆盖相邻进程的栈。第~\ref{chap:process} 章会介绍内核栈的动态分配。

\subsection{第一个用户进程的页表}

所有用户地址空间都由同一组函数建立。\code{uvmcreate()} 分配一页并清零，作为空的 L1 页表——全 0 意味着全部表项无效，这一点很重要：任何从 \code{kalloc()} 拿来做页表的页都必须先清零，否则 \code{0x05} 填充的垃圾会被 MMU 当成有效描述符。第一个进程没有 ELF 文件可读，内核把一小段编进内核镜像的 \code{initcode} 直接拷进一页：

\begin{ccode}
void uvminit(pagetable_t pagetable, uchar *src, uint sz) {
  if(sz >= PGSIZE) panic("inituvm: more than a page");
  char *mem = kalloc();
  memset(mem, 0, PGSIZE);
  mappages(pagetable, 0, PGSIZE, V2P(mem), PTE_NORMAL|PTE_U);
  memmove(mem, src, sz);              // 通过内核直接映射写入
  uvmsync_icache(pagetable, sz);      // 写入的是指令，要同步 I-cache
}
\end{ccode}

\code{userinit()} 再把 \code{trapframe} 的 \code{elr} 设为 0、\code{sp} 设为 \code{PGSIZE}：进程第一次“返回”用户态时，从虚拟地址 0 取第一条指令，栈从这一页的顶部向下长。同一个物理页此时有两个虚拟地址：内核通过 \code{KERNBASE+PA} 写它，用户通过 VA 0 执行它。拷贝时用的是前者，所以内核根本不需要先切换到这个进程的页表。

\subsection{\code{fork} 时复制地址空间}

\code{fork()} 要给子进程一份与父进程内容相同、但物理上独立的内存。\code{uvmcopy()} 逐页遍历父进程的页表，为每个已映射的页分配新物理页、拷贝内容、以相同的属性映射进子进程的页表：

\begin{ccode}
for(i = 0; i < sz; i += PGSIZE){
  if((pte = walk(old, i, 0)) == 0) panic("uvmcopy: pte should exist");
  if((*pte & PTE_AF) == 0)         panic("uvmcopy: page not present");
  pa = PTE2PA(*pte);
  flags = PTE_FLAGS(*pte);                     // 连同 U、RO、XN 一起复制
  if((mem = kalloc()) == 0) goto err;
  memmove(mem, (char*)P2V(pa), PGSIZE);        // 经内核直接映射读父进程的页
  if(mappages(new, i, PGSIZE, V2P(mem), flags) != 0){ kfree(mem); goto err; }
}
uvmsync_icache(new, sz);
\end{ccode}

源地址用的是 \code{P2V(pa)} 而不是用户地址 \code{i}。用户地址只有在 \reg{TTBR0} 恰好装着父进程页表时才有意义；走内核直接映射则与当前装的是谁的页表无关，在多核上也更不容易出错。失败时 \code{uvmunmap()} 把已经复制的页全部释放，子进程的页表保持“要么全有、要么全无”。

\subsection{TLB 与地址空间切换}

如果每次访存都要走三级页表，一次 \code{ldr} 就会变成四次内存访问。MMU 因此把最近用过的翻译结果缓存在 TLB（Translation Lookaside Buffer）里，命中时不再遍历页表。代价是：一旦软件改了页表，TLB 里可能还留着旧的翻译，必须由软件显式作废。

进程切换正是最大的一次“改页表”——\reg{TTBR0} 换成另一棵树，之前缓存的低地址翻译全都过期了：

\begin{ccode}
void switchuvm(struct proc *p) {
  if(p->vm == 0 || p->vm->pagetable == 0) panic("switchuvm: no pagetable");
  w_ttbr0_el1(V2P(p->vm->pagetable));
  flush_tlb();                 // dsb ishst; tlbi vmalle1is; dsb ish; isb
}
\end{ccode}

\code{flush\_tlb()} 作废全部 EL1 翻译，连内核的也一起丢掉了。ARMv8 提供了 ASID（地址空间标识）：把一个 8 或 16 位的进程编号写进 \reg{TTBR0} 的高位，TLB 表项会带上这个标签，切换时就不必全部作废。xv6-aarch64 没有使用 ASID，每次切换都整体作废，换来的是实现简单、不必管理 ASID 的回收；对于这个规模的系统，多出来的 TLB 缺失是可以接受的。注意内核页表（\reg{TTBR1}）在整个运行期间只装一次（\code{kvminithart()}），这也是把两个地址空间分开的好处。

\subsection{缺页与按需分配}

如果用户进程访问了一个没有映射的地址，页表遍历在某一级遇到无效项，CPU 产生同步异常进入 \code{usertrap()}。此时 \reg{ESR\_EL1} 的 EC 字段是 \code{0x24}（来自 EL0 的数据中止）或 \code{0x20}（来自 EL0 的指令中止）；ISS 的低 6 位 DFSC 说明原因，\code{0b0001xx} 是第 \code{xx} 级的翻译错误，\code{0b0011xx} 是权限错误；出错的虚拟地址在 \reg{FAR\_EL1} 里。

很多内核利用这个异常实现\emph{按需分配}（demand paging）：\code{sbrk} 只增大进程的 \code{sz}，并不真正分配物理页；进程第一次碰到新页时触发翻译错误，内核检查 \reg{FAR\_EL1} 落在 \code{[0, sz)} 之内，就当场 \code{kalloc()} 一页、\code{mappages()} 映射上去，然后 \code{eret} 重新执行那条出错的指令。进程看到的只是一条稍微慢一点的 \code{ldr}。

xv6-aarch64 目前走的是另一条路：\code{growproc()} 调用 \code{uvmalloc()} 立即分配并清零所有新页，所以合法的访问永远不会缺页。\code{usertrap()} 只处理 EC $=21$（\code{svc}），其余异常一律打印 \code{esr}、\code{elr}、\code{far}，调用 \code{uvmdump()} 打印这个进程的整张页表，然后杀掉进程：

\begin{ccode}
} else {
  printf("usertrap(): unexpected esr=%p ec=%p pid=%d\n", esr, ec, p->pid);
  printf("            elr=%p far=%p\n", r_elr_el1(), r_far_el1());
  uvmdump(p->vm->pagetable, p->pid, p->name, "usertrap");
  p->killed = 1;
}
\end{ccode}

这份输出正是 \ref{sec:caches} 节分析真机故障时用的线索。若要加上按需分配，只需在这个分支里先判断 \code{ec == 0x24 \&\& (esr \& 0x3c) == 0x04 \&\& r\_far\_el1() < p->vm->sz}，然后分配、映射，再像故障三那样广播作废 TLB 后返回；\code{growproc()} 则改为只增大 \code{sz}。

\section{MMIO、总线地址与 DMA}

\subsection{MMIO}

SoC 把外设寄存器和 RAM 放进同一个物理地址空间，CPU 用普通的 \code{ldr}/\code{str} 访问它们，由片内地址译码器决定一次访问是去 DDR 控制器还是去某个外设（表~\ref{tab:physmap}）。

\begin{table}[htbp]
  \centering
  \caption{Raspberry Pi 3 的物理地址布局（本内核使用的部分）}
  \label{tab:physmap}
  \small
  \begin{tabular}{@{}ll@{}}
    \toprule
    物理地址 & 用途 \\
    \midrule
    \code{0x00000000}～\code{0x06ffffff} & 内核管理的 RAM（\code{PHYSTOP} = 112\,MiB） \\
    \code{0x00080000} & 内核镜像装载地址 \\
    \code{0x07200000}～\code{0x072fffff} & 摄像头 DMA 区（\code{CAMDMA}） \\
    \code{0x3f003000} & BCM System Timer \\
    \code{0x3f00b000} & legacy 中断控制器；\code{+0x880} 处为 property mailbox \\
    \code{0x3f101000} & 时钟管理器（摄像头 \code{CM\_CAM1}） \\
    \code{0x3f200000} & GPIO \\
    \code{0x3f202000} & SDHOST（外置 SD 卡） \\
    \code{0x3f205000} & BSC0（I2C） \\
    \code{0x3f215000} & AUX / Mini UART \\
    \code{0x3f300000} & Arasan SDHCI（板载 Wi-Fi） \\
    \code{0x3f801000} & Unicam CSI1 \\
    \code{0x3f980000} & DWC2 USB 主机 \\
    \code{0x40000000} & ARM-local 中断与定时器路由 \\
    \bottomrule
  \end{tabular}
\end{table}

外设寄存器必须映射为 Device 内存。若映射成可缓存的普通内存，写操作可能停留在缓存里迟迟不到达设备，读操作可能拿到旧值，两次写还可能被合并或重排——而对设备来说，“写 FIFO”这个动作本身就有副作用。运行时可以用 \code{cat /proc/iomem} 查看内核登记的全部外设地址区间。

\subsection{三种地址视图}

Broadcom 的文档用 VideoCore 总线地址描述外设，例如 Mini UART 在 \code{0x7e215000}。同一组寄存器在不同观察者眼里地址不同：

\begin{txtcode}
VideoCore 总线地址   0x7e215000     (文档、设备树里的写法)
ARM 物理地址         0x3f215000     (CPU 发出的地址)
内核虚拟地址         0xffffff803f215000  (经 TTBR1 翻译)
\end{txtcode}

DMA 方向的地址转换更容易出错。DWC2 USB、Unicam 这些\emph{总线主设备}自己访问内存时，看到的是 VideoCore 总线地址，必须把 ARM 物理地址加上别名前缀：

\begin{ccode}
#define DWC2_DMA_BUS(a) (0xc0000000U | (uint32)V2P(a))
\end{ccode}

\code{0xC0000000} 别名表示“绕过 VideoCore 的 L2 缓存直接访问 SDRAM”。这不是新的内存，也不是 CPU 虚拟地址，而是同一片 RAM 在 GPU 总线上的另一种编址。

\section{缓存与 TLB：三个真机故障}
\label{sec:caches}

QEMU 不模拟缓存，也几乎不检查内存属性，下面三个问题都只在真机上出现。它们说明：页表不只是“虚拟地址到物理地址”的映射，还决定了 CPU 怎样访问这块内存。

\subsection{故障一：停在 \code{BPH0abde}}
\label{sec:fault1}

真机串口输出 \code{BPH0abde} 后不再有动静。对照表~\ref{tab:boot-marks}，\code{d} 表示进入了第一次 \code{kfree()}，\code{e} 表示 \code{memset} 已完成，缺少的 \code{f} 表示 \code{acquire(\&kmem.lock)} 没有返回。

当时只有 CPU0 在运行，不可能有别的核持锁。真正的原因在于 \code{acquire()} 的实现：

\begin{asmcode}
retry:  ldaxr w1, [x0]      // 独占加载，建立“独占监视”
        cbnz  w1, retry
        stxr  w2, w3, [x0]  // 独占存储，失败时 w2 != 0
        cbnz  w2, retry
\end{asmcode}

\code{ldaxr} 读到锁为 0，但 \code{stxr} 一直返回失败——这是活锁，不是死锁。当时普通 RAM 被映射成 \emph{Normal Non-cacheable}，也没有 Inner Shareable 属性，数据缓存也没有打开。在 Cortex-A53 上，独占监视器依赖缓存一致性机制工作，对这种属性的内存，独占存储可能永远不成功。

修复包括四处：增加 \code{MT\_NORMAL}（Write-Back）类型并让 RAM 使用它；页表项加 \code{PTE\_SH\_INNER}；\reg{TCR\_EL1} 的页表遍历属性设为 Write-Back、Inner Shareable；\reg{SCTLR\_EL1} 打开 C 和 I 位。之后串口变成 \code{BPH0abdefc1234}，系统顺利进入 shell。

\subsection{故障二：\code{ls} 跳到了不属于它的代码}

打开缓存之后，在 shell 里执行 \code{ls} 出现用户异常：

\begin{txtcode}
usertrap(): unexpected esr=0x92000047 ec=0x24 pid=3
            elr=0x13dc far=0x13f48
\end{txtcode}

EC$=$\code{0x24} 是来自 EL0 的数据中止，ISS 表示 L3 翻译错误、写访问。奇怪的是 \code{ls} 的代码只到 \code{0xbd7}，PC 却在 \code{0x13dc}。

原因在于 D-cache 与 I-cache 不一致。\code{exec()} 通过内核高地址把新程序的指令写进物理页，新指令此时还在 D-cache 里；用户态从低地址取指，走的是 I-cache，里面可能还留着同一虚拟地址上一个程序（\code{sh}）的指令。AArch64 不保证数据写入对取指自动可见，必须显式维护：

\begin{ccode}
void uvmsync_icache(pagetable_t pagetable, uint64 sz) {
  for(uint64 va = 0; va < PGROUNDUP(sz); va += PGSIZE){
    uint64 kva = uva2ka(pagetable, va);
    if(kva == 0) continue;
    for(uint64 p = kva; p < kva + PGSIZE; p += 64)
      asm volatile("dc cvau, %0" :: "r"(p) : "memory");   // 数据清到统一点
  }
  asm volatile("dsb ish" ::: "memory");
  asm volatile("ic ialluis\n\tdsb ish\n\tisb" ::: "memory"); // 广播作废所有核的 I-cache
}
\end{ccode}

第一版用 \code{ic ivau} 按内核高地址作废 I-cache，仍然失败：同一物理页有内核高地址和用户低地址两个别名，按其中一个别名维护不能保证覆盖另一个。改用 \code{ic iallu} 作废整个 I-cache 后问题消失。但打开四核以后它又以另一种形式回来了：\code{ic iallu} 只作废\emph{执行它的那个核}的 I-cache，\code{exec()} 在一个核上装好新程序，进程迁移到另一个核后，那里仍能命中旧程序的指令——表现为 \code{ELR} 指向一条根本不访存的指令却报数据中止（第~\ref{chap:drivers} 章摄像头一节）。最终版本用广播到 Inner Shareable 域的 \code{ic ialluis}。代价是每次 \code{exec()}/\code{fork()} 清一次 I-cache，对 xv6 的规模完全可以接受。\code{exec()}、\code{uvminit()}（第一个进程）、\code{uvmcopy()}（\code{fork}）三处都会调用它。

\subsection{故障三：\code{sbrk} 之后的旧 TLB}

进一步的日志显示，上面那次异常还有另一层原因：\code{sh} 通过 \code{sbrk()} 新增了 \code{0x13000} 这一页，但 \code{growproc()} 修改页表后没有作废 TLB，CPU 继续使用之前缓存的“该地址无效”的翻译结果。修复是在页表扩展或收缩后执行广播式 TLB 失效：

\begin{asmcode}
        dsb  ishst          // 确保页表写入完成
        tlbi vmalle1is      // 作废 Inner Shareable 域内所有核的 EL1 TLB
        dsb  ish
        isb
\end{asmcode}

\code{vmalle1is} 中的 \code{is} 很关键：多核以后，同一进程的线程可能正在别的核上运行，只作废本核 TLB 是不够的。收缩地址空间时顺序必须是“清页表项 → 广播作废 TLB → 释放物理页”，否则别的核可能通过陈旧的 TLB 访问一页已经被重新分配的内存。

\begin{note}
  经验：在 AArch64 上，改了“将要执行的内存”要做 D/I-cache 同步，改了“页表”要做 TLB 失效，两者不能互相替代。QEMU 两件事都不检查，所以必须上真机验证。
\end{note}

\chapter{进程调度}
\label{chap:process}

\begin{introduction}[本章内容]
  \item 进程的数据结构
  \item 动态进程表与按需分配的内核栈
  \item 第一个用户进程
  \item \code{fork()} 逐行解析
  \item 上下文切换与调度器
  \item \code{exec}、\code{exit} 与 \code{wait}
  \item 睡眠、唤醒与同步原语
  \item 多核软实时调度
\end{introduction}

\section{进程由什么组成}

在 xv6 里，一个进程就是一个 \code{struct proc}，加上它引用的几样东西：一张用户页表、一个内核栈、一个 trapframe、一组打开的文件。图~\ref{fig:proc} 画出了它们的关系。

\begin{figure}[htbp]
  \centering
  \begin{tikzpicture}[node distance=5mm and 10mm]
    \node[hbox, text width=30mm, minimum height=34mm] (p) {\code{struct proc}\\[1mm]\footnotesize state、pid、tgid\\chan、killed\\policy、cpumask\\context\\ofile[]、cwd};
    \node[box, right=of p, yshift=12mm, text width=34mm] (vm) {\code{struct vmspace}\\\footnotesize TTBR0 页表、sz、引用计数};
    \node[box, right=of p, yshift=-2mm, text width=34mm] (ks) {内核栈（一页）\\\footnotesize 顶部放 trapframe};
    \node[box, right=of p, yshift=-15mm, text width=34mm] (f) {\code{struct file}…\\\footnotesize 管道、inode、设备、socket};
    \node[gbox, right=of vm, text width=22mm] (pg) {用户物理页};
    \draw[arr] (p.east) ++(0,12mm) -- (vm.west);
    \draw[arr] (p.east) ++(0,-2mm) -- (ks.west);
    \draw[arr] (p.east) ++(0,-15mm) -- (f.west);
    \draw[arr] (vm) -- (pg);
  \end{tikzpicture}
  \caption{一个进程的组成}
  \label{fig:proc}
\end{figure}

\subsection{\code{struct proc}}

\begin{ccode}
enum procstate { UNUSED, USED, SLEEPING, RUNNABLE, RUNNING, ZOMBIE };

struct proc {
  struct spinlock lock;
  // 以下字段需持有 p->lock
  enum procstate state;
  void *chan;                  // 非 0 表示睡在这个“通道”上
  int killed, xstate;
  int pid, tgid, is_thread;    // tid、线程组 id
  int sid, pgid, ctty;         // 会话、进程组、控制终端
  int policy, rt_priority, nice;
  uint cpumask;                // CPU 亲和性
  int slice; uint64 rq_seq;    // 时间片、入队序号
  // 以下字段需持有 wait_lock
  struct proc *parent;
  // 以下字段只由进程自己访问
  uint64 kstack, kstack_pa, kstack_slot;
  struct proc *next;           // 动态进程表链表
  struct vmspace *vm;          // 用户地址空间（线程间共享）
  struct trapframe *trapframe; // 用户态现场，位于内核栈顶
  struct context context;      // 内核态现场，swtch() 用
  struct file *ofile[NOFILE];
  struct inode *cwd;
  char name[16];
};
\end{ccode}

字段按“谁在什么锁下访问”分组，这是 xv6 的一贯风格。与原版相比，这里多了会话和进程组（TTY 作业控制要用）、调度属性（软实时调度）、线程组 id，并且用户地址空间被抽成了可共享的 \code{struct vmspace}，为第~\ref{chap:user} 章的用户线程做准备。

\subsection{两种现场}

进程有两份寄存器现场，初学时很容易混淆：

\begin{itemize}
  \item \textbf{trapframe}：用户态被打断时（系统调用、中断、异常）由 \code{usavereg} 保存，包含全部 31 个通用寄存器、\code{elr}、\code{spsr} 和用户 \code{sp}。它描述“回到用户态时从哪里继续”。
  \item \textbf{context}：内核态主动让出 CPU 时由 \code{swtch()} 保存，只包含被调用者保存寄存器 \code{x18}～\code{x30} 和 \code{sp}。它描述“这个进程的内核线程停在了哪里”。
\end{itemize}

AArch64 的调用约定规定 \code{x19}～\code{x28}、\code{x29}（帧指针）和 \code{sp} 由被调用者保存，\code{x30} 是返回地址。\code{swtch()} 是一次普通的函数调用，调用者保存寄存器已经由编译器在调用前存好，所以 context 只需要这 14 个值。

\section{动态进程表与内核栈}

原版 xv6 用静态数组 \code{proc[NPROC]}，并在启动时为 64 个槽位全部分配内核栈。本项目改成了按需增长的链表：

\begin{ccode}
static struct proc* allocproc(void) {
  acquire(&proc_table_lock);
  for(p = proc_head; p; p = p->next) {        // 先找一个 UNUSED 节点复用
    acquire(&p->lock);
    if(p->state == UNUSED) {
      if(p->kstack_pa == 0 && alloc_kstack(p) < 0) { ... return 0; }
      goto found;
    }
    release(&p->lock);
  }
  p = newprocslot();                           // 没有就新建一个节点
  acquire(&p->lock);
found:
  release(&proc_table_lock);
  p->pid = allocpid();  p->tgid = p->pid;
  sched_defaults(p);
  p->state = USED;
  sp = (char*)p->kstack + PGSIZE;
  sp -= sizeof(*p->trapframe);                 // trapframe 放在内核栈顶
  p->trapframe = (struct trapframe*)sp;
  p->vm = vmspace_alloc();                     // 空的用户地址空间
  memset(&p->context, 0, sizeof(p->context));
  p->context.x30 = (uint64)forkret;            // 第一次被调度时“返回”到 forkret
  p->context.sp  = (uint64)sp;
  return p;                                    // 持有 p->lock 返回
}
\end{ccode}

几个设计要点：

\begin{itemize}
  \item 每个节点分到一个永不改变的 \code{kstack\_slot}，内核栈虚拟地址为 \code{KSTACK(slot)}，相邻栈之间隔一页保护页；
  \item 节点空闲时释放内核栈的\emph{物理页}，但节点本身永远留在链表里。调度器、\code{wait()}、\code{wakeup()} 都无锁遍历这个链表，保留节点就不会出现悬空的 \code{next} 指针，省去了 RCU 一类的复杂机制；
  \item 释放内核栈时先清页表项，再广播作废 TLB，最后才 \code{kfree} 物理页（第~\ref{chap:mmu} 章的规则）。
\end{itemize}

\code{forktest} 同时创建并回收 128 个子进程，验证已经越过了旧的 64 个槽位上限。

\section{第一个用户进程}

所有用户进程都是 \code{fork()} 出来的，但第一个进程只能由内核手工构造。\code{userinit()} 分配一个进程，把一段内置的机器码 \code{initcode} 拷到它的第 0 页：

\begin{ccode}
void userinit(void) {
  struct proc *p = allocproc();
  initproc = p;
  uvminit(p->vm->pagetable, initcode, sizeof(initcode));  // 拷贝并同步 I-cache
  p->vm->sz = PGSIZE;
  p->trapframe->elr  = 0;          // 用户态 PC
  p->trapframe->spsr = 0;          // M[3:0]=0：EL0t
  p->trapframe->sp   = PGSIZE;     // 用户栈顶
  safestrcpy(p->name, "initcode", sizeof(p->name));
  p->cwd = namei("/");
  make_runnable(p, 0);
  release(&p->lock);
}
\end{ccode}

\code{initcode} 由 \file{user/initcode.S} 汇编而来，只做一件事：调用 \code{exec("/init")}。

\begin{asmcode}
start:  ldr x0, =init          // 路径
        ldr x1, =argv          // 参数表
        mov x7, #SYS_exec
        svc #0
exit:   mov x7, #SYS_exit      // exec 失败才会到这里
        svc #0
        b exit
init:   .ascii "/init\0"
\end{asmcode}

这个进程第一次被调度时，\code{swtch()} 恢复的 \code{x30} 是 \code{forkret}，于是“返回”到 \code{forkret()}。它释放调度器替它拿着的 \code{p->lock}；如果是系统中第一个进程，还要在进程上下文中初始化文件系统（\code{fsinit} 读 superblock、恢复日志，期间可能睡眠，所以不能放在 \code{main()} 里）。然后调用 \code{usertrapret(trapframe)}，从 trapframe 恢复寄存器，\code{eret} 到 EL0 的地址 0 开始执行用户代码。

\section{\code{fork()} 逐行解析}

\begin{ccode}
int fork(void) {
  struct proc *np, *p = myproc();
  if((np = allocproc()) == 0)                               // (1)
    return -1;

  acquire(&p->vm->lock);                                    // (2)
  if(uvmcopy(p->vm->pagetable, np->vm->pagetable, p->vm->sz) < 0){
    release(&p->vm->lock);
    freeproc(np); release(&np->lock);
    return -1;
  }
  np->vm->sz = p->vm->sz;
  release(&p->vm->lock);

  *(np->trapframe) = *(p->trapframe);                       // (3)
  np->trapframe->x0 = 0;                                    // (4)

  for(int i = 0; i < NOFILE; i++)                           // (5)
    if(p->ofile[i]) np->ofile[i] = filedup(p->ofile[i]);
  np->cwd = idup(p->cwd);
  np->sid = p->sid;  np->pgid = p->pgid;  np->ctty = p->ctty;
  safestrcpy(np->name, p->name, sizeof(p->name));

  acquire(&p->lock);                                        // (6)
  np->policy = p->policy;  np->rt_priority = p->rt_priority;
  np->nice = p->nice;      np->cpumask = p->cpumask;
  release(&p->lock);
  np->slice = timeslice(np);

  int pid = np->pid;
  release(&np->lock);
  acquire(&wait_lock);  np->parent = p;  release(&wait_lock);   // (7)
  acquire(&np->lock);   make_runnable(np, 0);  release(&np->lock); // (8)
  return pid;                                               // (9)
}
\end{ccode}

\begin{enumerate}
  \item 分配进程描述符、内核栈和空的地址空间，context 设为“从 \code{forkret} 开始”；
  \item 逐页复制父进程的用户内存：为每一页新分配物理页，拷贝内容，用相同的权限映射到子进程页表。复制完调用 \code{uvmsync\_icache()}，避免子进程的新物理页命中旧的指令缓存；
  \item 复制整个 trapframe——子进程回到用户态时，PC、栈指针、全部寄存器都和父进程在 \code{svc} 之后一模一样；
  \item 唯一的区别：子进程的 \code{x0}（系统调用返回值）改成 0；
  \item 复制文件描述符表。注意复制的是 \code{struct file} 的\emph{引用}，父子共享同一个文件偏移，这正是 \code{sh} 实现重定向和管道的基础；
  \item 像 Linux 一样，子进程继承调度策略、优先级、nice 和 CPU 亲和性；
  \item 设置父进程指针，受 \code{wait\_lock} 保护，与 \code{wait()}、\code{exit()} 同步；
  \item 把子进程放入运行队列，从这一刻起别的 CPU 就可能开始运行它；
  \item 父进程返回子进程的 pid。
\end{enumerate}

图~\ref{fig:fork} 对比了父子两条返回路径。

\begin{figure}[htbp]
  \centering
  \begin{tikzpicture}[node distance=4mm and 6mm]
    \node[box] (svc) {用户态 \code{svc \#0}\\\code{x7=SYS\_fork}};
    \node[box, below=of svc] (ut) {\code{el0sync} → \code{usertrap()}\\→ \code{sys\_fork()} → \code{fork()}};
    \node[hbox, below left=8mm and -6mm of ut, text width=34mm] (par) {父进程：\code{fork()} 返回 pid\\写入 \code{trapframe->x0}\\\code{urestorereg}; \code{eret}};
    \node[hbox, below right=8mm and -6mm of ut, text width=40mm] (ch) {子进程：被调度\\\code{swtch} 恢复 \code{x30=forkret}\\\code{forkret()} → \code{usertrapret()}\\\code{eret}，\code{x0 = 0}};
    \node[box, below=of par] (pu) {父：\code{fork()} 返回 $>0$};
    \node[box, below=of ch] (cu) {子：\code{fork()} 返回 0};
    \draw[arr] (svc) -- (ut);
    \draw[arr] (ut) -- (par);
    \draw[darr] (ut) -- node[lab, right]{\code{make\_runnable}} (ch);
    \draw[arr] (par) -- (pu);
    \draw[arr] (ch) -- (cu);
  \end{tikzpicture}
  \caption{\code{fork()} 的两条返回路径}
  \label{fig:fork}
\end{figure}

子进程从来没有执行过 \code{fork()} 的前半段，它的内核栈上只有一个 trapframe。调度器第一次运行它时，\code{swtch()} 恢复的返回地址是 \code{forkret}，\code{forkret()} 再调用 \code{usertrapret()}，从复制来的 trapframe “返回”到用户态——在用户程序看来，就好像子进程也刚刚从 \code{fork()} 返回，只是返回值为 0。

\subsection{新进程的第一次运行}

\code{allocproc()} 为新进程伪造了一个“刚刚调用过 \code{swtch()}”的上下文：\code{context.sp} 指向内核栈上 trapframe 的下方，\code{context.x30}（返回地址）指向 \code{forkret}。调度器第一次 \code{swtch()} 到它时，\code{ret} 指令就跳进了一个从未被调用过的函数。这是所有操作系统创建新任务的共同技巧：教学内核里常见的写法是把返回地址设成一段叫 \code{ret\_from\_fork} 的汇编，Linux 也是如此。

这里有一个细节容易忽略：调度器在 \code{swtch()} 之前持有新进程的 \code{p->lock}，而持锁期间本 CPU 的中断是关闭的，所以新进程从 \code{forkret} 开始执行的那一刻不会被抢占。\code{forkret()} 的第一件事就是释放这把锁，相当于“完成交接、重新允许抢占”。若这时被定时器打断，也只是正常的抢占。然后它调用 \code{usertrapret()}，把父进程复制来的 trapframe 装回寄存器，\code{eret} 到用户态：

\begin{ccode}
void forkret(void)
{
  static int first = 1;
  struct proc *p = myproc();
  struct trapframe *tf = p->trapframe;

  release(&p->lock);            // 调度器替我们拿着的锁
  if (first) {                  // 第一个进程：在进程上下文中初始化文件系统
    first = 0;
    fsinit(ROOTDEV);
  }
  usertrapret(tf);              // 从 trapframe 返回用户态
}
\end{ccode}

\section{上下文切换与调度器}

\subsection{\code{swtch()}}

\begin{asmcode}
swtch:                          // void swtch(struct context *old, struct context *new)
        mov x9, x0
        mov x10, sp
        stp x10, x18, [x9, #16 * 0]   // 保存 sp 和 x18..x30
        stp x19, x20, [x9, #16 * 1]
        ...
        stp x29, x30, [x9, #16 * 6]
        mov x9, x1
        ldp x10, x18, [x9, #16 * 0]   // 恢复新的 sp 和 x18..x30
        ...
        ldp x29, x30, [x9, #16 * 6]
        mov sp, x10
        ret                           // 跳到新上下文的 x30
\end{asmcode}

\code{ret} 跳到恢复出来的 \code{x30}：对于一个曾经让出过 CPU 的进程，那是它当年调用 \code{swtch()} 的下一条指令（在 \code{sched()} 里）；对于新进程，就是 \code{forkret}。

\subsection{调度循环}

每个 CPU 都有一个调度器线程，使用启动时的 \code{stack0} 栈：

\begin{ccode}
void scheduler(void) {
  struct cpu *c = mycpu();
  int id = cpuid();
  c->proc = 0;  c->online = 1;
  for(;;){
    intr_on();                                  // 让设备有机会中断
    if((p = pick_next(c, id)) == 0) continue;   // 选优先级最高的可运行进程
    acquire(&p->lock);
    if(p->state == RUNNABLE && (p->cpumask & (1U << id))){
      p->state = RUNNING;  p->last_cpu = id;
      c->proc = p;  c->cur_prio = effective_prio(p, c);
      switchuvm(p);                             // TTBR0 = 进程页表
      swtch(&c->context, &p->context);          // 运行进程，直到它 sched()
      switchkvm();                              // TTBR0 = 0
      c->proc = 0;  c->cur_prio = PRIO_IDLE;
    }
    release(&p->lock);
  }
}
\end{ccode}

进程让出 CPU 的唯一入口是 \code{sched()}：调用者必须持有且只持有 \code{p->lock}，并已把状态改成非 RUNNING。锁在两次 \code{swtch()} 之间“传递”——进程持锁切到调度器，调度器在 \code{swtch()} 返回后释放；反方向同理，新进程在 \code{forkret()} 或 \code{sched()} 返回后释放。这保证了“设置状态”和“切换上下文”之间不会有别的 CPU 插进来运行同一个进程。

\subsection{调度对进程是透明的}

调度的目标是让多个进程分享 CPU 时间，而每个进程都以为自己独占 CPU。能做到这一点，靠的正是上面两份现场：被中断的那一刻，用户寄存器全部存在 trapframe 里；让出 CPU 的那一刻，内核寄存器存在 context 里。从进程自己的角度看，它只是调用了一次 \code{swtch()}，过了一段（也许很长的）时间后函数返回了，期间别的进程运行过这件事它察觉不到。也正因为切换只发生在一次函数调用之内，按调用约定可以被调用者破坏的 \code{x0}～\code{x17} 不需要保存。

\subsection{谁触发调度}

进程让出 CPU 有两种情形：

\begin{itemize}
  \item \textbf{主动让出}：进程在内核中没事可做了，例如 \code{sleep()} 等待磁盘、管道或网络，\code{wait()} 等待子进程，\code{exit()} 结束自己。它们把状态改成 SLEEPING 或 ZOMBIE，然后调用 \code{sched()}；
  \item \textbf{被动抢占}：每次 10\,ms 定时器中断，\code{sched\_tick()} 给当前进程的时间片减一，减到 0 就设置本 CPU 的 \code{need\_resched}；另一个 CPU 唤醒了更高优先级的进程时，会通过核间中断（IPI）设置这个标志。中断返回前，\code{userirq()}/\code{kernelirq()} 检查 \code{resched\_pending()}，若为真就调用 \code{yield()}，把自己放回运行队列再 \code{sched()}。
\end{itemize}

\begin{ccode}
void sched_tick(void)                   // 每个 CPU 每 10 ms 调用一次
{
  struct cpu *c = mycpu();
  struct proc *p = c->proc;
  if(p){
    acquire(&p->lock);
    p->run_ticks++;
    if(p->policy != SCHED_FIFO && --p->slice <= 0)
      c->need_resched = 1;              // 时间片用完：下次中断返回时让出
    release(&p->lock);
  }
  ...                                   // 实时任务节流计数
}
\end{ccode}

\subsection{关中断与禁止抢占}

教学内核里常见一个计数器 \code{preempt\_count}：它不为 0 时，定时器中断不会触发重新调度。xv6 没有单独的计数器，但效果相同：调度相关的代码总是持有 \code{p->lock}，而持有任何自旋锁时本 CPU 的中断都是关着的（\code{push\_off}），也就不可能被定时器抢占。\code{sched()} 开头的检查把这个约定写成了断言：必须持有 \code{p->lock}、只持这一把锁、状态已不是 RUNNING、中断是关着的。

“关中断”与“禁止抢占”并不总是一回事。调度器循环在每一轮开头都要 \code{intr\_on()}：如果此时没有任何可运行的进程，CPU 就在循环里空转，而让某个进程重新变成 RUNNABLE 的，往往正是一个设备中断（磁盘读完、网络包到达）。如果调度器关着中断空转，就永远等不到这个中断。早期 Linux 的调度函数也有同样的要求：它在运行期间禁止嵌套调度，但必须允许中断发生。

\subsection{如何选择下一个进程}

最早的 Linux（0.01 版）用一个很巧妙的计数器算法：每个任务有优先级 \code{priority} 和剩余时间 \code{counter}，每个定时器节拍当前任务的 \code{counter} 减一；调度时选 \code{counter} 最大的可运行任务；若所有可运行任务的 \code{counter} 都为 0，就给\emph{所有}任务（包括在睡眠的）执行 $\text{counter} \leftarrow \text{counter}/2 + \text{priority}$。睡得越久的任务累积得越多，但不会超过 $2\times\text{priority}$（这是迭代的不动点），所以交互式任务醒来后能优先得到 CPU。

xv6-aarch64 的 \code{pick\_next()} 把“剩余时间”和“谁先运行”分开：时间片只决定什么时候让出，选择则按有效优先级，同优先级按入队先后（\code{rq\_seq}）：

\begin{ccode}
static struct proc* pick_next(struct cpu *c, int id)
{
  struct proc *best = 0;
  int best_prio = 0;
  uint64 best_seq = 0;

  for(struct proc *p = proc_head; p; p = p->next){
    acquire(&p->lock);
    if(p->state == RUNNABLE && (p->cpumask & (1U << id))){
      int prio = effective_prio(p, c);
      if(best == 0 || prio > best_prio ||
         (prio == best_prio && p->rq_seq < best_seq)){
        best = p;
        best_prio = prio;
        best_seq = p->rq_seq;
      }
    }
    release(&p->lock);
  }
  return best;
}
\end{ccode}

时间片用完的进程在 \code{yield()} 中重新入队、拿到更大的 \code{rq\_seq}，于是排到同优先级的末尾——这就是轮转（round robin）。完整的优先级规则见第~\ref{sec:sched-rt} 节。

\subsection{跟踪一次时间片轮转}
\label{sec:sched-trace}

调度代码里有大量异步交互，单看函数很难想清楚整个系统的状态如何随时间变化。下面跟踪进程 A 在用户态运行时被定时器抢占、换成进程 B、之后又换回 A 的全过程，注意每一步发生在哪个栈上：

\begin{enumerate}
  \item A 在 EL0 运行，使用自己的用户栈（\reg{SP\_EL0}）；
  \item 定时器中断到来。硬件切到 \reg{SP\_EL1}，此时它指向 A 内核栈的顶端；\code{el0irq} 的 \code{usavereg} 把 A 的全部用户寄存器存进 A 的 trapframe；
  \item \code{userirq()} $\to$ \code{devintr()} $\to$ \code{sched\_tick()}：A 的时间片用完，设置 \code{need\_resched}；
  \item \code{userirq()} 发现 \code{resched\_pending()}，调用 \code{yield()}：拿 \code{A->lock}，状态改为 RUNNABLE，\code{sched()} $\to$ \code{swtch(\&A->context, \&cpu->context)}。A 的内核寄存器存进 \code{A->context}，\code{sp} 换成本 CPU 的调度器栈；
  \item 调度器从上次的 \code{swtch()} 之后继续：\code{switchkvm()}，释放 \code{A->lock}，下一轮 \code{pick\_next()} 选中 B；
  \item 调度器拿 \code{B->lock}，状态改为 RUNNING，\code{switchuvm(B)} 把 B 的页表装进 \reg{TTBR0\_EL1}，\code{swtch(\&cpu->context, \&B->context)}；
  \item B 从它自己上次的 \code{swtch()} 之后继续——如果 B 当初也是被定时器抢占的，那就是它的 \code{sched()} 返回到 \code{yield()}，释放 \code{B->lock}，再一路返回到 \code{userirq()}、\code{el0irq}，\code{urestorereg} 从 B 的 trapframe 恢复寄存器，\code{eret} 回到 B 的用户态；如果 B 是新进程，则 \code{ret} 进入 \code{forkret}；
  \item 若干时间片之后，B 以同样的方式让出，调度器再次选中 A，\code{swtch()} 恢复 \code{A->context}：A 回到第 4 步的 \code{sched()} 里面，沿 \code{yield()} $\to$ \code{userirq()} $\to$ \code{el0irq} 返回，\code{eret} 之后 A 从被打断的那条指令继续执行。
\end{enumerate}

这一过程中，A 在内核里的调用链 \code{el0irq} $\to$ \code{userirq} $\to$ \code{yield} $\to$ \code{sched} $\to$ \code{swtch} 一直“冻结”在 A 自己的内核栈上：

\begin{txtcode}
A 的内核栈（高地址在下）            调度器栈（stack0）
+-----------------------------+    +---------------------------+
| swtch/sched/yield 的栈帧    |    | scheduler() 的栈帧        |
| userirq 的栈帧              |    | （每个 CPU 一个）         |
+-----------------------------+    +---------------------------+
| trapframe：A 的用户寄存器   |
+-----------------------------+ <- p->kstack + PGSIZE
\end{txtcode}

这就是“以 A 的身份进入中断、以 B 的身份离开中断”之所以合法的原因：每个进程的中断现场（trapframe 中的 \code{elr}、\code{spsr}）和内核调用链都保存在它自己的内核栈上，\code{eret} 使用的永远是当前进程自己的那一份。与把任务结构放在栈页底部的教学内核相比，xv6 的每个进程在调度器之外还多了一个独立的调度器上下文，切换总是“进程 $\to$ 调度器 $\to$ 进程”两步，而不是进程之间直接切换；代价是多一次 \code{swtch()}，好处是调度器有自己干净的栈，进程退出时也不必在自己的栈上释放自己。

\section{\code{exec}、\code{exit} 与 \code{wait}}

\code{exec(path, argv)} 读取 ELF 文件头，为每个 \code{PT\_LOAD} 段分配并装入内存，然后在其上方分配一页保护页和 4 页用户栈，把参数字符串和 \code{argv[]} 指针数组压栈。全部就绪后才替换旧的地址空间——任何一步失败，原进程都不受影响。新页表生效前调用 \code{uvmsync\_icache()}。

\code{exit()} 关闭所有文件，把子进程过继给 \code{init}，把自己设为 ZOMBIE 并唤醒父进程，然后 \code{sched()}，再也不会回来。进程的内存和内核栈此时还不能释放（它正在这个内核栈上运行），要等父进程在 \code{wait()} 中找到 ZOMBIE 子进程后调用 \code{freeproc()} 回收。\code{init} 进程在一个循环里不停地 \code{wait()}，负责回收所有孤儿。

\section{睡眠、唤醒与同步原语}

\subsection{\code{sleep}/\code{wakeup}}

xv6 的阻塞机制只有两个函数。\code{sleep(chan, lk)} 原子地释放条件锁 \code{lk} 并在 \code{chan} 上睡眠；\code{wakeup(chan)} 把所有睡在 \code{chan} 上的进程置为 RUNNABLE。“通道”只是一个地址，通常就是被等待对象的地址。

关键是“原子”二字：\code{sleep()} 先拿到 \code{p->lock} 才释放 \code{lk}，而 \code{wakeup()} 必须拿到 \code{p->lock} 才能修改进程状态，所以唤醒不可能落在“检查完条件、还没睡着”的间隙里，不会丢失唤醒。被唤醒的进程总要在 \code{while} 循环中重新检查条件，因为 \code{wakeup} 会唤醒同一通道上的所有进程。

\subsection{用户态同步原语}

在 \code{sleep}/\code{wakeup} 之上，\file{kernel/sync.c} 提供了 64 个内核同步对象槽，用户通过带代数（generation）的整数句柄访问：

\begin{itemize}
  \item \code{mutex}：记录持有者 pid，只有持有者能解锁；
  \item \code{semaphore}：计数信号量，\code{sem\_wait} 在计数为 0 时阻塞；
  \item \code{event}：手动复位事件，\code{notify} 唤醒全部等待者并保持触发状态；
  \item \code{katomic}：跨进程的原子整数，支持 load/store/fetch\_add/compare\_exchange。
\end{itemize}

对象在 \code{fork()} 之前创建，父子进程就持有相同的句柄，可以跨越独立的地址空间同步。句柄中的代数防止对象销毁、重新分配之后旧句柄误用新对象。\code{prodcons} 示例用两个信号量（空槽、满槽）加一个互斥锁实现了容量为 4 的有界缓冲区。

\section{多核软实时调度}
\label{sec:sched-rt}

原版 xv6 的调度器按数组下标轮询，找到第一个 RUNNABLE 的进程就运行。为了让高优先级任务被唤醒后能尽快抢占，项目加入了与 Linux 语义一致的三种策略（表~\ref{tab:policy}）。

\begin{table}[htbp]
  \centering
  \caption{调度策略}
  \label{tab:policy}
  \small
  \begin{tabularx}{\linewidth}{@{}llX@{}}
    \toprule
    策略 & 优先级 & 行为 \\
    \midrule
    \code{SCHED\_OTHER} & nice $-20$～$19$ & 分时；时间片 $=(20-\text{nice})/2$ 个节拍（10\,ms～200\,ms） \\
    \code{SCHED\_FIFO} & 1～99 & 一直运行到阻塞、主动让出或被更高优先级抢占，没有时间片 \\
    \code{SCHED\_RR} & 1～99 & 同 FIFO，但同优先级之间每 100\,ms 轮转 \\
    \bottomrule
  \end{tabularx}
\end{table}

\paragraph{选择.} \code{pick\_next()} 在允许当前 CPU 运行的 RUNNABLE 进程中选\emph{有效优先级}最高的；同优先级选 \code{rq\_seq} 最小的（最早入队）。实时任务的有效优先级是 $100+\text{rt\_priority}$，普通任务是 0，所以任何实时任务都排在普通任务前面。被更高优先级抢占、时间片还有剩余的实时任务保留原来的 \code{rq\_seq}，回到队首。

\paragraph{抢占.} 进程变成 RUNNABLE 时，\code{resched\_hint()} 检查它允许运行的各个 CPU：如果有空闲 CPU，什么都不用做；否则找到当前优先级最低的 CPU，若新进程优先级更高，就设置该 CPU 的 \code{need\_resched}，并通过 ARM-local mailbox 0 向它发送核间中断（IPI）。目标 CPU 在中断返回、系统调用返回或内核态中断返回三处检查 \code{need\_resched}，调用 \code{yield()}。

\paragraph{节流.} 为防止失控的 FIFO 任务卡死整个系统，每个 CPU 每秒最多给实时任务 950\,ms；用完后实时任务的有效优先级降到 $-1$，让 shell 和内核工作线程在剩下的 50\,ms 里运行——与 Linux 默认的 \code{sched\_rt\_runtime\_us} 相同。

用户可以用与 Linux 同名的 \code{chrt}、\code{taskset}、\code{nice}、\code{renice} 命令设置这些属性，用 \code{rtlat} 测量唤醒延迟（第~\ref{chap:user} 章）。

\chapter{文件系统}
\label{chap:fs}

\begin{introduction}[本章内容]
  \item xv6 文件系统的分层
  \item 磁盘布局与 \code{mkfs}
  \item 块缓存、日志与 inode
  \item 块设备：从 ramdisk 到 SD 卡
  \item 根文件系统放在哪里
  \item VFS 中间层：procfs、ext2、FAT32
  \item FAT32 与启动分区：bootfs、BPB、FAT 表、目录项与长文件名
  \item xv6 根文件系统：superblock、inode、目录、路径查找与在线更新
  \item 烧录 SD 卡
\end{introduction}

\section{分层结构}

xv6 的文件系统是一叠职责单一的层（图~\ref{fig:fslayers}）。本项目保留了中间各层不动，在上面加了一层 VFS，在下面把块设备从 VirtIO 磁盘换成了真正的 SD 卡。

\begin{figure}[htbp]
  \centering
  \begin{tikzpicture}[node distance=2mm]
    \tikzset{layer/.style={box, minimum width=78mm, text width=74mm, align=left}}
    \node[layer] (sys) {\textbf{系统调用}\quad \code{open/read/write/mkdir/mount}};
    \node[layer, below=of sys] (file) {\textbf{文件描述符}\quad \code{struct file}：inode、设备、管道、vnode、socket、PTY};
    \node[layer, below=of file, fill=structurecolor!22] (vfs) {\textbf{VFS}\quad 挂载表、\code{vnode\_ops}：procfs、ext2、FAT32};
    \node[layer, below=of vfs] (path) {\textbf{路径名}\quad \code{namei}、目录项};
    \node[layer, below=of path] (inode) {\textbf{inode}\quad \code{ilock/readi/writei}、inode 缓存};
    \node[layer, below=of inode] (log) {\textbf{日志}\quad \code{begin\_op/log\_write/end\_op}};
    \node[layer, below=of log] (bio) {\textbf{块缓存}\quad \code{bread/bwrite}，\code{NBUF}=30};
    \node[layer, below=of bio, fill=structurecolor!22] (map) {\textbf{根块设备}\quad \file{rootdev.c}，\code{rootdev\_rw()}：原始分区 / 裸镜像};
    \node[layer, below=of map, fill=structurecolor!22] (sd) {\textbf{SD 卡驱动}\quad \code{sdsector()}：SDHOST，CMD17/CMD24};
  \end{tikzpicture}
  \caption{文件系统分层（深色为本项目新增或替换的层）}
  \label{fig:fslayers}
\end{figure}

\section{磁盘布局}
\label{sec:layout}

xv6 原生文件系统的块大小是 1\,KiB，磁盘依次为：

\begin{txtcode}
[ 启动块 | superblock | 日志 | inode 块 | 空闲位图 | 数据块 ]
   0          1          2..
\end{txtcode}

superblock 记录各区的起始块号和大小，魔数为 \code{0x10203040}。几个参数与原版不同：

\begin{itemize}
  \item \code{FSSIZE}$=32768$：文件系统 32\,MiB。原来只有 1000 块，加入网络、SSH 等程序后，打包就用掉了 986 块，首次启动创建 \code{/etc} 配置时 panic “balloc: out of blocks”；
  \item \code{DIRSIZ}$=30$：文件名最长 30 字节，目录项从 16 字节变为 32 字节，为了让 TFTP 下载的 \code{kernel8-xv6\_wifi.img} 这样的长文件名能通过普通 \code{ls} 显示；
  \item inode 有 12 个直接块和 1 个一级间接块，单个文件最大 $12+256=268$ 块，即 268\,KiB。
\end{itemize}

\begin{note}
  分区再大也不会改变可分配的块数——它由 superblock 里的 \code{size} 字段决定。修改 \code{FSSIZE} 或 \code{DIRSIZ} 之后必须重新生成并整体烧录 \code{fs.img}，只替换内核会让新内核去解释旧格式的磁盘。
\end{note}

\section{块缓存、日志与 inode}

这三层基本沿用 xv6 原样，这里只说明本项目中与它们相关的两点。

\paragraph{日志.} 一次文件系统操作可能修改多个块（例如创建文件要改 inode、目录、位图）。\code{begin\_op()}/\code{end\_op()} 之间的所有修改先写进日志区，提交时一次性写回原位置；启动时 \code{fsinit()} 检查日志，把已提交但未写完的事务重做一遍。这样即使在写的过程中断电，文件系统也只会处于“操作前”或“操作后”两种状态之一。对 SD 卡来说这很重要：真机经常是直接拔电源关机的。

\paragraph{inode 睡眠锁.} inode 操作可能触发 SD 卡读写，耗时很长，不能持有自旋锁，所以每个内存 inode 有一把睡眠锁。项目中出现过一次典型的锁配对错误：在为文件描述符层加入 \code{FD\_DEVICE}、\code{FD\_VNODE}、\code{FD\_SOCKET} 等分支时，\code{fileread()} 的 \code{FD\_INODE} 分支丢了 \code{iunlock()}。

\begin{ccode}
} else if(f->type == FD_INODE){
  ilock(f->ip);
  if((r = readi(f->ip, 1, addr, f->off, n)) > 0)
    f->off += r;
  iunlock(f->ip);       // 这一行被误删
}
\end{ccode}

现象非常迷惑：启动卡在 “exec-login: inode found” 之后，调试输出显示 \code{/bin/login} 的 inode 被 PID 1 锁着，而 \code{valid=0}。原因是 \code{init} 读 \code{/etc/fstab} 时没有解锁，关闭文件后引用计数归零，这个 inode 缓存槽被 \code{iget()} 复用给了 \code{/bin/login}——槽换了主人，锁却还“记得”上一任的持有者。

\begin{note}
  \code{iput()} 不能代替 \code{iunlock()}：前者管“有多少引用”，后者管“是否正在独占访问内容”。一个有用的调试断言是：\code{iget()} 复用 \code{ref==0} 的槽时，若发现锁仍被持有，立即 panic 并打印旧的 dev/inum/owner。
\end{note}

\section{块设备：从 ramdisk 到 SD 卡}

\subsection{三个阶段}

块设备驱动经历了三次替换：

\begin{enumerate}
  \item \textbf{ramdisk}：刚移植到 \code{raspi3b} 时，用 QEMU 的 loader 把 \code{fs.img} 装到物理地址 \code{0x07000000}，\code{bread/bwrite} 直接 \code{memmove}。简单，但改动不会写回宿主机；
  \item \textbf{Arasan SDHCI}（\file{kernel/sd.c}）：BCM2837 的 EMMC 控制器，GPIO48～53 设为 ALT3。QEMU 用 \code{-drive if=sd} 把 \code{fs.img} 接到 SD 总线，数据终于可以持久保存；
  \item \textbf{SDHOST}（\file{kernel/sdhost.c}）：后来要驱动板载 Wi-Fi，而 Wi-Fi 芯片只能接在 Arasan 控制器上（GPIO34～39）。于是外置 SD 卡改由 Broadcom 自己的 SDHOST 控制器（\code{0x3f202000}）驱动，GPIO48～53 改为 ALT0。
\end{enumerate}

两种 SD 驱动对上都提供同一个接口：

\begin{ccode}
int sdsector(uint32 sector, void *buffer, int write);   // 读写一个 512 字节扇区
\end{ccode}

xv6 的块是 1\,KiB，因此一个块对应两个扇区：块 $N$ 是扇区 $2N$ 和 $2N+1$。

\subsection{SD 卡初始化}

无论哪个控制器，SD 卡协议层的初始化顺序都一样（SD 物理层规范\cite{sdspec}）：

\begin{txtcode}
400 kHz 识别时钟
CMD0                复位，卡进入 idle
CMD8                检查电压范围，确认 SD 2.0
CMD55 + ACMD41      反复发送直到卡完成上电；OCR.CCS=1 表示 SDHC（按扇区寻址）
CMD2                读 CID
CMD3                取得相对地址 RCA
CMD7                选中卡，进入 transfer 状态
(SDSC) CMD16        设置块长 512
切换到 25 MHz 工作时钟
\end{txtcode}

读一个扇区发 CMD17，写一个扇区发 CMD24，数据通过数据寄存器以轮询 PIO 方式搬运 128 个 32 位字。驱动用一把自旋锁串行化对控制器的访问。

\subsection{命令超时与恢复}

真机在 shell 里执行 \code{ls} 时偶尔 panic “fat32: read FS.IMG”。在驱动里加入事务阶段编号后得到：

\begin{txtcode}
sd: I/O failure lba=295956 write=0 phase=1 status=17f0207 irq=18000 control1=e0807
\end{txtcode}

\code{phase=1} 表示失败发生在 CMD17 的命令阶段；中断状态 \code{0x18000} 是 ERROR 加 COMMAND TIMEOUT。更麻烦的是，超时后控制器的命令/数据状态机停在 inhibit 状态，之后的每条命令都会失败。仅仅清中断位不够，还要用 \reg{CONTROL1} 的软件复位位分别复位命令线和数据线，等硬件自动清零这两位。修复后的 \code{sdsector()} 对同一扇区最多重试三次，每次失败后先 \code{recover\_io()}。这只复位主机控制器内部的状态机，卡本身仍处于 transfer 状态，不需要重新初始化。

\section{根文件系统放在哪里}
\label{sec:rootfs}

启动时 \code{rootdev\_init()} 读 SD 卡第 0 扇区，判断根文件系统用哪种布局（表~\ref{tab:rootfs}），此后由 \code{rootdev\_rw()} 把 xv6 块号翻译成 SD 卡的绝对扇区号。这两个函数都在 \file{kernel/rootdev.c} 里。

\paragraph{根块设备不是 FAT32 的一部分.} \file{rootdev.c} 不解释任何文件系统格式。它只读 MBR，用 xv6 superblock 的魔数认出根分区，再把 FAT32 启动分区交给 \file{fat32.c} 挂起来，供 \code{/boot} 和固件加载使用。根分区的块里装的始终是 xv6 自己的 superblock、inode 和目录（\ref{sec:xv6fs} 节），与 FAT32 毫无关系。块缓存的读写入口 \code{rootdev\_rw()} 也只做地址换算，不调用任何 FAT32 函数。

\begin{table}[htbp]
  \centering
  \caption{根文件系统的两种存放方式}
  \label{tab:rootfs}
  \small
  \begin{tabularx}{\linewidth}{@{}lXl@{}}
    \toprule
    方式 & 识别条件与地址换算 & 适用场景 \\
    \midrule
    原始分区 & MBR 中某个非 FAT32 主分区的第 1 块是 xv6 superblock（魔数与大小匹配）；扇区 $=$ 分区起点 $+$ 块号$\times 2$ & 真机 \\
    裸镜像 & 没有 MBR；扇区 $=$ 块号$\times 2$ & QEMU \\
    \bottomrule
  \end{tabularx}
\end{table}

\paragraph{原始分区模式}是真机上的做法：SD 卡分成 4 个区，xv6 文件系统独占一个分区，与 FAT32 启动分区互不干扰。识别时只认 superblock 魔数，而不看 MBR 类型（macOS 只能建 \code{0x83} 类型的分区，与 ext2 相同）。启动日志会显示：

\begin{txtcode}
rootdev: raw xv6 root p2 lba=208896 sectors=819200 size=400 MiB
rootdev: bootfs p1 lba=2048 mounted as FAT32
\end{txtcode}

\section{VFS 中间层}
\label{sec:vfs}

\subsection{动机}

最初的 ext2 驱动只能通过两个专用系统调用 \code{ext2read}、\code{ext2readdir} 访问，普通的 \code{ls}、\code{cat} 完全不知道它的存在。VFS 层把非原生文件系统包装成 \emph{vnode}，让它们通过统一的文件描述符接口工作：

\begin{ccode}
struct vnode_ops {                       // kernel/vfs.c，参数是挂载点内的相对路径
  int (*stat)(char*, struct stat*);
  int (*read)(char*, uint64 off, void*, int n);
  int (*write)(char*, struct fat32_file*, uint64 off, int user_src, uint64 src, int n);
  int (*readdir)(char*, int index, void*);
  int (*create)(char*);
  int (*truncate)(char*, struct fat32_file*);
  int (*rename)(char*, char*);
};
\end{ccode}

这些操作以路径为参数而不是以 inode 为参数，后端每次按路径重新查找。这是一个极简的 VFS：没有 dentry 缓存，也没有把原生 xv6 inode 纳入 vnode，但足以让 \code{ls}、\code{cat}、\code{mv} 等普通命令透明地访问其他文件系统。

\code{struct file} 增加了 \code{FD\_VNODE} 类型，\code{fileread()}、\code{filestat()}、\code{fileclose()} 遇到它就转发给对应的 \code{vnode\_ops}。

\subsection{挂载表与路径路由}

\code{sys\_open()} 先调用 \code{vfsopen()}，它在挂载表里做\emph{最长前缀匹配}：

\begin{txtcode}
open("/boot/config.txt")
  -> findmount(): 绝对路径 /boot/config.txt，匹配挂载点 /boot，相对路径 /config.txt
  -> fat32_ops.stat("/config.txt")
\end{txtcode}

返回值约定是这一层设计的关键：\code{1} 表示 VFS 已打开；\code{-1} 表示路径属于某个挂载点但打开失败（不能退回原生根文件系统，否则会打开挂载点下面被遮住的同名文件）；\code{0} 表示不属于任何挂载点，继续走原生的 \code{namei()}。

\subsection{三个后端}

\begin{itemize}
  \item \textbf{procfs}（\code{/proc}）：没有磁盘数据，打开目录时遍历进程表生成每个 PID 的目录，读 \code{status} 时持锁生成快照。另有 \code{/proc/meminfo}、\code{/proc/iomem}、\code{/proc/devices} 等诊断文件；
  \item \textbf{ext2}（\code{/mnt/ext2}）：支持 1/2/4\,KiB 块，遇到 ext4 extent 等不兼容特性就拒绝挂载，可以和 Linux 交换文件。最初只能读、只支持一级间接块；后来补成了可读写、带外置事务日志的后端，见第~\ref{chap:optional} 章；
  \item \textbf{FAT32}（\code{/boot}）：挂载启动分区，支持根目录的读、创建、截断、顺序写和重命名，支持 ASCII 长文件名。这让系统可以在运行时更新自己的内核：\code{tftp} 下载 \code{kernel8-xv6\_wifi.img} 到 \code{/boot}，重启即生效。
\end{itemize}

\subsection{由用户态决定挂载}

挂载策略仿照嵌入式 Linux：内核只负责探测驱动，\code{init} 读取 \code{/etc/fstab} 决定挂载什么。首次启动时 \code{init} 创建：

\begin{txtcode}
proc    /proc      procfs ro 0 0
bootfs  /boot      fat32  rw 0 0
ext2    /mnt/ext2  ext2   rw 0 0
\end{txtcode}

然后逐行调用 \code{mount(source, target, fstype, flags)}。启动日志出现 “vfs: mounted fat32 at /boot read-write” 才说明 FAT32 真的挂上了；否则 \code{/boot} 只是 \code{init} 创建的空目录。

\section{FAT32 与启动分区}
\label{sec:fat32}

前面几节讲的是 xv6 自己的文件系统。SD 卡上还有一个 xv6 无法回避的 FAT32 分区：树莓派的固件只认 FAT32，内核镜像、\file{config.txt}、Wi-Fi 固件都必须放在这里。本节先看这个分区里放了什么、谁来读，再看 FAT32 在磁盘上的组织，最后逐个讲 \file{kernel/fat32.c} 里的函数。

\subsection{bootfs 里有什么，谁来读}

图~\ref{fig:bootfs} 按读者把启动分区的文件分成三组。

\begin{figure}[htbp]
  \centering
  \begin{tikzpicture}[x=1mm,y=1mm, font=\footnotesize,
      rd/.style={-{Stealth[length=2mm]}, thick, draw=structurecolor!70!black},
      wr/.style={{Stealth[length=2mm]}-{Stealth[length=2mm]}, thick, draw=orange!80!black},
      f/.style={font=\scriptsize\ttfamily, anchor=west, inner sep=1pt}]
    \draw[rounded corners=3pt, draw=gray!60, fill=gray!4] (48,0) rectangle (100,74);
    \node[lab, anchor=north] at (74,73.5) {SD 卡 p1：BOOTFS（FAT32 根目录）};
    % group A: firmware boot
    \draw[fill=gray!10, draw=gray!50] (51,52) rectangle (97,69);
    \node[f] at (52,66.5) {bootcode.bin};
    \node[f] at (52,63) {start.elf  fixup.dat};
    \node[f] at (52,59.5) {config.txt};
    \node[f] at (52,56) {kernel8-xv6\_wifi.img};
    \node[lab, anchor=east] at (96.5,53.8) {启动阶段};
    % group B: drivers
    \draw[fill=structurecolor!8, draw=structurecolor!50] (51,24) rectangle (97,49);
    \node[f] at (52,46) {BCM43430.BIN/.TXT/.CLM};
    \node[f] at (52,42.5) {BCM43455.BIN/.TXT/.CLM};
    \node[f] at (52,39) {MT7601U.BIN};
    \node[lab, anchor=east] at (96.5,26) {驱动固件};
    % group C: runtime
    \draw[fill=orange!8, draw=orange!50] (51,3) rectangle (97,21);
    \node[f] at (52,18) {camera.bmp  camera.raw};
    \node[f] at (52,14.5) {kernel8-xv6\_wifi.img（新）};
    \node[f] at (52,11) {（早期）FS.IMG};
    \node[lab, anchor=east] at (96.5,5) {运行时读写};
    % readers
    \node[gbox, text width=36mm, minimum height=16mm] (gpu) at (20,61) {VideoCore GPU\\ROM → \code{bootcode.bin}\\→ \code{start.elf}\\（ARM 尚未运行）};
    \node[hbox, text width=36mm, minimum height=14mm] (drv) at (20,37) {xv6 内核驱动\\\code{fat32openroot()}\\\code{fat32readnext()}};
    \node[box, text width=36mm, minimum height=14mm] (usr) at (20,12) {用户程序经 \code{/boot}\\\code{ls cat mv tftpclient}\\\code{camshot}};
    \draw[rd] (51,61) -- node[lab, above]{读} (gpu.east);
    \draw[rd] (51,37) -- node[lab, above]{读} (drv.east);
    \draw[wr] (51,12) -- node[lab, above]{读写} (usr.east);
    \node[lab, align=left, anchor=west] at (102,61) {把 \code{kernel8-xv6\_wifi.img}\\装到物理 \code{0x80000}，\\然后放开 ARM};
    \node[lab, align=left, anchor=west] at (102,37) {\file{brcmfmac.c}、\file{mt7601u.c}\\每次上电下载给网卡};
    \node[lab, align=left, anchor=west] at (102,12) {VFS 挂载 \code{/boot}，\\按 \code{/etc/fstab} 读写};
  \end{tikzpicture}
  \caption{启动分区里的文件和它们的三类读者。同一个 FAT32 分区，先被 GPU 固件读（ARM 还没开始执行），再被 xv6 内核驱动读，最后作为 \code{/boot} 被用户程序读写}
  \label{fig:bootfs}
\end{figure}

\paragraph{第一组由 GPU 读，在 ARM 执行第一条指令之前.} 树莓派上电后先运行的是 VideoCore GPU：片内 ROM 读 \file{bootcode.bin}，它再装入 \file{start.elf} 和 \file{fixup.dat}。\file{start.elf} 解析 \file{config.txt}：

\begin{txtcode}
arm_64bit=1                   # 以 AArch64 启动 ARM 核
enable_uart=1                 # 打开 Mini UART，固定 core 时钟
uart_2ndstage=1               # 固件自己的启动日志也从串口输出
disable_commandline_tags=1    # 不往内存里写 ATAG/设备树
core_freq=250                 # VPU 时钟 250 MHz（Unicam 的下限，串口波特率也依赖它）
core_freq_min=250
kernel=kernel8-xv6_wifi.img   # 内核镜像文件名
camera_auto_detect=0          # 不让固件占用摄像头
\end{txtcode}

然后把 \file{kernel8-xv6\_wifi.img} 原样拷到物理地址 \code{0x80000}，才放开 ARM 核（第~\ref{chap:boot} 章）。所以内核镜像必须是裸二进制而不是 ELF，固件也从不理会 xv6 的根文件系统。

\paragraph{第二组由 xv6 内核驱动读.} 两块 Wi-Fi 网卡每次上电都要由主机下载固件：BCM43430/43455 需要 \file{.BIN}（R4 的执行代码）、\file{.TXT}（NVRAM）和 \file{.CLM}（法规数据库），MT7601U 需要 \file{MT7601U.BIN}。驱动在 \code{main()} 里初始化，这时还没有任何进程，更谈不上挂载点，所以它们直接调用 \code{fat32openroot()} 按 8.3 短名在根目录里找文件，再用 \code{fat32readnext()} 顺序读出。这也是挂起 bootfs 的 \code{rootdev\_init()} 必须排在 \code{brcmfmac\_driver\_init()} 之前的原因。

\paragraph{第三组由用户程序经 \code{/boot} 读写.} 进入用户态以后，\code{init} 按 \file{/etc/fstab} 把同一个分区挂到 \code{/boot}（\ref{sec:vfs} 节）。于是 \code{ls /boot}、\code{cat /boot/config.txt} 都走普通的 \code{open/read}；\code{tftpclient} 把新内核下载成 \file{/boot/kernel8-xv6\_wifi.img}，重启即生效；\code{camshot} 把超过 xv6 单文件上限的照片写成 \file{/boot/camera.bmp}，关机后插到电脑上就能看到。

\subsection{磁盘上的 FAT32}

FAT32 的结构可以用一句话概括：\emph{文件内容按簇存放，簇与簇之间的先后关系记在一张表（FAT）里，目录本身也是一个文件}。图~\ref{fig:fat-layout} 是它在 SD 卡上的位置和内部划分。

\begin{figure}[htbp]
  \centering
  \begin{tikzpicture}[x=1mm,y=1mm, font=\scriptsize]
    \tikzset{seg/.style={draw=structurecolor!70!black, minimum height=10mm, inner sep=1pt, align=center, anchor=west}}
    % whole card
    \node[lab, anchor=east] at (-1,30) {整张卡};
    \node[seg, fill=gray!20, minimum width=12mm] (mbr) at (0,30) {MBR\\扇区 0};
    \node[seg, fill=gray!5, minimum width=8mm] (gap) at (mbr.east) {空};
    \node[seg, fill=structurecolor!20, minimum width=40mm] (p1) at (gap.east) {p1 BOOTFS（FAT32）\\LBA 2048 起};
    \node[seg, fill=orange!15, minimum width=34mm] (p2) at (p1.east) {p2 xv6 原始分区\\LBA 208896 起};
    \node[seg, fill=gray!10, minimum width=26mm] (p3) at (p2.east) {p3 ext2 …};
    % zoom p1
    \node[lab, anchor=east] at (-1,8) {p1 内部};
    \node[seg, fill=structurecolor!30, minimum width=10mm] (bpb) at (0,8) {BPB\\扇区 0};
    \node[seg, fill=structurecolor!20, minimum width=10mm] (fsi) at (bpb.east) {FSInfo\\扇区 1};
    \node[seg, fill=structurecolor!10, minimum width=12mm] (rsv) at (fsi.east) {保留区\\其余扇区};
    \node[seg, fill=orange!25, minimum width=20mm] (fat1) at (rsv.east) {FAT \#1};
    \node[seg, fill=orange!25, minimum width=20mm] (fat2) at (fat1.east) {FAT \#2（副本）};
    \node[seg, fill=gray!8, minimum width=48mm] (data) at (fat2.east) {数据区：簇 2、3、4 …\\根目录也是一条簇链};
    \draw[gray!60, dashed] (p1.south west) -- (bpb.north west);
    \draw[gray!60, dashed] (p1.south east) -- (data.north east);
    % labels below
    \draw[decorate, decoration={brace, amplitude=3pt, mirror}] ([yshift=-1mm]bpb.south west) -- node[lab, below=3pt]{\code{reserved} 个扇区} ([yshift=-1mm]rsv.south east);
    \draw[decorate, decoration={brace, amplitude=3pt, mirror}] ([yshift=-1mm]fat1.south west) -- node[lab, below=3pt]{\code{fats} $\times$ \code{fat\_sectors}} ([yshift=-1mm]fat2.south east);
    \node[lab, anchor=north] at (fat1.south west |- 0,-4) {\code{fat\_lba}};
    \node[lab, anchor=north] at (data.south west |- 0,-4) {\code{data\_lba}};
    \draw[gray!60] (fat1.south west) -- ++(0,-6.5);
    \draw[gray!60] (data.south west) -- ++(0,-6.5);
  \end{tikzpicture}
  \caption{SD 卡与 FAT32 分区的磁盘布局。\code{fat32mount()} 从 BPB 读出各区的大小，算出 \code{fat\_lba} 和 \code{data\_lba}；簇号从 2 开始编号}
  \label{fig:fat-layout}
\end{figure}

\paragraph{MBR 和分区表.} 整张卡的第 0 扇区是 MBR：偏移 446 起是 4 个 16 字节的分区表项，最后两个字节是签名 \code{0x55 0xAA}。每个表项里第 4 字节是分区类型（\code{0x0B}/\code{0x0C} 是 FAT32），第 8、12 字节是起始 LBA 和扇区数。读分区表的是 \file{rootdev.c} 的 \code{rootdev\_init()}，而不是 FAT32 驱动。它先把四个表项全部打印出来，再分两轮挑选：第一轮找 block 1 带 xv6 superblock 魔数的分区做根文件系统（\code{setup\_raw\_root()}），第二轮把第一个 FAT32 分区的起点交给 \code{fat32mount()}。两轮互不依赖，所以根文件系统和启动分区可以各在各的分区里。FAT32 驱动的工作从 \code{fat32mount()} 拿到分区起点开始。

\paragraph{BPB：分区的第一个扇区.} FAT32 分区的第 0 扇区叫 BPB（BIOS Parameter Block），记录了整个分区的几何参数。\code{fat32mount()} 用到的字段如表~\ref{tab:bpb}。

\begin{table}[htbp]
  \centering
  \caption{\code{fat32mount()} 读取的 BPB 字段}
  \label{tab:bpb}
  \small
  \begin{tabularx}{\linewidth}{@{}llX@{}}
    \toprule
    偏移 & 字段 & 用途 \\
    \midrule
    11 & 每扇区字节数 & 必须是 512，否则拒绝 \\
    13 & 每簇扇区数 \code{spc} & 必须是 2 的幂；一簇通常 4\,KiB 或更大 \\
    14 & 保留扇区数 \code{reserved} & FAT 表从 \code{partition\_lba + reserved} 开始 \\
    16 & FAT 份数 \code{fats} & 通常为 2，第二份是备份，写入时要同步 \\
    32 & 总扇区数 & 与 FAT 大小一起算出总簇数 \\
    36 & 每份 FAT 的扇区数 & 数据区从 \code{fat\_lba + fats×fat\_sectors} 开始 \\
    44 & 根目录首簇 & 通常为 2；根目录也是一条簇链 \\
    48 & FSInfo 扇区号 & 内有“下一个空闲簇”的提示 \\
    510 & \code{0x55 0xAA} & 签名 \\
    \bottomrule
  \end{tabularx}
\end{table}

有了这些字段，任何簇号都能换算成 SD 卡上的绝对扇区号：

\begin{ccode}
fat_lba  = partition_lba + reserved;
data_lba = fat_lba + fats * fat_sectors;
static uint32 cluster_lba(uint32 cluster) {     // 簇号从 2 开始
  return diskmap.data_lba + (cluster - 2) * diskmap.sectors_per_cluster;
}
\end{ccode}

\code{fat32mount()} 对每个字段都做了合法性检查（2 的幂、不越界、不溢出）——SD 卡上的数据是外部输入，一个损坏的 BPB 不应该让内核算出一个指向别处的 LBA 然后去写它。

\subsection{FAT 表与簇链}

FAT 表是一个巨大的 32 位整数数组，下标就是簇号，第 $c$ 项的值是“簇 $c$ 之后的下一个簇”。只有低 28 位有效，高 4 位保留且写入时要保持原值；0 表示空闲，\code{0x0FFFFFF8} 及以上表示链尾（EOC）。一个 512 字节的扇区装 128 项，所以查第 $c$ 项只需读一个扇区：

\begin{ccode}
static uint32 fat_next(uint32 cluster) {
  uint32 sector = diskmap.fat_lba + cluster / 128;   // 128 = 512 / 4
  uint32 pos    = (cluster % 128) * 4;
  if(sdsector(sector, scratch, 0) < 0) panic("fat32: read FAT");
  return le32(scratch + pos) & 0x0fffffffU;
}
\end{ccode}

\code{fat\_set\_next()} 是对应的写操作，它要把同一个扇区在每一份 FAT 里都读改写一遍，保证两份副本一致。图~\ref{fig:fat-chain} 下半部分画的就是一条簇链：从目录项里拿到首簇，然后反复 \code{fat\_next()} 直到 EOC。

\subsection{目录项与长文件名}

\begin{figure}[htbp]
  \centering
  \begin{tikzpicture}[x=1mm,y=1mm, font=\scriptsize]
    \tikzset{cell/.style={draw=gray!60, minimum height=7mm, inner sep=1pt, align=center, anchor=west}}
    % directory entries
    \node[lab, anchor=west] at (0,66) {根目录簇中的连续 32 字节目录项};
    \node[cell, fill=structurecolor!10, minimum width=118mm] (l2) at (0,60) {LFN \#2（序号 \code{0x42}＝最后一片 | 第 2 片）：\code{"ifi.img"}，属性 \code{0x0F}，校验和 \code{c}};
    \node[cell, fill=structurecolor!10, minimum width=118mm] (l1) at (0,53) {LFN \#1（序号 \code{0x01}）：\code{"kernel8-xv6\_w"}，前 13 个 UTF-16 字符，属性 \code{0x0F}，校验和 \code{c}};
    \node[cell, fill=orange!15, minimum width=22mm] (sn) at (0,45) {\code{KERNEL\~{}1IMG}\\名字 [0:11)};
    \node[cell, fill=orange!8, minimum width=14mm] (at) at (sn.east) {属性\\{}[11]＝\code{0x20}};
    \node[cell, fill=gray!5, minimum width=30mm] (mi) at (at.east) {时间、日期等};
    \node[cell, fill=orange!25, minimum width=22mm] (hi) at (mi.east) {首簇高 16 位\\{}[20:22)};
    \node[cell, fill=orange!25, minimum width=16mm] (lo) at (hi.east) {首簇低 16 位\\{}[26:28)};
    \node[cell, fill=orange!8, minimum width=14mm] (sz) at (lo.east) {大小\\{}[28:32)};
    \node[lab, anchor=west] at (119,56.5) {LFN 片倒序存放};
    \node[lab, anchor=west] at (119,45) {短名目录项};
    % FAT table
    \node[lab, anchor=west, align=left] at (109,22) {FAT 表\\每项 4 字节，低 28 位有效};
    \foreach \i/\v/\c in {2/EOC/gray!8, 3/4/gray!8, 4/EOC/gray!8, 5/6/orange!25, 6/9/orange!25, 7/0/gray!3, 8/0/gray!3, 9/EOC/orange!25, 10/0/gray!3}{
      \node[cell, minimum width=11mm, fill=\c] (fat\i) at ({(\i-2)*12},22) {[\i]\enspace\code{\v}};
    }
    % data clusters
    
    \foreach \i/\c in {2/gray!8, 3/gray!8, 4/gray!8, 5/structurecolor!20, 6/structurecolor!20, 7/gray!3, 8/gray!3, 9/structurecolor!20, 10/gray!3}{
      \node[cell, minimum width=11mm, fill=\c] (dc\i) at ({(\i-2)*12},6) {簇 \i};
    }
    \node[lab, anchor=west, align=left] at (109,6) {数据区的簇\\簇 2 是根目录};
    % arrows
    \draw[-{Stealth[length=2mm]}, thick, orange!80!black, rounded corners] (lo.south) -- (lo.south |- 0,37.5) -- node[lab, below, pos=0.5]{首簇号（高 16 位拼低 16 位）＝ 5} (fat5.north |- 0,37.5) -- (fat5.north);
    \draw[-{Stealth[length=1.6mm]}, thick, orange!80!black] (fat5.south east) to[out=-30, in=-150] (fat6.south west);
    \draw[-{Stealth[length=1.6mm]}, thick, orange!80!black] (fat6.south east) to[out=-25, in=-155] (fat9.south west);
    \foreach \i in {5,6,9}{ \draw[gray!60, dashed] (fat\i.south) ++(0,-0.2) -- (dc\i.north); }
  \end{tikzpicture}
  \caption{长文件名、目录项与簇链。文件 \code{kernel8-xv6\_wifi.img} 占两片 LFN 加一个短名目录项；短名项里的首簇号是 5，沿 FAT 表 5→6→9 直到 EOC 就是它的全部数据簇}
  \label{fig:fat-chain}
\end{figure}

目录的内容是一串 32 字节的目录项（图~\ref{fig:fat-chain} 上半部分）。短名目录项里有 8.3 格式的名字（11 字节，大写、空格填充）、属性字节、首簇号（拆成高低两个 16 位字段，这是 FAT16 向 FAT32 扩展时留下的痕迹）和文件大小。扫描目录时，首字节有两个特殊值：\code{0x00} 表示后面再没有目录项，\code{0xE5} 表示这一项已被删除。属性字节的 \code{0x10} 表示子目录，\code{0x08} 表示卷标。

长文件名（LFN）是后来加的：属性等于 \code{0x0F} 的目录项不是文件，而是一片名字，每片装 13 个 UTF-16 字符，倒序排在它所属的短名目录项前面。第一片（物理上最靠前的）序号带 \code{0x40} 标志，表示“最后一片”；每片都存着短名的校验和，以便发现“LFN 和短名不配套”（例如老系统删了短名但留下了 LFN）。\code{fat\_find\_path()} 扫描目录时把遇到的 LFN 片按序号拼进一个缓冲区，遇到短名项时检查三件事——片序号连续到 1、校验和与短名一致、拼出的名字与要找的名字相同（不区分大小写）——满足才算长名匹配；否则再按短名比较。

创建文件时反过来：\code{fat32createfile()} 先用 \code{fat\_make\_short\_alias()} 造一个 \code{KERNEL\~{}1.IMG} 这样的短名（取前 6 个字母数字、加 \code{\~{}1}～\code{\~{}9} 避免重名），算出需要几片 LFN，用 \code{fat\_reserve\_root\_slots()} 在根目录里找到足够的连续空位（不够就给根目录分配新簇），按倒序写入 LFN 片和短名项，最后补一个全 0 的目录项作为结束标记。

本项目的 FAT32 驱动只处理\emph{根目录}：查找、创建、枚举、重命名都在根目录簇链里进行，不支持子目录，LFN 也只支持 ASCII 字符。对启动分区来说这已经足够——上面列出的所有文件都在根目录里。

\subsection{函数一览与调用关系}

\begin{figure}[htbp]
  \centering
  \begin{tikzpicture}[x=1mm,y=1mm, font=\scriptsize,
      fn/.style={box, font=\scriptsize, minimum height=6.5mm, inner sep=1.5pt},
      call/.style={-{Stealth[length=1.6mm]}, semithick, draw=structurecolor!70!black}]
    \foreach \y/\t in {64/{调用者}, 48/{\file{fat32.c} 接口}, 30/{目录与 FAT}, 14/{扇区}, 2/{设备}}{
      \fill[gray!5] (0,\y-6.5) rectangle (150,\y+6.5);
      \node[lab, align=center, text width=14mm] at (7,\y) {\t};
    }
    % callers
    \node[fn, text width=34mm] (c1) at (34,64) {Wi-Fi 驱动下载固件\\\file{brcmfmac.c}、\file{mt7601u.c}};
    \node[fn, text width=34mm] (c2) at (80,64) {VFS \code{/boot}：\code{fat\_vread}\\\code{fat\_vwrite}、\code{fat\_vreaddir}…};
    \node[fn, text width=30mm] (c3) at (128,64) {根块设备 \file{rootdev.c}\\（仅 \code{FS.IMG} 兼容模式）};
    % interface
    \node[fn, text width=30mm] (i1) at (30,48) {\code{fat32openroot}\\\code{fat32readerinit/readnext}};
    \node[fn, text width=40mm] (i2) at (80,48) {\code{fat32statpath} \code{fat32pread}\\\code{fat32createfile} \code{fat32writefile}\\\code{fat32truncatefile} \code{fat32renamefile}};
    \node[fn, text width=30mm] (i3) at (128,48) {\code{fat32mapfile}\\\code{fat32mapsector}};
    % helpers
    \node[fn, text width=37mm] (h1) at (36,30) {\code{fat\_find\_path}/\code{fat\_find\_root}\\\code{fat\_reserve\_root\_slots}\\\code{fat\_make\_lfn\_entry}};
    \node[fn, text width=40mm] (h2) at (86,30) {\code{fat\_next} \code{fat\_set\_next}\\\code{fat\_alloc\_cluster} \code{fat\_free\_chain}\\\code{fat\_mark\_fsinfo\_unknown}};
    % sector
    \node[fn, text width=70mm] (s1) at (64,14) {\code{cluster\_lba(c) = data\_lba + (c-2)×spc}\\全部经 \code{sdsector(lba, scratch, write)}，持 \code{fat32\_lock}};
    \node[fn, text width=70mm, fill=gray!10, draw=gray!60] (d1) at (64,2) {\file{sdhost.c}：CMD17/CMD24，PIO 搬运 512 字节};
    \draw[call] (c1) -- (i1); \draw[call] (c2) -- (i2); \draw[call] (c3) -- (i3);
    \draw[call] (i1) -- (h1); \draw[call] (i2) -- (h1); \draw[call] (i2) -- (h2); \draw[call] (i1.south east) -- (h2.north west);
    \draw[call] (h1) -- (s1); \draw[call] (h2) -- (s1);
    \draw[call] (i3.east) -- ++(3,0) |- (s1.east);
    \draw[call] (s1) -- (d1);
    \node[fn, text width=30mm, fill=gray!10, draw=gray!60] (boot) at (127,30) {启动时 \code{rootdev\_init}\\→ \code{fat32mount}};
  \end{tikzpicture}
  \caption{FAT32 驱动的调用关系。三类调用者共用同一组目录与 FAT 操作，最终都落到一个 512 字节的 \code{scratch} 扇区缓冲和 \code{sdsector()} 上}
  \label{fig:fat-calls}
\end{figure}

\file{fat32.c} 的函数可以按图~\ref{fig:fat-calls} 分成四层：

\begin{itemize}
  \item \textbf{挂载}：\code{fat32mount(lba, sectors)} 解析 BPB、读 FSInfo 里的空闲簇提示；\code{fat32\_is\_boot\_sector()} 判断一个扇区是不是 FAT32 引导扇区，\file{rootdev.c} 用它识别“没有分区表、整张卡就是一个 FAT32”的 superfloppy 格式；\code{fat32init()} 只初始化 \code{fat32\_lock}。分区表扫描和 xv6 根分区识别都不在这个文件里。
  \item \textbf{对外接口}：只读的 \code{fat32openroot()}、\code{fat32openpath()}、\code{fat32statpath()}、\code{fat32readdirroot()}、\code{fat32pread()}、\code{fat32readerinit()}/\code{fat32readnext()}；写入的 \code{fat32createfile()}、\code{fat32openwrite()}、\code{fat32writefile()}、\code{fat32truncatefile()}、\code{fat32renamefile()}。VFS 的 \code{fat32\_ops} 把它们包装成 \code{/boot} 上的 vnode 操作。
  \item \textbf{目录与 FAT}：上一小节的查找和 LFN 函数，以及 \code{fat\_next()}、\code{fat\_set\_next()}、\code{fat\_alloc\_cluster()}、\code{fat\_free\_chain()}、\code{fat\_mark\_fsinfo\_unknown()}。
  \item \textbf{文件扇区映射}：\code{fat32mapfile()} 记下根目录中某个文件的整条簇链，\code{fat32mapsector()} 把“文件内第 $s$ 个扇区”换成卡上的扇区号。它们只为兼容早期把根镜像存成 \code{FS.IMG} 文件的旧卡而保留，调用者是 \file{rootdev.c}。FAT32 只回答“这个扇区在哪里”，从不解释扇区里的 xv6 数据。
\end{itemize}

两个读接口的区别值得一提。\code{fat32pread(file, offset, …)} 是无状态的：每次调用都从首簇沿链走 $\lfloor\text{offset}/\text{簇大小}\rfloor$ 步，适合 VFS 偶尔的随机读。如果用它按 512 字节分块读一个 600\,KiB 的 Wi-Fi 固件，总共要走 $O(n^2)$ 次 \code{fat\_next()}。\code{fat32reader} 则记住当前簇和簇内偏移，顺序读时每跨一个簇只查一次 FAT——固件下载就用它。

\subsection{写入：分配、追加与截断}

\code{fat32writefile()} 只支持追加：要求写入位置正好等于文件当前大小。写入过程是：

\begin{enumerate}
  \item 重新按路径查找目录项，核对首簇和大小与调用者手里的 \code{fat32\_file} 一致，防止文件已被别人改过；
  \item 写入跨进新簇时调用 \code{fat\_alloc\_cluster()}：从 FSInfo 提示的位置开始找 FAT 值为 0 的项，把它标成 EOC（每份 FAT 都写），把这一簇清零；再用 \code{fat\_set\_next()} 把它挂到链尾；
  \item 按扇区写数据：整扇区直接写，不满一个扇区的先读出来再改再写回；
  \item 最后一次性更新目录项里的首簇和大小。
\end{enumerate}

\code{fat32truncatefile()}（\code{O\_TRUNC}）的顺序恰好相反：先把目录项的首簇和大小清零并写回，再释放旧簇链。这个顺序的意义在于断电时的后果：先改目录项，最坏情况是旧簇成了没人引用的“丢失簇”，浪费一点空间；如果先释放簇链，断电后目录项会指向已被标为空闲、随时可能分给别人的簇，那才是真正的损坏。分配和释放之后，\code{fat\_mark\_fsinfo\_unknown()} 把 FSInfo 里的空闲簇计数标成“未知”（\code{0xFFFFFFFF}），只更新“下一个空闲簇”提示，避免留下一个过时的计数让别的系统误信。

两处性能问题都是真机上暴露出来的，修法相同——缓存当前 FAT 扇区：
\begin{itemize}
  \item \code{fat\_alloc\_cluster()} 最初对每个候选簇调用一次 \code{fat\_next()}，同一个扇区要重复读 128 次，根目录一长就像卡死了一样；现在扫描时只在跨扇区时才读；
  \item \code{fat\_free\_chain()} 最初对每个簇做一次读和两份 FAT 的读改写，截断一张 320$\times$240 的 BMP 要五百多次 SD 事务（第~\ref{chap:drivers} 章 \ref{sec:cam-bugs} 小节）；现在在一个扇区内批量清除链项，换扇区时才把整扇区写到每份 FAT。
\end{itemize}

\paragraph{并发与局限.} 所有接口都持同一把自旋锁 \code{fat32\_lock}，共用一个 512 字节的 \code{scratch} 扇区缓冲，没有块缓存，也没有日志。路径解析、LFN 字符、目录槽这些较大的临时数组也放在一个静态工作区 \code{fat\_work} 里，而不是放在栈上——每个进程的内核栈只有 4\,KiB，它们曾经把创建文件的调用链撑爆（第~\ref{chap:drivers} 章 \ref{sec:cam-bugs} 小节）；既然有 \code{fat32\_lock} 串行化，一个工作区就够用。这让实现很简单，代价也很清楚：持锁期间 \code{sdsector()} 以 PIO 方式等待 SD 卡，一次大文件写会长时间占着这把锁；同一时间只能有一个写者；写入途中断电可能损坏分区。启动分区平时几乎只读，偶尔更新内核或写一张照片，这些代价是可以接受的，但修改它之前最好备份。

\section{xv6 根文件系统：superblock、inode 与目录}
\label{sec:xv6fs}

前一节分析了启动分区上的 FAT32，这一节用同样的方法看 p2 上的 xv6 根文件系统：它怎么被认出来，磁盘上每一块放的是什么，文件和目录怎样表示，一个路径怎样一步步变成 SD 卡上的扇区。\ref{sec:rootfs} 节讲过根文件系统可以放在哪里；这里只讨论现在的“原始分区”方式。

\subsection{根文件系统怎样被认出来}

\code{rootdev\_init()} 扫描 MBR 时，对每个既不是 FAT32 也不是扩展分区的主分区调用 \code{setup\_raw\_root()}（两者都在 \file{kernel/rootdev.c}）：

\begin{ccode}
static int setup_raw_root(uint32 start, uint32 sectors) {
  // xv6 的第 1 块是 superblock；块大小 1024，所以它在分区起点之后第 2 个扇区
  if(start == 0 || sectors < FSSIZE * SECTORS_PER_BLOCK ||
     sdsector(start + SECTORS_PER_BLOCK, sector, 0) < 0 ||
     le32(sector) != FSMAGIC || le32(sector + 4) != FSSIZE)
    return -1;
  root.mode = ROOT_RAW;
  root.lba = start;
  root.sectors = sectors;
  return 0;
}
\end{ccode}

判断只看两样东西：superblock 开头的魔数 \code{0x10203040}，以及紧跟着的 \code{size} 是否等于内核编译时的 \code{FSSIZE}。MBR 里的分区类型不作数——macOS 的 \code{diskutil} 只能把这个分区建成 Linux 的 \code{0x83} 类型，和真正的 ext2 分区一样，靠魔数才能分开。\code{size} 的检查同样重要：如果有人改了 \code{FSSIZE} 重新编译内核却没重新烧 \code{fs.img}，这里会拒绝，而不是让新内核去解释旧格式的磁盘。识别成功后启动日志打印：

\begin{txtcode}
rootdev: raw xv6 root p2 lba=208896 sectors=819200 size=400 MiB
\end{txtcode}

此后块缓存的每一次读写都经过 \code{rootdev\_rw()}：

\begin{ccode}
void rootdev_rw(struct buf *b, int write) {         // kernel/rootdev.c
  uint32 first = b->blockno * SECTORS_PER_BLOCK;    // 1 块 = 2 扇区
  gate_enter();                                     // /dev/sdroot 更新时在此挡住
  for(int i = 0; i < SECTORS_PER_BLOCK; i++)
    if(sdsector(root_sector_lba(first + i),
                b->data + i * SECTOR_SIZE, write) < 0)
      panic(write ? "rootdev: write block" : "rootdev: read block");
  gate_exit();
}
\end{ccode}

原始分区模式下 \code{root\_sector\_lba()} 只做一次加法：扇区号 $=$ \code{root.lba} $+$ 块号 $\times 2 + i$。整条路径不经过 \file{fat32.c} 的任何函数。分区有 400\,MiB，文件系统只用前 32\,MiB；可分配多少块由 superblock 决定，与分区大小无关。

\subsection{superblock 与磁盘布局}

\begin{figure}[htbp]
  \centering
  \begin{tikzpicture}[x=1mm,y=1mm, font=\scriptsize]
    \tikzset{seg/.style={draw=structurecolor!70!black, minimum height=11mm, inner sep=1pt, align=center, anchor=west}}
    \node[seg, fill=gray!15, minimum width=10mm] (b0) at (0,10) {启动块\\0};
    \node[seg, fill=orange!30, minimum width=13mm] (sb) at (b0.east) {superblock\\1};
    \node[seg, fill=structurecolor!25, minimum width=24mm] (lg) at (sb.east) {日志\\2～31（30 块）};
    \node[seg, fill=structurecolor!15, minimum width=26mm] (in) at (lg.east) {inode 区\\32～44（13 块）};
    \node[seg, fill=orange!15, minimum width=20mm] (bm) at (in.east) {位图\\45～49（5 块）};
    \node[seg, fill=gray!5, minimum width=45mm] (dt) at (bm.east) {数据区\\50～32767（32718 块）};
    % zooms
    \node[lab, anchor=north, align=center] at (lg.south) {第一块是日志头：\\\code{n} + 30 个块号};
    \node[lab, anchor=north, align=center] at (in.south) {每块 16 个 \code{dinode}\\共 200 个 inode};
    \node[lab, anchor=north, align=center] at (bm.south) {每块 8192 位\\1 位管 1 块};
    \node[lab, anchor=north, align=center] at (dt.south) {文件内容、目录内容、\\间接块};
    \draw[-{Stealth[length=1.6mm]}, orange!70!black] (sb.north) to[out=70, in=110] node[lab, above]{\code{logstart}} (lg.north);
    \draw[-{Stealth[length=1.6mm]}, orange!70!black] (sb.north) to[out=60, in=115] node[lab, above, pos=0.6]{\code{inodestart}} (in.north);
    \draw[-{Stealth[length=1.6mm]}, orange!70!black] (sb.north) to[out=55, in=120] node[lab, above, pos=0.7]{\code{bmapstart}} (bm.north);
    \node[lab, anchor=south west] at (0,36) {块号（1\,KiB 一块）；分区内的扇区号 $=$ 块号 $\times 2$};
  \end{tikzpicture}
  \caption{\code{mkfs} 生成的 32\,MiB 根文件系统布局。所有数字都由 \code{FSSIZE}、\code{NINODES}、\code{LOGSIZE} 算出，并写进第 1 块的 superblock}
  \label{fig:xv6-layout}
\end{figure}

superblock 是第 1 块（第 0 块是从 Unix 继承下来的“启动块”，xv6 不用），内容就是 8 个 32 位整数：

\begin{table}[htbp]
  \centering
  \caption{\code{fs.img} 的 superblock（\code{struct superblock}，\file{kernel/fs.h}）}
  \label{tab:xv6-sb}
  \small
  \begin{tabularx}{\linewidth}{@{}llX@{}}
    \toprule
    字段 & 值 & 含义与来源 \\
    \midrule
    \code{magic} & \code{0x10203040} & \code{FSMAGIC}；\code{fsinit()} 和 \code{setup\_raw\_root()} 都检查它 \\
    \code{size} & 32768 & 总块数，\code{FSSIZE}；32\,MiB \\
    \code{nblocks} & 32718 & 数据块数 $=$ \code{size} $-$ 元数据块数 50 \\
    \code{ninodes} & 200 & \code{mkfs} 里的 \code{NINODES} \\
    \code{nlog} & 30 & \code{LOGSIZE} $=$ \code{MAXOPBLOCKS}$\times 3$ \\
    \code{logstart} & 2 & 紧跟 superblock \\
    \code{inodestart} & 32 & $2 + 30$；inode 区 $200/16+1=13$ 块 \\
    \code{bmapstart} & 45 & $32 + 13$；位图 $32768/8192+1=5$ 块 \\
    \bottomrule
  \end{tabularx}
\end{table}

这些数都是 \code{mkfs} 在宿主机上算出来写进去的（运行 \code{make fs.img} 时会打印 \code{nmeta 50 (boot, super, log blocks 30 inode blocks 13, bitmap blocks 5) blocks 32718 total 32768}）。内核启动时 \code{fsinit()} 读第 1 块、核对魔数、交给日志层，之后所有“某个东西在哪一块”的问题都由 superblock 加两个宏回答：

\begin{ccode}
#define IPB           (BSIZE / sizeof(struct dinode))   // 每块 16 个 inode
#define IBLOCK(i, sb) ((i) / IPB + sb.inodestart)       // inode i 所在的块
#define BPB           (BSIZE * 8)                       // 每个位图块管 8192 块
#define BBLOCK(b, sb) ((b) / BPB + sb.bmapstart)        // 块 b 的位所在的位图块
\end{ccode}

\paragraph{日志区.} 30 块中第一块是日志头 \code{struct logheader}：一个计数 \code{n} 加上 \code{n} 个“这些块属于哪个目标块号”。其余是被修改的块的副本。提交时先写副本、再写日志头（这一写就是提交点）、再把副本拷回原位、最后把 \code{n} 清零；\code{fsinit()} → \code{initlog()} 在启动时检查日志头，\code{n} 不为 0 就重做一遍。

\paragraph{位图.} 每一位对应一块，1 表示已用。\code{mkfs} 把前 \code{freeblock} 块（元数据加上镜像里已经写入的全部文件内容）一次性标成已用；运行时 \code{balloc()} 从头扫描位图找第一个 0 位，置 1、\code{log\_write}、再把新块清零。早期 \code{FSSIZE} 只有 1000 块时，光是打包进镜像的程序就用掉 986 块，首次启动创建 \code{/etc} 配置时 \code{balloc} 找不到 0 位，panic “out of blocks”——这就是 \code{FSSIZE} 改成 32768 的原因。

\subsection{inode：文件的元数据与块映射}

\begin{figure}[htbp]
  \centering
  \begin{tikzpicture}[x=1mm,y=1mm, font=\scriptsize]
    \tikzset{f/.style={draw=gray!60, minimum height=6mm, inner sep=1pt, align=center, anchor=west}}
    \node[lab, anchor=west] at (0,47) {\code{struct dinode}（64 字节，inode 区每块 16 个）};
    \node[f, fill=orange!20, minimum width=12mm] (t) at (0,40) {\code{type}\\2 B};
    \node[f, fill=gray!8, minimum width=12mm] (ma) at (t.east) {\code{major}\\2 B};
    \node[f, fill=gray!8, minimum width=12mm] (mi) at (ma.east) {\code{minor}\\2 B};
    \node[f, fill=orange!10, minimum width=12mm] (nl) at (mi.east) {\code{nlink}\\2 B};
    \node[f, fill=orange!10, minimum width=12mm] (sz) at (nl.east) {\code{size}\\4 B};
    \node[f, fill=structurecolor!15, minimum width=60mm] (ad) at (sz.east) {\code{addrs[0..11]}：12 个直接块号\\48 B};
    \node[f, fill=structurecolor!30, minimum width=16mm] (ind) at (ad.east) {\code{addrs[12]}\\间接块号};
    % data blocks
    \foreach \i in {0,...,4}{
      \draw[fill=gray!5, draw=gray!60] (62+\i*7,22) rectangle ++(6,6);
    }
    \node[lab] at (98,25) {…};
    \node[lab, anchor=west] at (62,18) {直接块：文件第 0～11 KiB};
    \foreach \i in {0,...,3}{ \draw[-{Stealth[length=1.2mm]}, gray!60] (64+\i*10,37) -- (65+\i*7,28); }
    \draw[fill=structurecolor!20, draw=structurecolor!70!black] (122,18) rectangle ++(18,10) node[midway, align=center]{间接块\\256 个块号};
    \draw[-{Stealth[length=1.6mm]}, thick, structurecolor!70!black] (ind.south) -- (131,28);
    \foreach \i in {0,...,4}{
      \draw[fill=gray!5, draw=gray!60] (110+\i*7,2) rectangle ++(6,6);
    }
    \node[lab] at (147,5) {…};
    \draw[-{Stealth[length=1.6mm]}, gray!60] (131,18) -- (124,8);
    \draw[-{Stealth[length=1.6mm]}, gray!60] (131,18) -- (138,8);
    \node[lab, anchor=east] at (108,5) {文件第 12～267 KiB};
    \node[lab, anchor=west, align=left] at (0,24) {\code{type}：0 空闲，1 目录，\\2 普通文件，3 设备\\[2pt]\code{nlink}：有几个目录项指向它\\[2pt]最大文件 $12+256=268$ 块};
  \end{tikzpicture}
  \caption{磁盘上的 inode 和它的块映射。\code{bmap(ip, bn)} 把文件内第 \code{bn} 块翻译成磁盘块号：小于 12 直接查 \code{addrs[bn]}，否则读间接块再查第 \code{bn-12} 项；遇到 0 就当场 \code{balloc()} 一块}
  \label{fig:xv6-inode}
\end{figure}

每个文件、目录和设备节点都由一个 64 字节的磁盘 inode 描述（图~\ref{fig:xv6-inode}）。它不包含文件名——名字在目录里——只有类型、链接数、大小和数据块的位置。内存里的 \code{struct inode} 在此之外再加上 \code{dev}、\code{inum}、引用计数 \code{ref}、睡眠锁和 \code{valid} 标志（是否已从磁盘读入）。\code{iget()} 只在 inode 缓存里占一个槽、不读盘；\code{ilock()} 第一次加锁时才 \code{bread(IBLOCK(...))} 把磁盘 inode 拷进来；修改后 \code{iupdate()} 写回。

块映射由 \code{bmap()} 完成：

\begin{ccode}
static uint bmap(struct inode *ip, uint bn) {
  uint addr, *a;
  struct buf *bp;
  if(bn < NDIRECT){                                   // 前 12 块：直接块
    if((addr = ip->addrs[bn]) == 0)
      ip->addrs[bn] = addr = balloc(ip->dev);
    return addr;
  }
  bn -= NDIRECT;
  if(bn < NINDIRECT){                                 // 之后 256 块：经间接块
    if((addr = ip->addrs[NDIRECT]) == 0)
      ip->addrs[NDIRECT] = addr = balloc(ip->dev);
    bp = bread(ip->dev, addr);
    a = (uint*)bp->data;
    if((addr = a[bn]) == 0){
      a[bn] = addr = balloc(ip->dev);
      log_write(bp);
    }
    brelse(bp);
    return addr;
  }
  panic("bmap: out of range");
}
\end{ccode}

所以单个文件最大 $12+256=268$ 块，即 268\,KiB。这是 xv6 根文件系统最明显的限制：Wi-Fi 固件、640$\times$480 的 BMP（900\,KB）、原始帧（375\,KB）都放不下，只能放在 FAT32 的 \code{/boot}（\ref{sec:fat32} 节）。

根文件系统里几个特殊的 inode 号：1 号永远是根目录（\code{ROOTINO}），\code{namex()} 遇到以 \code{/} 开头的路径就从它开始；2 号是 \code{mkfs} 建的 \code{/bin}；之后是按 \code{Makefile} 顺序装入的每个用户程序。\code{init} 是唯一一个在根目录下也有名字的程序（启动代码执行的是 \code{exec("init")}），所以 \code{/init} 和 \code{/bin/init} 是同一个 inode 的两个目录项，\code{mkfs} 把它的 \code{nlink} 设为 2。

\subsection{目录与路径查找}

\begin{figure}[htbp]
  \centering
  \begin{tikzpicture}[x=1mm,y=1mm, font=\scriptsize,
      ino/.style={draw=orange!70!black, fill=orange!12, rounded corners=1pt, align=left, inner sep=2pt, text width=34mm, anchor=north west},
      de/.style={draw=gray!60, fill=gray!4, minimum height=4.6mm, minimum width=34mm, inner sep=1pt, anchor=north west, font=\scriptsize\ttfamily},
      hit/.style={de, fill=structurecolor!20, draw=structurecolor!70!black},
      go/.style={-{Stealth[length=1.8mm]}, thick, draw=structurecolor!80!black},
      lk/.style={-{Stealth[length=1.6mm]}, semithick, draw=orange!70!black}]
    % column 1: root
    \node[lab, anchor=west] at (0,62) {\ding{172} \code{iget(ROOTDEV, 1)}};
    \node[ino] (i1) at (0,58) {inode 1：\code{type=1} 目录\\\code{size=1024}，\code{addrs[0]=50}};
    \node[lab, anchor=west] at (0,44) {块 50（根目录内容）};
    \node[de] (r0) at (0,41) {\ \ 1  .};
    \node[de] (r1) at (r0.south west) {\ \ 1  ..};
    \node[hit] (r2) at (r1.south west) {\ \ 2  bin};
    \node[de] (r3) at (r2.south west) {\ \ k  init};
    \node[de] (r4) at (r3.south west) {\ \ …  dev etc boot…（init 建）};
    \node[de] (r5) at (r4.south west) {\ \ 0  （空槽）};
    \draw[lk] (i1.south) ++(-6,0) -- ++(0,-3.5);
    % column 2: /bin
    \node[lab, anchor=west] at (52,62) {\ding{173} \code{dirlookup(root,"bin")}};
    \node[ino] (i2) at (52,58) {inode 2：\code{type=1} 目录\\\code{size=}若干 KiB，\code{addrs[0]=51}…};
    \node[lab, anchor=west] at (52,44) {块 51…（\code{/bin} 的内容）};
    \node[de] (b0) at (52,41) {\ \ 2  .};
    \node[de] (b1) at (b0.south west) {\ \ 1  ..};
    \node[de] (b2) at (b1.south west) {\ \ 3  cat};
    \node[de] (b3) at (b2.south west) {\ \ …};
    \node[hit] (b4) at (b3.south west) {\ \ m  ls};
    \node[de] (b5) at (b4.south west) {\ \ …};
    \draw[lk] (i2.south) ++(-6,0) -- ++(0,-3.5);
    % column 3: ls
    \node[lab, anchor=west] at (104,62) {\ding{174} \code{dirlookup(bin,"ls")}};
    \node[ino] (i3) at (104,58) {inode m：\code{type=2} 文件\\\code{size}＝ELF 大小\\\code{addrs[0..11]}、\code{addrs[12]}};
    \node[draw=gray!60, fill=gray!5, align=center, text width=34mm, anchor=north west, inner sep=2pt] (dl) at (104,41) {数据块：ELF 头、\\程序段……\\（\code{exec} 用 \code{readi} 读取）};
    \draw[lk] (i3.south) ++(-6,0) -- ++(0,-6.5);
    % lookups
    \draw[go] (r2.east) -- ++(4,0) |- (i2.west);
    \draw[go] (b4.east) -- ++(4,0) |- (i3.west);
    \node[lab, anchor=west, align=left] at (0,6) {每个目录项 32 字节：\code{ushort inum} + \code{char name[30]}，每块 32 项；\code{inum=0} 表示空槽。\\\code{/init} 与 \code{/bin/init} 是同一个 inode k 的两个名字，所以它的 \code{nlink=2}。};
  \end{tikzpicture}
  \caption{\code{namei("/bin/ls")} 的查找过程。目录本身也是文件，内容是一串目录项；\code{namex()} 每次用 \code{skipelem()} 取出一个路径分量，在当前目录的内容里 \code{dirlookup()}，得到下一个 inode。块号 50、51 是 \code{mkfs} 的分配顺序决定的，k、m 取决于 \code{Makefile} 里程序的排列顺序}
  \label{fig:xv6-namei}
\end{figure}

目录就是 \code{type=1} 的文件，内容是一串 32 字节的 \code{struct dirent}：2 字节 inode 号加 30 字节文件名（\code{DIRSIZ} 从原版的 14 扩到 30，见 \ref{sec:layout} 节），一块恰好 32 项，\code{inum=0} 表示空槽。\code{mkfs} 把目录大小向上取整到整块。目录操作只有两个基本函数：\code{dirlookup(dp, name)} 顺序读目录内容、比较名字，命中就 \code{iget()} 对应的 inode；\code{dirlink(dp, name, inum)} 先确认不重名，再找第一个空槽（没有就追加在末尾）写入新目录项。

路径查找是 \code{namex()}，图~\ref{fig:xv6-namei} 画的是 \code{namei("/bin/ls")}：

\begin{ccode}
if(*path == '/') ip = iget(ROOTDEV, ROOTINO);        // 绝对路径从 inode 1 开始
else             ip = idup(myproc()->cwd);            // 相对路径从当前目录开始
while((path = skipelem(path, name)) != 0){           // 依次取出 "bin"、"ls"
  ilock(ip);
  if(ip->type != T_DIR){ iunlockput(ip); return 0; }
  if(nameiparent && *path == '\0'){ iunlock(ip); return ip; }   // 停在父目录
  if((next = dirlookup(ip, name, 0)) == 0){ iunlockput(ip); return 0; }
  iunlockput(ip);
  ip = next;
}
\end{ccode}

注意加锁的顺序：每一步只锁当前目录，查到下一级以后先 \code{iunlockput()} 当前目录再进入下一级，任何时刻最多持有一把 inode 锁，所以不会因为两个进程以相反顺序查找路径而死锁。\code{nameiparent()} 在最后一个分量之前停下，返回父目录并把最后一个名字留给调用者——\code{create}、\code{unlink}、\code{link} 都需要它。

根目录的内容分两部分：\code{mkfs} 写入的 \code{.}、\code{..}、\code{bin}、\code{init}；以及 \code{init} 首次启动时创建的 \code{dev}（含 \code{console}、\code{ttyS0}、\code{input/event0}、\code{video0}、\code{sdroot} 等设备节点）、\code{etc}（\code{fstab} 等配置）、\code{proc}、\code{boot}、\code{mnt/ext2}、\code{root}、\code{usr/bin}。其中 \code{proc}、\code{boot}、\code{mnt/ext2} 只是空目录，挂载以后会被 VFS 的挂载表“盖住”（\ref{sec:vfs} 节）。设备节点是 \code{type=3} 的 inode，不占数据块，只在 \code{major}/\code{minor} 里记下交给哪个驱动。

\subsection{从路径到 SD 扇区}

把这一节和前面几节串起来，读 \code{/bin/ls} 的第 0 字节要经过这些层：

\begin{txtcode}
namei("/bin/ls")             -> inode m（经 inode 1、inode 2 的目录内容）
readi(ip, ..., off=0, n)     -> bn = off / 1024 = 0
bmap(ip, 0)                  -> 块号 B = ip->addrs[0]
bread(ROOTDEV, B)            -> 块缓存未命中，调用 rootdev_rw(b, 0)
root_sector_lba(2B), (2B+1)  -> 扇区 208896 + 2B 和 208896 + 2B + 1
sdsector(lba, buf, 0)        -> SDHOST CMD17，PIO 读 512 字节
\end{txtcode}

一个 1\,KiB 的 xv6 块对应两个 512 字节扇区。写入的路径与此相同，只是中间多了日志：\code{writei()} 修改的块经 \code{log\_write()} 记入日志，\code{end\_op()} 提交时才真正经 \code{rootdev\_rw()} 写到 SD 卡。

\subsection{\code{mkfs} 怎样生成 \code{fs.img}}

\code{mkfs} 是一个在宿主机上运行的普通 C 程序，它和内核共用 \file{kernel/fs.h}，因此磁盘格式只有一份定义：

\begin{enumerate}
  \item 按 \code{FSSIZE}、\code{NINODES}、\code{LOGSIZE} 算出各区大小，把 32768 块全部写 0，再在第 1 块写入 superblock；
  \item \code{ialloc(T\_DIR)} 得到 inode 1，写入 \code{.} 和 \code{..}（都指向自己）；再建 inode 2 作为 \code{/bin}，在根目录里加 \code{bin} 项；
  \item 对命令行上的每个程序（\code{\_ls}、\code{\_cat}…去掉前缀下划线）分配一个 inode、在 \code{/bin} 里加目录项、用 \code{iappend()} 把文件内容逐块追加——它和内核的 \code{bmap()} 一样，前 12 块直接放、之后经间接块；\code{init} 额外在根目录里加一项；
  \item 把两个目录的大小取整到整块；最后 \code{balloc(freeblock)} 把已经用掉的前 \code{freeblock} 块在位图里标成 1。
\end{enumerate}

\code{mkfs} 只按顺序分配块和 inode，不需要日志，也不处理删除；所有整数都用 \code{xint()}/\code{xshort()} 写成小端序，与 AArch64 内核一致。

\subsection{在系统里更新根文件系统}

根文件系统的格式一旦改动（例如 \code{DIRSIZ} 从 14 改成 30），就必须整体重写 p2。最初只能把 SD 卡拔到电脑上 \code{dd}；后来内核增加了一个受保护的字符设备 \code{/dev/sdroot}，它只允许写启动时 \code{setup\_raw\_root()} 选中的那个分区，绝不暴露整张卡。更新流程是：

\begin{shcode}
/bin/tftpclient 192.168.0.195 fs.img          # 下载到 /boot（FAT32）
/bin/dd if=/boot/fs.img of=/dev/sdroot bs=32768
\end{shcode}

\code{dd} 在写之前先检查镜像大小恰好是 \code{FSSIZE}$\times$\code{BSIZE}、把整个镜像读进内存并校验 superblock 的魔数、大小和各区起点——因为一旦开始覆盖 p2，就再也不能从旧的根文件系统读任何东西，连 \code{dd} 自己的代码也不行（用户程序是静态链接的，已全部在内存里）。打开 \code{/dev/sdroot} 时内核冻结根文件系统：不再接受新的事务，等已有事务提交、读者退出；然后每写一个扇区就读回比较一次。结束时打印：

\begin{txtcode}
sdroot: verified 32/32 MiB
sdroot: update complete and verified; reboot now
dd: 33554432 bytes installed and read-back verified
dd: root filesystem is frozen; power-cycle Raspberry Pi now
\end{txtcode}

冻结是单向的，无论成败都必须断电重启——内存里的块缓存和 inode 缓存描述的是旧磁盘，不能再拿来解释新磁盘。

\section{烧录 SD 卡}

\subsection{分区布局}

一张为 xv6 准备的 SD 卡分成四个区（表~\ref{tab:sdparts}）。

\begin{table}[htbp]
  \centering
  \caption{SD 卡分区（63.8\,GB 卡的实例）}
  \label{tab:sdparts}
  \small
  \begin{tabular}{@{}lllrl@{}}
    \toprule
    分区 & 名称 & 格式 & 起始 LBA & 用途 \\
    \midrule
    p1 & BOOTFS & FAT32 & 2048 & 固件、\code{config.txt}、内核、Wi-Fi 固件 \\
    p2 & XV6FS & 原始 & 208896 & xv6 根文件系统（400\,MiB，用到 32\,MiB） \\
    p3 & EXT2 & ext2 & 1030144 & Linux 文件交换 \\
    p5 & SPFS & 原始 & 123654144 & 预留给其他教学文件系统 \\
    \bottomrule
  \end{tabular}
\end{table}

\subsection{步骤}

以 macOS 为例（\code{disk4} 必须先用 \code{diskutil list} 确认，写错磁盘会毁掉数据）：

\begin{shcode}
# 1. 卸载整张卡并重新分区
diskutil unmountDisk force /dev/disk4
diskutil partitionDisk /dev/disk4 4 MBR \
  FAT32 BOOTFS 100MiB  UFSD_EXTFS XV6FS 400MiB \
  UFSD_EXTFS EXT2 R    UFSD_EXTFS SPFS 500MiB

# 2. 写入 xv6 根文件系统（只能写分区 2，绝不能写整卡 /dev/rdisk4）
make fs.img
make install-rpi3-rawfs RPI3_XV6_DEV=/dev/rdisk4s2

# 3. 把分区 3 格式化成不带 ext4 特性的 ext2
sudo mke2fs -F -t ext2 -L EXT2 /dev/rdisk4s3

# 4. 在 BOOTFS 中放入树莓派官方固件（bootcode.bin、start.elf、fixup.dat 等），
#    然后安装内核、config.txt 和 Wi-Fi 固件
make install-rpi3 RPI3_XV6_DEV=/dev/rdisk4s2
\end{shcode}

\code{Makefile} 对 \code{RPI3\_XV6\_DEV} 做了检查：变量为空时只更新内核、配置和固件，并给出警告；不是 \code{/dev/diskNsM} 形式的分区设备则拒绝执行。这些防护是吃过亏之后加上的——\code{dd} 写错一个字母就会覆盖整张卡的分区表。

\subsection{日常更新}

改了内核之后不必每次拔卡：运行 \code{make}，内核会被拷到宿主机的 TFTP 目录；在树莓派上执行

\begin{shcode}
tftp 192.168.0.195 kernel8-xv6_wifi.img
\end{shcode}

文件就会下载到 \code{/boot}（FAT32 启动分区），重启后生效。只有用户程序或文件系统格式变化时，才需要重新烧录分区 2。

\chapter{驱动架构}
\label{chap:drivers}

\begin{introduction}[本章内容]
  \item Linux 2.6 风格的总线/设备/驱动模型
  \item 中断上半部与 workqueue 下半部
  \item 控制台、Mini UART 与 TTY
  \item USB 与 HID 键盘：描述符、枚举、split 事务、URB、真机排错
  \item I2C（BSC）
  \item 网络协议栈
  \item SDIO 无线网卡 BCM43430/43455
  \item USB 无线网卡 MT7601U
  \item CSI-2 摄像头 OV5647
\end{introduction}

\section{设备模型}

\subsection{为什么不再只用 \code{devsw[]}}

原版 xv6 只有一张字符设备函数表：\code{inode.major} 直接索引 \code{devsw[major].read/write}。它够用来驱动一个控制台，但表达不了“硬件设备”“驱动程序”“总线”以及“二者何时绑定”。随着驱动从一个变成十几个，项目引入了 Linux 2.6 设备模型的最小子集（\file{kernel/device.h}）：

\begin{ccode}
struct bus_type {
  char *name;
  int  (*match)(struct device *dev, struct device_driver *drv);
  int  (*probe)(struct device *dev);
  void (*remove)(struct device *dev);
};
struct device_driver {
  char *name;  struct bus_type *bus;
  int  (*probe)(struct device *dev);
  void (*remove)(struct device *dev);
};
struct device {
  char *name;  int id;  struct device *parent;  struct bus_type *bus;
  struct resource resource[DEVICE_RES_MAX];  int nresource;  // MMIO、IRQ
  struct device_driver *driver;  void *driver_data;
};
struct file_operations {
  int  (*read)(int user_dst, uint64 dst, int n);
  int  (*write)(int user_src, uint64 src, int n);
  int  (*open)(struct file *f);                 // 可选：每次 open 的私有状态
  void (*release)(struct file *f);
  int  (*fread)(struct file *f, int user_dst, uint64 dst, int n);
};
\end{ccode}

设备和驱动可以按任意顺序注册：先注册的一方等着，后注册的一方触发总线扫描，\code{match()} 成功就调用驱动的 \code{probe()}。驱动在 \code{probe()} 中初始化硬件，然后用 \code{register\_chrdev(major, name, fops)} 注册字符设备，或者 \code{register\_netdev()} 注册网卡。

\begin{figure}[htbp]
  \centering
  \begin{tikzpicture}[node distance=6mm and 10mm]
    \node[box] (dev) {\code{device\_register()}\\板级代码 / 枚举};
    \node[hbox, right=of dev, minimum width=26mm] (bus) {\code{bus\_type}\\\code{match()}};
    \node[box, right=of bus] (drv) {\code{driver\_register()}\\驱动模块};
    \node[box, below=of bus] (probe) {\code{driver->probe(dev)}};
    \node[box, below left=6mm and -6mm of probe] (chr) {\code{register\_chrdev()}\\\code{/dev/*}};
    \node[box, below right=6mm and -6mm of probe] (net) {\code{register\_netdev()}\\\code{wlan0}};
    \draw[arr] (dev) -- (bus); \draw[arr] (drv) -- (bus);
    \draw[arr] (bus) -- (probe); \draw[arr] (probe) -- (chr); \draw[arr] (probe) -- (net);
  \end{tikzpicture}
  \caption{设备与驱动的绑定}
  \label{fig:devmodel}
\end{figure}

项目中有三条总线：

\begin{itemize}
  \item \textbf{platform}：SoC 内固定的设备（Mini UART、SDHOST、Arasan、DWC2），无法枚举，由板级代码静态描述，按名字匹配；
  \item \textbf{usb}：DWC2 枚举出的 USB 设备，按 VID/PID 或接口类别匹配；
  \item \textbf{sdio}：Arasan 枚举出的 SDIO function，按 vendor/device/function 匹配。
\end{itemize}

字符设备的主设备号见表~\ref{tab:majors}。xv6 没有 devtmpfs/udev，设备节点由 \code{init} 用 \code{mknod} 创建；\code{cat /proc/devices} 可以看到全部总线、设备、资源和已绑定的驱动。

\begin{table}[htbp]
  \centering
  \caption{字符设备主设备号}
  \label{tab:majors}
  \small
  \begin{tabular}{@{}cll@{}}
    \toprule
    major & 名称 & 节点 \\
    \midrule
    1 & console & \code{/dev/console} \\
    2 & tty（当前进程的控制终端） & \code{/dev/tty} \\
    3 & ttyS0（Mini UART） & \code{/dev/ttyS0} \\
    4 & input（evdev） & \code{/dev/input/event0} \\
    5 & camera & \code{/dev/video0} \\
    \bottomrule
  \end{tabular}
\end{table}

\section{上半部与下半部}

中断处理程序运行时，本 CPU 的中断是关着的，而且它打断的可能是任何代码，所以不能睡眠、不能长时间运行。真正耗时的工作——解析 USB 报告、执行几十次 SDIO 命令、处理一个 TCP 段——要推迟到可以调度、可以睡眠的上下文中去做。项目实现了与 Linux 同名的 workqueue（\file{kernel/workqueue.h}）：

\begin{ccode}
struct work_struct {
  void (*func)(struct work_struct *work);
  struct work_struct *next;
  int pending, running, rerun;
  struct workqueue *wq;
};
int  queue_work(struct workqueue *wq, struct work_struct *work);
int  queue_delayed_work(struct workqueue *wq, struct delayed_work *dw, uint64 delay);
void cancel_work_sync(struct work_struct *work);
\end{ccode}

每个 workqueue 有自己的内核线程（\code{kworker}），在 \code{ps} 中可以看到。几条原则贯穿所有驱动：

\begin{itemize}
  \item 上半部只做三件事：确认中断来源、屏蔽或清除该中断、\code{queue\_work()}；
  \item 每个设备使用\emph{私有}的 work 和 workqueue，慢设备（如 SDIO 的 CMD53）不会阻塞别的下半部；
  \item 同一个 work 在执行期间再次被排队，只记一个 \code{rerun} 标志，不会无限堆积；
  \item 移除设备时先 \code{cancel\_work\_sync()}，等正在运行的 work 返回后才释放设备内存。
\end{itemize}

定时器中断每 10\,ms 推进一次 delayed-work 时钟，TCP 重传、Wi-Fi 连接超时等都用 delayed work 实现，而不是在定时器中断里直接运行协议代码。网络接收还在 workqueue 之上实现了 NAPI 式的预算轮询：中断只调度一次，下半部每轮最多处理 16 个报文，没处理完就重新排队，处理完才重新打开中断。

\section{控制台、UART 与 TTY}

\subsection{Mini UART}

Mini UART 驱动（\file{kernel/uart.c}）在 \code{probe()} 中完成与早期汇编相同的初始化，再打开接收中断：

\begin{ccode}
REGS_AUX->enables |= 1U;               // 使能 AUX 中的 Mini UART
REGS_AUX->mu_control = 0;              // 暂停收发
REGS_AUX->mu_ier = 0;
REGS_AUX->mu_lcr = 3;                  // 8 位数据
REGS_AUX->mu_iir = 0xc6;               // 清空 FIFO
REGS_AUX->mu_baud_rate = 270;          // 115200 @ core_freq 250 MHz
REGS_AUX->mu_ier = MU_IER_RX_ENABLE;   // 接收中断
REGS_AUX->mu_control = 3;              // 打开收发
\end{ccode}

发送使用一个 32 字节的软件环形缓冲区：\code{uartputc()} 把字符放进缓冲区，缓冲区满时睡眠；\code{uartstart()} 在发送 FIFO 有空位时搬字符过去，并打开“发送空”中断；\code{uartintr()} 读空接收 FIFO、把字符交给 \code{consoleintr()}，再继续发送。缓冲区空了就关掉发送中断，否则会一直触发。内核 \code{printf} 在 panic 等场合使用同步的 \code{uartputc\_sync()}，不依赖中断。

\subsection{控制台与行规程}

\code{consoleintr()} 实现了最基本的行规程：回显、退格、整行缓冲、\code{Ctrl-D} 结束输入、\code{Ctrl-P} 打印进程表、\code{Ctrl-C} 向前台进程组发送 SIGINT。读进程在 \code{consoleread()} 里睡眠，直到收到一整行。

\subsection{TTY 层}

\code{/dev/console}、\code{/dev/ttyS0}、\code{/dev/tty} 共用 Mini UART 这一个硬件后端，但含义不同：

\begin{itemize}
  \item \code{/dev/console} 是系统控制台，内核日志直接写 UART，不经过文件系统；
  \item \code{/dev/ttyS0} 是具体的串口终端，\code{init} 在上面启动 \code{login}；
  \item \code{/dev/tty} 是“当前进程的控制终端”，按调用者的 \code{ctty} 动态转发：串口会话转到 \code{ttyS0}，SSH 会话转到对应的 PTY 从端。
\end{itemize}

\code{login} 调用 \code{tty\_attach()} 成为新会话和进程组的组长，并把 \code{ttyS0} 设为控制终端；之后 fork 出的 shell 和命令继承会话、进程组和控制终端。shell 运行外部命令时把它设为前台进程组，\code{Ctrl-C} 产生的 SIGINT 只发给这个组，所以能终止 \code{ping}、\code{tftp} 或整条管道，而不会杀掉 shell 和 \code{login}。

一个由此暴露的问题是：旧版 \code{ps} 通过系统调用让内核 \code{printf} 直接打印进程表，结果在 SSH 会话里执行 \code{ps}，输出却出现在串口上。内核 \code{printf} 写的是系统控制台，不看调用者的 fd 1。现在 \code{sys\_ps()} 把快照格式化到用户缓冲区，由 \code{ps} 自己 \code{write(1, ...)}。

\section{USB 与 HID 键盘}
\label{sec:usb}

USB 键盘是本项目里协议层次最多的驱动：从主机控制器的寄存器，到 Hub、事务转换器、描述符、HID 报告，再到输入子系统和控制台。本节沿着“插上键盘—识别—配置—读按键”的顺序，把每一层讲清楚。

\subsection{三层软件与硬件路径}

USB 驱动分三层：

\begin{itemize}
  \item \textbf{HCD}（\file{kernel/dwc2.c}）：复位 DWC2 控制器、强制主机模式、开 DMA；给根端口和 Hub 端口上电、复位；用 host channel 发起 control、bulk、interrupt 传输；
  \item \textbf{USB core}（\file{kernel/usb.c}）：\code{struct usb\_device}、按 ID 或接口类别匹配驱动、Linux 风格的 URB 接口；
  \item \textbf{类驱动}：HID 键盘（\file{usbkbd.c}）、CDC-ECM 网卡（\file{usbnet.c}，用于 QEMU）、MT7601U 无线网卡（\file{mt7601u.c}）。
\end{itemize}

\begin{figure}[htbp]
  \centering
  \begin{tikzpicture}[x=1mm,y=1mm, font=\footnotesize,
      data/.style={{Stealth[length=2.4mm]}-{Stealth[length=2.4mm]}, line width=1.6pt, draw=orange!85!black},
      ctl/.style={-{Stealth[length=2mm]}, thick, draw=structurecolor!70!black},
      evt/.style={-{Stealth[length=2mm]}, thick, dashed, draw=red!70!black},
      blk/.style={box, font=\footnotesize, minimum height=8mm}]
    \draw[rounded corners=3pt, draw=gray!60, fill=gray!4] (-2,-2) rectangle (70,74);
    \node[lab, anchor=north west] at (-1,73.5) {BCM2837 SoC};
    \draw[rounded corners=3pt, draw=gray!60, fill=gray!4] (76,20) rectangle (112,74);
    \node[lab, anchor=north west] at (77,73.5) {LAN9514（板载）};
    % SoC
    \node[hbox, minimum width=22mm, minimum height=36mm, text width=20mm] (cpu) at (11,52) {ARM\\Cortex-A53\\[2pt]\file{usbkbd.c}\\\file{usb.c}\\\file{dwc2.c}};
    \node[gbox, minimum width=30mm, minimum height=9mm, text width=30mm] (ram) at (48,6) {SDRAM\\8 字节报告缓冲};
    \node[hbox, text width=30mm, minimum height=44mm] (dwc) at (48,48) {};
    \node[font=\small\bfseries, anchor=north] at (48,69.5) {DWC2 主机控制器};
    \node[lab, anchor=north] at (48,65) {\code{0x3f980000}，IRQ 9};
    \node[box, font=\scriptsize, text width=26mm, minimum height=6mm] (c0) at (48,58) {ch0：control（枚举）};
    \node[box, font=\scriptsize, text width=26mm, minimum height=6mm, fill=orange!15] (c6) at (48,50) {ch6：interrupt-IN（键盘）};
    \node[box, font=\scriptsize, text width=26mm, minimum height=6mm] (c4) at (48,42) {ch4：bulk-IN（MT7601U）};
    \node[lab, align=center] at (48,32.5) {每个 channel：\code{HCCHAR HCSPLT}\\\code{HCTSIZ HCDMA HCINT}};
    \node[blk, text width=20mm] (ic) at (11,14) {中断控制器\\IRQ 9};
    % hub
    \node[hbox, text width=30mm, minimum height=14mm] (hub) at (94,60) {根 Hub（高速）\\\code{0424:2514}};
    \node[box, text width=30mm, minimum height=9mm, fill=orange!10] (tt) at (94,44) {事务转换器 TT\\高速 $\leftrightarrow$ 全速/低速};
    \node[box, text width=30mm, minimum height=9mm] (h2) at (94,28) {port 1：子 Hub\\→ 以太网 \code{0424:7800}};
    % keyboard
    \node[hbox, text width=26mm, minimum height=12mm] (kb) at (132,44) {USB 键盘接收器\\\code{2a7a:8a47}\\低速 1.5\,Mb/s};
    \node[lab] at (132,33) {hub port 3};
    % edges
    \draw[ctl] (cpu.east |- c0) -- node[lab, above]{MMIO} (c0.west -| 33,0);
    \draw[data] (dwc.east |- hub) -- node[lab, above, text=orange!50!black]{480 Mb/s} (hub.west);
    \draw[data] (hub) -- (tt);
    \draw[data] (tt.east) -- node[lab, below, text=orange!50!black]{1.5 Mb/s} (kb.west);
    \draw[gray!60, thick] (hub.west) ++(0,-3) -- ++(-1.5,0) |- (h2.west);
    \draw[data] (dwc.south) -- node[lab, right, text=orange!50!black]{DMA：\code{0xC0000000|PA}} (ram.north);
    \draw[evt] ([yshift=3mm]dwc.south west) -- node[lab, above]{HAINT} ++(-6,0) |- (ic.east);
    \draw[evt] (ic.north) -- (cpu.south);
  \end{tikzpicture}
  \caption{USB 键盘的硬件路径。DWC2 用 host channel 发起事务、用 DMA 读写内存；树莓派 3 的 USB 口都在板载 LAN9514 Hub 后面，低速键盘与 Hub 之间靠事务转换器（TT）完成速率转换}
  \label{fig:usb-hw}
\end{figure}

图~\ref{fig:usb-hw} 是硬件路径。有几个特点和前面的 Wi-Fi、摄像头都不一样：

\paragraph{主机发起一切.} USB 是严格的主从结构，设备永远不能主动说话。即使是名为“interrupt”的传输类型，也只是“主机承诺每隔 \code{bInterval} 毫秒来问一次”：主机发 IN 令牌，设备有数据就回 DATA，没有就回 NAK。所谓“键盘中断”，其实是主机控制器替 CPU 轮询，问到数据后才用 IRQ 9 通知 CPU。

\paragraph{host channel 与 DMA.} DWC2 有 8 个 host channel，每个 channel 是一组寄存器，描述“对哪个设备的哪个端点做一次什么传输”：\reg{HCCHAR}（设备地址、端点号、方向、传输类型、最大包长、是否低速）、\reg{HCSPLT}（split 事务参数）、\reg{HCTSIZ}（字节数、包数、DATA0/1）、\reg{HCDMA}（缓冲区的总线地址）、\reg{HCINT}（完成、ACK、NAK、NYET、错误等状态）。本项目固定分工：ch0 做枚举时的 control 传输，ch6 给键盘的 interrupt-IN，ch4 给 MT7601U 的 bulk-IN。和 Arasan 不同，DWC2 自己用 DMA 读写内存，所以缓冲区在提交前要 clean、完成后要 invalidate 缓存，地址也要换成 \code{0xC0000000|PA} 的总线地址。

\paragraph{一切都在 Hub 后面.} 树莓派 3 的 SoC 只有一个 USB 口，直接接板载 LAN9514，它内部是一个高速 Hub 加一块 USB 以太网卡（在 3B+ 上还有第二层 Hub）。外面插的键盘一般是低速（1.5\,Mb/s）设备，而 DWC2 和 Hub 之间跑的是高速（480\,Mb/s）。两种速率之间靠 Hub 里的\emph{事务转换器}（TT）衔接，这就是后面 split transaction 的由来。

\subsection{描述符：设备、配置、接口与端点}

\begin{figure}[htbp]
  \centering
  \begin{tikzpicture}[x=1mm,y=1mm, font=\scriptsize,
      d/.style={draw=structurecolor!70!black, rounded corners=1.5pt, fill=structurecolor!8, align=left, anchor=north west, inner sep=2pt, text width=#1},
      d/.default=50mm,
      sel/.style={fill=orange!15, draw=orange!70!black},
      tree/.style={gray!60!black, semithick}]
    \node[d=50mm] (dev) at (0,84) {\textbf{设备描述符}（18 B，类型 1）\\{}
      [4] \code{bDeviceClass}＝0（看接口）\\{}
      [7] \code{bMaxPacketSize0}＝8\\{}
      [8] \code{idVendor}＝\code{2a7a}\\{}[10] \code{idProduct}＝\code{8a47}};
    \node[d=52mm] (cfg) at (10,60) {\textbf{配置描述符}（9 B，类型 2）\\{}
      [2] \code{wTotalLength}：整串描述符的长度\\{}
      [4] \code{bNumInterfaces}＝2\\{}[5] \code{bConfigurationValue}＝1};
    \node[d=46mm, sel] (if0) at (20,39) {\textbf{接口 0}（9 B，类型 4）\\{}
      [5] class＝3 HID\\{}[6] subclass＝1 Boot\\{}
      [7] protocol＝1 键盘};
    \node[d=46mm] (hid0) at (30,20) {HID 描述符（类型 \code{0x21}）\\报告描述符长度等};
    \node[d=50mm, sel] (ep0) at (30,9) {\textbf{端点}（7 B，类型 5）\\{}
      [2] \code{0x81}：IN，端点 1\\{}
      [3] 属性＝3 interrupt\quad [4] MPS＝8\\{}[6] \code{bInterval}＝10（ms）};
    \node[d=44mm] (if1) at (88,39) {\textbf{接口 1}（9 B）\\class＝3，protocol＝2（鼠标）\\或 0（多媒体键等）};
    \node[d=40mm] (ep1) at (98,20) {端点 \code{0x82}：interrupt-IN};
    \draw[tree] (dev.south west) ++(3,0) |- (cfg.west);
    \draw[tree] (cfg.south west) ++(3,0) |- (if0.west);
    \draw[tree] (cfg.east) -| (if1.north);
    \draw[tree] (if0.south west) ++(3,0) |- (hid0.west);
    \draw[tree] (if0.south west) ++(3,0) |- (ep0.west);
    \draw[tree] (if1.south west) ++(3,0) |- (ep1.west);
    % byte stream
    \node[lab, anchor=west] at (0,-14) {\code{GET\_DESCRIPTOR(Configuration, wTotalLength)} 一次返回的字节串，逐项按 \code{bLength} 前进：};
    \foreach \x/\w/\t/\c in {0/14/配置 9/structurecolor!10, 14/14/接口0 9/orange!20, 28/14/HID 9/structurecolor!5, 42/12/端点 7/orange!20, 54/14/接口1 9/gray!10, 68/14/HID 9/gray!5, 82/12/端点 7/gray!10}{
      \draw[fill=\c, draw=gray!60] (\x,-23) rectangle ++(\w,6) node[midway]{\t};
    }
    \node[lab, anchor=west, align=left] at (96,-20) {只在接口 0 范围内\\挑 interrupt-IN 端点};
  \end{tikzpicture}
  \caption{一个 USB 键盘接收器的描述符层次（示意，方括号里是字节偏移；VID/PID 取自真机日志）。驱动按类别匹配，而不是按 VID/PID：先找 class/subclass/protocol＝3/1/1 的 Boot Keyboard 接口，再只在这个接口里找 interrupt-IN 端点}
  \label{fig:usb-desc}
\end{figure}

USB 设备用一组\emph{描述符}介绍自己（图~\ref{fig:usb-desc}）。每个描述符的第 0 字节是长度 \code{bLength}，第 1 字节是类型。层次是：一个设备有若干\emph{配置}（几乎总是 1 个）；一个配置包含若干\emph{接口}，每个接口是一项独立的功能，有自己的 class/subclass/protocol；一个接口有若干\emph{端点}，端点是真正收发数据的地方，有编号、方向、传输类型、最大包长和轮询间隔。端点 0 是所有设备都有的默认控制端点，枚举全靠它。

很多键盘接收器是复合设备：接口 0 是键盘，接口 1 是鼠标或多媒体键。所以驱动不能看设备描述符的 class（它通常是 0，意思是“看各个接口”），也不能凭 VID/PID 白名单，而要遍历配置描述符：

\begin{ccode}
// GET_DESCRIPTOR(Configuration) 先取 9 字节拿到 wTotalLength，再一次取全部
for(int off = 0; off + 9 <= total && ctrl_buf[off] >= 2; off += ctrl_buf[off]){
  if(ctrl_buf[off + 1] == 4 &&                   // 接口描述符
     ctrl_buf[off + 5] == 3 &&                   // class: HID
     ctrl_buf[off + 6] == 1 &&                   // subclass: Boot Interface
     ctrl_buf[off + 7] == 1){                    // protocol: Keyboard
    sel->interface_number = ctrl_buf[off + 2];
    return 1;
  }
}
\end{ccode}

\code{off += ctrl\_buf[off]} 按每个描述符自己的长度前进，所以中间夹着的 HID 描述符（类型 \code{0x21}）等不认识的类型会被自然跳过；\code{ctrl\_buf[off] >= 2} 防止一个长度为 0 的坏描述符造成死循环。找到键盘接口后，后面的端点只在这个接口的范围里挑（\code{current\_interface} 限定），否则复合设备里鼠标的端点会把键盘的覆盖掉。

\subsection{枚举：control 传输与请求序列}

端点 0 上的 control 传输分三个阶段：SETUP 阶段主机发出 8 字节的请求（\code{bmRequestType}、\code{bRequest}、\code{wValue}、\code{wIndex}、\code{wLength}），可选的 DATA 阶段按最大包长收发数据，STATUS 阶段反方向发一个空包确认。\code{control()} 就是在 ch0 上依次执行这三个阶段。表~\ref{tab:usb-enum} 是识别一个 Hub 后面的键盘所发出的请求。

\begin{table}[htbp]
  \centering
  \caption{从 Hub 端口到可用键盘的请求序列}
  \label{tab:usb-enum}
  \small
  \begin{tabularx}{\linewidth}{@{}lllX@{}}
    \toprule
    对象 & 请求 & \code{bmRequestType} & 作用 \\
    \midrule
    Hub & \code{SET\_FEATURE(PORT\_POWER)} & \code{0x23} & 给下游端口上电，之后最多等 500\,ms 连接去抖 \\
    Hub & \code{GET\_STATUS} & \code{0xa3} & 端口状态：bit 0 连接、bit 1 使能、bit 8 供电、bit 9 低速、bit 10 高速 \\
    Hub & \code{SET\_FEATURE(PORT\_RESET)} & \code{0x23} & 复位端口，设备进入默认状态、地址为 0 \\
    设备 0 & \code{GET\_DESCRIPTOR(Device, 8)} & \code{0x80} & 只取 8 字节，得到端点 0 的最大包长 \\
    设备 0 & \code{GET\_DESCRIPTOR(Device, 18)} & \code{0x80} & 完整设备描述符：class、VID、PID \\
    设备 0 & \code{SET\_ADDRESS(n)} & \code{0x00} & 分配唯一地址；把父 Hub、端口、速度记入 \code{usb\_route[n]} \\
    设备 $n$ & \code{GET\_DESCRIPTOR(Config, 9)} & \code{0x80} & 得到 \code{wTotalLength} \\
    设备 $n$ & \code{GET\_DESCRIPTOR(Config, total)} & \code{0x80} & 全部接口和端点，挑出 Boot Keyboard \\
    设备 $n$ & \code{SET\_CONFIGURATION} & \code{0x00} & 激活配置，端点开始可用 \\
    接口 & \code{SET\_PROTOCOL(boot)} & \code{0x21} & 让键盘用固定的 8 字节 Boot 报告格式 \\
    接口 & \code{SET\_IDLE(0)} & \code{0x21} & 只在按键状态变化时才回报告 \\
    \bottomrule
  \end{tabularx}
\end{table}

先取 8 字节再取 18 字节是 USB 枚举的惯例：在知道端点 0 的最大包长之前，只能假定它是最小的 8 字节，而设备描述符第 7 字节正好告诉我们真实值。端口状态的例子：真机日志里 Hub port 3 是 \code{0x301}（供电、低速、已连接），复位后变成 \code{0x303}（多了“使能”）。根 Hub 和嵌套的 Hub 用同一套请求递归扫描，深度限制为两层，防止一个异常的拓扑让启动无限递归。

\code{SET\_PROTOCOL(boot)} 很关键。HID 设备本来用“报告描述符”描述自己的报告格式，解析它需要一个不小的解释器；而 Boot 协议是 BIOS 时代为了不解析报告描述符而规定的固定格式，所有键盘都必须支持。用它，驱动只需理解图~\ref{fig:hid-report} 这 8 个字节。

\begin{figure}[htbp]
  \centering
  \begin{tikzpicture}[x=1mm,y=1mm, font=\scriptsize]
    \foreach \i/\t/\c in {0/{修饰键}/orange!20, 1/{保留}/gray!10, 2/{键 1}/structurecolor!15, 3/{键 2}/structurecolor!15, 4/{键 3}/structurecolor!15, 5/{键 4}/structurecolor!15, 6/{键 5}/structurecolor!15, 7/{键 6}/structurecolor!15}{
      \draw[fill=\c, draw=gray!60] (\i*16,10) rectangle ++(16,8) node[midway, align=center]{\t\\{}[\i]};
    }
    \foreach \b/\t in {0/{LCtrl},1/{LShift},2/{LAlt},3/{LGUI},4/{RCtrl},5/{RShift},6/{RAlt},7/{RGUI}}{
      \draw[fill=orange!8, draw=gray!50] (\b*16,0) rectangle ++(16,6) node[midway]{bit\,\b\ \t};
    }
    \draw[gray!60, dashed] (0,10) -- (0,6); \draw[gray!60, dashed] (16,10) -- (128,6);
    \node[lab, anchor=west, align=left] at (0,24) {例：\code{mod=0 keys=10,0,0,0,0,0} 表示只按着 HID usage \code{0x10}，即字母 \code{m}；\\同时按 Shift 时 byte 0 = \code{0x02}；按键超过 6 个时 6 个槽全是 \code{0x01}（rollover），驱动保持上一份状态。};
  \end{tikzpicture}
  \caption{HID Boot 键盘报告（8 字节）。byte 0 是 8 个修饰键的位图，byte 2～7 是当前按着的最多 6 个键的 usage 编号}
  \label{fig:hid-report}
\end{figure}

最后，\code{usb\_device\_register()} 把设备交给 USB 总线，总线按 class/subclass/protocol 匹配到 \code{usbhid-keyboard} 驱动，调用 \code{usbkbd\_probe()}：发送上面两个 HID 请求，分配每设备的 \code{usbkbd\_state} 和 URB，注册 \code{input\_dev}，然后第一次提交 URB。

\subsection{Split transaction：高速 Hub 后面的低速设备}

\begin{figure}[htbp]
  \centering
  \begin{tikzpicture}[x=1mm,y=1mm, font=\scriptsize,
      msg/.style={-{Stealth[length=1.8mm]}, semithick},
      hs/.style={msg, draw=structurecolor!80!black},
      ls/.style={msg, draw=orange!85!black},
      ev/.style={msg, dashed, draw=red!70!black},
      t/.style={font=\scriptsize, inner sep=0.5pt}]
    \foreach \i in {0,...,7}{
      \draw[gray!25] (20+\i*16,4) -- (20+\i*16,56);
      \node[lab] at (28+\i*16,58.5) {\i};
    }
    \draw[gray!25] (148,4) -- (148,56);
    \node[lab, anchor=east] at (19,58.5) {微帧};
    \node[lab, anchor=west] at (20,63) {一个 1\,ms 帧 = 8 个 125\,$\mu$s 微帧（\code{HFNUM} 低 3 位）；SSPLIT 发在微帧 0};
    \foreach \y/\l in {50/{DWC2 ch6}, 36/{LAN9514 TT}, 24/{低速键盘}, 10/{CPU}}{
      \node[lab, anchor=east] at (19,\y) {\l};
      \draw[gray!55] (20,\y) -- (148,\y);
    }
    % uframe 0
    \draw[hs] (22,50) -- (25,36); \node[t, text=structurecolor!80!black] at (23,52.5) {SSPLIT};
    \draw[hs] (28,36) -- (31,50); \node[t] at (32,43) {ACK};
    \draw[ev] (33,50) -- (33,10); \node[t, text=red!60!black, align=center, anchor=north] at (33,8) {\ding{172}\,记录 ACK\\所在微帧};
    % low speed transaction
    \fill[orange!12] (30,25) rectangle (64,35);
    \draw[ls] (31,36) -- (34,24); \draw[ls] (60,24) -- (63,36);
    \node[t, text=orange!50!black, align=center] at (47,30) {TT 在 1.5\,Mb/s 总线上\\替主机完成 IN / DATA / ACK};
    % uframe 2
    \draw[hs] (54,50) -- (56,36); \node[t, text=structurecolor!80!black] at (54,52.5) {CSPLIT};
    \draw[hs] (58,36) -- (60,50); \node[t] at (63,43) {NYET};
    \draw[ev] (61,50) -- (61,10); \node[t, text=red!60!black, align=center, anchor=north] at (61,8) {\ding{173}\,窗口内\\下一微帧重试};
    % uframe 3
    \draw[hs] (70,50) -- (72,36); \node[t, text=structurecolor!80!black] at (70,52.5) {CSPLIT};
    \draw[hs] (74,36) -- (76,50); \node[t] at (80,43) {DATA 8\,B};
    \draw[ev] (78,50) -- (78,10); \node[t, text=red!60!black, align=center, anchor=north] at (84,8) {\ding{174}\,\code{XFERCOMPL}\\交给 completion work};
    \node[draw=gray!60, fill=gray!5, rounded corners=1pt, align=left, anchor=north west, text width=50mm, font=\scriptsize] at (96,48) {键盘没有新按键时回 NAK；或者整个帧里一直 NYET。这两种情况本轮都结束，等 \code{bInterval}（10\,ms）后从新的 SSPLIT 开始。};
  \end{tikzpicture}
  \caption{低速键盘在高速 Hub 后面的一次 interrupt-IN 轮询。主机先向 TT 发 start-split，TT 在低速总线上替主机完成 IN 事务，主机再在随后的微帧里发 complete-split 取回结果。\ding{172}、\ding{173} 两种中间事件必须在中断里当场处理，晚一个帧就错过了窗口}
  \label{fig:usb-split}
\end{figure}

DWC2 不能直接对低速设备说话，它要先把事务交给 Hub 里的 TT（图~\ref{fig:usb-split}）：\emph{start-split}（SSPLIT）告诉 TT“请替我对 port 3 上地址 4 的端点 1 做一次 IN”，TT 回 ACK 表示接受；TT 随后在低速总线上完成真正的事务；主机在随后的微帧里发 \emph{complete-split}（CSPLIT）取结果，TT 要么回 NYET（还没完成），要么回数据，要么转达设备的 NAK。channel 寄存器里，这些信息体现为：

\begin{table}[htbp]
  \centering
  \caption{真机日志中的寄存器样本解码}
  \label{tab:usb-regs}
  \small
  \begin{tabularx}{\linewidth}{@{}lX@{}}
    \toprule
    样本 & 含义 \\
    \midrule
    \reg{HCCHAR}$=$\code{0x010e8808} & 设备地址 4、端点 1、IN、interrupt 类型、低速（\code{LSPDDEV}）、最大包长 8 \\
    \reg{HCSPLT}$=$\code{0x8001c083} & \code{SPLTENA}、\code{COMPSPLT}（当前是 CSPLIT）、XACTPOS$=$ALL、Hub 地址 1、端口 3 \\
    \reg{HCINT}$=$\code{0x22} & \code{CHHLTD | ACK}：SSPLIT 已被 TT 接受 \\
    \reg{HCINT}$=$\code{0x42} & \code{CHHLTD | NYET}：CSPLIT 暂时没有结果 \\
    \reg{HCINT}$=$\code{0x03} & \code{CHHLTD | XFERCOMPL}：DMA 缓冲里已经是完整的报告 \\
    \bottomrule
  \end{tabularx}
\end{table}

这里有三条规则，每一条都是在真机上踩过才写进代码的：低速设备的每一个事务都要在 \reg{HCCHAR} 里置 \code{LSPDDEV}；CSPLIT 收到 NAK 说明这次事务已经结束（设备没数据），必须从新的 SSPLIT 重来，而不是继续发 CSPLIT；CSPLIT 收到 NYET 只能在\emph{同一个帧}的窗口里重试，超出窗口就放弃本轮、等 \code{bInterval} 后重新开始。

\subsection{URB 与中断驱动的接收}

USB core 提供与 Linux 同名的异步接口：

\begin{ccode}
struct urb *usb_alloc_urb(void);
void usb_fill_int_urb(struct urb *urb, struct usb_device *dev, int ep, void *buf,
                      int len, void (*complete)(struct urb *), void *context, int interval);
int  usb_submit_urb(struct urb *urb);       // 提交给 HCD，立即返回
void usb_kill_urb(struct urb *urb);         // 同步取消
\end{ccode}

URB 的所有权在提交时交给 HCD，完成时由 \code{usb\_hcd\_giveback\_urb()} 交还类驱动并调用 \code{complete} 回调，回调通过 \code{urb->context} 找回所属的键盘。\code{dwc2\_submit\_urb()} 把 ch6 配好就返回：

\begin{ccode}
wr(HCINTMSK(6), HCINT_XFERCOMPL | HCINT_CHHLTD | HCINT_NAK |
                 HCINT_ACK | HCINT_NYET | HCINT_ERRORS);
if(usb_route[addr].hub)                        // 在 Hub 后面：split 事务
  split = HCSPLT_SPLTENA | HCSPLT_XACTPOS_ALL |
          HCSPLT_HUBADDR(usb_route[addr].hub) | HCSPLT_PRTADDR(usb_route[addr].port) |
          (req->complete_split ? HCSPLT_COMPSPLT : 0);
wr(HCSPLT(6), split);
wr(HCDMA(6),  DWC2_DMA_BUS(urb->transfer_buffer));
wr(HCTSIZ(6), HCTSIZ_XFERSIZE(8) | HCTSIZ_PKTCNT(1) | HCTSIZ_PID(pid));
wr(HAINTMSK, rd(HAINTMSK) | (1U << 6));
wr(HCCHAR(6), HCCHAR_DEVADDR(addr) | HCCHAR_EPNUM(ep) | HCCHAR_EPTYPE(EPTYPE_INTERRUPT) |
              HCCHAR_MPS(8) | HCCHAR_EPDIR_IN |
              (usb_route[addr].low_speed ? HCCHAR_LSPDDEV : 0) | HCCHAR_CHENA);
req->last_issue_frame = rd(HFNUM) & 0x3fff;    // 记下发出时的微帧号
\end{ccode}

\begin{figure}[htbp]
  \centering
  \begin{tikzpicture}[x=1mm,y=1mm, font=\scriptsize,
      fn/.style={box, font=\scriptsize, minimum height=8mm, inner sep=2pt, text width=#1},
      fn/.default=32mm,
      flow/.style={-{Stealth[length=2mm]}, line width=1.1pt, draw=structurecolor!80!black},
      loop/.style={-{Stealth[length=2mm]}, semithick, dashed, draw=orange!80!black}]
    \foreach \x/\w/\t in {0/36/{硬件}, 38/36/{硬中断}, 76/36/{HCD 完成 work}, 114/36/{\code{usbkbd\_wq}（进程上下文）}}{
      \fill[gray!5, rounded corners=2pt] (\x,0) rectangle ++(\w,62);
      \node[lab, anchor=north] at (\x+\w/2,61.5) {\t};
    }
    \node[fn] (h1) at (18,50) {键盘按下 \code{m}\\TT 取回 8 字节报告};
    \node[fn] (h2) at (18,32) {ch6 DMA 写入\\\code{kbd->report}};
    \node[fn] (h3) at (18,14) {\code{HCINT.XFERCOMPL}\\\code{HAINT} bit 6 → IRQ 9};
    \node[fn] (i1) at (56,50) {\code{dwc2\_irq()}};
    \node[fn] (i2) at (56,32) {快速路径不处理\\（不是 ACK/NYET）};
    \node[fn] (i3) at (56,14) {屏蔽 ch6\\\code{schedule\_work}};
    \node[fn] (w1) at (94,50) {\code{dwc2\_interrupt\_urb\_work}\\cache invalidate};
    \node[fn] (w2) at (94,32) {读 \code{HCTSIZ} 余量\\算出 \code{actual\_length}};
    \node[fn] (w3) at (94,14) {\code{usb\_hcd\_giveback\_urb}\\→ \code{usbkbd\_irq\_complete}};
    \node[fn] (k1) at (132,50) {\code{usbkbd\_rx\_work}\\与上一份报告比较};
    \node[fn] (k2) at (132,32) {\code{input\_report\_key}\\\code{input\_sync}\\\code{consoleintr('m')}};
    \node[fn] (k3) at (132,14) {\code{usb\_submit\_urb}\\重新 arm ch6};
    \draw[flow] (h1) -- (h2); \draw[flow] (h2) -- (h3);
    \draw[flow] (h3.east) -- ++(1,0) |- (i1.west);
    \draw[flow] (i1) -- (i2); \draw[flow] (i2) -- (i3);
    \draw[flow] (i3.east) -- ++(1,0) |- (w1.west);
    \draw[flow] (w1) -- (w2); \draw[flow] (w2) -- (w3);
    \draw[flow] (w3.east) -- ++(1,0) |- (k1.west);
    \draw[flow] (k1) -- (k2); \draw[flow] (k2) -- (k3);
    \draw[loop] (k3.south) -- ++(0,-8) -| node[lab, below, pos=0.25]{下一次 SSPLIT：回到硬件} (h3.south);
  \end{tikzpicture}
  \caption{一份键盘报告从硬件到 shell 的路径。每一列是一种执行上下文：硬中断只做与微帧时序有关的最少工作；HCD 的完成 work 负责缓存维护并把 URB 归还类驱动；HID 解析、输入事件和控制台都在键盘自己的 workqueue 里完成，最后重新提交 URB。键盘回 NAK（没有新按键）时，HCD 完成 work 等 \code{bInterval} 后自己重发，URB 不归还类驱动}
  \label{fig:usb-rx}
\end{figure}

之后的事全由中断推动（图~\ref{fig:usb-rx}）。\code{dwc2\_irq()} 先交给快速路径 \code{dwc2\_interrupt\_rx\_irq\_fast()}：SSPLIT 的 ACK 和窗口内的 NYET 在这里当场处理——记下当前微帧，忙等到“发出后第 2 个微帧”（最多 500\,$\mu$s），立即重发 CSPLIT——因为它们有微帧级的截止时间，交给任何可调度的线程都来不及。只有最终结果（完成、无数据或错误）才屏蔽 ch6 并 \code{schedule\_work}。HCD 的完成 work 做缓存维护、算实际长度并归还 URB；键盘的完成回调只 \code{queue\_work(kbd->rx\_work)}；\code{usbkbd\_rx\_work()} 在进程上下文里比较新旧报告、上报输入事件、调用 \code{consoleintr()}，最后重新 \code{usb\_submit\_urb()}。三种上下文各做各的：硬中断只管时序，HCD work 只管控制器，类驱动只管 HID。

拔出设备时顺序相反：\code{usbkbd\_remove()} 先标记断开，\code{usb\_kill\_urb()} 停止 ch6 并等正在运行的完成回调返回，\code{cancel\_work\_sync()} 等 \code{rx\_work} 返回，最后才释放 URB 和设备状态——保证不会有任何回调再碰到已释放的内存。

\subsection{HID 报告与输入子系统}

\code{usbkbd\_rx\_work()} 收到 8 字节报告后与上一份比较：修饰键位图里变化的位分别上报按下/松开；本次出现、上次没有的 usage 是新按下的键，查表转成字符；上次有、本次没有的是松开的键。每份报告最后 \code{input\_sync()}。按键多到键盘无法分辨时，6 个槽会全填错误码（rollover），这份报告不说明任何键的状态，驱动直接保留上一份，不能当成“全部松开”。

之后的事由输入子系统（\file{kernel/input.c}）完成，它把事件分发给两个“处理者”：

\begin{itemize}
  \item \textbf{evdev}：\code{/dev/input/event0}。每次 \code{open} 都分配一个独立的客户端环形缓冲区（\code{fops.open} 设置 \code{f->private\_data}），读者在自己的缓冲区上睡眠；溢出时插入 \code{SYN\_DROPPED}；
  \item \textbf{控制台}：把按键转成字符（处理 Shift、Caps Lock、Ctrl），送入 \code{consoleintr()}。所以 USB 键盘和串口一样可以登录、输入命令、按 \code{Ctrl-C}。
\end{itemize}

这正是 Linux 的分层：硬件驱动只关心“产生了什么事件”，不关心“谁在读”。

\subsection{真机排错：从 STALL 到第一份报告}
\label{sec:usb-debug}

这个驱动在 QEMU 上很快就能工作，到了真机（低速无线键盘接收器 \code{2a7a:8a47} 插在 LAN9514 的 port 3）却经历了九个阶段才读到第一个按键。表~\ref{tab:usb-bugs} 按时间顺序列出每一步的现象、根因和修复；每一步只把串口日志实际证明的东西当作结论——ACK 不等于下游数据已经返回，“ready”也不等于收到了按键。

\begin{table}[htbp]
  \centering
  \caption{USB 键盘真机排错的九个阶段}
  \label{tab:usb-bugs}
  \footnotesize
  \begin{tabularx}{\linewidth}{@{}p{2mm}>{\raggedright\arraybackslash}p{36mm}>{\raggedright\arraybackslash}X>{\raggedright\arraybackslash}X@{}}
    \toprule
    & 现象 & 根因 & 修复 \\
    \midrule
    1 & port 3 状态 \code{0x301}，复位后 \code{0x303} & 低速设备在高速 Hub 后面 & 记录父 Hub、端口、速度，走 split 事务 \\
    2 & \code{control setup failed req=6 addr=0}，\code{intr=0x0a} & SETUP 的 SSPLIT 没设 \code{LSPDDEV}，TT 回 STALL & 按 \code{usb\_route} 给所有低速事务置 \code{LSPDDEV} \\
    3 & GET\_DESCRIPTOR 数据阶段返回 $-2$（超时） & CSPLIT 收到 NAK 后仍反复发 CSPLIT & NAK 时清 \code{complete\_split}，下一帧重发 SSPLIT；NYET 才继续 CSPLIT \\
    4 & 设备描述符 \code{class=0 vid=0 pid=0} & 18 字节跨 3 个 8 字节包，被当成一个 split 请求，只收到前 8 字节 & 按最大包长分包，每包单独 SSPLIT/CSPLIT，完整收到一包后翻转 DATA0/1；检查长度、类型和 VID \\
    5 & “not a boot keyboard” & 复合设备的后一个接口覆盖了键盘接口的属性 & 先定位 Boot Keyboard 接口，端点只在该接口内选 \\
    6 & ready 之后又出现 \code{req=6 addr=4} STALL & 绑定键盘后仍继续 CDC-ECM 网卡探测，向键盘要配置描述符 & HID 交给 USB 总线后直接返回 \\
    7 & 按键无响应，反复 \code{hcint=0x42} & 把每个 NYET 都当作“TT 还在忙”，永久重发 CSPLIT & 记录 SSPLIT ACK 的帧号，同帧内最多重试 3 次；跨帧则结束本轮，等 \code{bInterval} \\
    8 & \code{ssack} 比 \code{issue} 晚几十个微帧 & ACK 交给 kworker 处理，调度延迟毫秒级，CSPLIT 窗口早过了 & SSPLIT ACK 和窗口内 NYET 改在 IRQ 快速路径里处理 \\
    9 & 快速路径启用后仍偶尔跨帧 & 快速路径之前有一行 \code{printf}；115200 波特率下一行要几毫秒 & 中间状态不打印，只在本轮结束后限量打印 \code{ch6 final} \\
    \bottomrule
  \end{tabularx}
\end{table}

其中几个阶段值得多说几句。

\paragraph{阶段 4：“读到了”但其实没读全.} 低速设备端点 0 的最大包长只有 8 字节，18 字节的设备描述符要分 3 个数据包传。经过 split 事务时，每个包都要单独走一遍 SSPLIT/CSPLIT，而且收完一个完整的包后要把 DATA0/DATA1 翻转。旧代码把整个 18 字节当成一次 split 请求，只取回了第一个包，后面 10 字节还是缓冲区里的 0——于是日志说 \code{vid=0 pid=0}。现在 \code{channel\_xfer()} 对 split 请求按最大包长切片，\code{get\_device\_descriptor()} 还检查长度、类型和 VID 不为 0，宁可失败也不接受半截数据。

\paragraph{阶段 7～8：时间窗口.} 周期性的 split 事务有严格的时序：CSPLIT 应在 SSPLIT 之后第 2 个微帧开始，并且只能在同一个帧内重试。早期代码对 NYET 无条件重试，永远等不到结果；改成按帧号判断窗口后，又发现记录下来的“ACK 帧号”比实际晚了几十个微帧——因为 ACK 是由 kworker 处理的，而 kworker 要等调度器给它 CPU，延迟是毫秒级的，比 125\,$\mu$s 的微帧粗了一个数量级。解决办法是把这两种有截止时间的中间事件移到中断上下文里（\code{dwc2\_interrupt\_rx\_irq\_fast()}），只把“这一轮已经结束”的结果交给线程。这是“上半部只做最少的事”这条原则的一个反例：对有硬实时要求的状态推进，上半部恰恰是唯一来得及的地方。

\paragraph{阶段 9：观察者效应.} 快速路径加上以后，问题偶尔还在。原因是 \code{dwc2\_irq()} 开头还留着一行调试 \code{printf}：串口 115200 波特率下每个字符约 87\,$\mu$s，一行几十个字符就是好几个毫秒，足够错过整个窗口。这是典型的“加日志就改变时序”的 Heisenbug。最终版本在有截止时间的分支里一行都不打印，只在本轮结束之后限量打印 \code{dwc2 ch6 final}；如果确实需要跟踪中间事件，应该写进内存里的 trace 环形缓冲，事后由进程上下文打印。

\paragraph{诊断顺序.} 整理下来，按键没反应时按这个顺序看日志：有没有 \code{ch6 armed}（channel 配置并启动了）；按键时有没有 \code{ch6 final}（host channel 产生了中断）；有没有 \code{usbkbd: channel 6 IRQ received}（URB 已归还、键盘 work 已排队）；有没有 \code{first HID report}（DMA 里读到了报告）；报告里的 usage 是否非 0（键盘真的报了键）；最后才看键码映射和控制台。真机第一次成功时打印的是 \code{usbkbd: first HID report mod=0 keys=10,0,0,0,0,0}——usage \code{0x10}，字母 \code{m}。

\paragraph{从两层轮询到中断驱动.} 最早的实现其实有两层轮询：定时器每个 tick 唤醒一次键盘 work，work 里启动一次传输再忙等 \code{HCINT}。改成 URB 加中断以后，CPU 只在真的有报告（或本轮结束）时才参与，空闲键盘不再占用任何 CPU 时间。

\section{I2C}

BCM2837 的 I2C 控制器叫 BSC（Broadcom Serial Controller）。项目目前只用到 BSC0，用来配置摄像头传感器。在树莓派 3B+ 上，摄像头排线口的 I2C 走 GPIO44/45 的 ALT1 功能，而不是 40 针排针上的 GPIO0/1。

一次传输的过程是：写从机地址寄存器 \code{A}、数据长度 \code{DLEN}，清 FIFO，向控制寄存器 \code{C} 写入 \code{I2CEN | ST}（读还要加 \code{READ}），然后轮询状态寄存器 \code{S}：\code{TXW} 表示可以继续写 FIFO，\code{RXD} 表示 FIFO 有数据可读，\code{DONE} 表示传输完成，\code{ERR} 表示从机没有应答，\code{CLKT} 表示时钟拉伸超时。分频寄存器 \code{DIV} $=$ 核心时钟 $/$ 100\,kHz。在真机上，BSC0 对摄像头传感器始终返回 NACK，驱动最后改用 GPIO44/45 上的软件 SCCB，排查过程见 \ref{sec:cam-probe} 小节。

\begin{ccode}
int i2c_write(int bus, uint8 addr, const uint8 *buf, int len);
int i2c_read(int bus, uint8 addr, uint8 *buf, int len);
int i2c_write_read(int bus, uint8 addr, const uint8 *wbuf, int wlen,
                   uint8 *rbuf, int rlen);     // 先写寄存器地址，再读
\end{ccode}

摄像头寄存器的访问很短（写是 2 字节地址加 1 字节数据，读是 2 字节地址后读 1 字节），所以全部用轮询加超时，不用中断。为了在真机上区分“控制器没配好”和“线上没有设备”，驱动还带了一个用 GPIO 模拟 I2C 的诊断路径：直接把 GPIO44/45 当开漏输出，按位产生起始、地址、应答和停止条件。

\section{网络协议栈}

在讲两块无线网卡之前，先看看它们的上层。\file{kernel/net.c} 参考 MIT xv6 网络实验的分层，逐步扩展成了一个小而完整的协议栈：

\begin{itemize}
  \item \code{struct net\_device}：每块网卡一个，保存 MAC、IP、掩码、网关、DNS 和发送函数 \code{start\_xmit}；驱动收到帧后调用 \code{net\_rx\_dev(dev, frame, len)}；
  \item Ethernet、ARP（表项以 \code{(ip, dev)} 为键，多网卡互不干扰）、IPv4、ICMP、UDP、DHCP；
  \item 路由表：最长前缀匹配，同长度按 metric 选择，DHCP 成功后自动加直连路由和默认路由；
  \item TCP：被动打开（\code{listen/accept}），每连接 8\,KiB 异步发送队列，单段停等式重传；
  \item socket 文件描述符和 epoll：socket、PTY、管道都是 \code{struct file}，可以统一 \code{read/write/epoll}。
\end{itemize}

在这之上，用户态实现了 \code{ping}、\code{dhcp}、\code{tftp}、\code{netdns}，以及一个基于 TweetNaCl 的 SSH 服务器（第~\ref{chap:user} 章）。

\section{SDIO 无线网卡：BCM43430/43455}

\subsection{硬件路径}

树莓派 3B 板载 BCM43430，3B+ 板载 BCM43455。它们通过 SDIO 接在 Arasan SDHCI 控制器（\code{0x3f300000}，GPIO34～39 ALT3，IRQ 62）上。这里有两个容易混淆的概念：Arasan 是 SDIO \emph{主机控制器}，负责产生 SDIO 命令和时钟、搬运数据；BCM43455 是 \emph{FullMAC} Wi-Fi 芯片，片内有一颗 Cortex-R4 运行厂商固件，处理大部分 802.11 MAC 工作。

\subsection{硬件结构：Arasan、SDIO 总线与 FullMAC 芯片}
\label{sec:wifi-hw}

和摄像头一样，先把参与的部件和它们之间的通路理清楚（图~\ref{fig:wifi-hw}）。

\begin{figure}[htbp]
  \centering
  \begin{tikzpicture}[x=1mm,y=1mm, font=\footnotesize,
      data/.style={{Stealth[length=2.4mm]}-{Stealth[length=2.4mm]}, line width=1.6pt, draw=orange!85!black},
      ctl/.style={-{Stealth[length=2mm]}, thick, draw=structurecolor!70!black},
      evt/.style={-{Stealth[length=2mm]}, thick, dashed, draw=red!70!black},
      blk/.style={box, font=\footnotesize, minimum height=8mm}]
    \draw[rounded corners=3pt, draw=gray!60, fill=gray!4] (-2,-2) rectangle (66,82);
    \node[lab, anchor=north west] at (-1,81.5) {BCM2837 SoC};
    \draw[rounded corners=3pt, draw=gray!60, fill=gray!4] (86,-2) rectangle (152,82);
    \node[lab, anchor=north west] at (87,81.5) {BCM43455 Wi-Fi 芯片};

    % SoC side
    \node[hbox, minimum width=23mm, minimum height=62mm, text width=23mm, inner sep=1pt] (cpu) at (11,40) {ARM\\Cortex-A53\\[2pt]\file{brcmfmac.c}\\\file{sdio.c}\\\file{arasan\_sdio.c}};
    \node[blk, text width=30mm] (mb) at (48,73) {mailbox → 固件\\查询 EMMC 时钟};
    \node[hbox, text width=30mm, minimum height=22mm] (ar) at (48,48) {\textbf{Arasan SDHCI}\\\code{0x3f300000}\\\code{ARG1 CMDTM}\\\code{INTERRUPT DATA}};
    \node[blk, text width=30mm] (ic) at (48,18) {legacy 中断控制器\\IRQ 62};

    % SDIO device side
    \draw[rounded corners=2pt, draw=structurecolor!60!black, fill=structurecolor!5] (89,26) rectangle (111,78);
    \node[lab, anchor=south] at (100,26.5) {SDIO 设备接口};
    \node[blk, text width=17mm, minimum height=9mm] (f0) at (100,71) {F0\\CCCR};
    \node[blk, text width=17mm, minimum height=9mm] (f1) at (100,57) {F1\\背板窗口};
    \node[blk, text width=17mm, minimum height=9mm, fill=orange!15] (f2) at (100,44) {F2 SDPCM\\FIFO};
    \fill[gray!25] (115,4) rectangle (119,78);
    \node[rotate=90, font=\scriptsize] at (117,16) {芯片内部背板（AXI）};
    \node[blk, text width=26mm] (cc) at (135,70) {ChipCommon\\chip ID、EROM};
    \node[hbox, text width=26mm, minimum height=16mm] (cr4) at (135,50) {ARM Cortex-R4\\TCM RAM：\\固件 + NVRAM};
    \node[blk, text width=26mm] (d11) at (135,22) {D11 MAC + PHY\\射频};
    \node[lab] (air) at (135,5) {802.11 无线 $\leftrightarrow$ AP};

    % control
    \draw[ctl] (cpu.east |- mb) -- (mb);
    \draw[ctl] (cpu.east |- 0,57) -- (ar.west |- 0,57);
    \draw[ctl] (ar.east |- 0,57) -- node[lab, above]{CMD 线} node[lab, below]{CMD52/53} (89,57) ;
    \draw[ctl] (f1.east) -- (115,58);
    \draw[ctl] (119,70) -- (cc.west);
    \draw[ctl] (119,58) -- (cr4.west |- 0,58);
    % data
    \draw[data] (cpu.east |- 0,44) -- node[lab, below, text=orange!50!black]{PIO} (ar.west |- 0,44);
    \draw[data] (ar.east |- 0,44) -- node[lab, above, text=orange!50!black]{DAT0–3} (89,44);
    \draw[data] (f2.east) -- (cr4.west |- f2.east);
    \draw[data] (cr4) -- node[lab, right]{802.3 $\leftrightarrow$ 802.11} (d11);
    \draw[data] (d11) -- (air);
    % events
    \draw[evt] (125,41) |- (111,31);
    \draw[evt] (89,31) -- node[lab, above]{DAT1 中断} (ar.east |- 0,31) -- (ar.south east);
    \draw[evt] (ar.south) -- node[lab, right]{\code{CARD\_INT}} (ic.north);
    \draw[evt] (ic.west) -- (cpu.east |- ic.west);
    \node[lab, align=center] at (76,12) {SDIO 总线\\CLK、CMD、\\DAT0–3\\GPIO34–39};
  \end{tikzpicture}
  \caption{SDIO Wi-Fi 的硬件拓扑。细实线是控制面：CMD 线上的 CMD52/CMD53 经 F0 配置总线、经 F1 访问芯片内部背板（读 chip ID、下载固件、控制 CR4 复位）；粗线是数据面：CPU 以 PIO 方式经 Arasan 的 \code{DATA} 寄存器读写 F2 的 SDPCM FIFO，芯片固件再与 802.11 帧互转；虚线是事件：固件经 DAT1 发出卡中断，成为 BCM2837 的 IRQ 62}
  \label{fig:wifi-hw}
\end{figure}

\paragraph{两颗处理器.} 这条链路两端各有一颗 CPU：一端是运行 xv6 的 Cortex-A53，另一端是 BCM43455 片内的 Cortex-R4。R4 执行主机下载进去的厂商固件，负责扫描、认证、关联、管理帧、ACK 与重传、用主机装入的密钥做 CCMP 加解密，以及 802.11 帧与以太网帧之间的转换——这就是 \emph{FullMAC} 的含义。主机看到的是一个“收发以太网帧、另外接受若干控制命令”的设备，802.11 的细节全藏在固件里。\ref{sec:mt7601u} 节的 MT7601U 恰好相反，是 SoftMAC：802.11 帧由主机驱动自己组装。

\paragraph{SDIO 总线.} 两者之间是 SDIO 总线：一根时钟线 CLK、一根命令线 CMD、四根数据线 DAT0～3，在 3B+ 上走 GPIO34～39 的 ALT3。CMD52 只用命令线，命令里带地址和一个字节，响应里带回一个字节，用来读写寄存器；CMD53 在命令线上发出地址和长度，数据走四根 DAT 线，用来搬运固件和网络帧。DAT1 还有第二个用途：总线空闲时，卡可以把它拉低来请求中断，所以 Wi-Fi 芯片不需要单独的中断引脚。

\paragraph{一个 SDIO 设备，三个 function.} SDIO 卡在协议上分成若干 function，各有独立的 17 位寄存器地址空间：
\begin{itemize}
  \item \textbf{F0} 是公共寄存器区 CCCR，设置总线宽度、使能其他 function、打开中断；
  \item \textbf{F1} 是进入芯片内部的“窗口”。BCM43455 内部的 ChipCommon、SDIO 核、R4 和 TCM RAM 挂在一条背板总线上，地址空间远大于 17 位，所以 F1 只开一个 32\,KiB 的窗口：先把目标地址的高位写进窗口寄存器，再用低 15 位访问。读 chip ID、解析 EROM、让 R4 停机、下载固件和 NVRAM、写复位向量，全都经过这个窗口；
  \item \textbf{F2} 是固件启动之后才就绪的 FIFO，运行时所有控制命令、事件和网络帧都从这里进出。
\end{itemize}

\paragraph{Arasan 是主机控制器，而且工作在 PIO 模式.} Arasan SDHCI 负责产生 SDIO 时钟、把命令编码到 CMD 线上、按 4 位总线收发数据。驱动发一条命令的方式是写 \code{ARG1}（参数）和 \code{CMDTM}（命令号与传输方式），然后等 \code{INTERRUPT} 里的完成位。有数据阶段时，本项目没有使用控制器的 DMA，而是 CPU 等 \code{READ\_RDY}/\code{WRITE\_RDY}，再一个 32 位字一个 32 位字地读写 \code{DATA} 寄存器。这和摄像头形成鲜明对比：Unicam 的数据完全不经过 CPU，而这里每个字都要 CPU 亲手搬运。这样实现简单，也没有缓存一致性问题，代价是吞吐受 CPU 和单次传输长度限制——实机曾在 512、256、128 字节的 CMD53 上出现数据 CRC 错误，驱动逐级减半，最后稳定在 64 字节。这个上限是运行时探测出来的，完成多块传输和时序调优之后应当重新评估。

\paragraph{中断要过三道闸门.} 固件有帧要交给主机时，DAT1 上的卡中断要穿过三层开关才能变成 CPU 的 IRQ：芯片内部 SDIO 核的 \code{HOSTINTMASK} 决定哪些事件送到 DAT1；F0 的 CCCR 里的 \code{INT\_ENABLE} 打开 master 位和 F1 位（\code{sdio\_claim\_irq()}）；Arasan 的 \code{IRPT\_EN} 打开 \code{CARD\_INT}，BCM2837 的 legacy 控制器再打开 IRQ 62。任何一层没开，接收都只能靠轮询。

\paragraph{固件从哪里来.} 芯片每次上电都是空的，固件、NVRAM、CLM 三个文件放在 SD 卡的 FAT32 启动分区里，由另一个控制器 SDHOST（\code{0x3f202000}）读出。两个 SD 控制器各管一条总线：SDHOST 接外置 SD 卡，Arasan 接板载 Wi-Fi，混用会造成 GPIO 和时钟的所有权冲突。


\subsection{SDIO 枚举}

Arasan 驱动 probe 后创建 \code{mmc\_host}，SDIO core 依次发送 CMD0（复位）、CMD5（协商 I/O 电压，等待就绪）、CMD3（取 RCA）、CMD7（选中），再用 CMD52 读 CCCR、FBR 和 CIS，切换到 4 位总线，最后把三个 function 注册到 \code{sdio} 总线。\code{brcmfmac} 驱动按 vendor \code{0x02d0}、device \code{0xa9a6}/\code{0xa9bf}、function 1 匹配。三条 SDIO 命令的分工是：CMD5 识别和上电；CMD52 单字节读写寄存器；CMD53 批量传输。

\subsection{固件装载}

FullMAC 芯片每次上电都要由主机下载固件。驱动从 FAT32 启动分区读取三个文件（所以必须排在 \code{fat32init()} 之后）：

\begin{itemize}
  \item \code{BCM43455.BIN}：芯片 CR4 的执行代码；
  \item \code{BCM43455.TXT}：板级 NVRAM 参数（天线、校准）；
  \item \code{BCM43455.CLM}：各国信道和功率的法规数据库。
\end{itemize}

Function 1 是“背板窗口”：设置窗口寄存器后，可以用 CMD52/CMD53 访问芯片内部的 ChipCommon、SDIO core、CR4 和片内 RAM。驱动读 chip ID、解析 EROM 找到各个核和 RAM，让 CR4 停机，用 CMD53 把固件和 NVRAM 写进 RAM，写入复位向量后放开 CR4，等待固件就绪的 mailbox 中断，最后才打开 Function 2。Function 2 是运行时的数据 FIFO。

\subsection{SDPCM 与 BCDC}

运行时的每一次 F2 传输都是分层封装的（图~\ref{fig:wifi-frame}）。

\begin{figure}[htbp]
  \centering
  \begin{tikzpicture}[x=1mm,y=1mm, font=\scriptsize]
    \tikzset{fld/.style={draw=structurecolor!70!black, minimum height=8mm, inner sep=1pt, align=center}}
    % data frame
    \node[lab, anchor=east] at (-1,24) {数据通道};
    \node[fld, fill=structurecolor!25, minimum width=36mm, anchor=west] (sd) at (0,24) {SDPCM 头 12\,B\\长度 | \textasciitilde 长度 | 序号 | 通道 2 | … | 偏移};
    \node[fld, fill=structurecolor!12, minimum width=16mm, anchor=west] (bc) at (sd.east) {BCDC 头\\4\,B};
    \node[fld, fill=orange!20, minimum width=30mm, anchor=west] (eh) at (bc.east) {以太网头 14\,B\\目的 | 源 | 类型};
    \node[fld, fill=orange!8, minimum width=40mm, anchor=west] (pl) at (eh.east) {载荷\\IPv4 / ARP / EAPOL（\code{0x888e}）};
    \node[fld, fill=gray!10, minimum width=10mm, anchor=west] (pd) at (pl.east) {填充\\到 4\,B};
    % control frame
    \node[lab, anchor=east] at (-1,4) {控制通道};
    \node[fld, fill=structurecolor!25, minimum width=36mm, anchor=west] (sd2) at (0,4) {SDPCM 头 12\,B\\… | 通道 0 | …};
    \node[fld, fill=structurecolor!12, minimum width=46mm, anchor=west] (dc) at (sd2.east) {BCDC 命令头 16\,B\\\code{cmd} | \code{len} | \code{flags}（含请求号） | \code{status}};
    \node[fld, fill=gray!6, minimum width=50mm, anchor=west] (iv) at (dc.east) {参数\\如 \code{"bsscfg:wpa\_auth\textbackslash0"} + 索引 + 值};
    % braces
    \draw[decorate, decoration={brace, amplitude=4pt}] (0,29) -- node[lab, above=4pt]{\code{brcmf\_f2\_transfer} 一次写出，再按传输上限（实测 64\,B）拆成多次 CMD53} (132,29);
    \draw[decorate, decoration={brace, amplitude=4pt, mirror}] (bc.south east) ++(0,-1.5) -- node[lab, below=4pt]{交给 \code{net\_rx\_dev} 的部分} ([yshift=-1.5mm]pl.south east);
  \end{tikzpicture}
  \caption{F2 上一帧的封装。SDPCM 是总线一层的帧格式，用通道号区分控制、事件和数据；控制通道里是 BCDC 命令（ioctl/iovar），数据通道里是 BCDC 数据头加一个完整的以太网帧}
  \label{fig:wifi-frame}
\end{figure}

发送控制命令后，下一个到达的帧不一定就是响应——事件和数据帧可能插在前面。接收分发函数必须按通道分流，只有请求号相同的控制帧才是对应的响应。

\subsection{WPA2 四次握手}

这版固件不处理 WPA2 握手，所以主机自己完成四次握手（\file{kernel/wpa\_crypto.c}）：用 PBKDF2 从口令和 SSID 算出 PMK；收到 M1 后生成 SNonce，派生 PTK，计算 MIC，发 M2；收到 M3 后验证 MIC，用 AES key unwrap 解出组密钥 GTK；把成对密钥和组密钥装进芯片，发 M4。握手完成后才放行普通数据帧，随后执行 DHCP。

\subsection{数据流与并发}
\label{sec:wifi-flow}

图~\ref{fig:wifi-flow} 把一帧数据在两个方向上经过的函数排成了层。两个方向的触发方式不同：

\begin{figure}[htbp]
  \centering
  \begin{tikzpicture}[x=1mm,y=1mm, font=\scriptsize,
      fn/.style={box, font=\scriptsize, minimum height=8mm, inner sep=2pt, text width=40mm},
      tx/.style={-{Stealth[length=2mm]}, line width=1.2pt, draw=orange!85!black},
      rx/.style={-{Stealth[length=2mm]}, line width=1.2pt, draw=structurecolor!80!black},
      evt/.style={-{Stealth[length=1.8mm]}, semithick, dashed, draw=red!70!black},
      call/.style={-{Stealth[length=1.8mm]}, semithick, draw=gray!60!black}]
    % layer bands
    \foreach \y/\h/\t in {84/{协议栈}, 70/{网卡接口}, 56/{\file{brcmfmac}\\封装/分流}, 42/{\file{brcmfmac}\\F2 收发}, 28/{SDIO core}, 14/{Arasan host}, 1/{芯片}}{
      \fill[gray!5] (0,\y-6) rectangle (150,\y+6);
      \node[lab, align=center, text width=14mm] at (7,\y) {\h};
    }
    \node[font=\footnotesize\bfseries, text=orange!60!black] at (38,93) {发送（TX）};
    \node[font=\footnotesize\bfseries, text=structurecolor!70!black] at (85,93) {接收（RX）};
    \node[font=\footnotesize\bfseries, text=red!60!black] at (131,93) {接收的触发};
    % TX
    \node[fn] (t1) at (38,84) {UDP/TCP/ICMP → IPv4\\路由 → ARP 解析下一跳 MAC};
    \node[fn] (t2) at (38,70) {\code{start\_xmit = brcmfmac\_xmit}\\握手未完成则拒绝；持 \code{bus->lock}};
    \node[fn] (t3) at (38,56) {\code{brcmf\_send\_ether\_locked}\\加 BCDC 头 + SDPCM 头（通道 2）};
    \node[fn] (t4) at (38,42) {\code{brcmf\_f2\_transfer(write)}\\按 \code{brcmf\_transfer\_limit} 分块};
    \node[fn] (t5) at (38,28) {\code{sdio\_cmd53(F2)}\\\code{mmc\_request\_sync}：持 \code{request\_lock}};
    \node[fn] (t6) at (38,14) {\code{arasan\_request}\\写 \code{ARG1/CMDTM}，逐字写 \code{DATA}};
    \node[gbox, font=\scriptsize, text width=40mm, minimum height=6mm] (t7) at (38,1) {F2 FIFO → 固件 → 802.11 发出};
    \foreach \a/\b in {t1/t2,t2/t3,t3/t4,t4/t5,t5/t6,t6/t7}{\draw[tx] (\a) -- (\b);}
    % RX
    \node[gbox, font=\scriptsize, text width=40mm, minimum height=6mm] (r7) at (85,1) {802.11 收到 → 固件 → F2 FIFO};
    \node[fn] (r6) at (85,14) {\code{arasan\_request}\\逐字读 \code{DATA}};
    \node[fn] (r5) at (85,28) {\code{sdio\_cmd53(F2)}\\\code{mmc\_request\_sync}};
    \node[fn] (r4) at (85,42) {\code{brcmf\_rx\_frame\_locked}\\经 F1 读 \code{INTSTATUS}；读 64\,B 头，再读余下};
    \node[fn] (r3) at (85,56) {\code{brcmf\_rx\_dispatch\_locked}\\控制 → \code{ctl} 缓冲；事件 → 打印\\数据：EAPOL → WPA；其余 ↑};
    \node[fn] (r2) at (85,70) {\code{net\_rx\_dev(wlan0, frame)}};
    \node[fn] (r1) at (85,84) {Ethernet → ARP / IPv4 →\\ICMP / UDP / TCP};
    \foreach \a/\b in {r7/r6,r6/r5,r5/r4,r4/r3}{\draw[rx] (\a) -- (\b);}
    \draw[rx] (r3) -- node[lab, right]{数据（握手后）} (r2);
    \draw[rx] (r2) -- (r1);
    % trigger column
    \node[fn, text width=26mm, draw=red!60!black] (i1) at (131,1) {DAT1 卡中断};
    \node[fn, text width=26mm, draw=red!60!black] (i2) at (131,14) {\code{arasan\_sdio\_irq}\\屏蔽 \code{CARD\_INT}};
    \node[fn, text width=26mm, draw=red!60!black] (i3) at (131,28) {\code{napi\_schedule}\\唤醒 \code{brcmf0\_wq}};
    \node[fn, text width=26mm] (i4) at (131,42) {kworker：\\\code{brcmf\_napi\_poll(16)}};
    \draw[evt] (i1) -- node[lab, right]{IRQ 62} (i2);
    \draw[evt] (i2) -- node[lab, right]{上半部} (i3);
    \draw[evt] (i3) -- node[lab, right]{下半部} (i4);
    \draw[call] (i4.north) |- node[lab, above, pos=0.75]{每轮至多 16 帧} (r3.east);
    \node[fn, text width=33mm] (i5) at (130,70) {读空：\code{napi\_complete\_done}\\\code{arasan\_sdio\_irq\_complete}};
    \draw[call] (i4.east) -- ++(2.5,0) |- (i5.east);
  \end{tikzpicture}
  \caption{SDIO Wi-Fi 的数据流。发送在调用者的上下文里一路同步向下；接收由 DAT1 卡中断触发，上半部只屏蔽中断并调度 NAPI，真正的 CMD53 读取和分流在 \code{brcmf0\_wq} 内核线程里完成，读空后才重新打开卡中断。两个方向都要经过 \code{mmc\_request\_sync()} 的 host 锁，且每个字都由 CPU 经 \code{DATA} 寄存器搬运}
  \label{fig:wifi-flow}
\end{figure}

\paragraph{发送是同步的.} 协议栈选好路由、解析出下一跳 MAC 以后，调用 \code{wlan0} 的发送函数 \code{brcmfmac\_xmit()}。握手完成之前它直接拒绝普通数据。之后在调用者自己的上下文里（可能是系统调用，也可能是 TCP 重传 work）持 \code{bus->lock} 加上 BCDC 头和 SDPCM 头，经 \code{brcmf\_f2\_transfer()} 拆成若干次 CMD53 写进 F2，函数返回时数据已经交给了芯片。

\paragraph{接收是中断触发、线程执行的.} DAT1 卡中断到来时，上半部 \code{arasan\_sdio\_irq()} 只做三件事：在 Arasan 上屏蔽 \code{CARD\_INT}、清掉状态位、调用 \code{napi\_schedule()} 唤醒设备私有的 \code{brcmf0\_wq} 内核线程。CMD53 要等总线、可能睡眠，绝不能在中断里做。下半部 \code{brcmf\_napi\_poll()} 在 kworker 里循环：经 F1 读 SDIO 核的 \code{INTSTATUS}，看到“有帧”就从 F2 读 64 字节的头、校验 SDPCM 长度和它的反码、再读剩余部分，然后按通道分流。每轮最多处理 16 帧；没读空就让 NAPI 重新排队，读空了才 \code{napi\_complete\_done()} 并重新打开 \code{CARD\_INT}。卡中断是电平触发的，在芯片里的帧被取走之前一直有效，所以必须先屏蔽、读空再打开，否则 CPU 会被同一个中断反复打断。

\paragraph{同步的控制命令.} 控制面是请求/应答式的：\code{brcmf\_dcmd()} 发出一个带递增请求号的 BCDC 命令，然后反复调用分流函数，直到读到请求号相同的控制帧。事件和数据帧可能插在应答前面，所以分流函数不能假定“下一帧就是回复”，必须按通道放到各自的去处：控制帧存进 \code{ctl} 缓冲，事件帧打印出来，数据帧中的 EAPOL 交给 WPA2 握手逻辑，其余在握手完成后交给 \code{net\_rx\_dev()}。

\paragraph{两把锁，两个粒度.} \code{bus->lock} 以\emph{设备}为粒度，保证同一块网卡的 SDPCM 序号、收发缓冲和连接状态不会被发送路径、NAPI 和连接 watchdog 同时改动。\code{host->request\_lock} 以\emph{控制器}为粒度，保证 Arasan 那一套 \code{ARG1/CMDTM/INTERRUPT/DATA} 寄存器在一条 CMD52/CMD53 完成之前不被另一条覆盖。\ref{sec:cam-bugs} 小节的 \code{CMD52 failed status=df0001} 就是缺了后一把锁造成的。两把都是睡眠锁或在进程上下文中持有，因为一次 SDIO 事务可能要等上毫秒级的超时。

\paragraph{从轮询到中断的演进.} 这条接收路径经历过三个版本：最早 CPU0 在每次定时器中断里直接轮询 F2，完整的接收和协议处理都跑在中断上下文里；后来有了 DAT1 中断，但仍等下一次定时器才去读；现在中断直接调度 NAPI，定时器只负责推进 delayed-work 的时钟（连接 watchdog、TCP 重传）。早期每 100\,ms 轮询一次时，SSH 输入一个字符要等 100～300\,ms 才回显；改成中断驱动后，这个延迟基本消失。

\section{USB 无线网卡：MT7601U}
\label{sec:mt7601u}

MT7601U 是一块常见的 USB 无线网卡（VID \code{148f}，PID \code{7601}）。与 BCM43455 最大的不同在于它是 \emph{SoftMAC}：芯片只负责射频、基带和收发帧，扫描、认证、关联、加密握手等 802.11 逻辑都要由主机驱动完成（表~\ref{tab:fullsoft}）。

\begin{table}[htbp]
  \centering
  \caption{两块无线网卡的对比}
  \label{tab:fullsoft}
  \small
  \begin{tabularx}{\linewidth}{@{}lXX@{}}
    \toprule
    & BCM43430/43455 & MT7601U \\
    \midrule
    总线 & SDIO（Arasan） & USB（DWC2 + LAN951x 集线器） \\
    类型 & FullMAC：固件完成扫描、关联 & SoftMAC：主机构造 802.11 帧 \\
    固件 & BIN + NVRAM + CLM，经背板写入 RAM & \code{MT7601U.BIN}，经 Bulk OUT 写入 MCU 的 ILM/DLM \\
    控制 & BCDC 命令 & MCU 命令（EP8 发、EP5 收）和寄存器读写 \\
    数据 & SDPCM 数据通道里的以太网帧 & TXWI 描述符 + 802.11 帧，EP8 发、EP4 收 \\
    网卡名 & \code{wlan0} & \code{wlan1} \\
    \bottomrule
  \end{tabularx}
\end{table}

驱动初始化的主要步骤是：

\begin{enumerate}
  \item 解析固件头，把 IVB、ILM、DLM 三段通过 Bulk OUT 写入 MCU，等待 \code{MT\_MCU\_COM\_REG0} 变成 1；
  \item 读 eFuse 得到永久 MAC 地址、频偏、RSSI 和温度校准参数；
  \item 初始化 MAC、BBP（数字基带）、RF（模拟射频）三组寄存器表，执行各项校准；
  \item 扫描信道，解析 beacon/probe response，选择目标 BSSID；
  \item 发送 Open System 认证帧和带 RSN IE 的关联请求；
  \item 完成 WPA2 四次握手，安装密钥，注册 \code{wlan1}。
\end{enumerate}

发送一个 IP 包要经过：以太网帧 → LLC/SNAP 封装 → 802.11 数据帧头 → TXWI → USB DMA 头和填充 → EP8 Bulk OUT。接收端在 EP4 上长期挂着一个 Bulk IN 请求，DWC2 用 DMA 把数据写到缓冲区（总线地址要加 \code{0xC0000000} 别名，见第~\ref{chap:mmu} 章），完成中断后由下半部解析。网络层的路由表决定一个包从 \code{wlan0} 还是 \code{wlan1} 发出，两块网卡可以同时工作。

\section{CSI-2 摄像头：OV5647}

\subsection{路线选择}

树莓派官方的摄像头栈运行在 GPU 固件里（\code{start\_x.elf} + MMAL），ARM 端只是发命令。本项目选择与 Linux 主线相同的“ARM 直接驱动”路线：内核自己通过 I2C 配置传感器，打开电源域和时钟，用 Unicam 接收 MIPI CSI-2 数据并 DMA 到内存。硬件是 Camera Module v1.3（OmniVision OV5647），接在 3B+ 的 CSI1 口。

\subsection{硬件结构：I2C、CSI-2、Unicam 与 DMA}
\label{sec:cam-hw}

摄像头驱动之所以涉及这么多部件，是因为一帧图像从产生到进入内存要经过几条性质完全不同的通路（图~\ref{fig:cam-hw}）。弄清楚每条通路传什么、由谁发起，后面的代码就很好读了。

\begin{figure}[htbp]
  \centering
  \begin{tikzpicture}[x=1mm,y=1mm, font=\footnotesize,
      data/.style={-{Stealth[length=2.6mm]}, line width=1.8pt, draw=orange!85!black},
      ctl/.style={-{Stealth[length=2mm]}, thick, draw=structurecolor!70!black},
      evt/.style={-{Stealth[length=2mm]}, thick, dashed, draw=red!70!black},
      blk/.style={box, font=\footnotesize, minimum height=9mm, text width=24mm}]
    \draw[rounded corners=3pt, draw=gray!60, fill=gray!4] (-2,-3) rectangle (106,80);
    \node[lab, anchor=north west] at (-1,79.5) {BCM2837 SoC};
    \draw[rounded corners=3pt, draw=gray!60, fill=gray!4] (113,16) rectangle (147,80);
    \node[lab, anchor=north west] at (114,79.5) {Camera v1.3 模块};

    % column 1
    \node[hbox, minimum width=22mm, minimum height=56mm, text width=20mm] (cpu) at (11,46) {ARM\\Cortex-A53\\[2pt]xv6 内核\\[2pt]\file{unicam.c}\\\file{ov5647.c}\\\file{i2c.c}\\\file{mbox.c}};
    \node[gbox, minimum width=56mm, minimum height=10mm, text width=54mm] (ram) at (28,5) {SDRAM：\code{CAMDMA}，物理 \code{0x07200000}，非缓存映射};
    % column 2
    \node[blk] (mbox) at (44,70) {mailbox\\属性标签};
    \node[blk] (bsc)  at (44,54) {BSC0 或 GPIO44/45\\软件 SCCB};
    \node[blk] (cm)   at (44,36) {\code{CM\_CAM1}\\lp 时钟 100\,MHz};
    \node[blk] (ic)   at (44,20) {legacy 中断控制器\\GPU IRQ 39};
    % column 3
    \node[blk, fill=gray!8, draw=gray!70] (vc) at (88,70) {VideoCore 固件};
    \node[hbox, minimum width=30mm, minimum height=40mm] (uni) at (88,27) {};
    \node[font=\small\bfseries, anchor=north] at (88,46.5) {Unicam CSI1};
    \node[box, font=\footnotesize, text width=24mm, minimum height=7mm] (phy) at (88,36) {D-PHY\\CLK + DAT0/1};
    \node[box, font=\footnotesize, text width=24mm, minimum height=7mm] (idi) at (88,26) {包过滤 \code{IDI0}\\VC 0、DT \code{0x2B}};
    \node[box, font=\footnotesize, text width=24mm, minimum height=7mm] (dma) at (88,15.5) {DMA \code{IBSA0/IBEA0}};
    % column 4
    \node[gbox, text width=27mm, minimum height=9mm] (reg) at (130,70) {稳压器使能\\\code{CAM\_GPIO0}};
    \node[hbox, text width=24mm, minimum height=12mm] (sen) at (130,50) {OV5647\\SCCB 地址 \code{0x36}};
    \node[gbox, text width=24mm, minimum height=7mm] (osc) at (130,24) {25\,MHz 晶振};
    \draw[gray!70] (reg) -- node[lab, right]{供电} (sen);
    \draw[gray!70] (osc) -- node[lab, right]{XCLK} (sen);

    % control plane
    \draw[ctl] (cpu.east |- mbox) -- (mbox);
    \draw[ctl] (cpu.east |- bsc) -- (bsc);
    \draw[ctl] (cpu.east |- cm) -- (cm);
    \draw[ctl] (mbox) -- (vc);
    \draw[ctl] (vc) -- node[lab, above]{GPIO 133} (reg);
    \draw[ctl] ([xshift=-6mm]vc.south) -- node[lab, left, pos=0.35, align=right]{电源域\\UNICAM1} ([xshift=-6mm]uni.north);
    \draw[ctl] (bsc.east) -- node[lab, above, pos=0.88]{SCCB} (sen.west |- bsc.east);
    \draw[ctl] (cm) -- (uni.west |- cm);
    % data plane
    \draw[data] (sen.west |- 0,47) -- (109,47) |- (uni.east |- phy);
    \node[lab, text=orange!50!black, anchor=west, align=left] at (110,36) {MIPI CSI-2\\2 lane};
    \draw[data] (phy) -- (idi);
    \draw[data] (idi) -- (dma);
    \draw[data] (dma.south) |- node[lab, below, pos=0.7, text=orange!50!black]{AXI 写，总线地址 \code{0xC0000000|PA}} (ram.east);
    \draw[data] ([xshift=-6mm]ram.north) -- node[lab, left]{拷给用户} ([xshift=-6mm]ram.north |- cpu.south);
    % events
    \draw[evt] (uni.west |- ic) -- node[lab, above]{FS/FE} (ic);
    \draw[evt] (ic) -- (cpu.east |- ic);
  \end{tikzpicture}
  \caption{摄像头的硬件拓扑。细实线是控制面：mailbox 上电和电源域、SCCB 写传感器寄存器、lp 时钟；粗线是数据面：CSI-2 → Unicam → DMA → 内存；虚线是事件：Unicam 的帧开始/结束中断，目前屏蔽，由 CPU 轮询 \code{STA}/\code{ISTA}}
  \label{fig:cam-hw}
\end{figure}

\paragraph{传感器有两个接口.} OV5647 对外有两个互不相干的接口。一个是低速的 \emph{SCCB} 控制口（OmniVision 对 I2C 的兼容实现，两根线 SDA/SCL，约 100\,kHz），CPU 通过它读芯片 ID、写一百来个配置寄存器：分辨率、binning、帧长、曝光、增益、MIPI 输出方式、测试图案，最后写“退出待机”让传感器开始出图。另一个是高速的 \emph{MIPI CSI-2} 输出口：一对时钟线加两对数据线（lane），每对都是差分的 D-PHY 链路，空闲时处于低功耗的 LP-11 状态，传数据时切换到高速模式。图像只走 CSI-2，配置只走 SCCB——I2C 再慢也不影响帧率，它只决定“开机配置”要花多久，这也是驱动后来能退到软件 SCCB 的原因。

\paragraph{CSI-2 链路上传的是包.} 传感器把每一帧编码成一串 CSI-2 包：帧开始（FS）和帧结束（FE）是不带数据的短包；中间每一行像素是一个长包，包头里有虚拟通道号（VC）和数据类型（DT，RAW10 是 \code{0x2B}），包尾有 CRC。接收方据此知道一帧何时开始、何时结束、每个包该放到哪里，CRC 和 ECC 错误也是在这一层被发现的（\code{STA} 里的 CRCE、PBE 等位）。

\paragraph{Unicam 是接收器，也是 DMA 主设备.} SoC 一侧的 Unicam CSI1 把上面的事情全做了，它内部可以分成三段：
\begin{itemize}
  \item \textbf{D-PHY 前端}：接收差分信号，完成低功耗态和高速态之间的切换。判断“稳定进入高速态”要等若干个 lp 时钟周期（\code{CLT}/\code{DLT} 里的 settle 和终端使能计数），所以它需要时钟管理器 \code{CM\_CAM1} 提供的 100\,MHz lp 时钟——时钟频率不对，settle 时间就跟着错，表现为链路误码；
  \item \textbf{包处理}：解析包头、校验 CRC，FS/FE 短包变成 \code{ISTA} 里的帧开始/结束事件；长包经 \code{IDI0} 按 VC 和 DT 过滤，只有 VC 0、RAW10 的行数据才往下走。本项目不让 Unicam 解包（\code{IPIPE}$=0$），行数据保持 CSI-2 原样的 10 位打包格式；
  \item \textbf{DMA 写入}：Unicam 自己就是 AXI 总线上的主设备，把每行按 \code{IBLS} 的行跨度写到 \code{IBSA0}～\code{IBEA0} 之间的内存，写到哪里由 \code{IBWP} 指示。这里的“DMA”不是 BCM2837 通用的 DMA 控制器，而是 Unicam 内置的写引擎；它看到的是 VideoCore 总线地址，所以内核写进去的是 \code{0xC0000000 | PA}（第~\ref{chap:mmu} 章“三种地址视图”）。
\end{itemize}

\paragraph{内存一侧的约束.} Unicam 只认一段起止地址，所以缓冲区必须物理连续，不能用 \code{kalloc()} 逐页拼；它放在 \code{PHYSTOP} 以上的 \code{CAMDMA} 区，分配器永远碰不到。这段内存映射成非缓存的普通内存：DMA 写进 DRAM 的数据，CPU 直接读就是最新的，不需要在每帧前后作废缓存行。

\paragraph{电源和时钟由谁管.} 摄像头模块上的稳压器由 \code{CAM\_GPIO0} 使能，这根线接在 VideoCore 管理的 GPIO 扩展器上，Unicam 所在的电源域也归固件管，ARM 只能通过 mailbox 发属性标签请固件代办。\code{CM\_CAM1} 时钟则是普通的时钟管理器寄存器，内核直接写。传感器自己的像素时钟来自模块上的 25\,MHz 晶振，不需要 SoC 提供。

\paragraph{CPU 只在两头出现.} 把这些放在一起看，CPU 的角色就清楚了：开始时它是\emph{配置者}——经 mailbox 上电、写时钟、经 SCCB 配传感器、经 MMIO 配 Unicam；数据开始流动以后，传感器 → Unicam → 内存这条路完全不经过 CPU；CPU 只需要对 FS/FE 事件作出反应，在恰当的帧开始前换掉 DMA 地址（图~\ref{fig:cam-fsm}），最后把整帧拷给用户。事件本可以用中断（GPU IRQ 39）通知，目前为了安全改成 CPU 每 1\,ms 读一次 \code{STA}/\code{ISTA}。

\paragraph{顺序约束.} 几条通路之间有严格的先后关系，驱动的步骤顺序就是由它们决定的：先上电，SCCB 才可能应答；传感器先停在 LP-11，Unicam 才能在干净的空闲态上完成模拟复位和校准；接收器先使能、传感器后出流，否则第一个 FS 可能丢失；停止时反过来，先停传感器再停接收器，避免接收器在一帧中间被关掉。表~\ref{tab:cam-board} 列出了这些部件在 Pi 3B+ 上的具体资源。


\subsection{\code{sensor\_ops}：仿照 V4L2 子设备}

桥接驱动 \file{unicam.c} 不认识任何具体传感器，传感器驱动也不碰 Unicam。两者之间只有一层接口，对应 Linux 的 \code{v4l2\_subdev\_ops}：

\begin{ccode}
struct sensor_ops {
  int  (*power_on)(struct camsensor *s);             // 供电、读芯片 ID、进入 LP-11
  void (*power_off)(struct camsensor *s);
  int  (*stream_on)(struct camsensor *s, const struct cam_mode *mode);
  int  (*stream_off)(struct camsensor *s);
  int  (*set_ctrl)(struct camsensor *s, int id, int value);   // 如测试图案
};
struct camsensor {
  char *name;  const struct sensor_ops *ops;
  const struct cam_mode *modes;  int nmodes;          // 宽、高、行字节数、CSI-2 数据类型
  struct camsensor_bus bus;                           // lane 数、是否非连续时钟、虚拟通道
  int warmup_frames;                                  // 开流后丢弃的帧数
};
\end{ccode}

启动时 \code{ov5647\_init()} 调用 \code{camsensor\_register()}，\code{camera\_init()} 按注册顺序逐个调用 \code{power\_on}，第一个读到正确芯片 ID 的传感器被绑定到 \code{/dev/video0}。要支持 Camera v2（IMX219），只需新写一个实现了 \code{sensor\_ops} 的文件。

\subsection{板级资源}

表~\ref{tab:cam-board} 列出了从 Raspberry Pi 设备树和 Linux 驱动中查到的资源。

\begin{table}[htbp]
  \centering
  \caption{摄像头相关的板级资源（Pi 3B+）}
  \label{tab:cam-board}
  \small
  \begin{tabularx}{\linewidth}{@{}lX@{}}
    \toprule
    项目 & 值 \\
    \midrule
    传感器 I2C & BSC0 走 GPIO44/45（ALT1），从机地址 \code{0x36}，芯片 ID \code{0x5647} \\
    传感器供电 & 固件扩展 GPIO 5，mailbox 编号 $128+5=133$ \\
    传感器时钟 & 摄像头板上自带 25\,MHz 晶振 \\
    接收器 & Unicam CSI1，寄存器 \code{0x7e801000}，lane 时钟门控 \code{0x7e802004} \\
    中断 & GPU IRQ 39 \\
    电源域 & UNICAM1：设备树编号 13，发给固件时要加 1，即 14 \\
    lp 时钟 & 时钟管理器 \code{CM\_CAM1}，100\,MHz \\
    VPU 时钟 & Unicam 要求至少 250\,MHz，\file{config.txt} 中 \code{core\_freq=250} \\
    CSI-2 链路 & 2 条数据 lane，非连续时钟，RAW10（数据类型 \code{0x2B}） \\
    \bottomrule
  \end{tabularx}
\end{table}

供电和电源域都要通过 VideoCore 固件的 property mailbox：内核在一块 16 字节对齐的缓冲区里写好“标签”（如 \code{0x00038041} 设置扩展 GPIO、\code{0x00038030} 设置电源域），把地址写入 mailbox 寄存器，等固件把结果写回同一块缓冲区。项目把原来分散在 SD、USB 驱动里的私有 mailbox 代码收拢成了带锁的 \file{kernel/mbox.c}。电源域编号是一个容易出错的地方：设备树里的 UNICAM1 是 13，发给固件时要加 1。

\subsection{上电与探测}
\label{sec:cam-probe}

启动时 \code{main()} 先调用 \code{ov5647\_init()} 登记传感器驱动（此时不碰硬件），再由 \code{camera\_init()} 注册字符设备 \code{CAMERA}（主设备号 5，\code{/dev/video0}），然后对每个登记过的传感器先检查其默认模式 Unicam 能否接收，再调用 \code{power\_on}。OV5647 的做法是：经 mailbox 把固件扩展 GPIO 5 先拉低 20\,ms、再置 1，等 100\,ms，把 GPIO44/45 交给传感器的控制总线，读寄存器 \code{0x300a/0x300b}，应得 \code{0x56 0x47}。下面几段讲这条控制总线是怎样从硬件 I2C 一步步退到 25\,kHz 软件 SCCB 的。成功时打印：

\begin{txtcode}
camera: OV5647 on CSI1 (2 lanes) -> /dev/video0, 640x480 RAW10
\end{txtcode}

\paragraph{mailbox.} 供电走的是固件 property mailbox。电源控制现在与 Linux \file{gpio-raspberrypi-exp.c} 一致：先用 \code{GET\_GPIO\_CONFIG} 读出并保留 CAM\_GPIO0 的极性，再设置方向和电平，最后用 \code{GET\_GPIO\_STATE} 回读，标准 Pi 3 应显示 \code{direction=output polarity=normal state=1}。消息本身按 Linux \code{rpi\_firmware\_property()} 的约定组织：标签的请求/应答长度字填 0；发送前清空 FIFO 里其他驱动超时后遗留的旧回复（日志里的 \code{stale replies K}）；只接受地址与本次请求相同的回复；判断成败时不看单个标签的 response 位，而看整条消息成功后返回结构的第一个字。这些规则是怎么来的，见 \ref{sec:cam-bugs} 小节的问题 1。

\paragraph{控制总线：从 BSC0 到软件 SCCB.} 读 16 位寄存器地址需要“写地址、不发 STOP、repeated START、读数据”的组合事务，\code{i2c\_write\_read()} 照 Linux \file{i2c-bcm2835.c} 实现了它（BSC 的写 FIFO 不能预填，要先启动空写事务、等 \code{TXW} 再填）。但在真机上，BSC0 对 \code{0x36} 始终 NACK，而 GPIO bit-bang 探测能得到 ACK，降速也无济于事。最终的驱动（\code{gpio-sccb-v8}）因此从第一次访问起就用 GPIO44/45 上约 25\,kHz 的软件 SCCB：STOP/START 分离读法，拒绝任何不是 \code{0x5647} 的读数，失败就断电重来，不再扫描其他地址。BSC0 仍按 100\,kHz 初始化并保留两段式中止，只是不再发事务。这条路是怎么一步步走过来的，见 \ref{sec:cam-bugs} 小节的问题 2～6。

\paragraph{延迟探测.} \code{/dev/video0} 在探测之前就已注册，所以 \code{open} 成功不代表绑定了传感器。启动时只探测一次，避免没有摄像头时拖慢开机；如果启动探测偶然失败，第一次 \code{read()} 或 \code{pattern=N} 写入会在拍照睡眠锁 \code{cam.busy} 下重新探测，最多三个完整的断电/上电周期，每次失败等 100\,ms。

\subsection{拍一帧}

用户每 \code{read()} 一次 \code{/dev/video0} 就拍一帧，返回 40 字节的帧头（magic、宽、高、行字节数、格式、长度、序号、时间戳）加打包的 RAW10 数据。帧头里的 \code{format} 同时给出 Bayer 顺序（OV5647 是 BGGR，IMX219 是 RGGB），用户程序据此去马赛克。\code{cam\_capture()} 的顺序与 Linux 相同——开始时先启动接收器再让传感器出流，结束时先停传感器再停接收器；整条调用链见图~\ref{fig:cam-calls}：

\begin{figure}[htbp]
  \centering
  \begin{tikzpicture}[x=1mm,y=1mm, font=\footnotesize,
      fn/.style={box, font=\footnotesize, minimum height=6mm, inner sep=2pt},
      st/.style={fn, text width=40mm, align=left},
      call/.style={-{Stealth[length=1.8mm]}, semithick, draw=structurecolor!70!black},
      irq/.style={call, dashed, draw=red!70!black}]
    % column backgrounds
    \foreach \x/\w/\t in {0/24/{用户态与 VFS}, 27/46/{\file{unicam.c}（桥接驱动）}, 76/30/{\file{ov5647.c}（\code{sensor\_ops}）}, 109/41/{\file{ov5647.c} 内部 → \file{i2c.c}/\file{mbox.c}}}{
      \fill[gray!5, rounded corners=2pt] (\x,2) rectangle ++(\w,84);
      \node[lab, anchor=north] at (\x+\w/2,85.5) {\t};
    }
    % user / vfs chain
    \node[fn, text width=22mm] (u)  at (12,72) {\code{camshot}\\\code{read(fd,buf,n)}};
    \node[fn, text width=20mm] (sr) at (12,60) {\code{sys\_read}\\\code{fileread}};
    \node[fn, text width=20mm] (cr) at (12,48) {\code{camread}\\持 \code{cam.busy}};
    \node[fn, text width=20mm] (cb) at (12,34) {\code{camera\_bind}\\（未绑定时）};
    \node[fn, text width=20mm] (co) at (12,10) {\code{copyout}\\帧头 + RAW10};
    \draw[call] (u) -- (sr); \draw[call] (sr) -- node[lab, right]{\code{devsw[5]}} (cr);
    \draw[call] (cr) -- (cb); \draw[call] (cb) -- (co);
    % cam_capture steps
    \node[font=\footnotesize\bfseries] at (50,76) {\code{cam\_capture()}};
    \foreach \i/\y/\t in {1/71/{\code{mbox\_set\_domain(14,1)}},
                         2/64/{\code{cam1\_clock\_on()}：\code{CM\_CAM1}},
                         3/57/{\code{ops->power\_on}},
                         4/50/{\code{unicam\_start(bus,mode)}},
                         5/43/{\code{ops->stream\_on(mode)}},
                         6/36/{\code{unicam\_service(0)} $\times N$},
                         7/29/{\code{ops->stream\_off}},
                         8/22/{\code{unicam\_stop()}},
                         9/15/{\code{cam\_power\_off()}}}{
      \node[st] (s\i) at (50,\y) {\textbf{\i}\enspace\t};
    }
    \draw[call] (cr.east) -- ++(2,0) |- (s1.west);
    \foreach \a/\b in {1/2,2/3,3/4,4/5,5/6,6/7,7/8,8/9}{ \draw[call] (s\a) -- (s\b); }
    \draw[call] (s9.south) |- (co.east);
    \node[fn, text width=20mm, draw=red!60!black] (irqn) at (12,22) {\code{devintr}\\→ \code{unicam\_irq}};
    \draw[irq] (irqn.east) -- ++(4,0) |- node[lab, above, pos=0.75, text=red!60!black]{屏蔽} (s6.west);
    % ov5647 ops
    \node[fn, text width=28mm] (pon) at (91,57) {\code{ov5647\_power\_on}};
    \node[fn, text width=28mm] (son) at (91,43) {\code{ov5647\_stream\_on}};
    \node[fn, text width=28mm] (sof) at (91,29) {\code{ov5647\_stream\_off}};
    \node[fn, text width=28mm] (pof) at (91,15) {\code{ov5647\_power\_off}};
    \draw[call] (s3) -- (pon); \draw[call] (s5) -- (son); \draw[call] (s7) -- (sof); \draw[call] (s9) -- (pof);
    % bus layer
    \node[fn, text width=37mm] (mb) at (129.5,68) {\code{mbox\_set\_expgpio(5,·)}\\\code{mbox\_set\_domain}\\→ \code{mbox\_property}};
    \node[fn, text width=37mm] (rw) at (129.5,50) {\code{ov5647\_read}/\code{ov5647\_write}\\\code{ov5647\_write\_array}（重试 3 次）};
    \node[fn, text width=37mm] (bsc) at (129.5,34) {\code{i2c\_write\_read}/\code{i2c\_write}\\硬件 BSC0};
    \node[fn, text width=37mm] (bb) at (129.5,18) {\code{i2c0\_bitbang\_write\_read}\\GPIO44/45 软件 SCCB};
    \draw[call] (pon.north east) -- ([yshift=-3mm]mb.west);
    \draw[call] (s1.east) -- (s1.east -| mb.west);
    \draw[call] (pon.east) -- (rw.west);
    \draw[call] (son.east) -- (rw.west);
    \draw[call] (sof.east) -- (rw.west);
    \draw[call] (rw) -- node[lab, right]{默认} (bsc);
    \draw[call] (rw.east) -- ++(2.5,0) |- node[lab, right, pos=0.25, align=left]{BSC\\失败后} (bb.east);
  \end{tikzpicture}
  \caption{一次 \code{read("/dev/video0")} 的调用关系。\code{unicam.c} 只通过 \code{sensor\_ops} 调用传感器驱动；传感器驱动再经 mailbox 上电、经 BSC0 或软件 SCCB 写寄存器。第 6 步每 1\,ms 轮询一次 Unicam 状态，直到 \textsc{done} 或 4\,s 超时，是目前采集帧的唯一路径，中断入口已屏蔽；第 9 步的 \code{ov5647\_power\_off} 同样经 mailbox 断电}
  \label{fig:cam-calls}
\end{figure}


\begin{enumerate}
  \item \code{mbox\_set\_domain(14, on)} 打开 UNICAM1 电源域；
  \item \code{CM\_CAM1} 时钟：PLLD\_PER（500\,MHz）$\div 5 = 100$\,MHz；若 BUSY 位不置起，就退回 19.2\,MHz 晶振并打印提示；
  \item \code{power\_on}：上电、校验 ID、进入 LP-11；
  \item 按传感器的总线参数配置 Unicam（照 Linux \code{unicam\_start\_rx()}）：lane 数、虚拟通道、数据类型、行字节数（16 的倍数），DMA 先指向一块大小为 0 的“假缓冲区”；
  \item \code{stream\_on}：依次写公共寄存器表、640$\times$480 模式表、HTS/VTS/曝光/增益、缓存的测试图案，退出软件待机，开始发送；
  \item 帧状态机：\textsc{warmup} 丢弃前 30 帧（约 0.5\,s）让传感器自带的 AEC/AGC 收敛；\textsc{armed} 在下一个帧开始时把 DMA 地址换成真缓冲区；\textsc{capture} 一开始就写回假缓冲区；帧结束后进入 \textsc{done}；
  \item \code{stream\_off} → 停 Unicam → \code{power\_off} → 关时钟和电源域，把数据从 \code{CAMDMA}（非缓存映射，物理 \code{0x07200000}）拷给用户。
\end{enumerate}

第 6 步的关键（图~\ref{fig:cam-fsm}）是：写进 \code{IBSA0/IBEA0} 的新地址\emph{要到下一个帧开始才生效}。所以真缓冲那一帧刚开始就把假缓冲写回去，这一帧会继续完整地写进真缓冲，之后的帧则落进假缓冲，缓冲区里恰好留下一整帧。假缓冲的大小编程为 0，是照搬 Linux 绕开循环缓冲模式越界写缺陷的做法。

\begin{figure}[htbp]
  \centering
  \begin{tikzpicture}[x=1mm,y=1mm, font=\footnotesize]
    \def\fw{22}\def\gap{5}
    % row labels
    \node[lab, anchor=east] at (14,40) {传感器输出};
    \node[lab, anchor=east] at (14,30) {DMA 实际写入};
    \node[lab, anchor=east] at (14,20) {\code{IBSA0} 寄存器};
    \node[lab, anchor=east] at (14,10) {\code{cam.state}};
    \foreach \i/\n in {0/29,1/30,2/31,3/32}{
      \pgfmathsetmacro{\xs}{16+\i*(\fw+\gap)}
      \draw[fill=structurecolor!15, draw=structurecolor!70!black] (\xs,36.5) rectangle ++(\fw,7) node[midway]{第 \n 帧};
      \draw[red!70!black, thick] (\xs,34) -- (\xs,46) node[above, font=\scriptsize]{FS};
      \draw[gray!70!black, thick] (\xs+\fw,34) -- (\xs+\fw,46) node[above, font=\scriptsize]{FE};
      \draw[gray!40, dashed] (\xs,4) -- (\xs,34);
    }
    % DMA target row
    \foreach \i/\t/\c in {0/假缓冲/gray!12, 1/假缓冲/gray!12, 2/真缓冲 \code{FRAME\_PA}/orange!30, 3/假缓冲/gray!12}{
      \pgfmathsetmacro{\xs}{16+\i*(\fw+\gap)}
      \draw[fill=\c, draw=gray!60] (\xs,26.5) rectangle ++(\fw,7) node[midway]{\t};
    }
    % IBSA row
    \draw[fill=gray!8, draw=gray!60] (16,16.5) rectangle (43,23.5) node[midway]{\code{DUMMY\_PA}};
    \draw[fill=orange!20, draw=gray!60] (43,16.5) rectangle (70,23.5) node[midway]{\code{FRAME\_PA}};
    \draw[fill=gray!8, draw=gray!60] (70,16.5) rectangle (119,23.5) node[midway]{\code{DUMMY\_PA}};
    % state row
    \foreach \a/\b/\t in {16/43/{\textsc{warmup}}, 43/70/{\textsc{armed}}, 70/92/{\textsc{capture}}, 92/119/{\textsc{done} → 停流}}{
      \draw[fill=structurecolor!8, draw=structurecolor!60!black] (\a,6.5) rectangle (\b,13.5) node[midway]{\t};
    }
    % annotations
    \node[lab, align=left, anchor=north west, text width=44mm] (n1) at (2,-1) {\textbf{A} 第 30 个 FS（或 1.5\,s 期限）：写入 \code{FRAME\_PA}，进入 \textsc{armed}；本帧仍写假缓冲};
    \node[lab, align=left, anchor=north west, text width=44mm] (n2) at (50,-1) {\textbf{B} 下一个 FS：新地址生效，第 31 帧写进真缓冲；处理程序立刻写回 \code{DUMMY\_PA}，进入 \textsc{capture}};
    \node[lab, align=left, anchor=north west, text width=40mm] (n3) at (98,-1) {\textbf{C} 第 31 帧的 FE：真缓冲里正好一整帧，进入 \textsc{done}；第 32 帧落进假缓冲};
    \node[circle, draw=red!70!black, inner sep=0.8pt, font=\scriptsize\bfseries, fill=white] at (43,54) {A};
    \node[circle, draw=red!70!black, inner sep=0.8pt, font=\scriptsize\bfseries, fill=white] at (70,54) {B};
    \node[circle, draw=red!70!black, inner sep=0.8pt, font=\scriptsize\bfseries, fill=white] at (92,54) {C};
  \end{tikzpicture}
  \caption{帧状态机的时序。\code{IBSA0/IBEA0} 的新值要到下一个帧开始（FS）才被硬件锁存，所以“写真地址”和“写回假地址”都比它们生效的那一帧早一个 FS}
  \label{fig:cam-fsm}
\end{figure}

\paragraph{为什么用轮询而不是中断.} 设备树说 CSI1 是 GPU IRQ 39，但这个编号没有在真机上验证过；第一次打开它时 CPU0 陷入了清不掉的电平中断（\ref{sec:cam-bugs} 小节问题 7）。所以目前屏蔽 IRQ 39，用 \reg{CNTVCT\_EL0} 计一个 4 秒期限、每 1\,ms 轮询一次 Unicam 状态。现在真机上的日志是：

\begin{txtcode}
unicam: frame captured by polling; CSI1 IRQ 39 intentionally masked
\end{txtcode}

它表示 sensor → CSI-2 → Unicam → DMA 的整条路径已经打通，只是 IRQ 路由还留待单独验证。

\subsection{用户程序 \code{camshot}}

\code{camshot} 用纯整数运算完成图像处理（用户程序以 \code{-mcpu=cortex-a53+nofp} 编译）：解包 RAW10（每 4 个像素占 5 字节），减黑电平 16，BGGR 双线性去马赛克，灰度世界白平衡（增益限制在 0.5～4），把绿色通道的第 99 百分位拉到满幅，查表做 gamma 2.2 校正，写成 24 位 BMP。\code{-t N} 让传感器输出内部测试图案（1 彩条、2 色块、3 随机），这时改用固定的 1.0 白平衡和固定电平，不做统计——测试图案本身是校准过的，正好单独验证传感器、CSI-2、DMA、转换和写盘；\code{-s} 输出 320$\times$240 半尺寸。

640$\times$480 的 BMP 有 900\,KB，超过原生 xv6 文件系统单文件 268\,KB 的上限（\code{MAXFILE}$=12+256$ 块），所以默认写到以读写方式挂载的 FAT32 启动分区 \code{/boot/camera.bmp}，关机后把 SD 卡插到电脑上就能看到。为了让慢速写盘不被误判成“卡死”，\code{camshot} 在每个阶段打印一行（收到帧、解包、统计、打开输出、写头、转换），BMP 按 16 行一批写入（半尺寸图从 240 次 FAT32 写降到 15 次），启动时打印的版本串 \code{converter fixed-pattern-v2 fat-progress batch16} 用来确认 SD 卡上的程序确实已经更新。

\subsection{一次成功拍照的日志解读}
\label{sec:cam-log}

下面是 v8 驱动在真机上一次成功拍照的完整串口输出（\code{camshot -t 1 -s -o /boot/cam02.bmp}：传感器彩条、半尺寸、写到启动分区）。左边的 (1)～(6) 是为了讲解加的阶段编号，与下文 \ding{172}～\ding{177} 一一对应：

\begin{txtcode}
(1) camera: driver gpio-sccb-v8 unicam-poll-v2
    ov5647: CAM_GPIO0 direction=output polarity=normal state=1
    i2c0: GPIO44/45 ALT1, 100000 Hz (core 250000000), BCM2837 two-stage abort
    ov5647: probing with 25 kHz GPIO SCCB; BSC0 bypassed
    ov5647: software SCCB power-on chip-id=5647
    ov5647: sensor outputs enabled
    ov5647: power-on complete, lanes in LP-11
    camera: OV5647 on CSI1 (2 lanes) -> /dev/video0, 640x480 RAW10
$ /bin/camshot -t 1 -s -o /boot/cam02.bmp
(2) camshot: converter fixed-pattern-v2 fat-progress batch16
    camshot: capturing...
    ov5647: CAM_GPIO0 direction=output polarity=normal state=1
    ...（与 (1) 相同的上电五行）
    ov5647: power-on complete, lanes in LP-11
(3) unicam: sensor powered, receiver setup begins
    unicam: receiver enabled, starting sensor stream
    ov5647: programming common registers (51)
    ov5647: programming mode registers (29)
    ov5647: programming controls (11)
    ov5647: streaming started
    unicam: sensor streaming, waiting for frame
(4) unicam: frame captured by polling; CSI1 IRQ 39 intentionally masked
(5) camshot: frame received 640x480, 384040 bytes
    camshot: unpacking RAW10...
    camshot: sensor test pattern; fixed WB/level
(6) camshot: preparing output /bovfs: boot create begin path=/boot/cam02.bmp sub=/cam02.bmp mode=601
    ot/cam02.bmp (230454 bytes)...
    vfs: boot create stat complete exists=0
    vfs: boot create directory-entry begin
    fat32: create lookup begin path=/cam02.bmp
    fat32: create lookup complete exists=0
    fat32: create reserve begin slots=2
    fat32: create reserve complete
    vfs: boot create directory-entry complete
    vfs: boot create truncate begin / complete
    vfs: boot create verify-stat begin / complete
    vfs: boot create openwrite begin / complete
    camshot: output opened; writing BMP header...
    camshot: converting and writing /boot/cam02.bmp (230454 bytes)...
    camshot: 320x240 frame 31, WB R x1.0 B x1.0, level x1.0 -> /boot/cam02.bmp
\end{txtcode}

\paragraph{\ding{172} 启动时的探测与绑定.} 这是 \code{main()} → \code{camera\_init()} → \code{camera\_bind()} → \code{ov5647\_power\_on()} 的输出。第一行是驱动版本串，用来确认 SD 卡上跑的确实是新内核。第二行来自 mailbox 的回读：CAM\_GPIO0 是输出、极性正常、电平为 1，模块的稳压器已经打开。第三行由 \code{i2c0\_init\_camera\_pins()} 打印：BSC0 仍按 100\,kHz 初始化（分频 $250\,\text{MHz}/100\,\text{kHz}=2500$），出错时采用“两段式中止”——先保持 \code{I2CEN} 清 FIFO 等 \code{TA} 归零，再关控制器（\ref{sec:cam-probe} 小节）。但紧接着的一行说明这次根本不用它：v8 直接用 25\,kHz 软件 SCCB。读到 \code{5647} 之后，驱动写一小段寄存器打开传感器的输出引脚，再写停流序列让 MIPI 链路停在 LP-11 空闲态——Unicam 稍后要在这个干净的空闲态上做模拟复位和校准。最后一行表示绑定成功；绑定之后 \code{camera\_bind()} 立即 \code{power\_off()}，模块在两次拍照之间是断电的。

\paragraph{\ding{173} 每次拍照都重新上电.} 用户 \code{read()} 进入 \code{camread()}，持 \code{cam.busy} 调 \code{cam\_capture()}，后者再次 \code{ops->power\_on()}，所以上电五行原样再出现一次。这是“每次拍照完整上电、预热、断电”这个设计的直接体现，也是一次拍照要一秒左右的主要原因。

\paragraph{\ding{174} 先接收器、后出流.} \code{unicam: receiver enabled} 在 \code{starting sensor stream} 之前，顺序与 Linux 一致。随后 \code{ov5647\_stream\_on()} 写三张表：公共寄存器 51 项、640$\times$480 模式 29 项、曝光增益等控制 11 项，再加上虚拟通道、测试图案、退出待机、MIPI 控制等几次读改写。它们全部经 25\,kHz 软件 SCCB：一次 16 位地址的写是 4 个字节、约 36 个时钟，约 1.5\,ms，九十多次写加起来约 0.15 秒——慢，但只在开流时发生一次。

\paragraph{\ding{175} 一帧到手.} \code{cam\_capture()} 每 1\,ms 调一次 \code{unicam\_service(0)}，状态机经过 \textsc{warmup}（丢 30 帧）、\textsc{armed}、\textsc{capture}，在帧结束时进入 \textsc{done}（图~\ref{fig:cam-fsm}）。“\code{intentionally masked}”说明 IRQ 39 是有意屏蔽的，不是丢了中断。

\paragraph{\ding{176} 数字对得上.} \code{read()} 返回 \code{384040} 字节：40 字节帧头加 $640\times480\times10/8=384000$ 字节 RAW10。测试图案模式下 \code{camshot} 用固定的 1.0 白平衡和电平，所以下一行直接是 \code{fixed WB/level}。

\paragraph{\ding{177} 创建输出文件.} 这一段同时有用户态和内核两个来源的输出，所以第一行被“切开”了：\code{camshot} 的 \code{printf} 在用户态逐字符 \code{write(1, \&c, 1)}，打到 \code{"/bo"} 时进程陷入 \code{open()}，VFS 的内核 \code{printf} 插了一整行进来，返回用户态后才接着打出 \code{"ot/cam02.bmp (230454 bytes)..."}。传给内核的路径并没有坏，\code{path=/boot/cam02.bmp} 写得清清楚楚。之后各行对应 \code{vfsopen()} 的创建流程：

\begin{itemize}
  \item \code{mode=601}：$\code{0x001}\,|\,\code{0x200}\,|\,\code{0x400}$，即 \code{O\_WRONLY|O\_CREATE|O\_TRUNC}；
  \item \code{stat complete exists=0}：\code{fat32statpath()} 正常返回，文件还不存在；
  \item \code{fat32: create lookup}：\code{fat32createfile()} 内部再查一次，确认不重名；
  \item \code{reserve begin slots=2}：\code{cam02.bmp} 本身就是合法的 8.3 名字，不需要 LFN，所以只要 2 个目录项——一个短名项，加一个全 0 的结束标记；根目录尾部空位够用，不必扩展目录簇；
  \item \code{truncate}：对刚建的空文件截断，首簇为 0，什么都不用释放；
  \item \code{verify-stat}、\code{openwrite}：重新确认目录项，并把首簇、大小装进 vnode 的 \code{fat32\_file}，供之后的追加写使用。
\end{itemize}

最后一行的 \code{230454} 是 54 字节 BMP 头加 $320\times240\times3=230400$ 字节像素；\code{frame 31} 是帧头里的序号，即帧结束计数——30 帧预热之后的第 31 帧。输出是 320$\times$240，是因为 \code{-s} 把每个 2$\times$2 Bayer 块合成一个像素，不是 Unicam 丢了一半分辨率。

\paragraph{这份日志证明了什么.} 对照 \ref{sec:cam-hw} 小节的四条通路：mailbox 供电（CAM\_GPIO0 回读）、SCCB 控制（读到 \code{5647}、三张表写完）、CSI-2 数据面与 DMA（\code{frame captured}、长度正确）、事件（轮询完成，IRQ 有意屏蔽）全部打通；用户态转换和 FAT32 写盘也都完成。它同时也是 \ref{sec:cam-bugs} 小节问题 5、6（软件 SCCB 读 ID 不稳）和问题 11（新建文件停在 \code{directory-entry begin}）修好之后的验证记录。

\subsection{真机调试记录}
\label{sec:cam-bugs}

摄像头从“QEMU 上能探测到没有摄像头”到“真机写出第一张 BMP”，经历了下面十一个问题。前七个属于摄像头本身的控制面和数据面；后四个是摄像头作为第一个同时用到大块 DMA、长时间用户态计算、FAT32 大文件写和 Wi-Fi 并发的程序，逼出来的系统问题。每个问题都按同样的顺序记录：先是串口上看到的日志，再是分析，最后是解决方案。日志中凡是标注“格式”的，是按驱动源码里的 \code{printf} 格式写出的形式，其余是真机原样输出。

\subsubsection{问题 1：mailbox 打开摄像头电源失败}

\paragraph{日志输出.}
\begin{txtcode}
ov5647: firmware GPIO 133 (camera power) failed, firmware rev N          （格式）
ov5647:   SET_GPIO_CONFIG: <结果> (status .., tag response .. = done D len L,
          first word W, stale replies K)
ov5647:   SET_GPIO_STATE:  <结果> (...)
ov5647: no reply at i2c0 0x36 (GPIO44/45), ...
\end{txtcode}

\paragraph{问题分析.} 传感器的稳压器由固件扩展 GPIO 133（CAM\_GPIO0）使能，ARM 只能请 VideoCore 固件代办；固件说“不支持”，传感器就没电，I2C 上自然读不到 \code{0x36}。第一行的结果字段可以这样判读：\code{no reply} 是固件根本没回；\code{message rejected} 是整条消息被拒；\code{done 1 len 0} 是标签被应答但没有数据（QEMU 就是这样）；\code{len 8} 且 \code{first word 133} 是固件认识这个标签、但认为 GPIO 133 无效。对照 Linux 的 \code{rpi\_firmware\_property()} 和 \file{gpio-raspberrypi-exp.c}，旧的 mailbox 代码有三处不同：标签的请求长度字填了数据长度，Linux 填 0；发送前不清 FIFO，而 SD、USB 驱动各有一份私有 mailbox 代码，它们超时后留下的回复会被当成本次的回复读走；收到回复不核对地址。另外，旧代码在复制数据之前就因为单个标签的 response 位没置而判失败，Linux 并不这样判。最后一列 \code{stale replies} 就是为验证第二点加的：它统计发送前清掉了多少条别人的回复。

\paragraph{解决方案.} 把各处私有代码收拢成带锁的 \file{kernel/mbox.c}，按 Linux 的约定重写：请求长度字填 0；发送前清空 FIFO；只接受地址与本次请求相同的回复；以“整条消息成功”加“返回结构第一个字为 0”判断 GPIO 操作成败。电源控制也按 Linux 的做法：先 \code{GET\_GPIO\_CONFIG} 读出并保留极性，再设置，最后 \code{GET\_GPIO\_STATE} 回读。修复后的日志是 \code{ov5647: CAM\_GPIO0 direction=output polarity=normal state=1}。如果某个固件确实不处理这个标签，兜底方案是自制 \file{dt-blob.bin}，让固件启动时就把 p133 置为输出高。

\subsubsection{问题 2：BSC0 对 \code{0x36} 始终 NACK}

\paragraph{日志输出.}
\begin{txtcode}
ov5647: no reply at i2c0 0x36 (GPIO44/45), BSC raw=90000131 flags=131 \
        now=... div=2500 fsel=55 levels=3
（更早的内核把 raw 打成 -6ffffecf）
\end{txtcode}

\paragraph{问题分析.} 供电已经正常，BSC0 却在 repeated-START 和 STOP/START 两种读法下都报错。逐个字段看：\code{flags} 的 \code{0x100} 是 \code{ERR}（从机没应答），\code{0x200} 才是时钟拉伸超时；\code{div=2500} 是 250\,MHz$/$100\,kHz；\code{fsel=55} 表示 GPIO44/45 都是 ALT1；\code{levels=3} 表示 SDA、SCL 都是高电平，总线空闲、没有被谁拉住。控制器确实发出了事务，从机确实没应答——但仅凭这一行，还分不清是“传感器不在”还是“控制器这条路不通”。另外两个副产品：真 BCM2837 在 \code{TA}（transfer active）仍为 1 时撤掉 \code{I2CEN}，状态机可能冻结，QEMU 不模拟这一点；内核 \code{printf} 的 \code{\%x} 把最高位为 1 的值当有符号数打印，\code{0x90000131} 成了 \code{-6ffffecf}。

\paragraph{解决方案.} 出错时采用两段式中止：先保持 \code{I2CEN}、清 FIFO，等 NACK+STOP 让 \code{TA} 归零，再关控制器、清粘滞状态位；\code{\%x} 改为无符号。为了回答“传感器在不在”，增加一个与 BSC 完全无关的 bit-bang 探测（问题 3）。

\subsubsection{问题 3：bit-bang 探测时而 ACK、时而 NACK}

\paragraph{日志输出.}
\begin{txtcode}
camera: GPIO bitbang probe 0x36=ACK 0x10=NACK 0x1a=NACK        （某次启动，格式）
camera: GPIO bitbang probe 0x36=NACK 0x10=NACK 0x1a=NACK       （同一模块，另一次）
\end{txtcode}

\paragraph{问题分析.} bit-bang 探测把 GPIO44/45 临时切成开漏 GPIO，只发一个地址字节、读第九个时钟上的 ACK，不读写任何寄存器。它本该给出确定的答案，同一个模块的结果却不稳定，问题出在诊断代码自己：第九个时钟调用的是通用的 \code{bb\_release()}，它释放 SDA 之后会等 SDA 回到高电平——这对数据位是对的，对 ACK 却正好相反，ACK 的定义就是从设备把 SDA 拉低。于是每次都白等 100\,µs 超时，再去采样时从设备可能已经松开了 SDA。

\paragraph{解决方案.} ACK 周期只把 SDA 切成输入，在 SCL 高电平的中间直接采样。修正以后结果稳定为 \code{0x36=ACK}。这个结论很有分量：bit-bang ACK 而 BSC NACK，说明传感器、地址、供电、上拉和排线都没问题，故障只在 BSC0 的复用或时序这一段。同时规定无论 bit-bang 结果如何，都继续执行问题 4 的恢复扫描。

\subsubsection{问题 4：BSC 四档降速全部失败}

\paragraph{日志输出.}
\begin{txtcode}
ov5647: BSC retry 100000 Hz failed status=... now=... div=...   （格式）
ov5647: BSC retry 50000 Hz failed status=...
ov5647: BSC retry 25000 Hz failed status=...
ov5647: BSC retry 10000 Hz failed status=...
\end{txtcode}

\paragraph{问题分析.} 既然传感器在，就怀疑速率或寄存器写入顺序：依次用 100、50、25、10\,kHz 重读芯片 ID，每档都按 Linux 的比例同时更新 \code{BSC\_DIV} 和边沿采样延迟 \code{BSC\_DEL}（不继承固件留下的值），并在写完地址、长度、FIFO 后加 \code{dsb sy}，保证配置在 START 之前到达控制器。四档全部是 \code{ERR}，速率和 MMIO 写序都被排除了；剩下的只能是 BSC0 到 GPIO44/45 这条路由本身，或者固件对 BSC0 的占用。继续在这里花时间，收益已经很低。

\paragraph{解决方案.} 绕开 BSC0：用 GPIO 完成完整的软件 SCCB 事务（写寄存器地址、STOP、START、读一个字节），读到 \code{0x5647} 后打印 \code{using GPIO software SCCB; BSC0 bypassed}，此后传感器的所有读写都走这条路。寄存器表里 \code{0x0103=1} 会让传感器软件复位、期间不应答，写完这一项等 10\,ms；普通寄存器写最多重试三次。软件 SCCB 慢，但只用来写一百来个寄存器，图像走 CSI-2，不受影响。

\subsubsection{问题 5：软件 SCCB 偶发读出 \code{0x3631}，随后 BUS-STUCK}

\paragraph{日志输出.}
\begin{txtcode}
ov5647: software SCCB chip-id=3631 attempt=1                   （格式）
...
camera: full GPIO I2C scan bus stuck at 0x..
\end{txtcode}

\paragraph{问题分析.} 地址被 ACK 了，读回的却是 \code{0x3631}——它不是任何传感器的 ID，事务层面成功、数据位却采错了，说明时序处于边缘。旧代码随后的两个动作让情况更糟：每次软件事务前无条件送 9 个 SCL 恢复脉冲；读到错误 ID 后立即扫描整个 7 位地址空间。对一个状态已经不稳的模块，这些都是额外的扰动，扫描最终把总线卡住（BUS-STUCK）。

\paragraph{解决方案.}（\code{gpio-sccb-v7}）软件时钟降到约 25\,kHz（半周期约 20\,µs）；任何不是 \code{0x5647} 的完整读数都被拒绝重读（日志 \code{rejected unstable chip-id=..}）；恢复脉冲只在发现 SDA 被拉低时才送，最多 9 个，然后统一发 STOP；只要 \code{0x36} 已经完成过寄存器事务，就不再做全地址扫描，而是打印 \code{address 0x36 responds but chip-id is unstable; power-cycle retry required} 后直接断电重来；上电前 CAM\_GPIO0 先拉低 20\,ms 让模块放电，置高后等 100\,ms 再访问。

\subsubsection{问题 6：\code{0x36} ACK，但三次读 ID 都失败}

\paragraph{日志输出.}
\begin{txtcode}
ov5647: software SCCB ID read retry 1/3 (probe=ACK)            （格式）
ov5647: software SCCB ID read retry 2/3 (probe=ACK)
ov5647: software SCCB ID read retry 3/3 (probe=ACK)
ov5647: address 0x36 ACK but chip-id read failed; power-cycle retry required (fsel=.. levels=..)
\end{txtcode}

\paragraph{问题分析.} v7 在真机上仍会出现这种情况。看调用顺序就明白了：每次上电，驱动仍然先用 BSC0 读 ID，失败几次之后才切到 GPIO。问题 4 已经证明这些 BSC 事务注定失败，而它们留下的半截事务会让传感器的 SCCB 接口进入不确定状态，紧接着的软件读取就跟着失败。

\paragraph{解决方案.}（\code{gpio-sccb-v8}）从第一次读 \code{0x300a/0x300b} 起就直接用 25\,kHz 软件 SCCB，BSC0 只初始化、不再发事务，日志变为 \code{probing with 25 kHz GPIO SCCB; BSC0 bypassed} 和 \code{software SCCB power-on chip-id=5647}。如果 \code{0x36} 已 ACK 但读寄存器仍失败，本轮立即断电，交给 \code{camera\_bind()} 做下一次完整的电源周期（用户首次 \code{read()} 时最多三次，每次间隔 100\,ms），不再探测其他地址。\ref{sec:cam-log} 小节的日志就是 v8 的样子。

\subsubsection{问题 7：开始出流后整机无响应}

\paragraph{日志输出.}
\begin{txtcode}
unicam: receiver enabled, starting sensor stream
ov5647: streaming started
unicam: sensor streaming, waiting for frame
（此后再无任何输出，4 秒超时也没有触发）
\end{txtcode}

\paragraph{问题分析.} 第一次真机出流时打开了设备树给出的 CSI1 中断 GPU IRQ 39。“4 秒超时也没触发”是关键线索：超时是按 xv6 的 \code{ticks} 计的，而 \code{ticks} 由 CPU0 的通用定时器中断推进；如果它不走了，说明 CPU0 一直在处理别的中断。IRQ 39 是电平触发的，若这个编号实际不对应 Unicam，或者处理程序清不掉它的源，CPU0 一从中断返回就立刻再次进入，定时器中断永远得不到服务——超时机制本身依赖于它要监督的那条路径，所以失效了。

\paragraph{解决方案.} bring-up 阶段屏蔽 IRQ 39，改用与中断无关的 \reg{CNTVCT\_EL0} 计 4 秒期限，每 1\,ms 轮询一次 Unicam 的 \code{STA}/\code{ISTA}；无论成败都能返回，失败时打印 FS/FE 计数、\code{STA}、\code{IBWP}。成功时的日志是 \code{unicam: frame captured by polling; CSI1 IRQ 39 intentionally masked}——“intentionally”说明中断是有意屏蔽的，不是丢了。以后再验证中断时，也只在拍照期间打开，连续 100 次没有发现 Unicam 状态就自动关掉，并始终保留这个有界轮询作兜底。

\subsubsection{问题 8：转换期间偶发 \code{SDIO CMD52 failed}}

\paragraph{日志输出.}
\begin{txtcode}
camshot: calculating white balance and level...
arasan-sdio: request CMD52 failed status=df0001 irq=100 c1=e0807
\end{txtcode}

\paragraph{问题分析.} 两行挨在一起，很容易得出“用户程序抢占了 Wi-Fi 工作线程”的结论，但这站不住：\code{camshot} 和 \code{brcmf0\_wq} 都是 \code{SCHED\_OTHER}、nice 0，用户态并不自动低于内核线程，SDIO 中断上半部也随时能打断用户程序。真正的线索在数值里：\code{status} 的 bit~0 是 \code{CMD\_INHIBIT}，\code{irq=100} 只有卡中断、没有本次的 \code{CMD\_DONE}——上一个事务还没结束，另一个就开始写寄存器了。原来 \code{sdio\_cmd52()}/\code{sdio\_cmd53()} 直接调用 Arasan 的 \code{request()}，而 Wi-Fi 的状态 work、NAPI 收发和 BCDC 控制请求可能在不同核上同时运行，共用同一组 \code{ARG1/CMDTM/INTERRUPT/DATA} 寄存器。

\paragraph{解决方案.} 给 \code{struct mmc\_host} 加睡眠锁 \code{request\_lock}，所有 CMD52/CMD53 都经 \code{mmc\_request\_sync()} 串行化。锁的粒度是物理控制器而不是某个 function 或某块网卡，对应 Linux MMC core 的 host claim；用睡眠锁是因为一次请求可能要等命令、数据和超时（\ref{sec:wifi-flow} 小节）。启动枚举时还没有进程，这时只有启动 CPU 发请求，直接执行、不取锁。此前为 \code{camshot} 加的 \code{sched\_yield()} 撤回——它掩盖不了并发访问。

\subsubsection{问题 9：执行不访存的指令却报数据中止}

\paragraph{日志输出.}
\begin{txtcode}
usertrap(): unexpected esr=0x92000004 ec=0x24
            elr=0x0000000000000b40 far=0xf3910003fda9a57b
\end{txtcode}

\paragraph{问题分析.} 更新 \code{/bin/camshot} 后，程序偶尔在第一行输出之前就退出。EC$=$\code{0x24} 是来自 EL0 的数据中止，可用当前 ELF 反汇编，\code{ELR=0xb40} 是 \code{strncpy()} 里一条纯寄存器的 \code{sub}，不可能访存；\code{FAR} 又是一个随机指针。CPU 实际执行的并不是磁盘上这个 ELF 的指令。\code{exec()} 通过内核高地址把新程序写进物理页，\code{uvmsync\_icache()} 做了 \code{dc cvau}，最后却用 \code{ic iallu}——它只作废\emph{执行它的那个核}的指令缓存。进程在一个核上装好，迁移到另一个核运行，那里的 I-cache 仍命中同一低地址上旧程序的指令：执行的是旧代码，用的却是新程序的数据和栈。

\paragraph{解决方案.} 改为广播到 Inner Shareable 域的 \code{ic ialluis}，前面的逐行 \code{dc cvau} 和 \code{dsb ish} 保留（第~\ref{chap:mmu} 章故障二）。\code{exec()}、\code{uvminit()}、\code{uvmcopy()} 都经过这个函数。排查这类问题时，要用实际装进 SD 卡的那个 ELF 去反查 \code{ELR}，不能用另一次构建的符号文件。

\subsubsection{问题 10：输出到 \code{fixed WB/level} 后长时间停顿}

\paragraph{日志输出.}
\begin{txtcode}
camshot: frame received 640x480, 384040 bytes
camshot: unpacking RAW10...
camshot: sensor test pattern; fixed WB/level
（长时间没有新输出）
\end{txtcode}

\paragraph{问题分析.} 看起来像转换死循环，但测试图案模式下转换并不做统计，下一步其实是以 \code{O\_TRUNC} 打开已存在的 \code{/boot/camera.bmp}。截断要释放旧文件的簇链，旧的 \code{fat\_free\_chain()} 对每个簇都读一次 FAT、再对两份 FAT 各做一次读改写，覆盖一个 320$\times$240 的 BMP 就要五百多次 SD 扇区事务，而且每次都是 PIO 轮询。

\paragraph{解决方案.} \code{fat\_free\_chain()} 缓存当前 FAT 扇区，在其中批量清除属于这条链的项，换扇区时才把整扇区写到每份 FAT（\ref{sec:fat32} 节）。同时 \code{camshot} 在打开输出文件、写 BMP 头和开始转换之前各打印一行阶段日志，BMP 每 16 行写一次，并用版本串 \code{converter fixed-pattern-v2 fat-progress batch16} 确认 SD 卡上的程序已经更新——以后再停住，可以直接看出停在哪一步。

\subsubsection{问题 11：新建文件停在 \code{directory-entry begin}}

\paragraph{日志输出.}
\begin{txtcode}
camshot: sensor test pattern; fixed WB/level
camshot: preparing output /boot/cavfs: boot create begin path=/boot/cam02.bmp sub=/cam02.bmp mode=601
m02.bmp (230454 bytes)...
vfs: boot create stat complete exists=0
vfs: boot create directory-entry begin
（此后再无输出）
\end{txtcode}

\paragraph{问题分析.} 先排除两个干扰：\code{/boot/cavfs: ... m02.bmp} 是用户态逐字符输出与内核 \code{printf} 交错的结果，路径本身没坏（\ref{sec:cam-log} 小节）；\code{exists=0} 说明 \code{stat} 已经正常返回，阻塞点在 \code{fat32createfile()} 里。这和问题 10 不是同一个问题：\code{exists=1} 且停在 \code{truncate begin} 才是释放旧簇链慢。第一个猜测是“根目录要扩展，簇分配器扫描太慢”，于是先给 \code{fat\_alloc\_cluster()} 加了按扇区缓存和 FSInfo 提示，并加了 \code{fat32: extending root directory} 日志——真机上仍停在同一行，连新日志都没出现，说明根本没走到分配器。

真正的原因是内核栈。每个进程的内核栈只有 4\,KiB，旧的 \code{fat32createfile()} 在栈上放了一个路径匹配结构、22 个目录槽和 256 字节的文件名，又嵌套调用 \code{fat\_find\_path()}，后者再放 256 字节文件名、520 字节 LFN 字符数组和 20 个目录槽。光这些数组就接近 3.5\,KiB，加上 VFS、系统调用的栈帧和保存的寄存器，越过了 4\,KiB。\code{stat} 只用到其中一层，所以没事。至于为什么是“停住”而不是一条报错：内核栈之间隔着不映射的保护页（第~\ref{chap:mmu} 章），越界先落进保护页、触发 EL1 数据中止；而 EL1 的异常入口要在同一个 \code{SP\_EL1} 上压 272 字节保存现场（\ref{sec:irq} 节），这一压又落在保护页里，异常入口自己再次中止，CPU 在异常向量里反复进出，串口上什么也不会出现。

\paragraph{解决方案.} 把这些大数组移出栈，放进 FAT32 私有的静态工作区 \code{fat\_work}。所有 FAT32 操作本来就持 \code{fat32\_lock} 串行执行，一个工作区就够，不必加大每个进程的内核栈。反汇编确认栈帧降到 \code{fat32createfile} 464 字节、\code{fat\_find\_path} 176 字节。create 路径上加了更细的阶段日志（\code{create lookup}、\code{create reserve}、\code{extending root directory}），可以直接区分是查找、保留目录项还是分配簇出了问题；之前对分配器的优化也保留下来。修复后的完整日志见 \ref{sec:cam-log} 小节。

\paragraph{小结.} 回头看这十一个问题，有效的做法是一样的：让每一层都有自己的日志，日志里带上能判读的原始数值（\code{flags}、\code{fsel}、\code{levels}、\code{status}、\code{esr}、\code{exists}），然后只把日志实际证明了的东西当作结论。几次走弯路，都是因为把“两行日志挨在一起”当成了因果（问题 8），或者先修了猜测中的原因而没有确认程序是否真的走到那里（问题 11）。

\subsection{当前状态}

真机（Pi 3B+，Camera v1.3）上，v8 驱动从第一次访问起就用 25\,kHz 软件 SCCB，稳定读到芯片 ID \code{0x5647} 并绑定到 \code{/dev/video0}；Unicam 以轮询方式拿到完整帧，\code{camshot -t 1 -s -o /boot/cam02.bmp} 能把彩条写成 BMP（\ref{sec:cam-log} 小节）。还没有完成的工作有：

\begin{itemize}
  \item 验证 CSI1 中断号并打开中断路径（保留轮询兜底）；
  \item 查清 BSC0 硬件读不到芯片 ID 的原因（控制器路由还是固件占用），以便退回硬件 I2C；
  \item 更多模式：Linux 的 1296$\times$972 等模式每帧超过 1.5\,MB，需要扩大 \code{CAMDMA}，桥接驱动不用改；目前也还没有选择模式的接口；
  \item 每次拍照都完整地上电和预热（约 1 秒），没有连续取流和 \code{poll()}，也没有手动曝光；
  \item \code{arasan\_sdio.c}、\code{dwc2.c}、\code{sd.c} 里仍各有一份私有 mailbox 代码，应统一到 \file{mbox.c}；
  \item \file{config.txt} 不能打开固件的摄像头支持（\code{start\_x=1}、\code{camera\_auto\_detect=1}），否则固件会占用同一个传感器和 I2C。
\end{itemize}

\begin{remark}
  摄像头驱动是全书最能体现“分层”价值的例子：mailbox、I2C、时钟、电源域、CSI-2 接收器、DMA、帧状态机、字符设备、用户态图像处理，每一层都可以单独验证。真机调试也正是这样一层层推进的：先看 mailbox 回读，再用 bit-bang 判断 I2C 故障在控制器内外，再用传感器内部彩条（\code{camshot -t 1}）验证 CSI-2 和 DMA，最后才看真实画面。
\end{remark}

\chapter{用户进程、线程与 libc}
\label{chap:user}

\begin{introduction}[本章内容]
  \item 用户程序的形态与系统调用桩
  \item 从 \code{init} 到 shell
  \item 用户程序一览
  \item 用户线程：\code{clone}、\code{texit}、\code{tjoin}
  \item libc 与 POSIX 源码兼容层
  \item 在 xv6 里编译 C 程序
  \item 一个大型用户程序：\code{sshd}
\end{introduction}

\section{用户程序的形态}

\subsection{为什么需要系统调用}

进程隔离的第一步，是把所有用户程序放到特权最低的 EL0：在这一级，程序不能修改页表基址、不能屏蔽中断、不能访问没有映射给它的内存，也就无法改写内核为隔离所做的任何设置。否则，其他一切隔离手段都没有意义。

可是用户程序总要打印字符、读文件、发网络包，而这些都需要直接操作设备。我们不能让每个程序直接去写 UART 寄存器，而是由内核提供一组受控的入口——系统调用。它不能是普通的函数调用，因为每次调用都要把执行级别从 EL0 提升到 EL1，而提升级别只有一条路：产生异常。所以每个系统调用本质上都是一次\emph{同步异常}：用户程序准备好参数，执行 \code{svc} 指令；硬件把它带到 EL1 的异常向量；内核检查参数、完成请求，然后执行普通的异常返回，从 \code{svc} 的下一条指令继续。进程隔离的第三个方面——每个进程有自己独立的内存视图——由页表完成（第~\ref{chap:mmu} 章）。


\subsection{链接与装载}

xv6 的用户程序是静态链接的 ELF64 文件，链接在虚拟地址 0，入口是 \code{main}：

\begin{makecode}
_%: %.o $(ULIB)
	$(LD) $(LDFLAGS) -N -e main -Ttext 0 -o $@ $^
\end{makecode}

\code{-N} 让代码段和数据段放在同一个可读写可执行的段里，\code{exec()} 只需处理一个 \code{PT\_LOAD}。没有动态链接，也没有 C 运行时启动代码：\code{exec()} 把 \code{argc}、\code{argv} 放进 \code{x0}、\code{x1}，把 PC 设为入口地址，\code{main} 返回后由用户库调用 \code{exit()}——实际上 xv6 的程序习惯在 \code{main} 末尾显式调用 \code{exit(0)}。

\subsection{系统调用桩}

\file{user/usys.pl} 为每个系统调用生成三条指令：

\begin{asmcode}
.global write
write:
 mov x7, #SYS_write      // 系统调用号放 x7
 svc #0                  // 陷入 EL1
 ret                     // 返回值已在 x0
\end{asmcode}

参数按 AArch64 调用约定已经在 \code{x0}～\code{x5} 里，所以桩函数什么都不用搬。内核侧 \code{syscall()} 用 \code{trapframe->x7} 索引函数表，返回值写回 \code{trapframe->x0}。目前共有 70 多个系统调用：除了 xv6 原有的 21 个，还有同步对象、网络（UDP、ICMP、TCP socket、epoll）、挂载、TTY 会话、线程、调度属性、高精度时钟等。

\begin{note}
  系统调用号放在 \code{x7}，这是本项目沿用的约定；Linux AArch64 用的是 \code{x8}。所以 Linux 的二进制或汇编系统调用包装不能直接拿来用，第~\ref{sec:posix} 节的 POSIX 层是\emph{源码}兼容，不是二进制兼容。
\end{note}

\subsection{内核如何分发系统调用}
\label{sec:syscall-dispatch}

\code{svc} 进入 \code{el0sync}，\code{usavereg} 把全部用户寄存器存进 trapframe，\code{usertrap()} 检查 \reg{ESR\_EL1} 的异常类别是否为 \code{0x15}（来自 AArch64 的 SVC），然后打开中断、调用 \code{syscall()}（第~\ref{sec:irq} 节）。分发函数在 \file{kernel/syscall.c} 中：

\begin{ccode}
void syscall(void)
{
  int num;
  struct proc *p = myproc();

  num = p->trapframe->x7;
  if(num > 0 && num < NELEM(syscalls) && syscalls[num]) {
    p->trapframe->x0 = syscalls[num]();
  } else {
    printf("%d %s: unknown sys call %d\n", p->pid, p->name, num);
    p->trapframe->x0 = -1;
  }
}
\end{ccode}

\code{syscalls[]} 是一张以系统调用号为下标的函数指针表。处理函数先检查调用号是否越界，越界就返回 $-1$，而不是像教学内核那样直接停机——一个用户程序传错调用号，不应该让整个系统死掉。

参数不需要搬运：用户态执行 \code{svc} 时它们就在 \code{x0}～\code{x5} 中，现在原样躺在 trapframe 里，各个 \code{sys\_*} 函数用 \code{argint()}/\code{argaddr()}/\code{argstr()} 按编号取出：

\begin{ccode}
static uint64 argraw(int n)
{
  struct proc *p = myproc();
  switch (n) {
  case 0: return p->trapframe->x0;
  case 1: return p->trapframe->x1;
  ...
  case 5: return p->trapframe->x5;
  }
  panic("argraw");
  return -1;
}
\end{ccode}

指针参数只取出数值，不在这里检查合法性；真正访问用户内存时由 \code{copyin()}/\code{copyout()} 查用户页表，非法地址会使系统调用返回 $-1$，而不会让内核出错。

返回值写回 \code{p->trapframe->x0}，而不是直接放进 \code{x0} 寄存器。原因是返回用户态时 \code{urestorereg} 会从 trapframe 恢复\emph{所有}寄存器，直接写寄存器的值会被覆盖；只有改写保存的那一份，用户程序在 \code{svc} 之后才能在 \code{x0} 中看到它。

\subsection{两个栈指针与返回级别}

EL0 和 EL1 各有自己的栈指针。从用户态陷入时，硬件立刻把 \code{sp} 换成 \reg{SP\_EL1}（内核栈），用户栈指针原样留在 \reg{SP\_EL0} 中。即使内核从不使用 \reg{SP\_EL0}，也必须把它连同其他寄存器一起保存和恢复：陷入期间可能发生进程切换，下一个返回用户态的进程需要的是\emph{它自己}的用户栈指针。所以 \code{usavereg} 把 \reg{SP\_EL0} 存进 trapframe 的 \code{sp} 字段，\code{urestorereg} 把它写回。

内核态的 \code{savereg} 则不需要恢复 \code{sp}：内核栈是同一个，即使中间发生了进程切换，\code{swtch()} 也已经把 \code{sp} 换成了正确的值。

\code{eret} 回到哪一级，由 \reg{SPSR\_EL1} 的模式字段决定，不需要另外指定，所以总是回到异常发生时所在的那一级。利用这一点，内核可以“伪造”一次从未发生过的异常来进入用户态：\code{userinit()} 在第一个进程的 trapframe 中写入 \code{elr = 0}、\code{spsr = 0}（模式 EL0t）、\code{sp = PGSIZE}，\code{usertrapret()} 一执行 \code{eret}，CPU 就第一次来到 EL0（第~\ref{chap:process} 章）。启动时从 EL2 降到 EL1 用的也是同一个技巧（第~\ref{chap:boot} 章）。

\section{从 \code{init} 到 shell}

第一个用户进程 \code{initcode} 执行 \code{exec("/init")}（第~\ref{chap:process} 章）。\file{user/init.c} 做的事情相当于一个嵌入式 Linux 的启动脚本：

\begin{enumerate}
  \item 创建 \code{/dev}，用 \code{mknod} 建立 \code{console}、\code{tty}、\code{ttyS0}、\code{input/event0}、\code{video0} 设备节点；以 \code{/dev/ttyS0} 打开 fd 0，复制为 fd 1、2；
  \item 创建 \code{/proc}、\code{/boot}、\code{/mnt/ext2}、\code{/etc}、\code{/root}、\code{/bin} 等目录，为根目录下的程序在 \code{/bin} 中建立硬链接；
  \item 首次启动时生成 \code{/etc/hostname}、\code{/etc/fstab}、\code{/etc/passwd}（默认账号 \code{root}，密码 \code{xv6}）、\code{/etc/rc}；
  \item 按 \code{/etc/fstab} 逐行调用 \code{mount()}；
  \item 如果存在 \code{/etc/wifi.conf}，先运行 \code{/bin/wifi} 连接无线网络；
  \item 在一个循环里 fork 并 exec \code{/bin/login}；\code{login} 退出后回收它，再启动一个新的登录提示；同时回收所有过继来的孤儿进程。
\end{enumerate}

\code{login} 校验密码后调用 \code{tty\_attach()} 建立会话，切到 passwd 中指定的家目录，exec 指定的 shell。\code{sh} 在 xv6 原版基础上增加了：读取 \code{/etc/rc}、\code{PATH} 搜索（找不到时再回退到根目录，作为存储异常时的恢复路径）、\code{export}、\code{exit}/\code{logout} 内建命令，以及把外部命令设为前台进程组的作业控制。

\section{用户程序一览}

\begin{table}[htbp]
  \centering
  \caption{用户程序（\code{UPROGS}）}
  \label{tab:uprogs}
  \small
  \begin{tabularx}{\linewidth}{@{}lX@{}}
    \toprule
    类别 & 程序 \\
    \midrule
    基本工具 & \code{sh}、\code{init}、\code{login}、\code{ls}、\code{cat}、\code{echo}、\code{grep}、\code{wc}、\code{mkdir}、\code{rm}、\code{ln}、\code{mv}、\code{touch}、\code{kill}、\code{ps}、\code{edit}、\code{file} \\
    文件系统 & \code{ext2ls}、\code{ext2cat}、\code{stressfs} \\
    网络 & \code{wifi}、\code{dhcp}、\code{ping}、\code{netdns}、\code{nettest}、\code{tftp}、\code{tcpd}、\code{epollserver}、\code{sshd} \\
    调度与同步 & \code{chrt}、\code{taskset}、\code{nice}、\code{renice}、\code{rtlat}、\code{spin}、\code{syncdemo}、\code{prodcons} \\
    线程与 POSIX & \code{threadtest}、\code{posixtest} \\
    设备 & \code{camshot}、\code{ptytest}、\code{vmmap} \\
    开发 & \code{tcc} \\
    测试 & \code{usertests}、\code{forktest}、\code{grind}、\code{zombie} \\
    \bottomrule
  \end{tabularx}
\end{table}

几个调度工具的用法与 Linux 相同，它们都是“先修改自己的调度属性，再 exec 目标程序”，利用了 \code{fork}/\code{exec} 对调度属性的继承：

\begin{shcode}
chrt -f 80 rtlat          # 以 SCHED_FIFO 优先级 80 运行 rtlat
taskset 0x8 spin &        # 只允许在 CPU3 上运行
nice -n 10 grind          # 降低普通进程的时间片
chrt -f -p 60 6           # 修改 PID 6（某个 kworker）的策略
\end{shcode}

\section{用户线程}

\subsection{设计}

原来的模型是“一个 \code{struct proc} 对应一个独立的地址空间”。加入线程时，项目没有引入新的“线程控制块”，而是让每个线程仍然是一个完整的 \code{struct proc}——有自己的内核栈、trapframe、调度属性、TID 和 CPU 亲和性——只是同一线程组的成员共享同一个引用计数的 \code{struct vmspace}（图~\ref{fig:threads}）。

\begin{figure}[htbp]
  \centering
  \begin{tikzpicture}[node distance=4mm and 14mm]
    \node[box, text width=30mm] (p) {leader\\TID $=$ TGID};
    \node[box, below=of p, text width=30mm] (t1) {线程 1\\独立内核栈/trapframe};
    \node[box, below=of t1, text width=30mm] (t2) {线程 2\\独立内核栈/trapframe};
    \node[hbox, right=of t1, text width=36mm] (vm) {\code{struct vmspace}\\TTBR0 页表、sz\\refcount $=3$};
    \draw[arr] (p.east) -- (vm); \draw[arr] (t1) -- (vm); \draw[arr] (t2.east) -- (vm);
  \end{tikzpicture}
  \caption{线程组共享 \code{vmspace}}
  \label{fig:threads}
\end{figure}

调度器切换到任何一个线程时，\code{switchuvm()} 都写入 \code{p->vm->pagetable}，所以同组线程看到相同的代码、全局变量、堆和彼此的栈；而它们的内核栈和 trapframe 各自独立，可以同时在多个 Cortex-A53 核上运行。

\subsection{系统调用}

\begin{table}[htbp]
  \centering
  \caption{线程相关系统调用}
  \small
  \begin{tabularx}{\linewidth}{@{}lX@{}}
    \toprule
    调用 & 作用 \\
    \midrule
    \code{clone(entry, arg, stack\_top)} & 创建共享地址空间的新调度实体：\code{x0}=arg，\code{SP}=stack\_top，\code{ELR}=entry \\
    \code{texit(status)} & 只退出调用线程；leader 调用时等同于整个进程退出 \\
    \code{tjoin(tid, \&status)} & 等待并回收同组中的指定线程（0 表示任意） \\
    \code{gettid()} / \code{getpid()} & 返回线程 ID / 线程组 ID \\
    \bottomrule
  \end{tabularx}
\end{table}

用户库在其上提供更安全的包装，它负责分配 16 字节对齐的 16\,KiB 用户栈，并保证线程函数返回后自动调用 \code{texit()}：

\begin{ccode}
static void thread_boot(void *opaque) {
  struct thread_start start = *(struct thread_start *)opaque;
  free(opaque);                       // 启动记录用完即释放
  start.fn(start.arg);
  texit(0);                           // 线程函数返回后自动退出
}

int thread_create(thread_t *thread, void (*fn)(void *), void *arg) {
  thread->stack = malloc(THREAD_STACK_SIZE + 16);
  struct thread_start *start = malloc(sizeof(*start));
  start->fn = fn;  start->arg = arg;
  uint64 top = ((uint64)thread->stack + THREAD_STACK_SIZE + 16) & ~15ULL;
  thread->tid = clone(thread_boot, start, (void *)top);
  ...
}
\end{ccode}

用户栈要等 \code{thread\_join()} 中内核回收了线程之后才释放，否则线程可能还在用这块栈。

\subsection{\code{clone} 为什么“返回两次”}

\code{fork()} 会返回两次：父进程中返回子进程的 pid，子进程中返回 0，因为子进程的 trapframe 是父进程的完整拷贝，只把 \code{x0} 改成了 0。Linux 的 \code{clone} 系统调用语义相同，所以 glibc 的 \code{clone()} 包装函数要在汇编里处理两次返回：调用前把函数指针和参数存进不会被系统调用破坏的寄存器；系统调用返回非 0 说明还在原线程，直接返回调用者；返回 0 说明已经身处新线程，就取出函数指针调用它，函数结束后再调用 \code{exit}——新线程永远不会回到包装函数的调用者那里。

xv6-aarch64 的 \code{clone(entry, arg, stack\_top)} 换了一种更简单的语义：内核直接把新线程的入口、参数和栈写进它的 trapframe，新线程第一次回到用户态时就从 \code{entry} 开始执行，用户态不需要任何汇编包装：

\begin{ccode}
int threadclone(uint64 entry, uint64 arg, uint64 stack_top)
{
  ...
  if(entry == 0 || stack_top < 16 || (stack_top & 15) != 0)
    return -1;                          // 栈必须 16 字节对齐
  ...
  vm->refcount++;                       // 共享同一个地址空间
  vmspace_put(np->vm);
  np->vm = vm;
  np->is_thread = 1;
  np->tgid = p->tgid;

  *(np->trapframe) = *(p->trapframe);
  np->trapframe->elr = entry;           // eret 之后从这里开始
  np->trapframe->sp  = stack_top;       // 新的用户栈
  np->trapframe->x0  = arg;             // 第一个参数
  np->trapframe->x30 = 0;               // 不允许“返回”
  ...
  make_runnable(np, 0);
  return tid;
}
\end{ccode}

\code{x30} 被清零，线程函数如果直接 \code{ret} 就会跳到地址 0 出错，所以用户库的 \code{thread\_boot()} 包了一层，保证函数返回后调用 \code{texit()}。线程退出后变成 ZOMBIE，等 \code{tjoin()} 回收——与进程的 \code{exit}/\code{wait} 是同一种约定：先保留“尸体”，让等待者能取到退出状态，再由等待者释放资源。

\subsection{多核下的细节}

\begin{itemize}
  \item \code{sbrk()} 修改共享的页表时持有 \code{vmspace.lock}，改完执行 \code{tlbi vmalle1is} 广播失效，覆盖其他核上正在运行的同组线程；
  \item 用户态的 K\&R \code{malloc} 加了一把基于 \code{ldaxr}/\code{stxr} 的锁；
  \item \code{exec()} 只在线程组只剩调用者时才允许执行；\code{exit()} 先终止并回收同组线程；\code{kill(pid)} 以 TGID 为目标，标记并唤醒整组；
  \item \code{clone()} 像 \code{fork()} 一样复制文件描述符\emph{引用}——底层 \code{struct file} 和偏移共享，但描述符表独立。这与 Linux 的 \code{CLONE\_FILES} 不同，是有意保留的简化。
\end{itemize}

\code{threadtest} 创建两个线程，用内核同步对象保护共享计数器，允许它们被调度到不同 CPU，预期输出 \code{threadtest: PASS counter=2000}。

\section{libc 与 POSIX 源码兼容层}
\label{sec:posix}

\subsection{\code{libc.a}}

用户库被打包成一个静态库：

\begin{itemize}
  \item \file{ulib.c}：字符串、内存函数、\code{gets}、\code{stat}、\code{atoi}；
  \item \file{usys.S}：系统调用桩；
  \item \file{printf.c}：\code{printf}/\code{fprintf}，支持 \code{\%d \%x \%p \%s \%c}；
  \item \file{umalloc.c}：K\&R 风格的 \code{malloc}/\code{free}，用 \code{sbrk} 向内核要内存；
  \item \file{thread.c}：上一节的线程包装；
  \item \file{posix.c}：POSIX 包装层；
  \item \file{tweetnacl.c}、\file{ssh\_crypto.c}：密码学原语（\code{sshd} 使用）。
\end{itemize}

\subsection{POSIX 层}

目标是让只依赖常见 POSIX 接口的 C 源码能直接用 xv6 工具链重新编译。它由三部分组成：

\begin{enumerate}
  \item \file{include/} 提供标准名字的头文件：\code{stdio.h}、\code{stdlib.h}、\code{string.h}、\code{unistd.h}、\code{fcntl.h}、\code{errno.h}、\code{time.h}、\code{pthread.h}、\code{semaphore.h}、\code{sys/types.h}、\code{sys/stat.h}、\code{sys/socket.h}、\code{netinet/in.h}、\code{arpa/inet.h}、\code{sys/epoll.h}；
  \item \file{user/posix.c} 把标准 API 映射到 xv6 已有的能力；
  \item 内核补上无法只靠库函数实现的语义，如 \code{lseek()}，以及 \code{socket()} → \code{bind()} → \code{listen()} → \code{accept()} 的生命周期。
\end{enumerate}

\begin{table}[htbp]
  \centering
  \caption{POSIX 接口到 xv6 的映射}
  \small
  \begin{tabularx}{\linewidth}{@{}lX@{}}
    \toprule
    POSIX 接口 & 实现 \\
    \midrule
    \code{open/read/write/lseek/fstat} & 直接对应系统调用；\code{O\_CREAT} 是 \code{O\_CREATE} 的别名 \\
    \code{pthread\_create/join/exit} & \code{clone}/\code{tjoin}/\code{texit}；\code{pthread\_join} 同时回收用户栈 \\
    \code{pthread\_mutex\_*}、\code{sem\_*} & 内核同步对象；支持 \code{PTHREAD\_MUTEX\_INITIALIZER} 延迟初始化 \\
    \code{clock\_gettime(CLOCK\_MONOTONIC)} & \code{clock\_us()}，读 \reg{CNTVCT\_EL0} \\
    \code{sleep/usleep/nanosleep} & 秒语义；原来按 tick 计的 \code{sleep} 改名为 \code{sleep\_ticks} \\
    \code{socket/bind/listen/accept} & 统一的 socket 文件描述符，可用 \code{read/write/close/epoll} \\
    \bottomrule
  \end{tabularx}
\end{table}

\code{posixtest} 只包含标准头文件，验证文件读写与 \code{lseek}、POSIX 信号量、两个 pthread 并发更新共享变量、静态初始化的互斥锁以及 socket 生命周期。

这仍是第一阶段：系统调用失败只返回 $-1$，没有传回具体错误码，\code{errno} 也还不是线程局部的；缺少 \code{waitpid}、\code{dup2}、条件变量、目录流和 TCP \code{connect()}。

\section{在 xv6 里编译 C 程序}

\code{tcc} 是一个面向教学的极小 C 编译器，运行在 xv6 用户态，直接输出 \code{exec()} 能加载的 ELF64 文件，不需要汇编器和链接器：

\begin{shcode}
$ edit hello.c
int main() {
  int answer = 6 * 7;
  printf("hello from xv6 C\n");
  return answer - 42;
}
$ tcc hello.c -o hello
$ hello
hello from xv6 C
\end{shcode}

它支持整数局部变量、四则运算、\code{return}、只含字面量的 \code{printf}/\code{puts} 和注释。生成的代码用 \code{x29} 指向 256 字节的局部变量区，表达式结果放在 \code{x0}，字符串地址用 \code{adr} 指令取得，所以无需重定位；输出和退出直接用 \code{svc \#0} 调用 \code{SYS\_write} 和 \code{SYS\_exit}。文件布局是一个 ELF 头、一个程序头、机器码和字符串池。

一个能在自己身上编译并运行程序的系统，才算是有了一点“自举”的味道。

\section{一个大型用户程序：\code{sshd}}

\code{sshd} 有一千多行，是整个系统里最复杂的用户程序，也是对内核各子系统的一次综合考验。它用到了：

\begin{itemize}
  \item TCP socket 与 \code{accept}——网络协议栈；
  \item 密码学全部在用户态完成：基于 TweetNaCl\cite{tweetnacl} 的 X25519 密钥交换和 Ed25519 主机密钥签名，传输层用 AES-128-CTR 加密、HMAC-SHA256 校验；熵来自内核的硬件随机数发生器（\file{kernel/rng.c}）；
  \item \code{/etc/passwd} 口令认证；
  \item PTY：为每个会话分配一对伪终端，shell 的 fd 0/1/2 指向 PTY 从端，\code{sshd} 读写主端并加密转发；
  \item 会话与控制终端：SSH 会话中的 \code{/dev/tty} 指向该会话的 PTY。
\end{itemize}

调试 \code{sshd} 时发现的问题推动了内核的多处改进：字符回显慢暴露了“网卡每 100\,ms 轮询一次”和“TCP \code{write()} 同步等 ACK”两个瓶颈，促成了 SDIO 中断、workqueue 下半部、NAPI 和 TCP 异步发送队列；\code{ps} 输出跑到串口上，促成了 \code{sys\_ps()} 改为返回缓冲区。

\begin{remark}
  从 \code{\_entry} 的第一条 \code{mrs} 指令，到一个可以通过 Wi-Fi 被 SSH 登录、能在多核上运行实时任务、能用 \code{tcc} 编译程序的系统，中间经过的每一层——串口、页表、进程、文件系统、驱动、协议栈、libc——都写在这本书里了。它们都还很简单，但每一层都在真机上跑过、出过错、被修好过。这也许就是“从裸机到操作系统”最有价值的部分。
\end{remark}

\chapter{可选工作}
\label{chap:optional}

\begin{introduction}[本章内容]
  \item 为什么要在 xv6 里读写 ext2，分区与日志放在哪里
  \item ext2 的磁盘结构：superblock、块组、inode、目录
  \item 从文件偏移到块号：直接块与三级间接块
  \item 接入 VFS：按 inode 打开、孤儿 inode
  \item 持久化：\code{sdflush()} 写屏障
  \item 外置事务日志 xjournal：格式、提交协议、崩溃恢复
  \item 用 WAL 程序和随机断电做测试
\end{introduction}

前面八章是系统的主线：没有它们，内核就启动不了、跑不了程序、用不了网络。本章收录的是主线之外的工作。它们不影响系统运行，但能把某个子系统做得更完整，也适合作为读者自己动手的练习。第一项是 ext2：让 xv6 能可靠地读写一块 Linux 也认识的分区，并在断电后保持一致。代码在 \file{kernel/ext2.c}、\file{kernel/xjournal.c} 和 \file{user/waltest.c}，说明文档是 \file{readme/readme\_ext2.md}。

\section{ext2 读写与外置事务日志}
\label{sec:ext2}

\subsection{动机与分区布局}

xv6 的根文件系统（第~\ref{chap:fs} 章）只有 xv6 自己认识，开发机上的 macOS 或 Linux 无法直接打开它。要在两边交换文件，需要一种双方都懂的格式。FAT32 启动分区可以，但它只支持根目录、没有权限位，平时也不该频繁写。ext2 是 Linux 最经典的文件系统，结构简单、文档齐全，非常适合在教学内核里实现。

第~\ref{sec:vfs} 节介绍 VFS 时，ext2 后端还只能读：只支持直接块和一级间接块，默认按只读挂载。本节把它补成一个可写的后端，并解决可写之后立刻出现的问题：写到一半断电怎么办。

ext2 放在 SD 卡的第三个分区，xv6 根文件系统仍在 p2，二者互不替代（图~\ref{fig:ext2-card}）。ext2 本身没有日志，这里给它配了一个\emph{外置}的事务日志，放在 p2 尾部：p2 有 400\,MiB，xv6 镜像只用前 32\,MiB（\code{FSSIZE} $=$ 32768 块，即 65536 扇区），后面的空间本来就闲着。

\begin{figure}[htbp]
  \centering
  \begin{tikzpicture}[x=1mm,y=1mm, font=\scriptsize]
    \tikzset{seg/.style={draw=structurecolor!70!black, minimum height=10mm, inner sep=1pt, align=center, anchor=west}}
    \node[lab, anchor=east] at (-1,30) {整张卡};
    \node[seg, fill=gray!20, minimum width=10mm] (mbr) at (0,30) {MBR};
    \node[seg, fill=structurecolor!20, minimum width=30mm] (p1) at (mbr.east) {p1 FAT32\\\code{/boot}};
    \node[seg, fill=orange!15, minimum width=46mm] (p2) at (p1.east) {p2 xv6 原始分区\\\code{/}，LBA 208896 起};
    \node[seg, fill=green!12, minimum width=40mm] (p3) at (p2.east) {p3 ext2\\\code{/mnt/ext2}，LBA 1030144 起};
    \node[lab, anchor=east] at (-1,8) {p2 内部};
    \node[seg, fill=orange!25, minimum width=46mm] (img) at (0,8) {xv6 根文件系统\\65536 扇区（32\,MiB）};
    \node[seg, fill=structurecolor!30, minimum width=14mm] (desc) at (img.east) {描述符\\9 扇区};
    \node[seg, fill=structurecolor!15, minimum width=34mm] (pay) at (desc.east) {载荷\\$\le$1024 扇区};
    \node[seg, fill=red!15, minimum width=12mm] (cr) at (pay.east) {提交\\记录};
    \node[seg, fill=gray!5, minimum width=20mm] (un) at (cr.east) {未用};
    \draw[gray!60, dashed] (p2.south west) -- (img.north west);
    \draw[gray!60, dashed] (p2.south east) -- (un.north east);
    \node[lab, anchor=north] at (desc.south west) {274432};
    \node[lab, anchor=north] at (cr.south) {275465};
    \draw[decorate, decoration={brace, amplitude=3pt, mirror}] ([yshift=-5mm]desc.south west) -- node[lab, below=3pt]{xjournal：共 1034 扇区，只保护 p3 的 ext2} ([yshift=-5mm]cr.south east);
  \end{tikzpicture}
  \caption{ext2 与外置日志在 SD 卡上的位置。下排数字是真机上的绝对 LBA：日志起点 $=$ p2 起点 208896 $+$ 65536}
  \label{fig:ext2-card}
\end{figure}

日志的位置在启动时算出。\code{ext2init()} 先向根块设备（第~\ref{sec:rootfs} 节）要 p2 的位置，再跳过 xv6 镜像：

\begin{ccode}
uint32 rlba, rsec, jstart = FSSIZE * (BSIZE / SECTOR);     // 65536
if(rootdev_raw_info(&rlba, &rsec) == 0 && rsec > jstart &&
   xj_init(rlba + jstart, rsec - jstart) == 0)               // 先恢复日志
  e2.journal = 1;
else
  printf("ext2: no raw xv6 root partition; external journal disabled\n");
\end{ccode}

\code{xj\_init()} 在解析任何 ext2 元数据之前就检查并重放日志，这样后面读到的 superblock 和位图都已经是一致的。QEMU 裸镜像没有 p2，这时日志关闭，ext2 退回逐扇区同步写。更新根文件系统的 \code{/dev/sdroot}（第~\ref{sec:xv6fs} 节）只写前 \code{FSSIZE} 块，不会碰到日志区。

\subsection{ext2 的磁盘结构}

ext2 把分区切成大小相同的\emph{块}（1、2 或 4\,KiB），再把连续的块分成若干\emph{块组}。每个块组都有自己的位图和 inode 表，这样文件的数据块可以和它的 inode 放在同一个组里，减少磁头移动——这是为机械硬盘设计的，但结构本身对 SD 卡同样适用。图~\ref{fig:ext2-layout} 是 4\,KiB 块时第 0 组的布局。

\begin{figure}[htbp]
  \centering
  \begin{tikzpicture}[x=1mm,y=1mm, font=\scriptsize]
    \tikzset{seg/.style={draw=structurecolor!70!black, minimum height=11mm, inner sep=1pt, align=center, anchor=west}}
    \node[seg, fill=gray!15, minimum width=22mm] (b0) at (0,20) {块 0\\引导 1\,KiB +\\superblock};
    \node[seg, fill=structurecolor!25, minimum width=16mm] (gdt) at (b0.east) {块 1\\组描述符表};
    \node[seg, fill=orange!25, minimum width=16mm] (bb) at (gdt.east) {块位图\\1 块};
    \node[seg, fill=orange!15, minimum width=16mm] (ib) at (bb.east) {inode 位图\\1 块};
    \node[seg, fill=green!15, minimum width=28mm] (it) at (ib.east) {inode 表\\每项 256 字节};
    \node[seg, fill=gray!5, minimum width=38mm] (data) at (it.east) {数据块\\文件内容、目录、间接块};
    \draw[decorate, decoration={brace, amplitude=3pt}] ([yshift=1mm]b0.north west) -- node[lab, above=3pt]{第 0 组（\code{blocks\_per\_group} 块，4\,KiB 块时为 32768 块 $=$ 128\,MiB）} ([yshift=1mm]data.north east);
    \node[box, font=\scriptsize, text width=46mm, align=left] (sb) at (24,0) {superblock（偏移 1024）\\块数、inode 数、块大小、每组块数、魔数 \code{0xEF53}、特性位};
    \node[box, font=\scriptsize, text width=52mm, align=left] (gd) at (88,0) {组描述符（每组 32 字节）\\块位图、inode 位图、inode 表的块号；空闲块数、空闲 inode 数};
    \draw[arr] (sb.north) -- (b0.south);
    \draw[arr] (gd.north) -- (gdt.south);
    \draw[darr] (gd.north east) to[bend right=20] (it.south);
  \end{tikzpicture}
  \caption{ext2 第 0 块组的布局。其他块组结构相同；开启 \code{sparse\_super} 时只有第 0、1 和 3、5、7 的幂次组保存 superblock 与描述符表的备份}
  \label{fig:ext2-layout}
\end{figure}

\code{setup\_ext2()} 从 superblock 读取表~\ref{tab:ext2-sb} 中的字段，并在用之前逐个检查。最重要的是三组\emph{特性位}：\code{compat} 中的 \code{has\_journal} 说明这是 ext3/ext4，\code{incompat} 中除 \code{filetype} 以外的任何一位（如 extent、64 位块号）都意味着本驱动看不懂磁盘布局。遇到这些就拒绝挂载，而不是把一个 ext4 分区当成 ext2 去写坏它。

\begin{table}[htbp]
  \centering
  \caption{\code{setup\_ext2()} 读取的 superblock 字段}
  \label{tab:ext2-sb}
  \small
  \begin{tabular}{@{}rll@{}}
    \toprule
    偏移 & 字段 & 用途 \\
    \midrule
    0 / 4 & \code{inodes\_count} / \code{blocks\_count} & 越界检查；块数 $\times$ 块大小不能超过分区 \\
    12 / 16 & 空闲块数 / 空闲 inode 数 & 分配、释放时加减 \\
    20 & \code{first\_data\_block} & 1\,KiB 块时为 1，否则为 0 \\
    24 & \code{log\_block\_size} & 块大小 $= 1024 \ll$ 该值 \\
    32 / 40 & 每组块数 / 每组 inode 数 & 块号、inode 号到组号的换算 \\
    56 & \code{magic} & 必须是 \code{0xEF53} \\
    58 & \code{state} & 第一次写时清除 \code{VALID\_FS} \\
    88 & \code{inode\_size} & 动态版本里一般为 256 \\
    92 / 96 / 100 & \code{compat} / \code{incompat} / \code{ro\_compat} & 拒绝 journal、extent、64bit 等 \\
    \bottomrule
  \end{tabular}
\end{table}

\paragraph{inode 与目录.} inode 号从 1 开始，2 号是根目录。找 $n$ 号 inode 时，先算组号 $g = (n-1)/\text{每组 inode 数}$ 和组内序号 $i = (n-1) \bmod \text{每组 inode 数}$，读第 $g$ 个组描述符拿到 inode 表的块号，inode 就在表内偏移 $i \times \text{inode\_size}$ 处：

\begin{ccode}
gd_off = (uint64)e2.gd_block * e2.block_size + group * e2.desc_size;
diskbytes(gd_off, inodebuf, e2.desc_size);
table = le32(inodebuf + 8);                                // inode 表块号
off = (uint64)table * e2.block_size + (uint64)index * e2.inode_size;
diskbytes(off, inodebuf, e2.inode_size);
ip->mode  = le16(inodebuf);        ip->links = le16(inodebuf + 26);
ip->size  = le32(inodebuf + 4);
if((ip->mode & EXT2_S_IFMT) == EXT2_S_IFREG)               // 普通文件的高 32 位
  ip->size |= (uint64)le32(inodebuf + 108) << 32;
for(int i = 0; i < 15; i++) ip->block[i] = le32(inodebuf + 40 + i*4);
\end{ccode}

\code{off} 必须用 64 位计算：一个几 GiB 的分区里，后面块组的 inode 表很可能位于 4\,GiB 之后，32 位乘法会悄悄溢出，读到另一个位置的“inode”。文件大小同理，高 32 位放在 inode 偏移 108 处。

目录是普通的数据块，里面是一串变长目录项：inode 号（4 字节）、本项长度 \code{rec\_len}（2 字节）、名字长度和文件类型（各 1 字节），然后是名字。删除一项时不移动后面的数据，而是把它的 \code{rec\_len} 并入前一项；新增时找一个 \code{rec\_len} 比实际需要大的项把它切开。\code{e2dirlookup()}、\code{e2diradd()}、\code{e2dirremove()}、\code{e2dirreplace()} 分别实现查找、插入、删除和改指向。

\subsection{从文件偏移到块号}

inode 的 \code{block[15]} 是一棵最多四层的树（图~\ref{fig:ext2-blocks}）：前 12 项直接指向数据块，第 13 项指向一个装满块号的\emph{间接块}，第 14、15 项分别多套一层和两层。4\,KiB 块时每个间接块能放 $n = 1024$ 个块号，所以一个文件最多有
\[
  12 + n + n^2 + n^3 = 12 + 1024 + 1\,048\,576 + 1\,073\,741\,824 \text{ 块} \approx 4\,\text{TiB}.
\]

\begin{figure}[htbp]
  \centering
  \begin{tikzpicture}[x=1mm,y=1mm, font=\scriptsize,
      cell/.style={draw=structurecolor!70!black, minimum width=8mm, minimum height=5mm, inner sep=0pt, fill=structurecolor!10},
      blk/.style={box, font=\scriptsize, minimum height=6mm, inner sep=1.5pt}]
    \node[lab] at (-6,40) {\code{block[]}};
    \foreach \i/\t in {0/0, 1/1, 2/\ldots, 3/11}
      \node[cell] (c\i) at (\i*8+6,40) {\t};
    \foreach \i/\t in {4/12, 5/13, 6/14}
      \node[cell, fill=orange!20] (c\i) at (\i*8+6,40) {\t};
    \node[blk, fill=gray!8] (d0) at (6,26) {数据};
    \node[blk, fill=gray!8] (d11) at (30,26) {数据};
    \draw[arr] (c0.south) -- (d0.north);
    \draw[arr] (c3.south) -- (d11.north);
    \node[blk, fill=orange!15] (s1) at (46,26) {间接块};
    \node[blk, fill=gray!8] (s1d) at (46,12) {数据 $\times n$};
    \draw[arr] (c4.south) -- (s1.north); \draw[arr] (s1) -- (s1d);
    \node[blk, fill=orange!15] (dd1) at (68,26) {二级间接};
    \node[blk, fill=orange!15] (dd2) at (68,12) {间接 $\times n$};
    \node[blk, fill=gray!8] (dd3) at (68,0) {数据 $\times n^2$};
    \draw[arr] (c5.south) to[out=-60,in=120] (dd1.north); \draw[arr] (dd1) -- (dd2); \draw[arr] (dd2) -- (dd3);
    \node[blk, fill=orange!15] (t1) at (96,40) {三级间接};
    \node[blk, fill=orange!15] (t2) at (96,26) {二级 $\times n$};
    \node[blk, fill=orange!15] (t3) at (96,12) {间接 $\times n^2$};
    \node[blk, fill=gray!8] (t4) at (96,0) {数据 $\times n^3$};
    \draw[arr] (c6.east) -- (t1.west); \draw[arr] (t1) -- (t2); \draw[arr] (t2) -- (t3); \draw[arr] (t3) -- (t4);
    \node[lab, align=left, anchor=west] at (106,26) {逻辑块 $\ge 12+n+n^2$\\走这一支};
  \end{tikzpicture}
  \caption{ext2 inode 的块指针树。块号为 0 表示空洞：读时返回全 0，写时才分配}
  \label{fig:ext2-blocks}
\end{figure}

\code{fileblock()} 按逻辑块号依次减去每一层能覆盖的块数，落到哪一层就沿哪一支往下读：

\begin{ccode}
static uint32 fileblock(struct einode *ip, uint64 logical) {
  uint64 n = e2.block_size / 4;
  if(logical < 12) return ip->block[logical];
  logical -= 12;
  if(logical < n){                                          // 一级间接
    if(ip->block[12] == 0 || readblock(ip->block[12], ptrbuf[0]) < 0) return 0;
    return ptrbuf[0][logical];
  }
  logical -= n;
  if(logical < n*n){ ... 读 block[13]，再读 ptrbuf[0][logical / n] ... }
  logical -= n*n;                                           // 三级间接
  ... 读 block[14]，再依次取 logical / n²、(logical % n²) / n、logical % n ...
}
\end{ccode}

三层用三个独立的 4\,KiB 缓冲 \code{ptrbuf[0..2]}，它们是全局变量而不是局部数组——内核栈只有 4\,KiB，放不下其中任何一个（第~\ref{chap:drivers} 章的摄像头调试记录里有过同样的教训）。所有 ext2 操作都持着同一把锁，所以全局缓冲是安全的。

写的方向由 \code{set\_fileblock()} 和 \code{get\_or\_alloc\_fileblock()} 完成：路径上缺哪一层间接块就分配哪一层。截断时 \code{clear\_fileblock()} 释放数据块，再检查它所在的间接块是否已经全为 0（\code{ptr\_empty()}），是就把间接块也释放掉，逐层向上。\code{waltest large} 在偏移 5\,GiB 处写一个字节：$5\,\text{GiB}/4\,\text{KiB} = 1\,310\,720 > 12 + n + n^2$，正好走到三级间接这一支，文件只占 4 个块。

\subsection{分配与释放}

分配一个数据块要改三处：块位图里的一位、组描述符里的空闲块数、superblock 里的空闲块数。\code{alloc\_block()} 从首选的组开始找第一个空闲位：

\begin{ccode}
for(uint32 bit = 0; bit < count; bit++){
  if((blockbuf[bit >> 3] & (1U << (bit & 7))) != 0) continue;
  if(freed_in_tx(first + bit)) continue;        // 本事务刚释放的块，暂不复用
  blockbuf[bit >> 3] |= 1U << (bit & 7);
  writeblock(bitmap, blockbuf);                 // 位图
  put16(gd + 12, le16(gd + 12) - 1);
  writegd(group, gd);                           // 组描述符
  adjust_super_count(12, -1);                   // superblock
  *out = first + bit;
  memset(blockbuf, 0, e2.block_size);
  writeblock_data(*out, blockbuf);              // 新块清零，直接写
  return 0;
}
\end{ccode}

两处细节和后面的日志有关。第一，\code{freed\_in\_tx()}：一个块如果在\emph{同一个事务里}刚被释放，就不能立刻再分配出去并原地改写——事务一旦中止或者断电，释放会被撤销，这个块又回到原主人手里，内容却已经是别人的了。第二，新块直接清零而不经过日志：在事务提交之前，没有任何 inode 指向它，写什么都不会被看见。

inode 的分配（\code{alloc\_inode()}）和释放（\code{free\_inode()}）是同样的三步，只是改的是 inode 位图和空闲 inode 数；目录还要维护组描述符里的目录计数。

\subsection{接入 VFS：按 inode 打开}

第~\ref{sec:vfs} 节的 VFS 最初是“按路径”的：每次 \code{read}、\code{write} 都把路径传给后端重新查找。对 ext2 这会带来两个问题：每次读写都要从根目录走一遍；更糟的是，文件被重命名或删除以后，已经打开的描述符会突然指向别的文件，或者失效。Unix 的语义是：打开的是 inode，而不是名字。

为此 \code{vnode\_ops} 增加了一组可选的\emph{句柄}操作：

\begin{ccode}
struct vnode_ops {
  ...                                                       // 原来的按路径操作
  int  (*open)(char*, uint64*);                  // 解析一次路径，返回句柄
  void (*release)(uint64);                       // 最后一次 close
  int  (*hstat)(uint64, struct stat*);
  int  (*hread)(uint64, uint64, void*, int);     // 句柄、偏移、缓冲、长度
  int  (*hwrite)(uint64, uint64, int, uint64, int);
  int  (*htruncate)(uint64, uint64);
};
\end{ccode}

\code{vfsopen()} 打开普通文件时，如果后端提供了 \code{open}，就调用它得到句柄并记在 vnode 里，之后的 \code{read}/\code{write}/\code{fstat}/\code{ftruncate} 都走句柄版本。ext2 的句柄就是 inode 号；\code{ext2open()} 还在一个小表 \code{e2.open[]} 里给这个 inode 记一个引用计数。

有了引用计数，就能实现“删除仍被打开的文件”：\code{unlink} 或被 \code{rename} 覆盖时，目录项立即消失；如果 inode 的链接数降到 0 而它还被打开着，\code{drop\_link()} 只把它标成\emph{孤儿}，块不释放。旧描述符照常读写原来的数据，最后一次 \code{close} 时 \code{ext2release()} 才在一个事务里截断并释放 inode。

\code{rename} 也补全了 POSIX 语义：目标已存在时原子地替换它（类型必须相同，目录必须为空）；把目录移进它自己的子树会被 \code{is\_ancestor()} 拒绝——否则这棵子树会从根上脱落，形成一个环。

\paragraph{锁.} 所有 ext2 操作由睡眠锁 \code{e2.lock} 串行化。最初用的是自旋锁，但 SD 卡是轮询 PIO，一次 \code{fsync} 最长可能等 5 秒，持自旋锁期间 CPU 一直关着中断，键盘和网络都会卡住。换成睡眠锁后，等待的进程会让出 CPU。

\subsection{持久化：\code{sdflush()} 写屏障}

\code{sdsector(lba, buf, 1)} 返回时，CMD24 的数据已经传给卡，但卡可能还在内部把它编程进 NAND。这时断电，这个扇区可能是新的、旧的，甚至是坏的。要建立“之前的写都已落盘”这样一个边界，需要单独的屏障：

\begin{ccode}
int sdflush(void) {                                         // kernel/sd.c
  acquire(&sdlock);                                         // 与所有传输互斥
  asm volatile("dsb sy" ::: "memory");
  uint64 deadline = r_cntvct_el0() + 5ULL * r_cntfrq_el0();
  for(;;){
    uint32 status = 0;
    if(send_command(CMD_INDEX(13) | CMD_RSPNS_48 | CMD_CRCCHK_EN | CMD_IXCHK_EN,
                    sd_rca << 16, &status) == 0 &&
       (status & (1U << 8)) && ((status >> 9) & 0xf) == 4)  // READY_FOR_DATA，TRAN 态
      return 0;                                             // （省略解锁与超时）
    ...
  }
}
\end{ccode}

CMD13 读卡的状态寄存器。第 8 位 \code{READY\_FOR\_DATA} 置位、第 9～12 位的状态回到 TRAN（4），说明卡已经处理完之前的写，可以接受新命令。因为 \code{sdflush()} 持有块设备锁，屏障之后的写也不可能越过它。SD 6.x 规范还定义了一个可选的写缓存，需要专门的 Flush Cache 命令；本驱动没有启用它，所以 CMD13 就是当前支持的卡的持久化边界。

xv6 原生日志（第~\ref{chap:fs} 章）也在四个位置用上了这个屏障：日志内容 → flush → 提交头 → flush → 写回原位 → flush → 清提交头 → flush。下面的 xjournal 用的是同一个思路。

\subsection{外置事务日志 xjournal}

\paragraph{要解决的问题.} 一次 \code{rename} 要改新目录的数据块、旧目录的数据块，覆盖已有文件时还要改目标 inode、位图和计数，一共可能十几个扇区。这些写落在卡上的不同位置，断电可能发生在任意两次之间，留下“两个名字指向同一个 inode”或者“块已分配但没人用”的状态。ext2 本身对此无能为力，只能事后让 \code{e2fsck} 修。

xjournal 的思路与 xv6 原生日志相同：先把一次操作要写的所有元数据扇区完整地写到一块独立的区域，写一个“提交”标记，然后才写回原位。断电后只要看到提交标记就重放一遍，没看到就当这次操作没发生。区别在于它工作在 \emph{512 字节扇区}一层，不依赖 xv6 的块缓存，也不知道 ext2 的格式，所以叫“外置”。

\paragraph{接口.} 客户（这里是 ext2）在每个修改型操作前后调用 \code{xj\_begin()} 和 \code{xj\_commit()}；中间的写调用 \code{xj\_write()}，它不写卡，而是把扇区存进事务缓冲；读调用 \code{xj\_read()}，命中缓冲时返回未提交的新内容，这样同一个操作后面的步骤能看到前面的修改。缓冲是 128 个物理页（1024 个扇区），启动时分配。ext2 只在两个函数里接入它：

\begin{ccode}
static int e2_rsector(uint32 lba, void *buf) {
  if(e2.journal && xj_read(lba, buf))                       // 先查事务缓冲
    return 0;
  return sdsector(lba, buf, 0);
}

static int e2_wsector(uint32 lba, const void *buf) {
  if(e2.journal && xj_active() && (!e2.direct || xj_contains(lba)))
    return xj_write(lba, buf);                              // 元数据：进事务
  return sdsector(lba, (void*)buf, 1);                      // 文件数据：直接写
}
\end{ccode}

\paragraph{ordered 模式.} 文件内容不进日志，由 \code{writeblock\_data()} 置 \code{e2.direct} 后直接写回原位。这借鉴的是 ext3 的 \code{data=ordered}：数据先落盘，引用它的元数据（inode 里的块号、文件大小）后提交。新分配的块在提交前没人能看到，所以永远不会出现“文件里有一段垃圾”；覆盖已有数据则不保证原子——这与 Linux 上的语义一致，需要原子更新的程序应该自己用 WAL（见下一小节的 \code{waltest}）。这样做让日志只承载元数据，一个事务 1024 个扇区的容量对绝大多数操作都绰绰有余。

\paragraph{磁盘格式与提交协议.} 日志区分三部分（图~\ref{fig:ext2-card}）：9 个扇区的\emph{描述符}（魔数 \code{XJD1}、序号、扇区数和每个扇区的原位 LBA，$16 + 4\times1024$ 字节恰好需要 9 个扇区）、最多 1024 个扇区的\emph{载荷}，以及 1 个扇区的\emph{提交记录}（魔数 \code{XJC1}、序号、扇区数、整个事务的校验和，再加一个对这 16 字节自身的校验和）。校验和用 32 位 FNV-1a，覆盖描述符和全部载荷。提交分四步，每步后面都有一次 \code{sdflush()}（图~\ref{fig:xj-commit}）：

\begin{figure}[htbp]
  \centering
  \begin{tikzpicture}[x=1mm,y=1mm, font=\scriptsize,
      st/.style={box, font=\scriptsize, minimum height=9mm, text width=24mm, align=center},
      fl/.style={draw=red!60!black, fill=red!10, rounded corners=1pt, font=\scriptsize, inner sep=1.5pt, minimum height=9mm, text width=7mm, align=center}]
    \node[st] (s1) at (14,20) {\textbf{1} 写描述符\\和载荷};
    \node[fl] (f1) at (32,20) {flush};
    \node[st, fill=red!15] (s2) at (50,20) {\textbf{2} 写提交记录\\（提交点）};
    \node[fl] (f2) at (68,20) {flush};
    \node[st] (s3) at (86,20) {\textbf{3} 载荷写回\\原位 LBA};
    \node[fl] (f3) at (104,20) {flush};
    \node[st] (s4) at (122,20) {\textbf{4} 清除\\提交记录};
    \node[fl] (f4) at (140,20) {flush};
    \foreach \a/\b in {s1/f1, f1/s2, s2/f2, f2/s3, s3/f3, f3/s4, s4/f4} \draw[arr] (\a) -- (\b);
    \draw[decorate, decoration={brace, amplitude=3pt, mirror}] (2,13) -- node[lab, below=3pt, align=center]{此处断电：没有有效提交记录\\→ 丢弃，操作“没发生”} (58,13);
    \draw[decorate, decoration={brace, amplitude=3pt, mirror}] (62,13) -- node[lab, below=3pt, align=center]{此处断电：提交记录有效\\→ 启动时重放 3、4 两步} (130,13);
    \node[lab, align=left, anchor=west] at (2,30) {之前 ordered 模式下直接写的文件数据，被第 1 步后的 flush 一起“压”到卡上};
  \end{tikzpicture}
  \caption{xjournal 的提交协议。四次 \code{sdflush()} 把四步严格排序；提交点之前和之后的断电有两种确定的结局}
  \label{fig:xj-commit}
\end{figure}

\begin{ccode}
int xj_commit(void) {                                       // kernel/xjournal.c，节选
  ...
  uint32 sum = txsum(count);                               // 描述符 + 载荷的 FNV
  // 1. 描述符和载荷；这次 flush 也保证此前直接写的数据已经落盘
  for(uint32 i = 0; i < XJ_DESC; i++)  sdsector(xj.start + i, desc + i*512, 1);
  for(uint32 i = 0; i < count; i++)    sdsector(xj.start + XJ_DESC + i, slot(i), 1);
  if(sdflush() < 0) goto fail;
  // 2. 提交记录：这个扇区落盘的那一刻，事务就生效了
  make_commit(seq, count, sum);
  if(sdsector(xj.start + XJ_DESC + XJ_CAP, xj.commit, 1) < 0 || sdflush() < 0)
    goto fail;
  // 3./4. 写回原位并清除提交记录；失败也没关系，下次启动会重放
  if(write_home(count) < 0 || clear_commit() < 0)
    printf("xjournal: checkpoint of seq=%d incomplete; replay at boot\n", seq);
  return 0;
fail:
  printf("xjournal: commit of seq=%d failed before commit point\n", seq);
  return -1;
}
\end{ccode}

提交点是一个扇区的写入。SD 卡保证单个扇区的写要么完成、要么不生效吗？并不完全保证，所以提交记录自带一个校验和：写坏的提交记录校验不过，等同于没有。

\paragraph{崩溃恢复.} \code{recover()} 在 \code{xj\_init()} 里执行，顺序是：读提交记录并校验它自身；读描述符，确认魔数、序号和扇区数与提交记录一致；读载荷，重算整个事务的校验和并与提交记录比较；全部通过才把载荷写回原位、清除提交记录。任何一步不符都丢弃——描述符或载荷写到一半时断电，校验和必然对不上。表~\ref{tab:xj-crash} 列出了各个断电时刻的结局。

\begin{table}[htbp]
  \centering
  \caption{断电时刻与恢复结果}
  \label{tab:xj-crash}
  \small
  \begin{tabularx}{\linewidth}{@{}lXl@{}}
    \toprule
    断电发生在 & 磁盘状态 & 启动时 \\
    \midrule
    第 1 步期间或之后 & 原位元数据未动；日志里有半个或一整个事务，但没有提交记录 & 无事可做，操作没发生 \\
    第 2 步期间 & 提交记录可能是旧的 0 扇区或写坏的扇区 & 校验失败，丢弃 \\
    第 3 步期间 & 原位一部分是新的、一部分是旧的 & 重放：\code{xjournal: replayed seq=N} \\
    第 4 步期间 & 原位已全部更新，提交记录还在 & 再重放一次（幂等），结果相同 \\
    第 4 步之后 & 一切都是新的 & 无事可做 \\
    \bottomrule
  \end{tabularx}
\end{table}

重放是幂等的：载荷是扇区的完整新内容，而不是“加 1”“置位”这样的增量，写一次和写两次结果一样。这让恢复本身被打断也没有关系，下次启动再来一遍即可。

\paragraph{ext2 怎样使用事务.} 每个修改型操作的外形都一样，以 \code{rename} 为例：

\begin{ccode}
int ext2rename(char *oldpath, char *newpath) {
  acquiresleep(&e2.lock);
  tx_begin();                                               // xj_begin()，清空 freed[]
  int r = tx_end(rename_locked(oldpath, newpath) == 0);     // 成功则提交，失败则中止
  releasesleep(&e2.lock);
  return r;
}
\end{ccode}

\code{tx\_end(0)} 调用 \code{xj\_abort()}，事务缓冲被直接丢掉，这次操作对磁盘的元数据没有留下任何痕迹。以前 \code{rename} 中途失败（比如目标目录满了又分配不到新块）会留下“新名字加上了、旧名字没删”的状态，有了事务就不会了。

少数操作会超过一个事务的容量：截断一个几 GiB 的文件要改成千上万个位图扇区。\code{tx\_checkpoint()} 在剩余空间不足 96 个扇区时，先把 inode 以\emph{当前已截断到的大小}写入事务并提交，再开一个新事务继续。每个分段的终点都是一个一致的状态：inode 的大小与它还引用的块相符，释放掉的块已经回到位图，只是文件还没截到最终长度。

最后，第一次写 ext2 时，\code{mark\_write\_session\_dirty()} 绕过日志直接清除 superblock 的 \code{VALID\_FS} 位并 flush。驱动还没有 \code{umount}，从来不会把它置回去，所以开发机上的 \code{e2fsck} 总会做一次完整检查。

\subsection{测试：WAL 程序与随机断电}

\code{user/waltest.c} 是一个极小的“数据库”：数据文件 \code{data.db} 里存一个值，修改时先写一条 redo 记录到 \code{redo.wal}。每条记录带魔数、事务号和校验和，提交顺序是：

\begin{txtcode}
写 redo.wal -> fdatasync -> 写 data.db -> fsync -> 清空 redo.wal -> fsync
\end{txtcode}

它同时检验了两件事：\code{fsync} 是否真的建立了持久化边界，以及文件系统操作的语义是否正确。子命令如表~\ref{tab:waltest}。

\begin{table}[htbp]
  \centering
  \caption{\code{waltest} 的子命令}
  \label{tab:waltest}
  \small
  \begin{tabularx}{\linewidth}{@{}lX@{}}
    \toprule
    命令 & 检查什么 \\
    \midrule
    \code{init} / \code{commit 42} / \code{show} & 正常提交与读取 \\
    \code{prepare 123}，断电，\code{recover} & 只持久化 WAL 后断电，重启必须输出 \code{REDO tx=… value=123} \\
    \code{fsops} & \code{mkdir}、\code{rename}、\code{unlink}、\code{ftruncate}、\code{pread}、\code{pwrite} \\
    \code{semantics} & rename 覆盖已有文件；目录不能移入自己的子树；删除仍被打开的文件后，旧描述符仍可读写 \\
    \code{large} & 5\,GiB 处写一个字节，走三级间接块 \\
    \code{jstress} & 无限循环：写文件、\code{rename} 覆盖、\code{unlink}，供随机断电使用 \\
    \bottomrule
  \end{tabularx}
\end{table}

断电测试在 QEMU 里做：准备一个与真机分区布局相同的 256\,MiB SD 镜像，运行 \code{waltest jstress}，在随机时刻直接杀掉 QEMU 进程，再从同一个镜像启动。每次重启后，把 p3 从镜像里取出来，在宿主机上运行 \code{e2fsck -fn}。

\begin{shcode}
dd if=sd.img of=p3.img bs=512 skip=215040 count=309248
e2fsck -fn p3.img          # 退出码 0 表示没有任何错误
\end{shcode}

结果：第一轮随机断电 14 次，其中 3 次正好落在“已提交、未写回”的窗口——这时如果不重放，\code{e2fsck} 会报错（退出码 4）；三次启动都打印了 \code{xjournal: replayed seq=…}，重放后全部无错误。其余 11 次断电前后都没有错误。拆分出 \file{rootdev.c} 以后又断电 4 次，1 次重放（\code{seq=85}，27 个扇区），\code{e2fsck} 全部通过。

\begin{txtcode}
xjournal: replayed seq=85 (27 sectors)
xjournal: ready lba=198656 capacity=1024 sectors
ext2: p2 has no ext superblock
ext2: p3 ready rw start=215040 sectors=309248 block=4096 groups=2
\end{txtcode}

（QEMU 镜像的 p2 较小，日志起点是 198656；真机上应为 274432。）

\subsection{在真机上使用}

p3 要格式化成不带 ext4 特性的 ext2。比第~\ref{chap:fs} 章“烧录 SD 卡”一节里的命令多指定几个特性开关更稳妥：

\begin{shcode}
diskutil unmount /dev/disk4s3
sudo mke2fs -F -t ext2 -b 4096 \
  -O filetype,sparse_super,large_file,^has_journal,^extent,^64bit,^metadata_csum \
  -L xv6-linux /dev/rdisk4s3
sudo e2fsck -fn /dev/rdisk4s3
\end{shcode}

\code{/etc/fstab} 里 ext2 一行要从 \code{ro} 改成 \code{rw}（\code{init} 不会覆盖已有的 \code{fstab}）。正常的启动日志：

\begin{txtcode}
xjournal: ready lba=274432 capacity=1024 sectors
ext2: p3 ready rw start=1030144 sectors=... block=4096 groups=...
vfs: mounted ext2 at /mnt/ext2 read-write
\end{txtcode}

\subsection{当前边界}

这仍是一个小型的 ext2 后端，与 Linux 的实现相比还缺不少：

\begin{itemize}
  \item 不支持硬链接、符号链接、扩展属性、配额、HTree 索引目录和 GPT 分区表；
  \item 不支持 ext3 日志、ext4 extent、64 位块号和元数据校验和；
  \item 外置日志只有本内核认识。如果日志里还有已提交未写回的事务，要先在 xv6 里启动一次让它重放，再把卡插到 Linux 上；
  \item 孤儿 inode 只记在内存里，没有 ext3 那样的磁盘孤儿链表。“删除仍被打开的文件”之后断电，这个 inode 要靠 \code{e2fsck} 回收；
  \item 还没有 \code{umount}，接回开发机时应运行 \code{e2fsck -f}。
\end{itemize}

这些都是不错的后续练习。其中最有价值的大概是磁盘孤儿链表和 \code{umount}：前者只需要在 superblock 的 \code{s\_last\_orphan} 和 inode 的 \code{dtime} 字段里串一条链，在同一个事务里维护；后者让 \code{VALID\_FS} 能在干净卸载时置回去，开发机就不必每次都做完整检查。


\nocite{*}
\printbibliography[heading=bibintoc, title=参考文献]

\end{document}
